% ============================================================================
% KARADUL (BLACK WIDOW) — Teknik Dokuman
% Tersine Muhendislik ve Statik Analiz Araci
% Berke Tezgocen — Fizik Muhendisligi, ITU 2026
% ============================================================================
\documentclass[12pt,a4paper,oneside]{book}

% --- Encoding & Fonts ---
\usepackage[utf8]{inputenc}
\usepackage[T1]{fontenc}
\usepackage{lmodern}
\usepackage{microtype}

% --- Page Geometry ---
\usepackage[
  a4paper,
  top=2.5cm,
  bottom=2.5cm,
  left=2.5cm,
  right=2.5cm,
  headheight=14pt,
]{geometry}

% --- SIMSIYAH ARKA PLAN ---
\usepackage{pagecolor}
\usepackage{afterpage}
\pagecolor{black}
\color{white}

% --- Math ---
\usepackage{amsmath,amssymb,amsthm,mathtools}
\usepackage{bm}

% --- Graphics & Figures ---
\usepackage{graphicx}
\usepackage{tikz}
\usetikzlibrary{
  arrows.meta,
  positioning,
  shapes.geometric,
  shapes.misc,
  fit,
  calc,
  decorations.pathreplacing,
  backgrounds,
  chains,
  matrix,
}
\usepackage{pgfplots}
\pgfplotsset{compat=1.18}

% --- Tables ---
\usepackage{booktabs}
\usepackage{longtable}
\usepackage{array}
\usepackage{multirow}
\usepackage{colortbl}

% Booktabs cizgi renklerini beyaz yap (siyah arka planda gorunsun)
\arrayrulecolor{white!70!black}
\setlength{\heavyrulewidth}{0.12em}
\setlength{\lightrulewidth}{0.06em}

% TABLO ICINDEKI YAZILARI BEYAZ YAP
% rowcolors varsayilan yazi rengini siyah yapiyor — bunu override et
\usepackage{etoolbox}
\AtBeginEnvironment{tabular}{\color{white}}
\AtBeginEnvironment{longtable}{\color{white}}

% --- Code Listings ---
\usepackage{listings}
\lstset{
  language=Python,
  basicstyle=\ttfamily\small\color{white},
  keywordstyle=\color{cyan}\bfseries,
  stringstyle=\color{green!70!white},
  commentstyle=\color{gray!70!white}\itshape,
  numberstyle=\tiny\color{gray!50!white},
  numbers=left,
  numbersep=8pt,
  frame=single,
  rulecolor=\color{gray!50!white},
  backgroundcolor=\color{white!12!black},
  breaklines=true,
  breakatwhitespace=false,
  showstringspaces=false,
  tabsize=4,
  captionpos=b,
  xleftmargin=1.5em,
  framexleftmargin=1.5em,
  literate={ı}{{\i}}1 {ö}{{\"o}}1 {ü}{{\"u}}1 {ş}{{\c{s}}}1 {ç}{{\c{c}}}1 {ğ}{{\u{g}}}1 {İ}{\.{I}}1 {Ö}{{\"O}}1 {Ü}{{\"U}}1 {Ş}{{\c{S}}}1 {Ç}{{\c{C}}}1 {Ğ}{{\u{G}}}1,
}
\lstdefinelanguage{JavaScript}{
  keywords={break,case,catch,const,continue,debugger,default,delete,do,else,export,extends,finally,for,function,if,import,in,instanceof,let,new,of,return,switch,this,throw,try,typeof,var,void,while,with,yield,class,async,await},
  morekeywords={true,false,null,undefined},
  sensitive=true,
  comment=[l]{//},
  morecomment=[s]{/*}{*/},
  string=[b]',
  morestring=[b]",
  morestring=[b]`,
}

% --- Colors ---
\usepackage{xcolor}
\definecolor{chapterblue}{RGB}{100,149,237}
\definecolor{sectioncyan}{RGB}{0,200,220}
\definecolor{linkcolor}{RGB}{120,180,255}
\definecolor{citecolor}{RGB}{180,220,100}
\definecolor{boxbg}{RGB}{60,60,80}
\definecolor{boxborder}{RGB}{120,180,255}
\definecolor{warningbg}{RGB}{80,50,50}
\definecolor{warningborder}{RGB}{255,130,130}
\definecolor{successbg}{RGB}{45,75,45}
\definecolor{successborder}{RGB}{130,255,130}
\definecolor{formulabg}{RGB}{50,50,70}

% --- Hyperref ---
\usepackage[
  colorlinks=true,
  linkcolor=linkcolor,
  citecolor=citecolor,
  urlcolor=linkcolor,
  bookmarks=true,
  bookmarksnumbered=true,
  pdfauthor={Berke Tezgocen},
  pdftitle={Karadul (Black Widow) -- Teknik Dokuman},
  pdfsubject={Tersine Muhendislik ve Statik Analiz},
]{hyperref}

% --- Chapter & Section Formatting ---
\usepackage{titlesec}
\titleformat{\chapter}[display]
  {\normalfont\Huge\bfseries\color{chapterblue}}
  {\chaptertitlename\ \thechapter}{20pt}{\Huge}
\titleformat{\section}
  {\normalfont\Large\bfseries\color{sectioncyan}}
  {\thesection}{1em}{}
\titleformat{\subsection}
  {\normalfont\large\bfseries\color{white!90!blue}}
  {\thesubsection}{1em}{}
\titleformat{\subsubsection}
  {\normalfont\normalsize\bfseries\color{white!80!cyan}}
  {\thesubsubsection}{1em}{}

% --- Headers & Footers ---
\usepackage{fancyhdr}
\pagestyle{fancy}
\fancyhf{}
\fancyhead[L]{\color{gray!60!white}\small\leftmark}
\fancyhead[R]{\color{gray!60!white}\small\thepage}
\fancyfoot[C]{\color{gray!40!white}\tiny Karadul v1.0 --- Teknik Dokuman}
\renewcommand{\headrulewidth}{0.4pt}
\renewcommand{\headrule}{\hbox to\headwidth{\color{gray!30!white}\leaders\hrule height \headrulewidth\hfill}}

% --- Custom Environments ---
\usepackage{tcolorbox}
\tcbuselibrary{skins,breakable}

% Bolum Ozeti kutusu
\newtcolorbox{ozet}{
  colback=boxbg,
  colframe=boxborder,
  coltext=white,
  coltitle=white,
  fonttitle=\bfseries,
  title={Bu Bolumun Ozeti},
  breakable,
  sharp corners,
  boxrule=0.8pt,
}

% Tanim kutusu
\newtcolorbox{tanim}[1][]{
  colback=formulabg,
  colframe=sectioncyan,
  coltext=white,
  coltitle=white,
  fonttitle=\bfseries,
  title={#1},
  breakable,
  sharp corners,
  boxrule=0.6pt,
}

% Uyari kutusu
\newtcolorbox{uyari}[1][]{
  colback=warningbg,
  colframe=warningborder,
  coltext=white,
  coltitle=white,
  fonttitle=\bfseries,
  title={#1},
  breakable,
  sharp corners,
  boxrule=0.6pt,
}

% Basari kutusu
\newtcolorbox{basari}[1][]{
  colback=successbg,
  colframe=successborder,
  coltext=white,
  coltitle=white,
  fonttitle=\bfseries,
  title={#1},
  breakable,
  sharp corners,
  boxrule=0.6pt,
}

% --- Theorem Environments ---
\theoremstyle{definition}
\newtheorem{definition}{Tanim}[chapter]
\newtheorem{theorem}{Teorem}[chapter]
\newtheorem{lemma}[theorem]{Lemma}
\newtheorem{corollary}[theorem]{Sonuc}
\newtheorem{example}{Ornek}[chapter]

% --- Bracket Commands ---
\newcommand{\llbracket}{\ensuremath{[\![}}
\newcommand{\rrbracket}{\ensuremath{]\!]}}

% --- Custom Commands ---
\newcommand{\code}[1]{\texttt{\color{green!60!white}#1}}
\newcommand{\file}[1]{\texttt{\color{yellow!70!white}#1}}
\newcommand{\bigO}[1]{\mathcal{O}\!\left(#1\right)}
\newcommand{\prob}[1]{P\!\left(#1\right)}
\newcommand{\given}{\,|\,}
\newcommand{\sigmoid}{\sigma}
\newcommand{\logodds}{\text{log-odds}}
\DeclareMathOperator*{\argmax}{arg\,max}
\DeclareMathOperator*{\argmin}{arg\,min}

% --- Bibliography ---
\usepackage[numbers,sort&compress]{natbib}

% --- Index ---
\usepackage{makeidx}
\makeindex

% --- Spacing ---
\usepackage{setspace}
\onehalfspacing
\setlength{\parindent}{0pt}
\setlength{\parskip}{6pt}

% ============================================================================
\begin{document}

% --- Kapak Sayfasi ---
\begin{titlepage}
\centering
\vspace*{4cm}
{\fontsize{36}{42}\selectfont\bfseries\color{chapterblue}
Karadul\\[8pt]
\Large\color{white!80!blue}(Black Widow)}\\[1.5cm]
{\Large\itshape\color{gray!70!white} Tersine Muhendislik ve Statik Analiz Araci\\[4pt]
Kapsamli Teknik Dokuman}\\[2cm]
{\large\color{white}
\textbf{Berke Tezgocen}\\[8pt]
Fizik Muhendisligi --- ITU\\[4pt]
2026}\\[3cm]
{\small\color{gray!60!white}
Bu dokuman, Karadul (Black Widow) tersine muhendislik aracinin\\
tum algoritmalarini, matematiksel temellerini, tasarim kararlarini\\
ve implementasyon detaylarini A'dan Z'ye anlatir.}\\[4cm]
{\footnotesize\color{gray!40!white}
Hazirlayan: Codex AI Danismanligi + Claude Opus 4.6\\
Son guncelleme: 24 Mart 2026}
\end{titlepage}

% --- Icindekiler ---
\tableofcontents
\newpage

% ============================================================================
% BOLUM 1: GIRIS VE MOTIVASYON
% ============================================================================
\chapter{Giris ve Motivasyon}
\label{ch:giris}

\section{Tersine Muhendislik Nedir?}
\label{sec:re-nedir}

Tersine muhendislik (reverse engineering), bir yazilim urununu --- kaynak kodu olmadan --- analiz ederek ic yapisini, algoritmalarini, veri yapilarini ve davranisini anlama surecidir. Derlenm\-is bir binary dosyadan veya obfuscate edilmis bir JavaScript bundle'indan orijinal kaynak koda yaklasma cabasi, bilgisayar bilimlerinin en zorlayici alanlarindan biridir.

\begin{tanim}[Formal Tanim]
Bir $f: S \to T$ fonksiyonu verildigi\-nde, tersine muhendislik $f^{-1}: T \to \hat{S}$ yaklasik ters fonksiyonunu bulma problemidir. Burada:
\begin{itemize}
\item $S$: Orijinal kaynak kod uzayi
\item $T$: Derlenmis/obfuscate edilmis cikti uzayi
\item $f$: Derleyici veya obfuscator donusumu (genellikle surjective ama injective degil)
\item $\hat{S}$: Kaynak kodun yaklasik rekonstruksiyonu
\end{itemize}
Kritik nokta: $f$ genellikle bijective degil --- birden fazla farkli kaynak kod ayni binary'yi uretebilir. Bu nedenle $\hat{S} \neq S$ olabilir, ama $f(\hat{S}) \approx f(S)$ olma\-lidir.
\end{tanim}

\subsection{Neden Onemli?}

Tersine muhendislik bircok kritik alanda kullanilir:

\begin{enumerate}
\item \textbf{Guvenlik arastirmasi:} Malware analizi, zafiyet kesfetme (vulnerability research), exploit gelistirme ve savunma.
\item \textbf{Uyumluluk:} Clean-room reverse engineering ile uyumlu yazilim gelistirme (ornek: Wine projesi, ReactOS).
\item \textbf{Dijital adli bilisim:} Siber suc arastirmalarinda delil analizi.
\item \textbf{Legacy sistem bakimi:} Kaynak kodu kaybolmus yazilimlarin anlasil\-masi ve surdurul\-mesi.
\item \textbf{Yarisma (CTF):} Capture The Flag yarismalarinda binary exploitation.
\end{enumerate}

\subsection{Endustri Baglami}

Tersine muhendislik araci pazari buyuk ve buyuyor:

\begin{table}[h]
\centering
\begin{tabular}{lll}
\toprule
\textbf{\color{sectioncyan}Arac} & \textbf{\color{sectioncyan}Fiyat} & \textbf{\color{sectioncyan}Ozellik} \\
\midrule
IDA Pro & \$2,500+/yil & Endustri standardi, kapali kaynak \\
Ghidra & Ucretsiz (NSA) & Acik kaynak, Java tabanli \\
Binary Ninja & \$400/yil & Modern UI, Python API \\
radare2/Cutter & Ucretsiz & Acik kaynak, komut satiri odakli \\
\textbf{Karadul} & \textbf{Ucretsiz} & \textbf{Otonom pipeline, Bayesian isimlendirme} \\
\bottomrule
\end{tabular}
\caption{Tersine muhendislik arac karsilastirmasi.}
\label{tab:re-tools}
\end{table}

\subsection{Mevcut Cozumlerin Sinirliliklari}

Mevcut araclar tipik olarak su sorunlara sahiptir:

\begin{enumerate}
\item \textbf{Manuel is yueku:} IDA Pro ve Ghidra gueclue decompiler'lar saglar, ancak fonksiyon isimlendirme, tip cikari\-mi ve yapı tespiti buyuek olcuede analistin elindedir. 10.000 fonksiyonluk bir binary'de bu haftalarca suerer.
\item \textbf{Tek platform odakli:} Cogu arac ya binary ya da JS analizi yapar, ikisini birden degil.
\item \textbf{Otomasyon eksikligi:} Toplu analiz (batch processing), otomatik raporlama ve end-to-end pipeline mevcut araclarda sinirlidir.
\item \textbf{Isim kurtarma kalitesi:} Stripped binary'lerde fonksiyon isimlerini kurtarma orani genellikle duesuektuer.
\end{enumerate}

\subsection{Karadul'un Yaklasimi ve Farki}
\label{sec:karadul-fark}

Karadul (Black Widow), bu sinirlamalari su yeniliklerle asar:

\begin{enumerate}
\item \textbf{6 asamali otonom pipeline:} Hedef tespitinden rapora kadar tamamen otomatik.
\item \textbf{Coklu hedef destegi:} 12 farkli binary/bundle formati, 13 programlama dili.
\item \textbf{Bayesian isim fuzyonu:} 7 bagimsiz kaynaktan gelen isimleri korelasyon-aware Bayesian log-odds ile birlestirme --- enduestride benzersiz.
\item \textbf{Derin JS deobfuscation:} 9 fazli AST transform pipeline --- LLM'siz semantik degisken isimlendirme.
\item \textbf{Graceful degradation:} Herhangi bir asamanin basarisizligi pipeline'i durdurmaz.
\item \textbf{Olceklenebilirlik:} 224MB+ binary'ler ve 200MB+ JS dosyalari icin chunked processing.
\end{enumerate}

\section{Proje Kapsami}
\label{sec:kapsam}

Karadul v1.0 su bilesenlerden olusur:

\begin{itemize}
\item \textbf{110 Python modulu,} 65.000+ satir kod
\item \textbf{460+ JavaScript/MJS analiz scripti}
\item \textbf{1.393 test} (tumu gecen)
\item \textbf{1.5M+ kutuephane imzasi} (FLIRT veritabani)
\item \textbf{Docker destegi} + pip kurulabilir paket
\item \textbf{CLI arayuzu} (Click tabanli)
\end{itemize}

\section{Dogrulanan Gercek Hedefler}

\begin{table}[h]
\centering
\begin{tabular}{lllrrl}
\toprule
\textbf{\color{sectioncyan}Hedef} & \textbf{\color{sectioncyan}Tip} & \textbf{\color{sectioncyan}Dil} & \textbf{\color{sectioncyan}Boyut} & \textbf{\color{sectioncyan}Fonk.} & \textbf{\color{sectioncyan}Kurtarma} \\
\midrule
Rectangle & Mach-O & Swift & 1.8 MB & 2.852 & \%89.9 \\
Steam & Mach-O & C++ & 4.1 MB & 11.296 & \%74.2 \\
Claude Code & JS bundle & esbuild & 11.8 MB & 43.747 & \%57.3 \\
Cursor & JS bundle & webpack & 5.97 MB & 17.705 & \%65.4 \\
LM Studio & Electron & webpack & 52 MB & 150K+ & 4/4 bundle \\
\bottomrule
\end{tabular}
\caption{Gercek duenyada dogrulanan hedefler.}
\label{tab:targets}
\end{table}

\section{Bu Dokumanin Organizasyonu}

\begin{description}
\item[Bolum 2] --- \textbf{Pipeline Mimarisi:} 6 asamali pipeline, context pattern, error recovery.
\item[Bolum 3] --- \textbf{Hedef Tespit Sistemi:} Magic bytes, binary format analizi, dil tespiti.
\item[Bolum 4] --- \textbf{Statik Analiz:} Binary ve JS statik analiz algoritmalari.
\item[Bolum 5] --- \textbf{JavaScript Deobfuscation:} AST transform pipeline, 9 faz, webpack unpack.
\item[Bolum 6] --- \textbf{Bayesian Isim Fuzyonu:} Log-odds formulu, 7 kaynak, korelasyon duzeltmesi.
\item[Bolum 7] --- \textbf{Kaynak Kod Rekonstruksiyonu:} C naming, tip kurtarma, algoritma tespiti.
\item[Bolum 8] --- \textbf{Guvenlik Analiz Algoritmalari:} Entropy, YARA, CFG, opaque predicate.
\item[Bolum 9] --- \textbf{Performans ve Olceklendirme:} Buyuek dosya destegi, bellek yoenetimi, Apple Silicon.
\item[Bolum 10] --- \textbf{Tasarim Kararlari ve Retrospektif:} Neden bu yaklasimlar? Alternatifler.
\item[Ek A] --- \textbf{Parametre Tablosu:} Tuem konfiguerasyon parametreleri.
\item[Ek B] --- \textbf{Muuelakat Hazirlik:} 20 teknik soru ve detayli cevaplari.
\end{description}

\begin{ozet}
Karadul, derlenm\-is binary'ler ve obfuscate edilmis JavaScript bundle'lari icin otonom bir tersine muhendislik arac\-idir. 6 asamali pipeline, Bayesian isim fuzyonu ve 9 fazli derin deobfuscation ile enduestri standartlarinin oetesinde bir otomasyon seviyesi saglar. 5 gercek duenya hedefinde dogrulana\-rak v1.0 olarak tamamlandi.
\end{ozet}


% ============================================================================
% BOLUM 2: PIPELINE MIMARISI
% ============================================================================
\chapter{Pipeline Mimarisi}
\label{ch:pipeline}

Bu boluemde Karadul'un temelini olusturan 6 asamali pipeline mimarisini, yazilim muehendisligi prensiplerini ve hata kurtarma mekanizmalarini detayli olarak inceliyoruz.

\section{Pipeline Pattern --- Matematiksel Temel}
\label{sec:pipeline-pattern}

Karadul'un pipeline'i, \textbf{Orchestrated DAG Executor} patternidir --- akademik literaturde ``workflow engine'' siniflandirmasi. Bu, yazilim muehendisligindeki \textbf{Pipes and Filters} (Shaw \& Garlan, 1996) ile \textbf{Chain of Responsibility}'nin melez bir formudur.

\subsection{Formal Tanim}

Pipeline'i bir yoenlu dogrusiz cizge (Directed Acyclic Graph --- DAG) olarak modelleyebiliriz:

\begin{equation}
G = (V, E)
\label{eq:dag}
\end{equation}

burada:
\begin{itemize}
\item $V = \{v_1, v_2, \ldots, v_n\}$ stage kuemesi (her $v_i$ bir Stage nesnesi)
\item $E \subseteq V \times V$ bagimlilik kenarlari ($v_i \to v_j$ demek ``$v_i$, $v_j$'den oence calismali'')
\end{itemize}

\subsection{Gercek Dependency Grafi}

Koddan cikarilan gercek bagimlilik yapisi (\file{stages.py}):

\begin{center}
\begin{tikzpicture}[
  node/.style={rectangle, draw=sectioncyan, fill=boxbg, minimum width=2.8cm, minimum height=0.7cm, text=white, font=\small\bfseries, rounded corners=2pt},
  arr/.style={-{Stealth[length=2.5mm]}, thick, color=gray!60!white},
]
  \node[node] (id) at (0,0) {IDENTIFY};
  \node[node] (st) at (-4,-1.5) {STATIC};
  \node[node] (dy) at (-1.3,-1.5) {DYNAMIC};
  \node[node] (de) at (1.3,-1.5) {DEOBFUSCATE};
  \node[node] (rp) at (4,-1.5) {REPORT};
  \node[node] (rc) at (1.3,-3) {RECONSTRUCT};

  \draw[arr] (id) -- (st);
  \draw[arr] (id) -- (dy);
  \draw[arr] (id) -- (de);
  \draw[arr] (id) -- (rp);
  \draw[arr] (de) -- (rc);
\end{tikzpicture}
\end{center}

Bu bir \textbf{agac degil, DAG}. Identify koek dueguem, 4 dal ondan cikar ve reconstruct'in deobfuscate'e bagimliligi var. Lineer bir zincir degil --- static, dynamic ve deobfuscate teorik olarak paralel calisabilir (hepsi sadece identify'a bagimli).

\subsection{Topolojik Siralama}

Stage'lerin calisti\-rilma siras\-i, topolojik siralama (topological sort) ile belirlenir. Kahn's Algorithm:

\begin{enumerate}
\item Her dueguemuen in-degree'sini hesapla: $\text{deg}^-(v)$ ($v$'ye giren kenar sayisi)
\item $\text{deg}^-(v) = 0$ olan tuem dueguemleri kuyru\-ga ekle
\item Kuyruktan bir $v$ al, siralanan listeye ekle
\item $v$'nin tuem komsularinin in-degree'sini 1 azalt
\item Yeni $\text{deg}^- = 0$ olanlari kuyru\-ga ekle
\item Tekrarla (tuem dueguemler islene kadar)
\end{enumerate}

Hesaplama karmasikligi: $\bigO{|V| + |E|}$.

\begin{tanim}[Kritik Tasarim Karari: Kayit Sirasi vs Topolojik Siralama]
Karadul topolojik siralama \textbf{yapmiyor}. \file{pipeline.py:357}'deki \code{\_resolve\_run\_order()} metodu \textbf{kayit sirasi}ni kullanir:
\begin{lstlisting}[language=Python]
def _resolve_run_order(self, requested):
    if requested is None:
        return list(self._stage_order)  # kayit sirasi!
    return [s for s in self._stage_order if s in requested]
\end{lstlisting}
Stage'ler \code{register\_stage()} ile eklenme sirasinda calistirilir. CLI her zaman dogru sirada register ettigi icin pratikte sorun yok. Yanlis sirada ekleme durumunda \code{\_check\_dependencies()} runtime'da fail eder ve stage atlanir.
\end{tanim}

Bu yaklasim \textbf{static topological sort} yerine \textbf{runtime dependency validation} tercih etmistir:

\begin{table}[ht]
\centering
\begin{tabular}{p{5cm}p{4cm}p{4cm}}
\toprule
\textbf{\color{sectioncyan}Yaklasim} & \textbf{\color{sectioncyan}Avantaj} & \textbf{\color{sectioncyan}Dezavantaj} \\
\midrule
Topolojik siralama & Yanlis siralama imkansiz & $\bigO{V+E}$ karmasiklik, cycle detection \\
Kayit sirasi + runtime check (Karadul) & Basit implementasyon & Kullanicinin dogru eklemesi gerekir \\
\bottomrule
\end{tabular}
\end{table}

\subsection{Neden Sequential Pipeline?}

Dependency grafina bakildiginda static, dynamic ve deobfuscate paralel calisabilir. Ama sequential secilmis. Bunu Amdahl Yasasi ile analiz edelim:

\begin{tanim}[Amdahl Yasasi ile Sayisal Analiz]
Tipik stage suereleri:
\begin{align}
T_{\text{identify}} &= 2\text{s}, \quad T_{\text{static}} = 60\text{s (Ghidra)}, \quad T_{\text{dynamic}} = 30\text{s} \notag \\
T_{\text{deobfuscate}} &= 45\text{s}, \quad T_{\text{reconstruct}} = 20\text{s}, \quad T_{\text{report}} = 5\text{s} \notag
\end{align}

Sequential toplam: $T_{\text{seq}} = 2 + 60 + 30 + 45 + 20 + 5 = 162\text{s}$

Concurrent (static $\|$ dynamic $\|$ deobfuscate):
\begin{equation}
T_{\text{parallel}} = 2 + \max(60, 30, 45) + 20 + 5 = 87\text{s}
\end{equation}
\begin{equation}
\text{Speedup} = \frac{162}{87} = 1.86\times
\end{equation}

Amdahl formuelue ile: Paralel kisim orani $p = (60+30+45)/162 = 0.83$:
\begin{equation}
S_{\text{teorik}} = \frac{1}{(1-0.83) + 0.83/3} = \frac{1}{0.17 + 0.28} = 2.22\times
\end{equation}
\end{tanim}

\textbf{1.86x speedup icin concurrent pipeline karmasikligi (lock'lar, barrier'lar, partial failure handling) eklenmemis. Dogru karar} --- cuenkue:

\begin{enumerate}
\item \textbf{Kaynak catismasi:} Ghidra 2--4 GB RAM, Frida + Babel ayri ag\-ir surecler. Paralel calissa $8+$ GB --- OOM riski.
\item \textbf{Gizli bagimliliklar:} \code{stages.py}'de deobfuscate stage'i \code{requires=["identify"]} diyor ama gercekte static analiz sonuclarini da kullaniyor. ``Hidden coupling'' paralelde race condition yaratir.
\item \textbf{Debugging:} Sequential'da hata hangi stage'de oldugu acik. Concurrent'da cok daha karisik.
\end{enumerate}

\section{PipelineContext --- Shared State Pattern}
\label{sec:context}

\file{pipeline.py:41}'de tanimlanan \code{PipelineContext}, stage'ler arasi paylasilan veri yapisidir:

\begin{lstlisting}[language=Python,caption={PipelineContext veri yapisi}]
@dataclass
class PipelineContext:
    target: TargetInfo       # Hedef bilgileri
    workspace: Workspace     # Calisma dizini
    config: Config           # Konfiguerasyon
    results: dict[str, StageResult]  # Onceki sonuclar
    extra: dict[str, Any]    # Ek paylasilam veri
\end{lstlisting}

Bu tasarim, Martin Fowler'in \textbf{Introduce Parameter Object} refactoring kalibinin bir uygulamasidir. Avantajlari:

\begin{itemize}
\item Stage'ler arasinda explicit veri gecisi
\item Yeni veriler \code{extra} dict'e eklenebilir (genisletilebilirlik)
\item Tum state tek bir nesnede --- test edilebilirlik
\end{itemize}

\section{Error Recovery --- Circuit Breaker Pattern}
\label{sec:error-recovery}

\subsection{State Machine Modeli}

Circuit breaker uec durumlu bir state machine'dir:

\begin{equation}
\text{States} = \{\text{CLOSED}, \text{OPEN}, \text{HALF\_OPEN}\}
\end{equation}

\begin{center}
\begin{tikzpicture}[
  statenode/.style={circle, draw=sectioncyan, fill=boxbg, minimum size=2cm, text=white, font=\small\bfseries},
  trans/.style={-{Stealth[length=3mm]}, thick, color=gray!70!white},
  lbl/.style={font=\tiny, text=gray!80!white, align=center},
]
  \node[statenode] (closed) at (0,0) {CLOSED};
  \node[statenode] (open) at (5,0) {OPEN};
  \node[statenode] (half) at (2.5,-3.5) {\shortstack{HALF\\OPEN}};

  \draw[trans] (closed) -- node[above, lbl] {fail $\geq 5$} (open);
  \draw[trans] (open) -- node[right, lbl] {60s timeout} (half);
  \draw[trans] (half) -- node[left, lbl] {success} (closed);
  \draw[trans] (half) -- node[below right, lbl] {fail} (open);
  \draw[trans] (closed) to[loop above] node[above, lbl] {success} (closed);
\end{tikzpicture}
\end{center}

\subsection{Exponential Backoff}

Yeniden deneme gecikmesi:

\begin{equation}
\text{delay}(k) = \min\!\big(\text{base} \cdot 2^k,\; \text{max\_delay}\big)
\label{eq:backoff}
\end{equation}

burada $k$ = deneme sayisi, base $= 1.0$s, max\_delay $= 30$s.

Toplam beklenen bekleme sueresi ($n$ deneme icin):

\begin{equation}
E[T] = \sum_{k=0}^{n-1} \min\!\big(1.0 \cdot 2^k,\; 30\big) = 1 + 2 + 4 + 8 + 16 + 30 + \cdots
\end{equation}

$n = 3$ icin: $E[T] = 1 + 2 + 4 = 7$ saniye.

\subsection{Neden Exponential Backoff?}

Farkli backoff stratejilerinin karsilastirmasi:

\begin{table}[ht]
\centering
\begin{tabular}{lllll}
\toprule
\textbf{\color{sectioncyan}Strateji} & \textbf{\color{sectioncyan}Formuul} & \textbf{\color{sectioncyan}$n=3$} & \textbf{\color{sectioncyan}$n=5$} & \textbf{\color{sectioncyan}Bueyueme} \\
\midrule
Linear & $n(n+1)/2$ & 6s & 15s & $\bigO{n^2}$ \\
Quadratic & $n(n+1)(2n+1)/6$ & 14s & 55s & $\bigO{n^3}$ \\
Exponential & $2^n - 1$ & 7s & 31s & $\bigO{2^n}$ \\
\bottomrule
\end{tabular}
\caption{Backoff stratejilerinin toplam bekleme suereleri.}
\end{table}

Exponential tercih sebebi: \textbf{geri basinc (backpressure)}. Servis down ise ilk retry'lar hizli, sonrakiler dramatik artan --- servise ``nefes alma'' sueresi verir. Linear bunu yapmaz.

\subsection{Jitter}

Kodda \textbf{zaten var} (\file{error\_recovery.py:205}):
\begin{lstlisting}[language=Python]
if jitter:
    jitter_range = delay * 0.2
    delay += random.uniform(-jitter_range, jitter_range)
\end{lstlisting}

$\pm 20\%$ uniform jitter. Amazon AWS ``full jitter'' ($[0, \text{delay}]$) oenerir, ama Karadul'da tek kullanici var --- thundering herd problemi yok.

\subsection{Basari Olasiligi}

$k$ deneme sonrasi basari olasiligi (her denemenin bagimsiz basari olasiligi $p$):
\begin{equation}
\prob{\text{basari sonrasi } k \text{ deneme}} = 1 - (1-p)^k
\end{equation}

$p = 0.8$ (tipik gecici hata) ve $k = 3$: $P = 1 - 0.2^3 = 1 - 0.008 = 0.992$ --- \%99.2 basari.

\begin{ozet}
Karadul pipeline'i, 6 stage'li sequential bir DAG execution engine'dir. PipelineContext shared state pattern ile veri paylasimi saglar. ErrorRecovery, circuit breaker + exponential backoff ile tuem stage'leri korur. Graceful degradation sayesinde tek bir stage basarisizligi pipeline'i durdurmaz.
\end{ozet}


% ============================================================================
% BOLUM 3: HEDEF TESPIT SISTEMI
% ============================================================================
\chapter{Hedef Tespit Sistemi}
\label{ch:target}

Karadul'un ilk asamasi, analiz edilecek dosyanin tipini, programlama dilini ve format\-ini tespit etmektir.

\section{Magic Number Sistemi}
\label{sec:magic}

Her binary dosya formati, dosyanin bas\-inda benzersiz bir ``magic number'' tasir. Bu degerler formati tanimlar:

\begin{equation}
\text{magic}: \{0, 1\}^{32} \to \text{Format}
\end{equation}

\begin{table}[h]
\centering
\begin{tabular}{llll}
\toprule
\textbf{\color{sectioncyan}Format} & \textbf{\color{sectioncyan}Magic (hex)} & \textbf{\color{sectioncyan}ASCII} & \textbf{\color{sectioncyan}Oykue} \\
\midrule
Mach-O 64-bit & \code{0xFEEDFACF} & --- & ``feed face'' (Apple) \\
Mach-O 32-bit & \code{0xFEEDFACE} & --- & ``feed face'' (Apple) \\
Universal & \code{0xCAFEBABE} & --- & ``cafe babe'' (NeXT/Java) \\
ELF & \code{0x7F454C46} & \code{.ELF} & ELF standard header \\
PE/DOS & \code{0x4D5A} & \code{MZ} & Mark Zbikowski (DOS yaraticisi) \\
ZIP/JAR/APK & \code{0x504B0304} & \code{PK..} & Phil Katz (PKZIP yaraticisi) \\
\bottomrule
\end{tabular}
\caption{Binary dosya formatlarinin magic byte'lari.}
\label{tab:magic}
\end{table}

\subsection{Collision Olasiligi}

Rastgele bir dosyanin belirli bir 4-byte magic ile basla\-ma olasiligi:

\begin{equation}
\prob{\text{false positive}} = \frac{1}{2^{32}} \approx 2.33 \times 10^{-10}
\end{equation}

Bu son derece duseuk bir olasiliktir, ancak Karadul ek dogrulamalar da yapar (icerik analizi, dosya uzantisi).

\subsection{Big-Endian vs Little-Endian Okuma}

Magic byte'lar big-endian sirada okunur (\code{struct.unpack(">I")}, \file{target.py:607}). Mach-O formati ``cigam'' (magic'in tersi) varyantlarina da sahiptir:

\begin{equation}
\text{CIGAM}(m) = \text{bswap32}(m)
\end{equation}

Oernegin: \code{0xFEEDFACF} $\to$ \code{0xCFFAEDFE} (little-endian ARM islemcilerde).

\section{Dil Tespit Algoritmasi}
\label{sec:lang-detect}

Binary icindeki string pattern'lerden dil tespiti yapar (\file{target.py:470}):

\subsection{Puanlama Sistemi}

Her dil icin bilinen string marker'lar taranir:

\begin{equation}
\text{score}(L) = \sum_{s \in \text{markers}(L)} \mathbb{1}[s \in \text{binary\_text}]
\end{equation}

En yueksek skorlu dil secilir: $\hat{L} = \arg\max_L \text{score}(L)$.

\begin{tanim}[Bayes Siniflandirici Perspektifi]
Bu puanlama sistemi aslinda bir Naive Bayes siniflandiricisi\-nin basitlestirilmis hali\-dir:
\begin{equation}
\prob{L \given S_1, \ldots, S_n} \propto \prob{L} \prod_{i=1}^{n} \prob{S_i \given L}
\end{equation}
Mevcut implementasyonda $\prob{L}$ uniform (esit prior) ve $\prob{S_i \given L} \in \{0, 1\}$ (deterministik).
\end{tanim}

\section{JavaScript Bundler Tespiti}
\label{sec:bundler}

JS dosyalarinin ilk 64KB'i okunarak bundler tespiti yapilir:

\begin{table}[h]
\centering
\begin{tabular}{ll}
\toprule
\textbf{\color{sectioncyan}Signature} & \textbf{\color{sectioncyan}Bundler} \\
\midrule
\code{\_\_webpack\_require\_\_} & webpack \\
\code{parcelRequire} & parcel \\
\code{System.register} & systemjs \\
\code{\_rollupPluginBabelHelpers} & rollup \\
\code{define([} & amd/requirejs \\
\bottomrule
\end{tabular}
\caption{JavaScript bundler signature'lari.}
\label{tab:bundlers}
\end{table}

\section{SHA-256 Hashing}
\label{sec:hashing}

Dosya buetunluegue icin SHA-256 streaming hash:

\begin{equation}
H = \text{SHA-256}\!\left(\bigoplus_{i=1}^{\lceil n/8192 \rceil} \text{chunk}_i\right)
\end{equation}

8KB chunk boyutu disk I/O vs CPU tradeoff'u optimize eder. Collision olasiligi:

\begin{equation}
\prob{\text{collision}} \approx \frac{n^2}{2^{257}}
\end{equation}

$n = 10^{12}$ dosya icin bile: $\prob{\text{collision}} \approx 10^{-54}$.

\begin{ozet}
Hedef tespit sistemi, magic byte analizi, string-based dil tespiti, bundler signature eslestirmesi ve SHA-256 hashing kullanarak 12 farkli dosya tipini ve 13 programlama dilini otomatik olarak tespit eder. False positive olasiligi astronomik derecede duseuktueer.
\end{ozet}

% --- DETAYLI GENISLETME ---
% ============================================================================
% BOLUM 3 GENISLETME: HEDEF TESPIT SISTEMI — DETAYLI
% ============================================================================

% Bu dosya main.tex'teki Bolum 3'un SONUNA \input ile eklenir

\section{Mach-O Format Detaylari}
\label{sec:macho-detail}

Apple platformlarinda kullanilan Mach-O (Mach Object) formati, NeXTSTEP mirasinin urunudur. Carnegie Mellon'daki Mach kernel projesi (1985) icin tasarlanmistir.

\subsection{Header $\to$ Load Commands $\to$ Segments $\to$ Sections}

Mach-O dosyasi katmanli bir yapidir:

\begin{center}
\begin{tikzpicture}[
  block/.style={rectangle, draw=sectioncyan, fill=boxbg, minimum width=8cm, minimum height=0.7cm, text=white, font=\footnotesize},
  sep/.style={rectangle, draw=gray!30!white, fill=black!70!white, minimum width=8cm, minimum height=0.3cm},
]
  \node[block] (header) at (0,0) {Mach-O Header (32 byte): magic, cputype, filetype, ncmds};
  \node[block] (lc1) at (0,-0.9) {LC\_SEGMENT\_64: \_\_TEXT segment tanimlamasi};
  \node[block] (lc2) at (0,-1.8) {LC\_SYMTAB: Symbol table offset ve boyut};
  \node[block] (lc3) at (0,-2.7) {LC\_DYSYMTAB: Dynamic symbol bilgisi};
  \node[block] (lc4) at (0,-3.6) {LC\_DYLIB: Bagli kutuphane bilgisi};
  \node[sep] (s1) at (0,-4.3) {};
  \node[block] (text) at (0,-5.0) {\_\_TEXT Segment: derlenm\-is makine kodu, string sabitleri};
  \node[block] (data) at (0,-5.9) {\_\_DATA Segment: global degiskenler, BSS};
  \node[block] (link) at (0,-6.8) {\_\_LINKEDIT Segment: symbol table, string table};
\end{tikzpicture}
\end{center}

\textbf{Header} (mach\_header\_64, 32 byte):
\begin{itemize}
\item \code{magic} (4 byte): \code{0xFEEDFACF} (64-bit) veya \code{0xFEEDFACE} (32-bit)
\item \code{cputype}: \code{CPU\_TYPE\_ARM64 = 0x0100000C}, \code{CPU\_TYPE\_X86\_64 = 0x01000007}
\item \code{filetype}: \code{MH\_EXECUTE} (2), \code{MH\_DYLIB} (6), \code{MH\_BUNDLE} (8)
\item \code{ncmds}: Load command sayisi
\item \code{flags}: \code{MH\_PIE} (position independent), \code{MH\_TWOLEVEL} (iki seviyeli namespace)
\end{itemize}

\textbf{Segmentler:} Bellek yuekleme birimi. \code{\_\_TEXT} segmenti yurutuluebilir ama yazilamaz (W\^{}X, code signing icin kritik). \code{\_\_DATA} segmenti okunur-yazilir.

\subsection{Universal Binary (Fat Binary)}

Birden fazla mimari icin tek dosyada derlenmis binary:

\begin{center}
\begin{tikzpicture}[
  block/.style={rectangle, draw=sectioncyan, fill=boxbg, minimum width=7cm, minimum height=0.7cm, text=white, font=\footnotesize},
]
  \node[block] (fat) at (0,0) {Fat Header: magic=0xCAFEBABE, nfat\_arch=2};
  \node[block] (a0) at (0,-0.9) {fat\_arch[0]: arm64 -- offset, size, alignment};
  \node[block] (a1) at (0,-1.8) {fat\_arch[1]: x86\_64 -- offset, size, alignment};
  \node[block,fill=boxbg!80!blue] (m0) at (0,-3.0) {ARM64 Mach-O (tam dosya)};
  \node[block,fill=boxbg!80!green] (m1) at (0,-3.9) {x86\_64 Mach-O (tam dosya)};
\end{tikzpicture}
\end{center}

Karadul'un \file{macho.py}'sindaki \code{analyze\_static()} bunu dogru handle eder:
\begin{lstlisting}[language=Python]
if target.target_type == TargetType.UNIVERSAL_BINARY:
    # lipo -thin arm64 ile sadece ARM slice'i cikar
    proc = subprocess.run(["lipo", str(binary_path), "-thin", "arm64", ...])
\end{lstlisting}

\textbf{Neden universal binary hala kullanilir:} Apple 2020'de Intel'den ARM64'e gecti. Universal binary gecis doeneminde zorunlu --- M-serisi ARM64, eski Mac'ler x86\_64 calistirir. Rosetta 2 JIT translation katmani sadece ARM64 slice yoksa devreye girer.

\subsection{CAFEBABE Cakismasi}

\begin{uyari}[Potansiyel Sorun]
\code{0xCAFEBABE} magic hem Java class dosyalarinda hem Apple universal binary'lerinde kullanilir. Ayrim: Java'da magic'ten sonra minor/major version byte'lari gelir (oernegin \code{0x0000003D} = Java 17), fat binary'de ise mimari sayisi gelir (1--9 arasi kueceuk integer). Mevcut Karadul implementasyonu bu ayrimi \textbf{yapmiyor} --- potansiyel false positive kaynagi.
\end{uyari}


\section{ELF Format ve Go Binary Tespiti}
\label{sec:elf-go}

\subsection{ELF Yapisi}

ELF (Executable and Linkable Format), Unix System V Release 4 (1989) ile standartlastirilmistir:

\begin{itemize}
\item \textbf{e\_ident[16]:} Ilk 16 byte --- magic (\code{7F 45 4C 46}), class (32/64-bit), endianness, version, OS/ABI
\item \textbf{Program Headers:} Runtime yuekleme bilgisi --- \code{PT\_LOAD}, \code{PT\_INTERP}, \code{PT\_DYNAMIC}
\item \textbf{Section Headers:} Link-time bilgisi --- \code{.text}, \code{.rodata}, \code{.data}, \code{.bss}, \code{.symtab}
\end{itemize}

\textbf{Kritik fark:} Program headers = runtime view (loader kullanir), section headers = link-time view (linker kullanir). Stripped binary'de section headers \textbf{silinebilir} ama program headers silinmez.

\subsection{GOPCLNTAB (Go Program Counter Line Number Table)}

Go binary'lerinin en onemli ic yapisi. Her Go binary'de bulunur --- \textbf{stripped olsa bile} genellikle kalir, cueuenkue Go runtime panic/recover/stack trace icin buna bagimlidir.

GOPCLNTAB icerigi:
\begin{itemize}
\item Her fonksiyonun baslangic adresi (PC)
\item Fonksiyon adi (tam package path: \code{main.(*Server).HandleRequest})
\item Kaynak dosya yolu
\item Satir numarasi mapping'i (PC $\to$ satir)
\item Stack frame boyutu
\end{itemize}

\textbf{Magic byte evrimi} (Karadul 4 versiyonu da handle ediyor):

\begin{table}[ht]
\centering
\begin{tabular}{lll}
\toprule
\textbf{\color{sectioncyan}Magic} & \textbf{\color{sectioncyan}Go Versiyon} & \textbf{\color{sectioncyan}Neden Degisti} \\
\midrule
\code{FB FF FF FF 00 00} & Go 1.2--1.15 & Orijinal format \\
\code{FA FF FF FF 00 00} & Go 1.16--1.17 & Format yeniden tasarim \\
\code{F0 FF FF FF 00 00} & Go 1.18--1.19 & Register-based ABI \\
\code{F1 FF FF FF 00 00} & Go 1.20+ & Gueuencel format \\
\bottomrule
\end{tabular}
\caption{GOPCLNTAB magic byte evrimi.}
\end{table}


\section{PE Format ve .NET/Delphi Ayrimi}
\label{sec:pe-dotnet-delphi}

\subsection{PE Yapisi}

PE (Portable Executable) formati, her Windows .exe, .dll, .sys dosyasi:

\begin{itemize}
\item \textbf{DOS Header} (64 byte): \code{e\_magic="MZ"} (Mark Zbikowski, DOS yaraticisi), \code{e\_lfanew} = PE offset
\item \textbf{DOS Stub} ($\sim$64 byte): ``This program cannot be run in DOS mode''
\item \textbf{PE Signature} (4 byte): \code{"PE\textbackslash 0\textbackslash 0"}
\item \textbf{COFF Header} (20 byte): Machine, NumberOfSections, TimeDateStamp
\item \textbf{Optional Header}: Entry point, ImageBase, \textbf{Data Directories} (16 adet)
\item \textbf{Data Directory index 14} = CLR Runtime Header --- bu entry sifirdan bueyuekse \textbf{.NET assembly}
\end{itemize}

\subsection{.NET Tespiti: mscoree.dll}

Her .NET executable'da import table'da \code{mscoree.dll} referansi bulunur:

\begin{lstlisting}[language=Python]
# target.py satir 524-536
if b"mscoree.dll" in data or b"_CorExeMain" in data:
    return True
if b"mscorlib" in data or b"System.Runtime" in data:
    return True
\end{lstlisting}

\textbf{False positive neredeyse imkansiz:} \code{mscoree.dll} string'i yalnizca .NET binary'lerinde bulunur.

\begin{equation}
\prob{.NET \given \text{mscoree.dll bulundu}} \approx 0.999
\end{equation}

\subsection{Delphi Binary Tespiti}

Delphi binary'leri native x86/x64 makine kodu icerir. Ayirt edici ozellikler:

\textbf{RTTI (Runtime Type Information):} Her class icin otomatik uretilir, binary'nin \code{.rdata} section'inda goemuelueduer:

\begin{lstlisting}[language=Python]
# Delphi mangled name pattern:
# @Unit@Class@Method$qqr
pattern = rb"@([A-Z]\w+)@([A-Z]\w+)@([A-Za-z_]\w*)\$qqr"
\end{lstlisting}

\code{\$qqr} suffixi ``no parameters'' anlamina gelir (Delphi mangling convention).

\textbf{Tespit hiyerarsisi:} \file{target.py}'da oencelik sirasi dogru:
\begin{enumerate}
\item \code{\_is\_dotnet()} $\to$ DOTNET\_ASSEMBLY
\item \code{\_is\_delphi()} $\to$ DELPHI\_BINARY
\item Fallback $\to$ PE\_BINARY
\end{enumerate}

.NET kontrolu daha belirgin (mscoree.dll) bu yueuezden oence yapilmali.


\section{Dil Tespit Algoritmasi --- Bayes Perspektifi}
\label{sec:lang-bayes}

\subsection{Mevcut Puanlama Sistemi}

\file{target.py:470}'deki dil tespiti basit bir sayma yaklasimi kullaniyor:

\begin{equation}
\text{score}(L) = \sum_{s \in \text{markers}(L)} \mathbb{1}[s \in \text{binary\_text}]
\end{equation}

Her signature esit agirlikta --- \code{rust\_begin\_unwind} ile \code{std::} ayni puani verir.

\subsection{Naive Bayes Formulasyonu}

Bu problem bir Bayes siniflandiricidir:

\begin{equation}
\prob{L \given S_1, \ldots, S_n} \propto \prob{L} \prod_{i=1}^{n} \prob{S_i \given L}
\end{equation}

\textbf{Prior olasiliklar} platform-conditioned olabilir:
\begin{itemize}
\item Mach-O ise: Swift/Obj-C prior'u yueksek
\item ELF ise: C/Go prior'u yueksek
\item PE ise: C++/.NET/Delphi prior'u yueksek
\end{itemize}

\textbf{Signature likelihood'lari:}

\begin{table}[ht]
\centering
\begin{tabular}{lllll}
\toprule
\textbf{\color{sectioncyan}Signature} & \textbf{\color{sectioncyan}$\prob{S \given \text{Rust}}$} & \textbf{\color{sectioncyan}$\prob{S \given \text{C++}}$} & \textbf{\color{sectioncyan}$\prob{S \given \text{Go}}$} & \textbf{\color{sectioncyan}Gueuec} \\
\midrule
\code{rust\_begin\_unwind} & 0.98 & 0.001 & 0.001 & Cok gueueclue \\
\code{std::} & 0.15 & 0.95 & 0.02 & Cakisma riski \\
\code{runtime.gopanic} & 0.001 & 0.001 & 0.99 & Cok gueueclue \\
\code{objc\_msgSend} & 0.001 & 0.01 & 0.001 & Obj-C kesin \\
\bottomrule
\end{tabular}
\caption{Dil tespit signature likelihood degerleri.}
\end{table}

\begin{tanim}[Iyilestirme Oenerisi]
Mevcut esit agirlikli sayma yerine agirlikli puanlama daha dogru olur:
\begin{lstlisting}[language=Python]
_LANGUAGE_WEIGHTS = {
    "rust_begin_unwind": (Language.RUST, 3),  # Cok guclu
    "std::":             (Language.CPP, 1),   # Zayif (cakisma)
    "runtime.gopanic":   (Language.GO, 3),    # Cok guclu
}
\end{lstlisting}
\end{tanim}

\subsection{False Positive Analizi}

\textbf{En yueksek risk: \code{std::}} --- Rust binary'leri C++ standard library ile linklenebilir (FFI), Go binary'leri cgo ile C++ kodu icerebilir. Bu durumda \code{std::} yanlis dili goesterir.

\textbf{En yueksek false negative riski:} Stripped binary'ler (agresif \code{-s -w}) cogu string'i siler.

\begin{uyari}[Tutarsizlik]
\code{\_detect\_language\_from\_binary()} 1MB okurken, \code{\_is\_go\_binary()} 2MB okuyor. Go binary'lerin bueyuekluegueuenden mantikli ama genel dil tespiti icin de 2MB daha dogru sonuc verir.
\end{uyari}


\section{Bundler Tespit Heuristikleri}
\label{sec:bundler-detail}

\subsection{webpack Runtime Yapisi}

webpack duenyanin en yaygin JS bundler'idir. Temel yapi:

\begin{lstlisting}[language=JavaScript]
var __webpack_modules__ = {
  "./src/index.js": function(module, exports, __webpack_require__) {
    // Modul kodu
  }
};
function __webpack_require__(moduleId) {
  var cachedModule = __webpack_module_cache__[moduleId];
  if (cachedModule !== undefined) return cachedModule.exports;
  var module = __webpack_module_cache__[moduleId] = { exports: {} };
  __webpack_modules__[moduleId](module, module.exports, __webpack_require__);
  return module.exports;
}
\end{lstlisting}

\code{\_\_webpack\_require\_\_} fonksiyonu: (1) Cache'e bak, (2) modul fonksiyonunu CommonJS convention ile cagir, (3) sonucu cache'e yaz.

\subsection{esbuild --- Tespiti Zor}

esbuild (Go ile yazilmis) minimal output ueretir --- belirgin bir runtime signature'i \textbf{yoktur}. Bu kasitli tasarim karari. Karadul'un mevcut sistemi esbuild ciktisini tespit \textbf{edemez} --- gueuencel JS ekosisteminde giderek onemli bir eksiklik.

\subsection{Bundler Tespitinin Ooenemi}

Bundler tespiti deobfuscation stratejisini belirler:
\begin{itemize}
\item \textbf{webpack:} Module map parse edilebilir, her modul ayri dosyaya yazilabilir
\item \textbf{parcel:} Benzer module map yapisi
\item \textbf{rollup:} Tek dosya, modul sinirlari silindigi icin parcalamak zor
\item \textbf{esbuild:} Agresif inlining, moduelleri ayirmak neredeyse imkansiz
\end{itemize}


\section{SHA-256 Hashing --- Detayli Analiz}
\label{sec:sha256-detail}

\subsection{Merkle-Damg\r{a}rd Construction}

SHA-256 su pipeline'i izler:

\begin{equation}
M \xrightarrow{\text{padding}} M_1 | M_2 | \cdots | M_n \quad \text{(512-bit bloklar)}
\end{equation}

\begin{align}
H_0 &= \text{IV (Initial Value, 256-bit sabit)} \notag \\
H_i &= f(H_{i-1}, M_i) \quad i = 1, 2, \ldots, n \notag \\
\text{Hash} &= H_n \notag
\end{align}

Compression function $f$: 64 round, her round'da Ch, Maj, $\Sigma$ fonksiyonlari (bitwise operations).

\subsection{8KB Chunk Boyutu Optimalligi}

\begin{table}[ht]
\centering
\begin{tabular}{lll}
\toprule
\textbf{\color{sectioncyan}Chunk Boyutu} & \textbf{\color{sectioncyan}Throughput} & \textbf{\color{sectioncyan}Bellek} \\
\midrule
256 B & $\sim$200 MB/s & 256 B \\
4 KB & $\sim$800 MB/s & 4 KB \\
\textbf{8 KB} & $\sim$\textbf{850 MB/s} & \textbf{8 KB} \\
64 KB & $\sim$870 MB/s & 64 KB \\
1 MB & $\sim$860 MB/s & 1 MB \\
\bottomrule
\end{tabular}
\caption{Hash chunk boyutu vs throughput. 8KB ``sweet spot''.}
\end{table}

8KB = 128 $\times$ 64B (SHA-256 compression block size) --- mueuekemmel hizalanma. L1 cache'te ($\sim$32--64KB) rahatca sigar.

\subsection{Collision Olasiligi}

Birthday paradox'undan:

\begin{equation}
\prob{\text{collision}, n \text{ dosya}} = 1 - e^{-n^2 / (2 \cdot 2^{256})} \approx \frac{n^2}{2^{257}}
\end{equation}

$n = 10^{12}$ (1 trilyon dosya): $P \approx 10^{18} / 2^{257} \approx 10^{-59}$ --- astronomik derecede dueueseek.

$P = 0.5$ icin gerekli $n$: $n \approx 2^{128} \approx 3.4 \times 10^{38}$ dosya --- evrendeki atom sayisindan ($\sim 10^{80}$) bile kat kat azdir, yine de collision pratikte imkansiz.



% ============================================================================
% BOLUM 4: STATIK ANALIZ
% ============================================================================
\chapter{Statik Analiz}
\label{ch:static}

Pipeline'in ikinci asamasi olan statik analiz, dosyayi calisti\-rmadan yapisini analiz eder.

\section{Binary Statik Analiz}

\subsection{Mach-O Analizi}

Apple platformlarinda kullanilan Mach-O formatinin yapisi:

\begin{center}
\begin{tikzpicture}[
  block/.style={rectangle, draw=sectioncyan, fill=boxbg, minimum width=6cm, minimum height=0.8cm, text=white, font=\small},
]
  \node[block] (header) at (0,0) {Mach-O Header (magic, cputype, filetype)};
  \node[block] (lc1) at (0,-1) {Load Command 1 (LC\_SEGMENT\_64)};
  \node[block] (lc2) at (0,-2) {Load Command 2 (LC\_SYMTAB)};
  \node[block] (lcn) at (0,-3) {Load Command N (LC\_DYSYMTAB)};
  \node[block] (text) at (0,-4.5) {\_\_TEXT Segment (code)};
  \node[block] (data) at (0,-5.5) {\_\_DATA Segment (variables)};
  \node[block] (link) at (0,-6.5) {\_\_LINKEDIT Segment (symbols)};
\end{tikzpicture}
\end{center}

\subsection{Symbol Extraction}

Binary'den sembol cikarmak icin \code{nm} ve \code{otool} araclari kullanilir:

\begin{itemize}
\item \code{nm -gU binary}: Global tanimli semboller
\item \code{otool -tV binary}: Disassemble edilmis kod
\item \code{strings binary}: Okunabilir string'ler
\end{itemize}

\subsection{Ghidra Headless Analiz}

Buyuek binary'ler icin Ghidra headless mode:

\begin{lstlisting}[language=bash,caption={Ghidra headless invocation}]
analyzeHeadless /tmp/project ProjectName \
  -import /path/to/binary \
  -scriptPath /path/to/scripts \
  -postScript ExportDecompiled.java \
  -deleteProject
\end{lstlisting}

Ghidra, binary'yi P-code adli bir intermediate representation'a donuesteruer, sonra C koduna decompile eder.

\section{JavaScript Statik Analiz}

JS analizi Babel AST parser kullanir:

\begin{lstlisting}[language=JavaScript,caption={Babel AST parse ornegi}]
const ast = parser.parse(sourceCode, {
  sourceType: "module",
  plugins: ["jsx", "typescript", "decorators"]
});
// AST traversal
traverse(ast, {
  FunctionDeclaration(path) {
    functions.push(path.node.id.name);
  },
  StringLiteral(path) {
    strings.push(path.node.value);
  }
});
\end{lstlisting}

Cikarilan bilgiler:
\begin{itemize}
\item Fonksiyon sayisi ve isimleri
\item String literal'ler
\item Import/export ifadeleri
\item Webpack modul ID'leri
\end{itemize}

\begin{ozet}
Statik analiz, binary'ler icin sembol tablosu, Ghidra decompilation ve string extraction; JavaScript icin Babel AST parsing, fonksiyon/string/import cikartma islemleri gerceklestirir. Buyuek dosyalar icin chunked processing destegi vardir.
\end{ozet}

% --- DEVELOPER AGENT DETAYLI GENISLETME: BOLUM 3+4 ---
% ============================================================================
% BOLUM 3: HEDEF TESPIT SISTEMI
% BOLUM 4: STATIK ANALIZ
% Karadul (Black Widow) --- Teknik Dokuman
% ============================================================================
% Bu dosya \include veya \input ile main.tex'ten dahil edilir.
% \chapter{} main.tex'te tanimlidir; burada \section ve \subsection kullanilir.
% ============================================================================

% ============================================================================
%                       B O L U M   3
%              HEDEF TESPIT SISTEMI (Target Detection)
% ============================================================================

\section{Hedef Tespit Sistemine Genel Bakis}
\label{sec:target-overview}

Karadul pipeline'inin ilk asamasi olan \emph{hedef tespit sistemi}, analiz edilecek dosyanin turunu, dilini ve metadata'sini belirler. Bu asama, sonraki tum islemlerin temelini oluturur: yanlis bir tespit, yanlis analyzer zincirine yonlendirme ve basarisiz analiz demektir.

\begin{tanim}[Hedef Tespit Problemi]
Bir $b \in \mathcal{B}$ binary dosyasi verildiginde, hedef tespit fonksiyonu:
\begin{equation}
\label{eq:target-detect}
\mathcal{D}: \mathcal{B} \longrightarrow \mathcal{T} \times \mathcal{L} \times \mathcal{M}
\end{equation}
seklinde tanimlanir. Burada:
\begin{itemize}
\item $\mathcal{T} = \{\text{MACHO}, \text{ELF}, \text{PE}, \text{JS\_BUNDLE}, \text{ELECTRON}, \text{GO}, \text{UNIVERSAL}, \ldots\}$: 12 eleman\-li hedef tipi kumesi.
\item $\mathcal{L} = \{\text{C}, \text{C++}, \text{Rust}, \text{Go}, \text{Swift}, \text{ObjC}, \text{JavaScript}, \ldots\}$: 13 eleman\-li dil kumesi.
\item $\mathcal{M}$: metadata sozlugu (mimari, bundler, framework vb.).
\end{itemize}
\end{tanim}

Tespit sistemi \file{target.py} dosyasinda \code{TargetDetector} sinifi olarak implement edilmistir. Toplam 630 satir Python kodundan olusur.

% --- TikZ: Tespit Akis Diyagrami ---
\begin{figure}[htbp]
\centering
\begin{tikzpicture}[
    node distance=1.5cm and 2.5cm,
    every node/.style={font=\small},
    block/.style={rectangle, draw=sectioncyan, fill=boxbg, text=white,
                  minimum width=3.5cm, minimum height=1cm, rounded corners=3pt},
    decision/.style={diamond, draw=yellow!70!white, fill=boxbg, text=white,
                     aspect=2, inner sep=2pt},
    arrow/.style={-{Stealth[length=3mm]}, thick, white},
]
    % Nodes
    \node[block] (input) {Dosya Yolu Girisi};
    \node[decision, below=of input] (isdir) {Dizin mi?};
    \node[block, right=of isdir] (appbundle) {.app Bundle Analizi};
    \node[decision, below=of isdir] (ext) {Uzanti Kontrolu};
    \node[block, right=of ext] (jsdetect) {JS Bundle Analizi};
    \node[block, below=of ext] (magic) {Magic Bytes Okuma};
    \node[decision, below=of magic] (format) {Format Tespiti};
    \node[block, below left=1.2cm and 0.5cm of format] (macho) {Mach-O};
    \node[block, below=of format] (elf) {ELF};
    \node[block, below right=1.2cm and 0.5cm of format] (pe) {PE};
    \node[block, right=3cm of format] (zip) {ZIP (JAR/APK)};
    \node[block, below=2.5cm of format] (lang) {Dil Tespiti};
    \node[block, below=of lang] (output) {\textbf{TargetInfo}};

    % Arrows
    \draw[arrow] (input) -- (isdir);
    \draw[arrow] (isdir) -- node[above, font=\tiny]{.app} (appbundle);
    \draw[arrow] (isdir) -- node[right, font=\tiny]{dosya} (ext);
    \draw[arrow] (ext) -- node[above, font=\tiny]{.js/.mjs} (jsdetect);
    \draw[arrow] (ext) -- node[right, font=\tiny]{diger} (magic);
    \draw[arrow] (magic) -- (format);
    \draw[arrow] (format) -- (macho);
    \draw[arrow] (format) -- (elf);
    \draw[arrow] (format) -- (pe);
    \draw[arrow] (format) -- (zip);
    \draw[arrow] (macho) |- (lang);
    \draw[arrow] (elf) -- (lang);
    \draw[arrow] (pe) |- (lang);
    \draw[arrow] (lang) -- (output);
    \draw[arrow] (jsdetect) |- (output);
    \draw[arrow] (appbundle) |- (output);
\end{tikzpicture}
\caption{Hedef tespit akis diyagrami. Dosya yolundan TargetInfo'ya uzanan karar agaci.}
\label{fig:target-flow}
\end{figure}


% ===================================================================
\section{Magic Byte Sistemi}
\label{sec:magic-bytes}

Her binary dosya formati, dosyanin ilk birkac byte'inda kendine ozgu bir \emph{magic number} barindirir. Bu sabitler, dosya turunu hizli ve guvenilir sekilde tanimlama amaciyla kullanilir.

\subsection{Temel Magic Byte Tanimlari}

\file{target.py:L53--L66}'da tanimlanan sabitler:

\begin{table}[htbp]
\centering
\begin{tabular}{llll}
\toprule
\textbf{\color{sectioncyan}Format} & \textbf{\color{sectioncyan}Magic (hex)} & \textbf{\color{sectioncyan}Endianness} & \textbf{\color{sectioncyan}Python Sabiti} \\
\midrule
Mach-O 64-bit   & \texttt{0xFEEDFACF} & Big-endian    & \code{\_MACHO\_64\_MAGIC} \\
Mach-O 32-bit   & \texttt{0xFEEDFACE} & Big-endian    & \code{\_MACHO\_32\_MAGIC} \\
Mach-O 64 Cigam & \texttt{0xCFFAEDFE} & Little-endian & \code{\_MACHO\_64\_CIGAM} \\
Mach-O 32 Cigam & \texttt{0xCEFAEDFE} & Little-endian & \code{\_MACHO\_32\_CIGAM} \\
Universal       & \texttt{0xCAFEBABE} & Big-endian    & \code{\_UNIVERSAL\_MAGIC} \\
ELF             & \texttt{7F 45 4C 46} & ---           & \code{\_ELF\_MAGIC} \\
PE (DOS)        & \texttt{4D 5A}       & ---           & \code{\_PE\_MAGIC} \\
ZIP (JAR/APK)   & \texttt{50 4B 03 04} & ---           & \code{\_ZIP\_MAGIC} \\
\bottomrule
\end{tabular}
\caption{Desteklenen binary formatlarinin magic byte'lari.}
\label{tab:magic-bytes}
\end{table}

\subsection{Big-Endian ve Little-Endian (Cigam) Kavrami}
\label{sec:endianness}

Mach-O formati iki farkli byte sirasini destekler. ``Cigam'' sozcugu ``magic'' kelimesinin tersten yazilisidir ve little-endian kodlamayi isaret eder.

\begin{tanim}[Byte Order]
$n$ byte'lik bir $v$ tamsayisi icin:
\begin{equation}
\label{eq:big-endian}
v_{\text{BE}} = \sum_{i=0}^{n-1} b_i \cdot 256^{n-1-i}
\quad\text{(big-endian)}
\end{equation}
\begin{equation}
\label{eq:little-endian}
v_{\text{LE}} = \sum_{i=0}^{n-1} b_i \cdot 256^{i}
\quad\text{(little-endian)}
\end{equation}
Burada $b_i$, dosyadan okunan $i$'inci byte'dir.
\end{tanim}

\paragraph{Sayisal Ornek.}
\texttt{0xFEEDFACF} sabitinin bellekteki gorunumu:

\begin{center}
\begin{tikzpicture}[
    cell/.style={rectangle, draw=sectioncyan, fill=boxbg, text=white,
                 minimum width=1.8cm, minimum height=0.8cm},
]
    % Big-endian
    \node at (-3, 1.2) {\color{sectioncyan}\textbf{Big-endian (Magic):}};
    \node[cell] at (0, 0.5) {\texttt{0xFE}};
    \node[cell] at (2, 0.5) {\texttt{0xED}};
    \node[cell] at (4, 0.5) {\texttt{0xFA}};
    \node[cell] at (6, 0.5) {\texttt{0xCF}};
    \node at (8, 0.5) {\color{gray!70!white}\small addr 0..3};

    % Little-endian
    \node at (-3, -0.8) {\color{yellow!70!white}\textbf{Little-endian (Cigam):}};
    \node[cell] at (0, -1.5) {\texttt{0xCF}};
    \node[cell] at (2, -1.5) {\texttt{0xFA}};
    \node[cell] at (4, -1.5) {\texttt{0xED}};
    \node[cell] at (6, -1.5) {\texttt{0xFE}};
    \node at (8, -1.5) {\color{gray!70!white}\small addr 0..3};
\end{tikzpicture}
\end{center}

\code{\_read\_magic} metodu (\file{target.py:L604--L614}) dosyanin ilk 4 byte'ini her zaman big-endian olarak okur:

\begin{lstlisting}[language=Python, caption={struct.unpack ile magic byte okuma}]
@staticmethod
def _read_magic(path: Path) -> int | None:
    """Dosyanin ilk 4 byte'ini big-endian uint32 olarak oku."""
    try:
        with open(path, "rb") as f:
            data = f.read(4)
        if len(data) < 4:
            return None
        return struct.unpack(">I", data)[0]  # ">" = big-endian, "I" = uint32
    except OSError:
        return None
\end{lstlisting}

Bu tasarim sayesinde, big-endian Mach-O dosyalari dogrudan \texttt{0xFEEDFACF} ile eslesirken, little-endian dosyalar \texttt{0xCFFAEDFE} degerini dondurur ve \code{\_MACHO\_64\_CIGAM} sabiti ile eslestirilir.

\paragraph{struct.unpack Format Karakterleri.}

\begin{table}[htbp]
\centering
\begin{tabular}{cll}
\toprule
\textbf{\color{sectioncyan}Karakter} & \textbf{\color{sectioncyan}Anlam} & \textbf{\color{sectioncyan}Boyut} \\
\midrule
\texttt{>} & Big-endian byte order & --- \\
\texttt{<} & Little-endian byte order & --- \\
\texttt{I} & Unsigned 32-bit integer & 4 byte \\
\texttt{H} & Unsigned 16-bit integer & 2 byte \\
\texttt{Q} & Unsigned 64-bit integer & 8 byte \\
\texttt{B} & Unsigned 8-bit integer & 1 byte \\
\bottomrule
\end{tabular}
\caption{Python \texttt{struct.unpack} format karakterleri.}
\label{tab:struct-format}
\end{table}


% ===================================================================
\section{Mach-O Format Yapisi}
\label{sec:macho-format}

Apple'in Mach-O (Mach Object) formati, macOS, iOS, watchOS ve tvOS uzerinde kullanilan native binary formatidir. Karadul'un birincil hedef formatidir.

\subsection{Mach-O Header Yapisi}

Mach-O dosyasinin ilk 32 byte'i \texttt{mach\_header\_64} yapisini icerir:

\begin{table}[htbp]
\centering
\begin{tabular}{llrl}
\toprule
\textbf{\color{sectioncyan}Alan} & \textbf{\color{sectioncyan}Tip} & \textbf{\color{sectioncyan}Offset} & \textbf{\color{sectioncyan}Aciklama} \\
\midrule
\texttt{magic}      & uint32 & 0  & \texttt{0xFEEDFACF} (64-bit) \\
\texttt{cputype}    & int32  & 4  & CPU ailesi (ARM64: \texttt{0x100000C}) \\
\texttt{cpusubtype} & int32  & 8  & CPU alt tipi \\
\texttt{filetype}   & uint32 & 12 & Dosya tipi (2=executable, 6=dylib) \\
\texttt{ncmds}      & uint32 & 16 & Load command sayisi \\
\texttt{sizeofcmds} & uint32 & 20 & Tum load command'larin toplam boyutu \\
\texttt{flags}      & uint32 & 24 & Binary flag'leri (PIE, stripped vb.) \\
\texttt{reserved}   & uint32 & 28 & Rezerve (64-bit icin) \\
\bottomrule
\end{tabular}
\caption{Mach-O 64-bit header alanlari.}
\label{tab:macho-header}
\end{table}

% --- TikZ: Mach-O File Layout ---
\begin{figure}[htbp]
\centering
\begin{tikzpicture}[
    segment/.style={rectangle, draw=sectioncyan, fill=boxbg, text=white,
                    minimum width=6cm, minimum height=0.9cm},
    label/.style={font=\small\ttfamily, text=gray!70!white},
]
    \node[segment, fill=blue!20!black] (header) at (0, 0)
        {\textbf{Mach-O Header} (32 bytes)};

    \node[segment, fill=cyan!15!black] (lc1) at (0, -1.2)
        {LC\_SEGMENT\_64 (\texttt{\_\_TEXT})};
    \node[segment, fill=cyan!15!black] (lc2) at (0, -2.4)
        {LC\_SEGMENT\_64 (\texttt{\_\_DATA})};
    \node[segment, fill=cyan!10!black] (lc3) at (0, -3.6)
        {LC\_SYMTAB};
    \node[segment, fill=cyan!10!black] (lc4) at (0, -4.8)
        {LC\_DYSYMTAB};
    \node[segment, fill=cyan!8!black] (lcn) at (0, -6.0)
        {LC\_LOAD\_DYLIB $\times$ N};

    \node[segment, fill=green!15!black] (text) at (0, -7.5)
        {\textbf{\_\_TEXT Segment} (kod + string'ler)};
    \node[segment, fill=yellow!10!black] (data) at (0, -8.7)
        {\textbf{\_\_DATA Segment} (global veriler)};
    \node[segment, fill=red!10!black] (link) at (0, -9.9)
        {\textbf{\_\_LINKEDIT} (sembol tablolari)};

    % Labels
    \node[label, right] at (3.5, 0) {Offset 0x00};
    \node[label, right] at (3.5, -1.2) {Load Commands};
    \node[label, right] at (3.5, -6.0) {baslangi\c{c}};
    \node[label, right] at (3.5, -7.5) {Executable kod};
    \node[label, right] at (3.5, -8.7) {Read-write};
    \node[label, right] at (3.5, -9.9) {Read-only};

    % Brace for Load Commands
    \draw[decorate, decoration={brace, amplitude=8pt, mirror}, thick, white]
        (3.2, -0.6) -- (3.2, -6.5)
        node[midway, right=10pt, font=\small, text=white] {ncmds adet};
\end{tikzpicture}
\caption{Mach-O 64-bit dosya yapisi. Header + Load Commands + Segments.}
\label{fig:macho-layout}
\end{figure}

\subsection{Load Commands}

Load command'lar, isletim sisteminin binary'yi nasil yukleyecegini ve baglayacagini tarif eder. Karadul'un \code{MachOAnalyzer} sinifi (\file{macho.py:L389--L391}) \texttt{otool -l} komutuyla tum load command'lari cikarir.

Onemli load command turleri:

\begin{description}
\item[\texttt{LC\_SEGMENT\_64}:] Bir segment'in bellek adresini, boyutunu ve icerigini tanimlar. Her segment birden fazla \emph{section} icerebilir.
\item[\texttt{LC\_SYMTAB}:] Sembol tablosunun offset ve boyutunu belirtir.
\item[\texttt{LC\_DYSYMTAB}:] Dinamik sembol tablosu bilgilerini icerir.
\item[\texttt{LC\_LOAD\_DYLIB}:] Yuklenmesi gereken bir dinamik kutuphanenin yolunu belirtir.
\item[\texttt{LC\_UUID}:] Binary'nin benzersiz tanimlayicisi (build ID).
\item[\texttt{LC\_CODE\_SIGNATURE}:] Apple code signing bilgisi.
\end{description}

\subsection{Segment ve Section Yapisi}

Mach-O dosyasindaki temel segmentler:

\begin{table}[htbp]
\centering
\begin{tabular}{llp{7cm}}
\toprule
\textbf{\color{sectioncyan}Segment} & \textbf{\color{sectioncyan}Section'lar} & \textbf{\color{sectioncyan}Aciklama} \\
\midrule
\texttt{\_\_TEXT} & \texttt{\_\_text, \_\_stubs, \_\_stub\_helper, \_\_cstring, \_\_const} & Executable kod, sabit string'ler. Read-only + execute. \\
\texttt{\_\_DATA} & \texttt{\_\_data, \_\_bss, \_\_common, \_\_la\_symbol\_ptr, \_\_got} & Global degiskenler, BSS. Read-write. \\
\texttt{\_\_DATA\_CONST} & \texttt{\_\_const, \_\_cfstring, \_\_objc\_classlist} & Salt-okunur veri (runtime'da write-protect). \\
\texttt{\_\_LINKEDIT} & (monolitik) & Sembol tablolari, string tablolari, code signature. Read-only. \\
\bottomrule
\end{tabular}
\caption{Mach-O segmentleri ve section'lari.}
\label{tab:macho-segments}
\end{table}

Karadul'un \code{lief} entegrasyonu (\file{macho.py:L431--L495}) tum segment ve section bilgilerini programatik olarak cikarir:

\begin{lstlisting}[language=Python, caption={lief ile segment/section cikartma}]
binary = lief.parse(str(binary_path))
for seg in binary.segments:
    result["segments"].append({
        "name": seg.name,
        "virtual_address": seg.virtual_address,
        "virtual_size": seg.virtual_size,
        "file_offset": seg.file_offset,
        "file_size": seg.file_size,
    })
for sec in binary.sections:
    result["sections"].append({
        "name": sec.name,
        "segment": sec.segment.name,
        "size": sec.size,
        "offset": sec.offset,
    })
\end{lstlisting}


% ===================================================================
\section{ELF Format Yapisi}
\label{sec:elf-format}

ELF (Executable and Linkable Format), Linux ve diger Unix-benzeri sistemlerde kullanilan standart binary formatidir. Magic byte'i \texttt{7F 45 4C 46} ($\backslash$x7fELF) ile baslar (\file{target.py:L64}).

\subsection{ELF Header}

\begin{table}[htbp]
\centering
\begin{tabular}{llrl}
\toprule
\textbf{\color{sectioncyan}Alan} & \textbf{\color{sectioncyan}Tip} & \textbf{\color{sectioncyan}Offset} & \textbf{\color{sectioncyan}Aciklama} \\
\midrule
\texttt{e\_ident[0..3]}  & byte[4]  & 0   & Magic: \texttt{7F 45 4C 46} \\
\texttt{e\_ident[4]}     & uint8    & 4   & Class: 1=32-bit, 2=64-bit \\
\texttt{e\_ident[5]}     & uint8    & 5   & Data: 1=LE, 2=BE \\
\texttt{e\_type}         & uint16   & 16  & Tip: 2=ET\_EXEC, 3=ET\_DYN \\
\texttt{e\_machine}      & uint16   & 18  & Mimari: 0x3E=x86\_64, 0xB7=AArch64 \\
\texttt{e\_entry}        & uint64   & 24  & Entry point adresi \\
\texttt{e\_phoff}        & uint64   & 32  & Program header table offset \\
\texttt{e\_shoff}        & uint64   & 40  & Section header table offset \\
\texttt{e\_phnum}        & uint16   & 56  & Program header sayisi \\
\texttt{e\_shnum}        & uint16   & 60  & Section header sayisi \\
\bottomrule
\end{tabular}
\caption{ELF 64-bit header alanlari (secilmis).}
\label{tab:elf-header}
\end{table}

\subsection{ELF ve Mach-O Karsilastirmasi}

\begin{table}[htbp]
\centering
\begin{tabular}{lll}
\toprule
\textbf{\color{sectioncyan}Kavram} & \textbf{\color{sectioncyan}Mach-O} & \textbf{\color{sectioncyan}ELF} \\
\midrule
Kod bolumu    & \texttt{\_\_TEXT,\_\_text} & \texttt{.text} \\
Veri bolumu   & \texttt{\_\_DATA,\_\_data} & \texttt{.data} \\
Salt-okunur   & \texttt{\_\_DATA\_CONST}   & \texttt{.rodata} \\
Sembol tablosu & \texttt{LC\_SYMTAB}       & \texttt{.symtab} / \texttt{.dynsym} \\
Dinamik baglama & \texttt{LC\_LOAD\_DYLIB} & \texttt{.dynamic} + \texttt{DT\_NEEDED} \\
Magic number  & \texttt{0xFEEDFACF}        & \texttt{7F 45 4C 46} \\
\bottomrule
\end{tabular}
\caption{Mach-O ve ELF terminoloji eslestirmesi.}
\label{tab:macho-elf-compare}
\end{table}

Karadul, ELF binary'ler icin \texttt{otool} adimlarini atlar (\file{macho.py:L122--L143}), dogrudan \texttt{nm}, \texttt{strings} ve Ghidra headless analiz kullanir.


% ===================================================================
\section{PE Format Yapisi}
\label{sec:pe-format}

PE (Portable Executable), Microsoft Windows'un native binary formatidir. Magic byte'i \texttt{4D 5A} (ASCII ``MZ'') ile baslar ve \emph{DOS header} icerir.

\subsection{PE Dosya Yapisi}

\begin{figure}[htbp]
\centering
\begin{tikzpicture}[
    segment/.style={rectangle, draw=sectioncyan, fill=boxbg, text=white,
                    minimum width=5.5cm, minimum height=0.8cm},
]
    \node[segment, fill=red!15!black] (dos) at (0, 0)
        {\textbf{DOS Header} (\texttt{MZ})};
    \node[segment, fill=red!10!black] (stub) at (0, -1.1)
        {DOS Stub};
    \node[segment, fill=blue!15!black] (pe) at (0, -2.2)
        {\textbf{PE Signature} (\texttt{PE\textbackslash 0\textbackslash 0})};
    \node[segment, fill=cyan!15!black] (coff) at (0, -3.3)
        {COFF Header};
    \node[segment, fill=cyan!12!black] (opt) at (0, -4.4)
        {Optional Header (PE32+)};
    \node[segment, fill=green!12!black] (sec) at (0, -5.5)
        {Section Table};
    \node[segment, fill=green!15!black] (text) at (0, -6.8)
        {\texttt{.text} (kod)};
    \node[segment, fill=yellow!10!black] (rdata) at (0, -7.9)
        {\texttt{.rdata} (salt-okunur veri)};
    \node[segment, fill=yellow!8!black] (data) at (0, -9.0)
        {\texttt{.data} (veri)};
    \node[segment, fill=red!8!black] (rsrc) at (0, -10.1)
        {\texttt{.rsrc} (kaynaklar)};

    % Offset labels
    \node[font=\small\ttfamily, text=gray!70!white, right] at (3, 0) {0x00};
    \node[font=\small\ttfamily, text=gray!70!white, right] at (3, -2.2) {e\_lfanew};
    \node[font=\small\ttfamily, text=gray!70!white, right] at (3, -6.8) {Virtual sections};
\end{tikzpicture}
\caption{PE dosya yapisi. DOS header'daki \texttt{e\_lfanew} alani PE signature'a isaret eder.}
\label{fig:pe-layout}
\end{figure}

Karadul'un PE tespit mantigi (\file{target.py:L402--L430}):

\begin{enumerate}
\item Ilk 2 byte \texttt{MZ} mi? $\to$ PE adayi.
\item .NET kontrolu: \code{\_is\_dotnet()} metodu \texttt{mscoree.dll} veya \texttt{\_CorExeMain} arar.
\item Delphi kontrolu: \code{\_is\_delphi()} metodu ``Borland Delphi'', mangled name pattern'leri arar.
\item Hic biri degilse: genel \code{PE\_BINARY} olarak siniflandirir.
\end{enumerate}


% ===================================================================
\section{Universal Binary (Fat Binary)}
\label{sec:universal-binary}

Apple'in Universal Binary formati, birden fazla mimarinin tek dosyada birlestirilmesini saglar. Magic byte'i \texttt{0xCAFEBABE}'dir.

\begin{uyari}[Dikkat: Java Class File Cakismasi]
\texttt{0xCAFEBABE}, Java class dosyalarinin da magic byte'idir. Ancak pratikte bu bir sorun yaratmaz, cunku:
\begin{enumerate}
\item Java class dosyalari genellikle .jar (ZIP) icinde gelir ve \texttt{PK} magic'i ile baslar.
\item Tek basina \texttt{.class} dosyalari nadiren analiz hedefi olur.
\item Karadul, \texttt{0xCAFEBABE} goruldugunde dosya boyutu ve icerik heuristikleri uygular.
\end{enumerate}
\end{uyari}

Universal binary tespit edildikten sonra (\file{macho.py:L85--L110}), \texttt{lipo} komutu ile tercih edilen mimari cikarilir:

\begin{lstlisting}[language=Python, caption={lipo thin ile mimari cikartma}]
# Oncelik: arm64 (Apple Silicon) > x86_64
proc = subprocess.run(
    ["lipo", str(binary_path), "-thin", "arm64",
     "-output", str(thin_path)],
    capture_output=True, text=True, timeout=30,
)
if proc.returncode == 0:
    binary_path = thin_path  # arm64 slice ile devam et
\end{lstlisting}


% ===================================================================
\section{Dil Tespit Algoritmasi}
\label{sec:language-detection}

Binary icindeki string'lerden kaynak programlama dilini tespit etmek, Karadul'un en onemli heuristik adimlarindan biridir. \file{target.py:L470--L491}'de implement edilen bu algoritma, bir \emph{puanlama sistemi} kullanir.

\subsection{Puanlama Sistemi}

\begin{tanim}[Dil Puanlama]
$\mathcal{S} = \{(s_i, l_i)\}_{i=1}^{N}$ imza kumesi verildiginde, her dil $l \in \mathcal{L}$ icin puan:
\begin{equation}
\label{eq:lang-score}
\text{score}(l) = \sum_{i=1}^{N} \mathbb{1}[s_i \in \text{binary\_text}] \cdot \mathbb{1}[l_i = l]
\end{equation}
Tespit edilen dil:
\begin{equation}
\label{eq:lang-argmax}
\hat{l} = \arg\max_{l \in \mathcal{L}} \text{score}(l)
\end{equation}
Eger $\text{score}(\hat{l}) = 0$ ise, $\hat{l} = \text{UNKNOWN}$ dondurulur.
\end{tanim}

\subsection{Dil Imza Tablosu}

\file{target.py:L136--L154}'te tanimlanan \code{\_LANGUAGE\_SIGNATURES} sozlugu:

\begin{table}[htbp]
\centering
\begin{tabular}{lll}
\toprule
\textbf{\color{sectioncyan}String Imzasi} & \textbf{\color{sectioncyan}Dil} & \textbf{\color{sectioncyan}Aciklama} \\
\midrule
\texttt{rust\_begin\_unwind}  & Rust      & Panic handler baslangici \\
\texttt{rust\_panic}          & Rust      & Panic cagirma fonksiyonu \\
\texttt{\_\_rust\_alloc}      & Rust      & Bellek ayirici \\
\texttt{core::panicking}      & Rust      & Core kutuphane panic modulu \\
\texttt{runtime.gopanic}      & Go        & Go runtime panic \\
\texttt{runtime.goexit}       & Go        & Goroutine cikis \\
\texttt{runtime/internal}     & Go        & Runtime internal paket \\
\texttt{go.buildid}           & Go        & Build identifier \\
\texttt{\_swift\_}            & Swift     & Swift runtime prefix \\
\texttt{Swift.}               & Swift     & Swift modulleri \\
\texttt{\_\_cxa\_throw}       & C++       & C++ exception throwing \\
\texttt{std::}                & C++       & C++ standard kutuphane \\
\texttt{objc\_msgSend}        & Obj-C     & Objective-C mesaj gonderme \\
\texttt{@objc}                & Obj-C     & Objective-C dekoratoru \\
\texttt{NSObject}             & Obj-C     & Foundation sinifi \\
\bottomrule
\end{tabular}
\caption{Dil tespit imza tablosu (\code{\_LANGUAGE\_SIGNATURES}).}
\label{tab:lang-sigs}
\end{table}

\subsection{Naive Bayes Perspektifi}

Dil tespitini Bayesian bir cercevede formalize edebiliriz. Her imza $s_i$'nin binary'de bulunmasi, dilin $l$ olmasi olasiligini guncelleyen bir \emph{evidence} olarak gorulur.

\begin{equation}
\label{eq:naive-bayes-lang}
P(l \given s_1, s_2, \ldots, s_N) \propto P(l) \prod_{i=1}^{N} P(s_i \given l)
\end{equation}

Karadul'un mevcut implementasyonunda $P(l)$ uniform prior ve $P(s_i \given l) \in \{0, 1\}$ (deterministik) olarak basitlestirilmistir. Bu, formul~(\ref{eq:lang-score})'daki toplam sayma ile esdegerdir.

\paragraph{Karmasiklik Analizi.} Dosyanin ilk 1~MB'i okunur ($\bigO{1}$ --- sabit boyut), $N = 15$ imza uzerinde dogrusal tarama yapilir. Toplam:
\begin{equation}
\label{eq:lang-complexity}
T_{\text{lang}} = \bigO{|\text{1MB}| \times N} = \bigO{N} \approx \bigO{1}
\end{equation}


% ===================================================================
\section{Go Binary Tespiti}
\label{sec:go-detection}

Go dili, derlenmis binary'lere son derece zengin metadata gomer. Bu metadata, tersine muhendislik icin degerlidir ve Karadul bunu iki asamada tespit eder.

\subsection{GOPCLNTAB (Program Counter Line Table)}

Go derleyicisi, her binary'ye bir \texttt{pclntab} (program counter line table) ekler. Bu tablo fonksiyon adlarini, dosya yollarini ve satir numaralarini icerir. \file{go\_binary.py:L47--L60}'da tanimlanan magic byte'lar:

\begin{table}[htbp]
\centering
\begin{tabular}{ll}
\toprule
\textbf{\color{sectioncyan}Go Versiyonu} & \textbf{\color{sectioncyan}GOPCLNTAB Magic} \\
\midrule
Go 1.2+  & \texttt{0xFB 0xFF 0xFF 0xFF 0x00 0x00} \\
Go 1.16+ & \texttt{0xFA 0xFF 0xFF 0xFF 0x00 0x00} \\
Go 1.18+ & \texttt{0xF0 0xFF 0xFF 0xFF 0x00 0x00} \\
Go 1.20+ & \texttt{0xF1 0xFF 0xFF 0xFF 0x00 0x00} \\
\bottomrule
\end{tabular}
\caption{GOPCLNTAB magic byte'lari Go versiyonlarina gore.}
\label{tab:gopclntab-magic}
\end{table}

\file{target.py:L493--L516}'deki \code{\_is\_go\_binary} metodu, ilk 2~MB icinde bu magic byte'lari ve Go runtime string'lerini arar:

\begin{lstlisting}[language=Python, caption={Go binary tespit algoritmasi}]
def _is_go_binary(self, path: Path) -> bool:
    with open(path, "rb") as f:
        data = f.read(2_097_152)  # 2 MB
    # GOPCLNTAB magic bytes
    if b"GOPCLNTAB" in data or b"\xFB\xFF\xFF\xFF\x00\x00" in data:
        return True
    # go.buildid section
    if b"go.buildid" in data or b"Go build" in data:
        return True
    # Runtime string'leri
    text = data.decode("ascii", errors="replace")
    go_markers = ("runtime.gopanic", "runtime.goexit", "runtime/internal")
    return any(marker in text for marker in go_markers)
\end{lstlisting}

\subsection{Go Binary Isim Kurtarma Zenginligi}

Go binary'ler, strip edilse bile GOPCLNTAB genellikle korunur. Bu, Go'yu tersine muhendislik icin ``en kolay hedef'' yapar:

\begin{equation}
\label{eq:go-recovery}
\text{Kurtarma Orani}_{\text{Go}} \approx 70\%\text{--}95\%
\end{equation}

Karsilastirma olarak, tipik bir stripped C/C++ binary'de bu oran $<5\%$'tir.


% ===================================================================
\section{.NET ve Delphi Tespiti}
\label{sec:dotnet-delphi}

\subsection{.NET Assembly Tespiti}

.NET assembly'leri PE formati icinde paketlenmis IL (Intermediate Language) kodudur. Karadul, \file{target.py:L518--L536}'daki \code{\_is\_dotnet} metodu ile tespit eder:

\begin{itemize}
\item \texttt{mscoree.dll}: .NET CLR runtime yukleyicisi --- her .NET PE'de bulunur.
\item \texttt{\_CorExeMain}: .NET entry point fonksiyonu.
\item \texttt{mscorlib} veya \texttt{System.Runtime}: .NET framework referanslari.
\end{itemize}

Herhangi birinin binary icinde bulunmasi, .NET assembly oldugunu gosterir.

\subsection{Delphi Binary Tespiti}

Delphi compiler'in urettigi binary'lerde kendine ozgu \emph{mangled name} pattern'leri vardir. \file{target.py:L538--L571}'deki \code{\_is\_delphi} metodu iki asamali tespit uygular:

\textbf{Asama 1 --- Bilinen string'ler:} ``Borland Delphi'', ``CodeGear Delphi'', ``Embarcadero Delphi'', \texttt{@System@TObject@}, \texttt{System.SysUtils} string'lerinden en az 2 tanesinin bulunmasi yeterlidir.

\textbf{Asama 2 --- Mangled name regex:}
\begin{equation}
\label{eq:delphi-regex}
\texttt{@[A-Z][\ldots]@[A-Z][\ldots]@[A-Za-z\_][\ldots]\$qqr}
\end{equation}
Bu pattern, Delphi'nin \texttt{@Unit@Class@Method\$qqr} seklindeki mangled fonksiyon isimlerini yakalar. 5 veya daha fazla eslesme Delphi binary oldugunu gosterir.

\begin{equation}
\label{eq:delphi-confidence}
\text{confidence}_{\text{Delphi}} =
\begin{cases}
0.95 & \text{string\_hits} \geq 2 \\
0.85 & \text{regex\_matches} \geq 5 \\
0.0  & \text{aksi halde}
\end{cases}
\end{equation}


% ===================================================================
\section{JavaScript Bundle Tespiti}
\label{sec:js-detection}

JavaScript dosyalari uzanti tabanli (\texttt{.js}, \texttt{.mjs}, \texttt{.cjs}) tespit edilir ve ardindan bundler tespiti yapilir.

\subsection{Bundler Tespit Sistemi}

\file{target.py:L126--L133}'te tanimlanan imza tablosu:

\begin{table}[htbp]
\centering
\begin{tabular}{ll}
\toprule
\textbf{\color{sectioncyan}Imza String'i} & \textbf{\color{sectioncyan}Bundler} \\
\midrule
\texttt{\_\_webpack\_require\_\_} & webpack \\
\texttt{\_\_webpack\_modules\_\_} & webpack \\
\texttt{parcelRequire}            & Parcel \\
\texttt{System.register}          & SystemJS \\
\texttt{\_rollupPluginBabelHelpers} & Rollup \\
\texttt{define([}                  & AMD/RequireJS \\
\bottomrule
\end{tabular}
\caption{JS bundler imza tablosu.}
\label{tab:js-bundlers}
\end{table}

Tespit algoritmasi:
\begin{enumerate}
\item Dosyanin ilk 64~KB'i okunur (\file{target.py:L288--L291}).
\item Her imza string'i icin dogrusal arama yapilir.
\item Ilk eslesen bundler, metadata'ya kaydedilir.
\end{enumerate}

\paragraph{Karmasiklik.} $K = 6$ imza, $M = 65536$ byte okuma:
\begin{equation}
\label{eq:bundler-complexity}
T_{\text{bundler}} = \bigO{K \times M} = \bigO{1}
\end{equation}

\subsection{esbuild Tespiti}

esbuild, imza tablosunda yoktur cunku belirgin bir runtime marker'i birakmaz. Bu nedenle, bilinen bundler'lardan hicbiri eslesmediyse ve dosya minified gorunuyorsa, metadata'ya \texttt{bundler: null} olarak kaydedilir. Kullanici tarafindan manuel belirtilebilir.

\subsection{Webpack vs Parcel vs Rollup Farklari}

\begin{table}[htbp]
\centering
\begin{tabular}{lp{3.5cm}p{3.5cm}p{3.5cm}}
\toprule
\textbf{\color{sectioncyan}Ozellik} & \textbf{\color{sectioncyan}webpack} & \textbf{\color{sectioncyan}Parcel} & \textbf{\color{sectioncyan}Rollup} \\
\midrule
Runtime          & \texttt{\_\_webpack\_require\_\_} & \texttt{parcelRequire} & Yok (tree-shaken) \\
Modul ID         & Sayisal / string & Hash-based & Degisken adi \\
Unpack kolayligi & Kolay & Orta & Zor \\
\bottomrule
\end{tabular}
\caption{Bundler runtime farklari ve Karadul'un unpack stratejisi.}
\label{tab:bundler-compare}
\end{table}


% ===================================================================
\section{Electron .app Bundle Analizi}
\label{sec:electron-detection}

Electron uygulamalari, macOS'ta \texttt{.app} bundle formunda dagitilir. Karadul, \file{target.py:L215--L279} ve \file{electron.py} ile Electron tespiti yapar.

\subsection{Electron Tespit Algoritmasi}

\begin{enumerate}
\item \texttt{.app/Contents/Resources/app.asar} veya \texttt{.app/Contents/Resources/app/} dizini kontrol edilir.
\item Varsa: \texttt{ELECTRON\_APP} tipi atanir, dil \texttt{JAVASCRIPT} olarak ayarlanir.
\item Yoksa: standart macOS uygulamasi olarak siniflandirilir.
\end{enumerate}

\subsection{ASAR Arsiv Formati}

ASAR (Atom Shell Archive), Electron'un kullandigi ozel bir arsiv formatidir. JSON-tabanli bir header ve birlesmis dosya icerikleri icerir:

\begin{figure}[htbp]
\centering
\begin{tikzpicture}[
    segment/.style={rectangle, draw=sectioncyan, fill=boxbg, text=white,
                    minimum width=5cm, minimum height=0.7cm},
]
    \node[segment, fill=blue!15!black] (hdr) at (0, 0)
        {Header Size (4 byte, LE uint32)};
    \node[segment, fill=cyan!15!black] (json) at (0, -1.0)
        {JSON Header (dosya listesi + offset'ler)};
    \node[segment, fill=green!15!black] (f1) at (0, -2.0)
        {File Content 1};
    \node[segment, fill=green!12!black] (f2) at (0, -3.0)
        {File Content 2};
    \node[segment, fill=green!10!black] (fn) at (0, -4.0)
        {File Content N};
\end{tikzpicture}
\caption{ASAR arsiv yapisi.}
\label{fig:asar-format}
\end{figure}

ASAR extraction (\file{electron.py:L346--L397}):
\begin{lstlisting}[language=Python, caption={ASAR extraction komutu}]
cmd = [
    npx_path, "--yes", "@electron/asar",
    "extract", str(asar_path), str(output_dir),
]
result = self._runner.run_command(cmd, timeout=timeout)
\end{lstlisting}

\subsection{Ana JS Dosyasi Bulma Stratejisi}

Extract edilen ASAR icinde onlarca JS dosyasi olabilir. Karadul, ana dosyayi su sirayla arar (\file{electron.py:L453--L498}):

\begin{enumerate}
\item \texttt{package.json}'dan \texttt{"main"} field'ini oku.
\item Bilinen isim adaylarina bak: \texttt{dist/cli.js}, \texttt{dist/main.js}, \texttt{dist/index.js}, \texttt{app/main.js}, \ldots (12 aday; \file{electron.py:L32--L46}).
\item Hicbiri yoksa: en buyuk JS dosyasini sec.
\end{enumerate}


% ===================================================================
\section{SHA-256 Streaming Hash}
\label{sec:sha256-hash}

Her hedef icin benzersiz bir taniimlayici olarak SHA-256 hash hesaplanir. \file{target.py:L616--L629}'daki \code{\_compute\_sha256} metodu \emph{streaming} yaklasimiyla calisir.

\subsection{Streaming Hash Algoritmasi}

\begin{lstlisting}[language=Python, caption={8~KB chunk'larla SHA-256 hesaplama}]
_HASH_CHUNK_SIZE = 8192  # 8 KB

@staticmethod
def _compute_sha256(path: Path) -> str:
    h = hashlib.sha256()
    with open(path, "rb") as f:
        while True:
            chunk = f.read(8192)
            if not chunk:
                break
            h.update(chunk)
    return h.hexdigest()
\end{lstlisting}

\subsection{Neden 8~KB Chunk?}

\begin{itemize}
\item \textbf{Bellek verimliligi:} 224~MB binary tamamen RAM'e alinmaz; sabit 8~KB tampon kullanilir.
\item \textbf{Sayfa hizalanmasi:} Cogu isletim sistemi 4~KB sayfa boyutu kullanir; 8~KB = 2 sayfa, I/O kararliligi saglar.
\item \textbf{Performans:} Daha kucuk chunk'lar syscall overhead'i arttirir; daha buyuk chunk'lar cache'i kirletir. 8~KB iyi bir denge noktasidir.
\end{itemize}

\paragraph{Karmasiklik.}
\begin{equation}
\label{eq:sha256-complexity}
T_{\text{SHA256}} = \bigO{\frac{n}{8192}} = \bigO{n}
\end{equation}
burada $n$ dosya boyutudur (byte).

\subsection{Collision Olasiligi}

SHA-256, 256-bit cikti ureten bir kriptografik hash fonksiyonudur.

\begin{equation}
\label{eq:collision-birthday}
P(\text{collision} \given k \text{ dosya}) \approx 1 - e^{-k^2 / (2 \cdot 2^{256})}
\end{equation}

$k = 10^{9}$ dosya icin:
\begin{equation}
\label{eq:collision-numeric}
P \approx \frac{(10^9)^2}{2^{257}} \approx \frac{10^{18}}{2.31 \times 10^{77}} \approx 4.33 \times 10^{-59}
\end{equation}

Bu, pratik amaclar icin sifir kabul edilir. SHA-256'nin ikinci preimage direnci $2^{256}$ islem gerektirir ki bu evrenin yasi boyunca tum yeryuzu bilgisayarlariyla bile ulasilamaz.


% ===================================================================
\section{False Positive Analizi}
\label{sec:false-positive}

Magic byte tabanli tespit, false positive riski tasir. Bu bolum, her format icin cakisma olasiliklerini analiz eder.

\subsection{Magic Byte Collision Hesabi}

Rastgele bir dosyanin ilk $k$ byte'inin belirli bir magic byte ile eslesmesi olasiligi:

\begin{equation}
\label{eq:magic-fp}
P(\text{FP}_k) = \frac{1}{256^k}
\end{equation}

\begin{table}[htbp]
\centering
\begin{tabular}{lccc}
\toprule
\textbf{\color{sectioncyan}Format} & \textbf{\color{sectioncyan}$k$ (byte)} & \textbf{\color{sectioncyan}$P(\text{FP})$} & \textbf{\color{sectioncyan}1/N} \\
\midrule
Mach-O (4 byte magic)    & 4 & $2.33 \times 10^{-10}$ & 4.3 milyar dosyada 1 \\
ELF (4 byte magic)       & 4 & $2.33 \times 10^{-10}$ & 4.3 milyar dosyada 1 \\
PE (2 byte ``MZ'')       & 2 & $1.53 \times 10^{-5}$  & 65.536 dosyada 1 \\
ZIP (4 byte magic)       & 4 & $2.33 \times 10^{-10}$ & 4.3 milyar dosyada 1 \\
\bottomrule
\end{tabular}
\caption{Magic byte false positive olasiliklari.}
\label{tab:magic-fp}
\end{table}

PE formati icin 2 byte'lik magic, daha yuksek false positive riski tasir. Ancak Karadul, PE tespitinden sonra ek kontroller (PE signature, section table validity) uygulayarak bu riski azaltir.

\subsection{Coklu Kontrol Stratejisi}

Karadul'un false positive oranini azaltan katmanli yaklasimi:

\begin{equation}
\label{eq:multi-check-fp}
P(\text{FP}_{\text{toplam}}) = P(\text{FP}_{\text{magic}}) \times P(\text{FP}_{\text{icerik}}) \times P(\text{FP}_{\text{uzanti}})
\end{equation}

Ornegin, PE icin:
\begin{equation}
P(\text{FP}_{\text{PE}}) = \frac{1}{256^2} \times P(\text{valid section table}) \times P(\texttt{.exe}/\texttt{.dll}) \approx 10^{-12}
\end{equation}


\begin{ozet}
Hedef Tespit Sistemi, Karadul pipeline'inin temeli olarak 12 farkli binary formatini ve 13 programlama dilini tespit eder. Magic byte karsilastirmasi $\bigO{1}$ zamanda calisir, dil tespiti puanlama tabanlidir ve $\bigO{1}$ karmasikliga sahiptir. SHA-256 streaming hash $\bigO{n}$ ile dosya butunlugunu garanti eder. False positive orani, coklu kontrol stratejisi ile $10^{-10}$ mertebesinin altinda tutulur. Tum bu islemler \file{target.py}'daki 630 satirlik \code{TargetDetector} sinifinda birlestirilmistir.
\end{ozet}


% ============================================================================
%                       B O L U M   4
%                     STATIK ANALIZ
% ============================================================================

\newpage

\section{Statik Analize Genel Bakis}
\label{sec:static-overview}

Statik analiz, bir programi \emph{calistirmadan} --- yalnizca yapisinan inceleyerek --- bilgi cikarma surecidir. Karadul'un statik analiz katmani iki ana dal icindenir:

\begin{enumerate}
\item \textbf{Binary Statik Analiz} (\file{macho.py}): Mach-O, ELF, PE binary'ler icin otool, nm, strings, lief ve Ghidra headless.
\item \textbf{JavaScript Statik Analiz} (\file{javascript.py}): Babel AST ile fonksiyon/string/import/export cikartma.
\end{enumerate}

\begin{tanim}[Statik Analiz Formalizasyonu]
Bir hedef $t \in \mathcal{T}$ uzerinde statik analiz fonksiyonu:
\begin{equation}
\label{eq:static-analysis}
\mathcal{A}_{\text{static}}: \mathcal{T} \longrightarrow \mathcal{F} \times \mathcal{S} \times \mathcal{I} \times \mathcal{G}
\end{equation}
Burada:
\begin{itemize}
\item $\mathcal{F}$: Fonksiyon kumesi (isim, adres, boyut, decompiled kod)
\item $\mathcal{S}$: String kumesi (literal, referans adresi)
\item $\mathcal{I}$: Import/export kumesi (kutuplhane, sembol)
\item $\mathcal{G}$: Call graph (yonlu graf, $G = (V, E)$ burada $V$ fonksiyonlar, $E$ cagri iliskileri)
\end{itemize}
\end{tanim}

% --- TikZ: Statik Analiz Pipeline ---
\begin{figure}[htbp]
\centering
\begin{tikzpicture}[
    node distance=1.3cm,
    block/.style={rectangle, draw=sectioncyan, fill=boxbg, text=white,
                  minimum width=3.5cm, minimum height=0.8cm, rounded corners=2pt,
                  font=\small},
    arrow/.style={-{Stealth[length=2.5mm]}, thick, white},
]
    % Binary track
    \node[block, fill=blue!15!black] (bin_input) {Native Binary};
    \node[block, below=of bin_input] (otool) {otool -L, -l};
    \node[block, below=of otool] (strings) {strings extraction};
    \node[block, below=of strings] (nm) {nm -g (symbols)};
    \node[block, below=of nm] (lief) {lief header parse};
    \node[block, below=of lief] (ghidra) {Ghidra Headless};
    \node[block, below=of ghidra, fill=green!15!black] (bin_out) {StageResult};

    % JS track
    \node[block, fill=yellow!15!black, right=4cm of bin_input] (js_input) {JS Bundle};
    \node[block, below=of js_input] (copy) {Raw Copy};
    \node[block, below=of copy] (chunk) {ChunkedProcessor};
    \node[block, below=of chunk] (babel) {Babel AST Parse};
    \node[block, below=of babel] (extract) {Function/String Extract};
    \node[block, below=of extract] (webpack) {Webpack Module ID};
    \node[block, below=of webpack, fill=green!15!black] (js_out) {StageResult};

    % Arrows
    \draw[arrow] (bin_input) -- (otool);
    \draw[arrow] (otool) -- (strings);
    \draw[arrow] (strings) -- (nm);
    \draw[arrow] (nm) -- (lief);
    \draw[arrow] (lief) -- (ghidra);
    \draw[arrow] (ghidra) -- (bin_out);

    \draw[arrow] (js_input) -- (copy);
    \draw[arrow] (copy) -- (chunk);
    \draw[arrow] (chunk) -- (babel);
    \draw[arrow] (babel) -- (extract);
    \draw[arrow] (extract) -- (webpack);
    \draw[arrow] (webpack) -- (js_out);

    % Labels
    \node[font=\footnotesize\bfseries, text=sectioncyan] at ([yshift=0.7cm]bin_input.north) {Binary Track};
    \node[font=\footnotesize\bfseries, text=yellow!70!white] at ([yshift=0.7cm]js_input.north) {JavaScript Track};
\end{tikzpicture}
\caption{Statik analiz iki paralel track: binary ve JavaScript.}
\label{fig:static-pipeline}
\end{figure}


% ===================================================================
\section{Binary Statik Analiz Araclari}
\label{sec:binary-static-tools}

\code{MachOAnalyzer} sinifi (\file{macho.py:L34--L610}), native binary'ler icin 7 adimlik bir statik analiz pipeline'i isletir.

\subsection{otool: Dynamic Library ve Load Command Analizi}

macOS'a ozgu \texttt{otool} araci, Mach-O binary'lerin yapisal bilgilerini cikarir.

\paragraph{otool -L: Dynamic Libraries.}
\file{macho.py:L362--L387}'deki \code{\_run\_otool\_libs} metodu:

\begin{lstlisting}[language=Python, caption={otool -L cikti parse'i}]
output = self.runner.run_otool(binary_path, flags=["-L"])
for line in output.splitlines()[1:]:  # ilk satir dosya adi
    parts = line.strip().split(" (")
    lib_path = parts[0].strip()
    version_info = parts[1].rstrip(")") if len(parts) > 1 else ""
    libraries.append({"path": lib_path, "version_info": version_info})
\end{lstlisting}

\paragraph{otool -l: Load Commands.}
Tum load command'lari ham metin olarak kaydeder (\file{macho.py:L389--L391}). Bu cikti, segment offset'leri, UUID, code signature ve encryption bilgilerini icerir.

\paragraph{Karmasiklik.} Her iki komut da $\bigO{n}$ dosya boyutuyla orantili zamanda calisir. Tipik bir 10~MB binary icin $<1$~saniye.

\subsection{nm: Symbol Table Extraction}

\texttt{nm -g} komutu, binary'nin global (export) sembollerini cikarir. \file{macho.py:L393--L429}'daki \code{\_run\_nm} metodu:

\begin{lstlisting}[language=Python, caption={nm cikti formati}]
# nm cikti formati: "ADDRESS TYPE NAME"
# Ornekler:
# 0000000100003a40 T _main
# 0000000100004b20 T _process_data
#                  U _printf        (undefined/external)
for line in result.stdout.splitlines():
    parts = line.strip().split(None, 2)
    if len(parts) == 3:
        symbols.append({
            "address": parts[0],  # hex adres
            "type": parts[1],     # T=text, D=data, U=undefined
            "name": parts[2],     # sembol adi
        })
\end{lstlisting}

\paragraph{Symbol Type Kodlari.}

\begin{table}[htbp]
\centering
\begin{tabular}{cl}
\toprule
\textbf{\color{sectioncyan}Kod} & \textbf{\color{sectioncyan}Anlam} \\
\midrule
\texttt{T/t} & Text (kod) section'da tanimli / lokal \\
\texttt{D/d} & Data section'da tanimli / lokal \\
\texttt{B/b} & BSS (sifir-initialize) section / lokal \\
\texttt{S/s} & Diger section'da tanimli / lokal \\
\texttt{U}   & Undefined (baska kutuphane\-den import) \\
\bottomrule
\end{tabular}
\caption{nm sembol tip kodlari.}
\label{tab:nm-types}
\end{table}

\subsection{strings: String Extraction}
\label{sec:strings-extraction}

Binary icindeki okunabilir string'ler, tersine muhendisligin en degerli bilgi kaynaklarindan biridir. Karadul iki yontem kullanir:

\textbf{1. Subprocess yontemi} (kucuk binary'ler, $<$100~MB): Sistemin \texttt{strings} komutunu SubprocessRunner ile calistirir.

\textbf{2. mmap yontemi} (buyuk binary'ler, $\geq$100~MB): \file{macho.py:L497--L572}'deki \code{\_extract\_strings\_mmap} metodu, dosyayi \texttt{mmap} ile belleye map'leyerek byte-by-byte tarar:

\begin{lstlisting}[language=Python, caption={mmap ile string extraction}]
with mmap.mmap(f.fileno(), 0, access=mmap.ACCESS_READ) as mm:
    for i in range(mm.size()):
        byte = mm[i]
        if 0x20 <= byte <= 0x7E:  # Printable ASCII
            current.append(byte)
        else:
            if len(current) >= min_length:
                strings.append(bytes(current).decode("ascii"))
                if len(strings) >= max_strings:
                    return strings  # Erken cikis
            current = []
\end{lstlisting}

\paragraph{mmap Avantajlari.}
\begin{itemize}
\item OS seviyesinde sayfa sayfa okuma --- tum dosya RAM'e alinmaz.
\item 224~MB binary icin $\sim$15 saniye (subprocess ile $>$30 saniye).
\item \texttt{max\_strings = 50000} limiti ile bellek kullanimi kontrol altinda.
\end{itemize}

\paragraph{Karmasiklik.}
\begin{equation}
\label{eq:strings-complexity}
T_{\text{strings}} = \bigO{n}
\end{equation}
Tek gecis (single pass) algoritmasi. Her byte tam bir kez incelenir.


% ===================================================================
\section{Ghidra Headless Analiz}
\label{sec:ghidra-headless}

Ghidra, NSA tarafindan gelistirilen acik kaynakli bir tersine muhendislik framework'udur. Karadul, Ghidra'yi \emph{headless} modda (GUI olmadan) kullanarak derinlemesine binary analiz yapar.

\subsection{PyGhidra API vs CLI Modu}

\file{headless.py:L1--L68}'de iki calisma modu desteklenir:

\begin{enumerate}
\item \textbf{PyGhidra API (tercih edilen):} \texttt{pyghidra.open\_program} ile Ghidra JVM dogrudan Python 3 process'i icinde baslatilir. Hizli, stabil, ama Ghidra + PyGhidra kurulumu gerektirir.
\item \textbf{CLI fallback:} \texttt{analyzeHeadless} komutu ile ayri bir JVM process'i baslatilir. Jython gerektirir, daha yavas.
\end{enumerate}

\subsection{Decompilation Pipeline}

Ghidra'nin decompilation sureci P-code (Processor code) ara temsilini kullanir:

\begin{figure}[htbp]
\centering
\begin{tikzpicture}[
    node distance=1.2cm,
    block/.style={rectangle, draw=sectioncyan, fill=boxbg, text=white,
                  minimum width=3.2cm, minimum height=0.7cm, rounded corners=2pt,
                  font=\small},
    arrow/.style={-{Stealth[length=2.5mm]}, thick, white},
]
    \node[block] (bin) {Raw Binary};
    \node[block, below=of bin] (disasm) {Disassembly};
    \node[block, below=of disasm] (pcode) {P-code IR};
    \node[block, below=of pcode] (ssa) {SSA Form};
    \node[block, below=of ssa] (opt) {Optimization Passes};
    \node[block, below=of opt] (decomp) {C Pseudocode};

    \draw[arrow] (bin) -- node[right, font=\tiny] {SLEIGH} (disasm);
    \draw[arrow] (disasm) -- node[right, font=\tiny] {Lifting} (pcode);
    \draw[arrow] (pcode) -- node[right, font=\tiny] {Analysis} (ssa);
    \draw[arrow] (ssa) -- node[right, font=\tiny] {Dead code elim.} (opt);
    \draw[arrow] (opt) -- node[right, font=\tiny] {Struct recovery} (decomp);
\end{tikzpicture}
\caption{Ghidra decompilation pipeline: Binary $\to$ P-code $\to$ C pseudocode.}
\label{fig:ghidra-pipeline}
\end{figure}

\paragraph{P-code.}
P-code, Ghidra'nin mimariden bagimsiz ara dilidir. Her native instruction (ARM64, x86\_64 vb.) bir veya birden fazla P-code islemine donusturulur. Bu, Ghidra'nin tek bir decompiler'la 40+ mimariyi desteklemesini saglar.

\subsection{Ghidra Script'leri}

Karadul, Ghidra'ya 5 ozel script gonderir:

\begin{table}[htbp]
\centering
\begin{tabular}{lp{8cm}}
\toprule
\textbf{\color{sectioncyan}Script} & \textbf{\color{sectioncyan}Gorevi} \\
\midrule
\texttt{function\_lister.py}  & Tum fonksiyonlarin listesi (isim, adres, boyut) \\
\texttt{string\_extractor.py} & Cross-reference'li string cikartma \\
\texttt{call\_graph.py}       & Fonksiyonlar arasi cagri grafinin cikarilmasi \\
\texttt{decompile\_all.py}    & Her fonksiyonun C pseudocode'a cevirilmesi \\
\texttt{export\_results.py}   & Tum sonuclarin JSON olarak birlestirilmesi \\
\bottomrule
\end{tabular}
\caption{Karadul'un Ghidra script'leri.}
\label{tab:ghidra-scripts}
\end{table}


% ===================================================================
\section{Mach-O Segment/Section Analizi}
\label{sec:macho-sections}

\code{lief} kutuphanesi ile Mach-O binary'nin tum segment ve section bilgileri programatik olarak cikarilir (\file{macho.py:L431--L495}).

\subsection{\_\_TEXT Segmenti}

\begin{description}
\item[\texttt{\_\_text}:] Derlenm\-is makine kodu. ARM64 veya x86\_64 instruction'lar.
\item[\texttt{\_\_stubs}:] Lazy binding stub'lari --- dinamik kutuplhane fonksiyonlarina ilk cagridaki trampoline.
\item[\texttt{\_\_stub\_helper}:] Stub helper fonksiyonlari --- dyld'e geri cagri.
\item[\texttt{\_\_cstring}:] Null-terminated C string'leri.
\item[\texttt{\_\_const}:] Salt-okunur sabitler.
\item[\texttt{\_\_unwind\_info}:] Exception handling unwind bilgileri.
\end{description}

\subsection{\_\_DATA Segmenti}

\begin{description}
\item[\texttt{\_\_data}:] Initialize edilmis global degiskenler.
\item[\texttt{\_\_bss}:] Sifir-initialize edilmis degiskenler (dosyada yer kaplamaz).
\item[\texttt{\_\_la\_symbol\_ptr}:] Lazy symbol pointer'lar --- ilk cagri sirasinda doldurulur.
\item[\texttt{\_\_got}:] Global Offset Table --- non-lazy symbol pointer'lar.
\end{description}

\subsection{\_\_LINKEDIT Segmenti}

Monolitik bir segment olup sembol tablolari, string tablolari, code signature ve diger baglama bilgilerini icerir. Bu segmentin ic yapisi \texttt{LC\_SYMTAB} ve \texttt{LC\_DYSYMTAB} load command'lari ile tarif edilir.


% ===================================================================
\section{Symbol Extraction: SYMTAB ve DYSYMTAB}
\label{sec:symbol-extraction}

\subsection{SYMTAB (Symbol Table)}

\texttt{LC\_SYMTAB} load command'i, sembol tablosunun dosya icindeki konumunu belirtir:

\begin{equation}
\label{eq:symtab-struct}
\text{LC\_SYMTAB} =
\begin{cases}
\texttt{symoff}  & \text{sembol tablosu offset'i} \\
\texttt{nsyms}   & \text{sembol sayisi} \\
\texttt{stroff}  & \text{string tablosu offset'i} \\
\texttt{strsize} & \text{string tablosu boyutu}
\end{cases}
\end{equation}

Her sembol girisi (\texttt{nlist\_64}) 16 byte'tir:

\begin{table}[htbp]
\centering
\begin{tabular}{llr}
\toprule
\textbf{\color{sectioncyan}Alan} & \textbf{\color{sectioncyan}Tip} & \textbf{\color{sectioncyan}Boyut} \\
\midrule
\texttt{n\_strx}  & uint32 & 4 byte (string tablosu index'i) \\
\texttt{n\_type}  & uint8  & 1 byte (tip flag'leri) \\
\texttt{n\_sect}  & uint8  & 1 byte (section numarasi) \\
\texttt{n\_desc}  & uint16 & 2 byte (aciklama) \\
\texttt{n\_value} & uint64 & 8 byte (sembol adresi) \\
\bottomrule
\end{tabular}
\caption{\texttt{nlist\_64} sembol girisi yapisi.}
\label{tab:nlist64}
\end{table}

\subsection{DYSYMTAB (Dynamic Symbol Table)}

Dinamik baglayici (dyld) tarafindan kullanilan ek bilgileri icerir: hangi sembollerin local, hangilesinin external, hangilerinin undefined oldugunu belirten index'ler.

\paragraph{Karadul'un sembol cikartma stratejisi:}
\begin{enumerate}
\item \texttt{nm -g} ile global sembolleri al (\file{macho.py:L393}).
\item lief ile tum segment/section bilgisini al (\file{macho.py:L431}).
\item Ghidra ile fonksiyon listesini al (decompilation dahil).
\item Uc kaynagi birlestir: \texttt{nm} $\cup$ \texttt{lief} $\cup$ \texttt{Ghidra}.
\end{enumerate}


% ===================================================================
\section{JavaScript Babel AST Analizi}
\label{sec:babel-ast}

JavaScript statik analizi, \texttt{@babel/parser} ile kaynak kodu Abstract Syntax Tree (AST) agacina donusturur ve bu agac uzerinde gezinerek bilgi cikarir.

\subsection{AST Kavramsal Yapisi}

\begin{tanim}[Abstract Syntax Tree]
Bir kaynak kod $s$ icin AST, yaprak dugumleri token'lar, ic dugumleri sentaktik yapilar olan bir koklu agactir:
\begin{equation}
\label{eq:ast-def}
\text{AST}(s) = (V, E, r, \lambda)
\end{equation}
Burada:
\begin{itemize}
\item $V$: dugum kumesi
\item $E \subseteq V \times V$: ebeveyn-cocuk iliskileri
\item $r \in V$: kok dugum (Program)
\item $\lambda: V \to \text{NodeType}$: dugum tip etiketi
\end{itemize}
\end{tanim}

\subsection{Babel Parser Entegrasyonu}

Karadul, Babel'i dogrudan Python'dan degil, Node.js scriptleri uzerinden cagirir (\file{javascript.py:L49--L162}):

\begin{lstlisting}[language=Python, caption={deobfuscate.mjs cagirma}]
cmd = [
    str(self.config.tools.node),
    "--max-old-space-size=8192",
    str(deobfuscate_script),
    str(analysis_file),
]
sub_result = self._runner.run_command(
    cmd, timeout=effective_timeout, cwd=self._scripts_dir,
)
result_json = sub_result.parsed_json
\end{lstlisting}

\subsection{Cikarilan AST Node Tipleri}

\begin{table}[htbp]
\centering
\begin{tabular}{llp{5.5cm}}
\toprule
\textbf{\color{sectioncyan}Node Tipi} & \textbf{\color{sectioncyan}Kategori} & \textbf{\color{sectioncyan}Aciklama} \\
\midrule
\texttt{FunctionDeclaration}   & Fonksiyon & \texttt{function foo() \{...\}} \\
\texttt{ArrowFunctionExpression} & Fonksiyon & \texttt{(x) => x + 1} \\
\texttt{ClassMethod}           & Fonksiyon & \texttt{class Foo \{ bar() \{...\} \}} \\
\texttt{FunctionExpression}    & Fonksiyon & \texttt{const f = function() \{...\}} \\
\texttt{StringLiteral}         & String    & \texttt{"hello world"} \\
\texttt{TemplateLiteral}       & String    & \texttt{`Hello \$\{name\}`} \\
\texttt{ImportDeclaration}     & Import    & \texttt{import x from 'y'} \\
\texttt{ExportNamedDeclaration} & Export   & \texttt{export \{ foo \}} \\
\texttt{CallExpression}        & Cagri     & \texttt{foo(a, b)} \\
\bottomrule
\end{tabular}
\caption{Babel AST'den cikarilan node tipleri.}
\label{tab:ast-nodes}
\end{table}

\subsection{Fonksiyon Extraction Algoritmasi}

Babel AST uzerinde \texttt{traverse} ile derinlemesine gezinti yapilir. Her fonksiyon node'u icin:

\begin{enumerate}
\item \textbf{Isim:} \texttt{FunctionDeclaration.id.name}, \texttt{ClassMethod.key.name}, veya anonim fonksiyonlar icin \texttt{null}.
\item \textbf{Parametre sayisi:} \texttt{node.params.length}.
\item \textbf{Baslangic/bitis satir numaralari:} \texttt{node.loc.start.line}, \texttt{node.loc.end.line}.
\item \textbf{Govde boyutu:} Bitis - baslangic satir farki.
\end{enumerate}

\paragraph{Pseudocode:}
\begin{lstlisting}[language=JavaScript, caption={Fonksiyon extraction pseudocode}]
traverse(ast, {
  "FunctionDeclaration|ArrowFunctionExpression|ClassMethod"(path) {
    functions.push({
      name: path.node.id?.name || path.node.key?.name || null,
      params: path.node.params.length,
      loc: path.node.loc,
      body_size: path.node.loc.end.line - path.node.loc.start.line,
    });
  }
});
\end{lstlisting}

\subsection{String Extraction}

AST'den iki tipteki string cikarilir:

\textbf{1. StringLiteral:} Tek veya cift tirnak icindeki sabit string'ler. Degeri dogrudan \texttt{node.value}'dan alinir.

\textbf{2. TemplateLiteral:} Backtick icindeki sablonlu string'ler. \texttt{node.quasis} dizisindeki \texttt{cooked} degerleri birlestirilir.

\paragraph{Karmasiklik.} AST traversal $\bigO{|V|}$ --- dugum sayisiyla dogrusal.


% ===================================================================
\section{Webpack Modul ID Tespiti}
\label{sec:webpack-modules}

Webpack bundle'lari, her modulu bir ID ile tanimlar ve \texttt{\_\_webpack\_require\_\_(id)} ile yukler.

\subsection{Webpack Runtime Yapisi}

Tipik bir webpack bundle'in yapisi:

\begin{lstlisting}[language=JavaScript, caption={Webpack bundle yapisi (basitlestirilmis)}]
(function(modules) {
  // webpack bootstrap
  function __webpack_require__(moduleId) {
    // cache kontrolu
    if (installedModules[moduleId]) return installedModules[moduleId].exports;
    // yeni modul
    var module = installedModules[moduleId] = { exports: {} };
    modules[moduleId].call(module.exports, module, module.exports,
                           __webpack_require__);
    return module.exports;
  }
  __webpack_require__(0); // entry point
})({
  0: function(module, exports, __webpack_require__) { /* ... */ },
  1: function(module, exports, __webpack_require__) { /* ... */ },
  // ...
});
\end{lstlisting}

\subsection{Modul ID Cikartma Algoritmasi}

Babel AST uzerinde \texttt{CallExpression} node'lari taranir:

\begin{enumerate}
\item \texttt{callee.name === "\_\_webpack\_require\_\_"} mi kontrol edilir.
\item Argument olarak sayi veya string literal alinir --- bu modul ID'dir.
\item ID'ler toplanir ve \texttt{webpack-modules} artifact'i olarak kaydedilir.
\end{enumerate}

\file{javascript.py:L213--L220}'de kayit:
\begin{lstlisting}[language=Python, caption={Webpack modul kaydi}]
webpack_modules = result_json.get("webpack_modules", [])
if webpack_modules:
    wp_path = workspace.save_json("static", "webpack-modules", {
        "module_ids": webpack_modules,
        "count": len(webpack_modules),
    })
\end{lstlisting}

\subsection{Modul ID Mapping}

Webpack 4+ surumlerde modul ID'leri deterministik hash'lerdir:
\begin{equation}
\label{eq:webpack-id}
\text{module\_id} = \text{MD4}(\text{relative\_path})[:4] \quad \text{(webpack 4)}
\end{equation}
\begin{equation}
\label{eq:webpack5-id}
\text{module\_id} = \text{MD4}(\text{absolute\_path + loaders})[:4] \quad \text{(webpack 5)}
\end{equation}

Bu hash'ler, orijinal dosya yolunu geri kazanmayi olanaklsiz kilar; ancak modul sayisi ve bagimlilik yapisi analizlerinde kullanilir.


% ===================================================================
\section{Buyuk Dosya Destegi: ChunkedProcessor}
\label{sec:chunked-processor}

200~MB+ JavaScript dosyalari Babel parser'in bellek limitlerini asar. \file{chunked\_processor.py}'deki \code{ChunkedProcessor} sinifi bu problemi cozer.

\subsection{Chunk Bolme Algoritmasi}

\begin{tanim}[Brace-Aware Chunking]
Bir JS dosyasi $F = (l_1, l_2, \ldots, l_n)$ satirlarindan olusmaktadir. Brace derinligi:
\begin{equation}
\label{eq:brace-depth}
d(i) = d(i-1) + \#\{l_i\}^{\texttt{\{}} - \#\{l_i\}^{\texttt{\}}}
\end{equation}
Bir chunk siniri $i$'de olusur eger:
\begin{equation}
\label{eq:chunk-boundary}
d(i) = 0 \quad \wedge \quad \sum_{j=k}^{i} |l_j| \geq C_{\max}
\end{equation}
Burada $C_{\max}$ = \code{max\_chunk\_mb} $\times 1024^2$ (varsayilan 50~MB).
\end{tanim}

\file{chunked\_processor.py:L113--L223}'teki implementasyon:

\begin{lstlisting}[language=Python, caption={Brace-aware chunking}]
depth = 0
for line_num, line in enumerate(f, start=1):
    for ch in line:
        if ch == "{": depth += 1
        elif ch == "}":
            depth -= 1
            if depth < 0: depth = 0
    current_lines.append(line)
    current_size += len(line.encode("utf-8"))
    # Chunk boundary: depth=0 ve yeterince buyuk
    if depth == 0 and current_size >= max_chunk_bytes:
        write_chunk(current_lines)
        current_lines, current_size = [], 0
\end{lstlisting}

\subsection{Top-Level Statement Tespiti}

Chunk sinirlari, tercihen top-level statement baslangiclarina hizalanir. \file{chunked\_processor.py:L22--L36}'daki pattern'ler:

\begin{lstlisting}[language=Python, caption={Top-level statement pattern'leri}]
_TOP_LEVEL_STARTS = (
    "function ", "async function ", "class ",
    "const ", "let ", "var ",
    "export ", "module.exports", "exports.",
    "Object.defineProperty", "__webpack_require__",
    '"use strict"', "'use strict'",
)
\end{lstlisting}

\subsection{Paralel Isleme}

\code{ChunkedProcessor.process\_chunks} metodu (\file{chunked\_processor.py:L225--L258}) hem sirali hem paralel isleme destekler:

\begin{equation}
\label{eq:parallel-speedup}
S_p = \frac{T_{\text{sequential}}}{T_{\text{parallel}}} \leq \min(p, N_{\text{chunks}})
\end{equation}

Burada $p$ = \code{max\_workers} (varsayilan 4), $N_{\text{chunks}}$ toplam chunk sayisi. Python GIL nedeniyle I/O-bound islemlerde etkilidir, CPU-bound islemlerde sinirlidir.

\paragraph{Veri Yapilari.}

\begin{table}[htbp]
\centering
\begin{tabular}{ll}
\toprule
\textbf{\color{sectioncyan}Sinif} & \textbf{\color{sectioncyan}Icerik} \\
\midrule
\code{ChunkInfo}   & path, index, start\_line, end\_line, size\_bytes, line\_count \\
\code{ChunkResult} & chunk, success, data, error \\
\code{SplitResult}  & success, chunks[], total\_lines, errors[] \\
\bottomrule
\end{tabular}
\caption{ChunkedProcessor veri yapilari.}
\label{tab:chunk-dataclasses}
\end{table}


% ===================================================================
\section{FLIRT Signature Sistemi}
\label{sec:flirt-signatures}

FLIRT (Fast Library Identification and Recognition Technology), IDA Pro'nun kutuphane fonksiyon tanima sistemidir. Karadul, \file{flirt\_parser.py}'de kendi FLIRT-uyumlu parser'ini implement eder.

\subsection{.pat Dosya Formati}

FLIRT pattern dosyasi (\texttt{.pat}) su formatta satirlar icerir:

\begin{verbatim}
HEXBYTES CRC16 SIZE TOTAL :OFFSET NAME [REFERENCED...]
\end{verbatim}

\paragraph{Ornek:}
\begin{verbatim}
558BEC83EC10 0C 0025 003A :0000 _my_function
558BEC........8B4508 00 0000 001F :0000 _another_func ^0010 _ref
\end{verbatim}

Burada \texttt{..} wildcard byte'i temsil eder (herhangi bir deger eslesir).

\subsection{Pattern Matching Algoritmasi}

\file{flirt\_parser.py:L721--L779}'deki \code{match\_function\_bytes} metodu:

\begin{tanim}[Masked Byte Comparison]
Bir fonksiyonun ilk $n$ byte'i $f = (f_0, f_1, \ldots, f_{n-1})$ ve bir FLIRT pattern'i $(p, m)$ (pattern + mask) verildiginde:
\begin{equation}
\label{eq:flirt-match}
\text{match}(f, p, m) = \bigwedge_{i=0}^{|p|-1}
\begin{cases}
f_i = p_i & \text{eger } m_i = \texttt{0xFF} \\
\text{true} & \text{eger } m_i = \texttt{0x00} \text{ (wildcard)}
\end{cases}
\end{equation}
\end{tanim}

\paragraph{Confidence Hesabi.}
\begin{equation}
\label{eq:flirt-confidence}
\text{conf}(p, m) = \min\!\left(0.98,\; \text{base\_conf} \times \frac{|\{i : m_i = \texttt{0xFF}\}|}{|p|}\right)
\end{equation}
Burada \texttt{base\_conf} = 0.95 (FLIRT imzalarinin yuksek guvenilirligi).

\subsection{Hex-to-Bytes-with-Mask Donusumu}

\file{flirt\_parser.py:L218--L256}'deki \code{\_hex\_to\_bytes\_with\_mask} metodu:

\begin{lstlisting}[language=Python, caption={Hex string ile mask olusturma}]
# "558BEC..8B45" -> (pattern_bytes, mask_bytes)
i = 0
while i < len(hex_str):
    if hex_str[i:i+2] == "..":
        pattern_bytes.append(0x00)
        mask_bytes.append(0x00)   # Wildcard
    else:
        pattern_bytes.append(int(hex_str[i:i+2], 16))
        mask_bytes.append(0xFF)   # Kesin
    i += 2
\end{lstlisting}

\subsection{SignatureDB Entegrasyonu}

FLIRT parser, bagimsiz bir modul olarak tasarlanmistir. SignatureDB'ye enjeksiyon \code{inject\_into\_signature\_db} metodu ile yapilir (\file{flirt\_parser.py:L647--L715}):

\begin{equation}
\label{eq:flirt-inject}
\text{SignatureDB} = \text{SignatureDB}_{\text{builtin}} \cup \text{FLIRT}_{\text{external}}
\end{equation}

Duplikasyon kontrolu yapilir: builtin DB'deki imzalar onceliklidir, ayni isimli FLIRT imzasi atlanir.

\paragraph{Karadul'un toplam imza sayisi:}
\begin{equation}
|\text{SignatureDB}| = 6100+ \text{ (builtin)} + \text{FLIRT external} \approx 1.5\text{M+}
\end{equation}


% ===================================================================
\section{Assembly Instruction Analizi}
\label{sec:assembly-analysis}

\file{assembly\_analyzer.py}'deki \code{AssemblyAnalyzer} sinifi, Ghidra decompiler'in kacirdigi dusuk seviye bilgileri dogrudan assembly instruction'lardan cikarir.

\subsection{ARM64 (AArch64) vs x86\_64}

\begin{table}[htbp]
\centering
\begin{tabular}{lll}
\toprule
\textbf{\color{sectioncyan}Ozellik} & \textbf{\color{sectioncyan}ARM64 (AArch64)} & \textbf{\color{sectioncyan}x86\_64} \\
\midrule
Instruction boyutu  & Sabit 4 byte         & Degisken 1--15 byte \\
Register sayisi     & 31 genel amacli      & 16 genel amacli \\
Calling convention  & AAPCS64              & System V AMD64 / Win64 \\
Parametre register  & x0--x7 (8 adet)      & rdi,rsi,rdx,rcx,r8,r9 (6 adet) \\
Return register     & x0, x1               & rax, rdx \\
SIMD extension      & NEON                 & SSE/AVX \\
Endianness          & Little-endian        & Little-endian \\
\bottomrule
\end{tabular}
\caption{ARM64 ve x86\_64 mimari karsilastirmasi.}
\label{tab:arm-vs-x86}
\end{table}

\subsection{Calling Convention Tespiti}

\file{assembly\_analyzer.py:L207--L224}'deki \code{\_detect\_calling\_convention}:

\begin{lstlisting}[language=Python, caption={Calling convention tespiti}]
if arch == "aarch64":
    return "aapcs64"

first_lines = "\n".join(lines[:20]).lower()
# System V AMD64: rdi, rsi parametreler
if re.search(r"\brdi\b", first_lines) and re.search(r"\brsi\b", first_lines):
    return "sysv_amd64"
# Windows x64: rcx, rdx parametreler
if re.search(r"\brcx\b", first_lines) and not re.search(r"\brdi\b", first_lines):
    return "win64"
return "sysv_amd64"  # macOS/Linux default
\end{lstlisting}

\subsection{Stack Frame Analizi}

Fonksiyonun stack frame boyutu, prolog'daki \texttt{sub rsp, N} instruction'indan cikarilir:

\begin{equation}
\label{eq:stack-frame}
\text{frame\_size} =
\begin{cases}
N & \text{eger } \texttt{sub rsp, N} \text{ bulunduysa} \\
0 & \text{aksi halde (leaf function)}
\end{cases}
\end{equation}

Stack degiskenleri \texttt{[rbp-0xNN]} pattern'leri ile tespit edilir. Ardisik offset farklari degisken boyutunu verir:

\begin{equation}
\label{eq:stack-var-size}
\text{size}(v_i) = \text{offset}(v_{i+1}) - \text{offset}(v_i)
\end{equation}

\subsection{SIMD Pattern Tespiti}

\file{assembly\_analyzer.py:L82--L110}'da tanimlanan instruction setleri:

\begin{description}
\item[SSE:] \texttt{addps}, \texttt{mulps}, \texttt{movaps}, \ldots --- 128-bit SIMD (float32$\times$4)
\item[AVX:] \texttt{vaddps}, \texttt{vmulps}, \texttt{vfmadd231ps}, \ldots --- 256-bit SIMD (float32$\times$8)
\item[NEON:] \texttt{fmla}, \texttt{fadd}, \texttt{ld1}, \texttt{st1}, \ldots --- ARM 128-bit SIMD
\item[Crypto:] \texttt{aesenc}, \texttt{sha256rnds2}, \texttt{pclmulqdq}, \texttt{aese}, \ldots --- AES-NI / ARM Crypto Extensions
\end{description}

FMA (Fused Multiply-Add) instruction'larinin varligii, dot product veya matris carpiimi pattern'ini isaret eder.


% ===================================================================
\section{Entropy Analizi ile Packing Tespiti}
\label{sec:entropy-packing}

Paketlenm\-is (packed) binary'ler, normal binary'lerden ayirt edilebilir entropy profilleri gosterir.

\subsection{Shannon Entropisi}

\begin{equation}
\label{eq:shannon-entropy}
H(X) = -\sum_{x \in \mathcal{X}} p(x) \log_2 p(x)
\end{equation}

Burada $p(x) = \text{freq}(x) / n$, byte $x$'in frekans olasiligidir. $H$ degerleri:

\begin{itemize}
\item $H = 0$: Tamamen tekduze (tek byte tekrari).
\item $H \approx 4$--$5$: Tipik executable binary (kod + veri).
\item $H \approx 6$--$7$: Sikistirilmis veri veya rastgele-benzeri bolgeler.
\item $H \approx 8$: Tamamen rastgele / sifrelenmis / UPX-packed.
\end{itemize}

\file{packed\_binary.py:L137--L160}'daki implementasyon:
\begin{lstlisting}[language=Python, caption={Shannon entropy hesaplama}]
def calculate_entropy(data: bytes) -> float:
    if not data:
        return 0.0
    freq = [0] * 256
    for byte in data:
        freq[byte] += 1
    length = len(data)
    entropy = 0.0
    for f in freq:
        if f > 0:
            p = f / length
            entropy -= p * math.log2(p)
    return entropy
\end{lstlisting}

\paragraph{Karmasiklik.}
\begin{equation}
\label{eq:entropy-complexity}
T_{\text{entropy}} = \bigO{n + 256} = \bigO{n}
\end{equation}
Tek gecis frekans sayma + 256 eleman toplam.

\subsection{Packing Karar Algoritmasi}

\file{packed\_binary.py:L209--L388}'deki \code{PackingDetector.detect} metodu 5 adimli bir tespit uygular:

\begin{enumerate}
\item \textbf{UPX kontrolu:} \texttt{UPX!} magic byte'i aranir. Bulunursa: $\text{confidence} = 0.95$.
\item \textbf{PyInstaller kontrolu:} \texttt{MEI} magic byte'i (8 byte). Bulunursa: $\text{confidence} = 0.95$.
\item \textbf{Nuitka kontrolu:} \texttt{\_\_nuitka\_}, \texttt{nuitka-compiled} string'leri. Bulunursa: $\text{confidence} = 0.85$.
\item \textbf{Entropy analizi:} $H > 7.0$ esigi ve $>\!50\%$ packed section orani.
\item \textbf{Import heuristigi:} $<\!10$ import $\implies$ muhtemelen packed.
\end{enumerate}

Karar formulu:
\begin{equation}
\label{eq:packing-decision}
\text{is\_packed} =
\begin{cases}
\text{true, UPX} & \texttt{UPX!} \in \text{binary} \\
\text{true, PyInstaller} & \texttt{MEI} \in \text{binary} \\
\text{true, Nuitka} & \texttt{nuitka} \in \text{binary} \\
\text{true, unknown} & H > 7.0 \wedge \text{packed\_ratio} > 0.5 \\
\text{true, generic} & H > 6.5 \wedge \text{imports} < 10 \\
\text{false} & \text{aksi halde}
\end{cases}
\end{equation}

% --- TikZ: Entropy Distribution ---
\begin{figure}[htbp]
\centering
\begin{tikzpicture}
\begin{axis}[
    width=12cm, height=6cm,
    xlabel={Entropy ($H$)},
    ylabel={Binary sayisi (goreli)},
    xmin=0, xmax=8.5,
    ymin=0, ymax=1.1,
    xtick={0,1,2,3,4,5,6,7,8},
    grid=both,
    grid style={gray!20!black},
    axis background/.style={fill=black},
    axis line style={white},
    tick style={white},
    ticklabel style={white},
    xlabel style={white},
    ylabel style={white},
    legend style={fill=boxbg, draw=sectioncyan, text=white, at={(0.98,0.98)}, anchor=north east},
]
    % Normal binary dağılımı
    \addplot[domain=2:7, samples=100, thick, cyan, smooth]
        {exp(-0.5*((x-4.5)/0.8)^2)};
    \addlegendentry{Normal binary}

    % Packed binary dağılımı
    \addplot[domain=6:8.5, samples=100, thick, red, smooth]
        {0.9*exp(-0.5*((x-7.5)/0.3)^2)};
    \addlegendentry{Packed binary}

    % Threshold çizgisi
    \draw[yellow, dashed, thick] (axis cs:7,0) -- (axis cs:7,1.1);
    \node[yellow, font=\small] at (axis cs:7.3,1.0) {$H=7.0$};
\end{axis}
\end{tikzpicture}
\caption{Normal vs packed binary entropy dagilimi. $H = 7.0$ esigi optimal ayirim noktasidir.}
\label{fig:entropy-distribution}
\end{figure}


% ===================================================================
\section{Her Analyzer Icin Giris/Cikis Veri Yapilari}
\label{sec:analyzer-io}

Tum analyzer'lar ortak bir arayuz kullanir: \code{BaseAnalyzer} soyut sinifinden miras alir ve \code{StageResult} dondurur.

\subsection{Giris Veri Yapisi: TargetInfo}

\begin{lstlisting}[language=Python, caption={TargetInfo veri yapisi}]
@dataclass
class TargetInfo:
    path: Path              # Dosya yolu
    name: str               # Hedef adi
    target_type: TargetType # 12 tip
    language: Language       # 13 dil
    file_size: int          # Byte
    file_hash: str          # SHA-256
    metadata: dict[str, Any]# Ek bilgi
\end{lstlisting}

\subsection{Cikis Veri Yapisi: StageResult}

\begin{lstlisting}[language=Python, caption={StageResult veri yapisi}]
@dataclass
class StageResult:
    stage_name: str                  # "static", "deobfuscate", ...
    success: bool                    # Basarili mi
    duration_seconds: float          # Sure
    artifacts: dict[str, Path]       # Cikti dosyalari
    stats: dict[str, object]         # Istatistikler
    errors: list[str]                # Hata mesajlari
\end{lstlisting}

\subsection{Analyzer Giris/Cikis Ozet Tablosu}

\begin{table}[htbp]
\centering
\begin{tabular}{lp{4cm}p{6cm}}
\toprule
\textbf{\color{sectioncyan}Analyzer} & \textbf{\color{sectioncyan}Giris} & \textbf{\color{sectioncyan}Cikis Artifact'leri} \\
\midrule
MachOAnalyzer & Mach-O/ELF/PE binary + Workspace & dynamic\_libraries, load\_commands, strings\_raw, symbols, lief\_analysis, ghidra\_* \\
JavaScriptAnalyzer & JS dosya + Workspace & ast\_analysis, functions, strings, imports\_exports, webpack\_modules \\
ElectronAnalyzer & .app bundle + Workspace & asar\_extracted\_dir, js\_files, electron-analysis.json \\
GoBinaryAnalyzer & Go binary + Workspace & go\_functions, go\_buildinfo, go\_types \\
PackingDetector & Binary path & PackingInfo (is\_packed, type, entropy, evidence) \\
FLIRTParser & .pat/.json dosya & FLIRTSignature listesi (name, library, byte\_pattern, mask) \\
AssemblyAnalyzer & Assembly text + arch & AssemblyAnalysisResult (convention, params, stack, simd) \\
\bottomrule
\end{tabular}
\caption{Her analyzer'in giris ve cikis veri yapilari.}
\label{tab:analyzer-io}
\end{table}

\subsection{Graceful Degradation}

Karadul'un statik analiz katmaninin kritik tasarim prensibi: herhangi bir alt-aracin basarisizligi pipeline'i durdurmaz. \file{macho.py:L284--L291}'deki basari kriterr:

\begin{lstlisting}[language=Python, caption={Kismi basari kriteri}]
return StageResult(
    stage_name="static",
    success=len(errors) == 0 or len(artifacts) > 0,  # En az 1 artifact yeterli
    # ...
)
\end{lstlisting}

Bu, ``ya hep ya hic'' yerine ``ne kadar bilgi cikarabilirsen o kadar iyi'' yaklasimini benimser.


% ===================================================================
\section{Binary Intelligence: String Clustering ve Mimari Cikarim}
\label{sec:binary-intelligence}

\file{macho.py:L222--L280}'de cagirilan \code{BinaryIntelligence} modulu, cikarilan string, sembol ve dylib verilerini birlestirerek ust seviye mimarik bilgi cikarir:

\begin{itemize}
\item \textbf{App type tespiti:} CLI, GUI, daemon, library, framework.
\item \textbf{Subsystem tespiti:} networking, database, crypto, UI, logging.
\item \textbf{Algoritma tespiti:} AES, SHA, zlib, protobuf, JSON.
\item \textbf{Guvenlik mekanizmasi tespiti:} certificate pinning, code signing, anti-debug.
\item \textbf{Protokol tespiti:} HTTP/2, gRPC, WebSocket, MQTT.
\end{itemize}

Bu bilgiler, ileriki asamalarda fonksiyon isimlendirme ve rapor uretiminde kullanilir.


\begin{ozet}
Statik Analiz katmani, Karadul'un bilgi toplama omurgasidir. Binary track'te otool, nm, strings, lief ve Ghidra Headless 7 adimli bir analiz pipeline'i olusturur. JavaScript track'te Babel AST analizi fonksiyon, string, import/export ve webpack modul bilgilerini cikarir. Buyuk dosyalar (200~MB+) icin ChunkedProcessor brace-aware bolme stratejisi kullanir. FLIRT signature sistemi, 1.5~M+ imza ile kutuphane fonksiyon tanimlamasi yapar. Assembly analiz, calling convention, stack frame ve SIMD pattern tespiti ile derin bilgi saglar. Entropy analizi, packed binary tespitini $H > 7.0$ esigi ile yapar. Tum bu bilesenler, graceful degradation prensibiyle calisir: herhangi bir alt aracin basarisizligi pipeline'i durdurmaz.
\end{ozet}



% ============================================================================
% BOLUM 5: JAVASCRIPT DEOBFUSCATION
% ============================================================================
\chapter{JavaScript Deobfuscation}
\label{ch:deobf}

Bu bolum, Karadul'un en yenilikci bilesenlerinden birini --- 9 fazli derin JavaScript deobfuscation pipeline'ini --- detayli olarak aciklar.

\section{Obfuscation Nedir?}
\label{sec:obf-intro}

JavaScript obfuscation, kaynak kodu okunaksiz hale getirme suerecidir:

\begin{equation}
\text{obfuscate}: S_{\text{readable}} \to S_{\text{obfuscated}} \quad \text{s.t.} \quad \text{semantics}(S_{\text{readable}}) = \text{semantics}(S_{\text{obfuscated}})
\end{equation}

Yaygin teknikler:

\begin{enumerate}
\item \textbf{Minification:} Bosluk, yorum silme, kisa isimler (\code{a}, \code{b}, \code{c})
\item \textbf{String encoding:} \code{"Hello"} $\to$ \code{"\textbackslash x48\textbackslash x65\textbackslash x6c\textbackslash x6c\textbackslash x6f"}
\item \textbf{Control flow flattening:} if/else $\to$ switch-case state machine
\item \textbf{Opaque predicates:} Her zaman true/false olan kosullar ekleme
\item \textbf{Dead code injection:} Calismayan sahte kod ekleme
\item \textbf{String array rotation:} String'leri bir diziye koy, indeksle eris, diziyi doenduer
\end{enumerate}

\section{Standart 4 Adimli Zincir}
\label{sec:std-chain}

\file{manager.py} dosyasindaki standart deobfuscation zinciri:

\begin{center}
\begin{tikzpicture}[
  stepbox/.style={rectangle, draw=sectioncyan, fill=boxbg, minimum width=3.5cm, minimum height=1cm, text=white, font=\small\bfseries, rounded corners=3pt},
  arrow/.style={-{Stealth[length=3mm]}, thick, color=gray!70!white},
]
  \node[stepbox] (beautify) at (0,0) {Beautify};
  \node[stepbox] (synchrony) at (4.5,0) {Synchrony};
  \node[stepbox] (babel) at (9,0) {Babel AST};
  \node[stepbox] (webpack) at (13.5,0) {\shortstack{Webpack\\Unpack}};

  \draw[arrow] (beautify) -- (synchrony);
  \draw[arrow] (synchrony) -- (babel);
  \draw[arrow] (babel) -- (webpack);

  \node[below=0.3cm, font=\tiny, text=gray!60!white] at (beautify.south) {Format kurtarma};
  \node[below=0.3cm, font=\tiny, text=gray!60!white] at (synchrony.south) {obfuscator.io};
  \node[below=0.3cm, font=\tiny, text=gray!60!white] at (babel.south) {9 faz transform};
  \node[below=0.3cm, font=\tiny, text=gray!60!white] at (webpack.south) {Modul cikartma};
\end{tikzpicture}
\end{center}

\subsection{Graceful Degradation}

Her adim basarisiz olabilir --- bu durumda bir oenceki basarili cikti kullanilir:

\begin{equation}
\text{output}(i) = \begin{cases}
\text{step}_i(\text{output}(i-1)) & \text{eger basarili} \\
\text{output}(i-1) & \text{eger basarisiz}
\end{cases}
\end{equation}

\section{9 Fazli Deep Deobfuscation}
\label{sec:deep-deobf}

\subsection{Faz 1: Constant Folding (Sabit Katlama)}

Derleme zamaninda hesaplanabilir ifadeleri basitlestirme:

\begin{equation}
\text{fold}(e_1 \oplus e_2) = \begin{cases}
v_1 \oplus v_2 & \text{eger } e_1 = v_1, e_2 = v_2 \text{ (her ikisi sabit)} \\
e_1 \oplus e_2 & \text{aksi halde (degismez)}
\end{cases}
\end{equation}

\begin{example}
\begin{lstlisting}[language=JavaScript]
// Once:
var x = 1 + 2 * 3;
var y = "He" + "llo";
// Sonra:
var x = 7;
var y = "Hello";
\end{lstlisting}
\end{example}

\subsection{Faz 2: Dead Code Elimination}

Kullanilmayan kodu tespit ve silme. Liveness analizi:

\begin{equation}
\text{live\_out}(B) = \bigcup_{S \in \text{succ}(B)} \text{live\_in}(S)
\end{equation}
\begin{equation}
\text{live\_in}(B) = \text{use}(B) \cup (\text{live\_out}(B) \setminus \text{def}(B))
\end{equation}

Eger bir degisken hicbir noktada ``live'' degilse, atanmasi silinebilir.

\subsection{Faz 3: String Unhexing}

Hex-encoded string'leri orijinaline doenduerme:

\begin{equation}
\text{unhex}(\texttt{"\textbackslash xHH"}) = \text{char}(\text{parseInt}(\texttt{HH}, 16))
\end{equation}

\begin{example}
\code{\textbackslash x48\textbackslash x65\textbackslash x6c\textbackslash x6c\textbackslash x6f} $\to$ \code{"Hello"}
\end{example}

\subsection{Faz 4: Computed-to-Static Property}

\begin{equation}
\text{obj}[\texttt{"foo"}] \to \text{obj}.\texttt{foo} \quad \text{(eger key string literal ise)}
\end{equation}

AST duenusemuue: \code{MemberExpression \{ computed: true \}} $\to$ \code{MemberExpression \{ computed: false \}}.

\subsection{Faz 5--8: Ileri Basitlestirmeler}

\begin{itemize}
\item \textbf{Faz 5:} Comma-split: \code{a=1, b=2, c=3} $\to$ uec ayri ifade
\item \textbf{Faz 6:} Conditional simplification: \code{true ? a : b} $\to$ \code{a}
\item \textbf{Faz 7:} Expression flattening: icice ifadeleri duzlestir
\item \textbf{Faz 8:} Literal restoration: obfuscated literal'leri orijinaline doenduer
\end{itemize}

\subsection{Faz 9: Semantik Degisken Isimlendirme}

LLM kullanmadan, 25 context-aware heuristik kural ile anlamli isimler uereetme:

\begin{table}[h]
\centering
\small
\begin{tabular}{p{5cm}p{8cm}}
\toprule
\textbf{\color{sectioncyan}Kural} & \textbf{\color{sectioncyan}Oernek} \\
\midrule
Event listener callback & \code{addEventListener("click", a)} $\to$ \code{clickHandler} \\
Error constructor & \code{new Error(a)} $\to$ \code{errorMessage} \\
Fetch URL & \code{fetch(a)} $\to$ \code{fetchUrl} \\
Array index & \code{for(a=0; a<b.length)} $\to$ \code{index} \\
DOM element & \code{document.getElementById(a)} $\to$ \code{elementId} \\
\bottomrule
\end{tabular}
\caption{Semantik isimlendirme oernek kurallari.}
\label{tab:semantic-rules}
\end{table}

\begin{ozet}
JavaScript deobfuscation, 4 adimli standart zincir (beautify $\to$ synchrony $\to$ babel $\to$ webpack) ve 9 fazli derin transform pipeline'indan olusur. Her faz AST uezerinde calisir ve belirli bir obfuscation tekni\-gini geri alir. Graceful degradation her adimda basarisizliga tolerans saglar. Faz 9'daki LLM-siz semantik isimlendirme 25 heuristik kuralla etkileyici sonuclar ueretir.
\end{ozet}

% --- DETAYLI GENISLETME ---
% ============================================================================
% BOLUM 5 GENISLETME: JAVASCRIPT DEOBFUSCATION — COK DETAYLI
% ============================================================================

\section{Abstract Syntax Tree (AST) --- Derinlemesine}
\label{sec:ast-deep}

\subsection{Formal Tanim}

AST, kaynak kodun soezdizimsel yapisini soyutlastiran bir yoenlue, etiketli agactir:

\begin{definition}[Abstract Syntax Tree]
$T = (V, E, r, L)$ burada:
\begin{itemize}
\item $V$: dugum (node) kuemesi
\item $E \subseteq V \times V$: kenar kuemesi (parent $\to$ child)
\item $r \in V$: koek dugum (root)
\item $L: V \to \text{NodeTypes}$: etiketleme fonksiyonu
\end{itemize}
\end{definition}

\textbf{Parse tree vs AST farki:} Parse tree (Concrete Syntax Tree --- CST) gramerin her intermediate kuralini goesterir. AST sadece semantik olarak anlamli oegeleri tutar.

\begin{example}
\code{1 + 2 * 3} icin CST:
\begin{lstlisting}
ExpressionStatement
  Expression
    AdditiveExpression
      NumericLiteral("1") "+" MultiplicativeExpression
        NumericLiteral("2") "*" NumericLiteral("3")
\end{lstlisting}

AST (Babel):
\begin{lstlisting}
BinaryExpression { operator: "+",
  left: NumericLiteral(1),
  right: BinaryExpression { operator: "*",
    left: NumericLiteral(2),
    right: NumericLiteral(3)
  }
}
\end{lstlisting}
\end{example}

\subsection{Babel vs Acorn (ESTree)}

Karadul her iki parser'i da kullanir. Kritik farklar:

\begin{table}[ht]
\centering
\begin{tabular}{lll}
\toprule
\textbf{\color{sectioncyan}Kavram} & \textbf{\color{sectioncyan}Acorn (ESTree)} & \textbf{\color{sectioncyan}Babel AST} \\
\midrule
String literal & \code{Literal \{ value: "hello" \}} & \code{StringLiteral \{ value: "hello" \}} \\
Sayi literal & \code{Literal \{ value: 42 \}} & \code{NumericLiteral \{ value: 42 \}} \\
Object key & \code{Property} & \code{ObjectProperty} \\
Program & \code{Program \{ body: [...] \}} & \code{File \{ program: \{ body \} \}} \\
\bottomrule
\end{tabular}
\caption{Babel vs Acorn AST node tipleri.}
\end{table}

\file{deobfuscate.mjs:164}'te bu ayrim handle ediliyor:
\begin{lstlisting}[language=JavaScript]
const usedAcorn = ast.body !== undefined && !ast.program;
\end{lstlisting}

\subsection{Hesaplama Karmasikligi}

Parse etme: $\bigO{n}$ burada $n$ = token sayisi. Babel'in parser'i Pratt parsing (top-down operator precedence) kullanir. AST node sayisi $\sim \bigO{n}$ cueuenkue her token en fazla sabit sayida node ueretir.


\section{Constant Folding --- Lattice Teorisi}
\label{sec:const-fold-deep}

\subsection{Abstract Interpretation Framework}

Abstract interpretation'da her degisken degeri bir lattice uezerinde izlenir:

\begin{center}
\begin{tikzpicture}[
  nd/.style={circle, draw=sectioncyan, fill=boxbg, minimum size=1.2cm, text=white, font=\small},
]
  \node[nd] (top) at (0,2) {$\top$};
  \node[nd] (c1) at (-2,0) {$c_1$};
  \node[nd] (c2) at (0,0) {$c_2$};
  \node[nd] (c3) at (2,0) {$c_3$};
  \node[nd] (bot) at (0,-2) {$\bot$};

  \draw[gray!50!white] (top) -- (c1);
  \draw[gray!50!white] (top) -- (c2);
  \draw[gray!50!white] (top) -- (c3);
  \draw[gray!50!white] (c1) -- (bot);
  \draw[gray!50!white] (c2) -- (bot);
  \draw[gray!50!white] (c3) -- (bot);
\end{tikzpicture}
\end{center}

\begin{itemize}
\item $\top$ (top): Degiskenin degeri henueuez belli degil
\item $c_i$ (constant): Tam olarak bilinen sabit deger
\item $\bot$ (bottom): Birden fazla farkli deger atanmis, sabit degil
\end{itemize}

Transfer fonksiyonu $f$: \code{x = 3 + 4} ise $f(x) = 7$ (constant). \code{x = a + b} ve $a = 3$, $b = \top$ ise $f(x) = \top$.

\subsection{Karadul Implementasyonu}

\file{deep-deobfuscate.mjs} Phase 2'de \textbf{exit visitor} (bottom-up) kullanilir:

\begin{lstlisting}[language=JavaScript]
// BinaryExpression.exit:
if (t.isNumericLiteral(left) && t.isNumericLiteral(right)) {
    let result;
    switch (operator) {
        case "+": result = left.value + right.value; break;
        case "-": result = left.value - right.value; break;
        case "*": result = left.value * right.value; break;
        // ... 11 operator toplam
    }
    if (result !== undefined && Number.isFinite(result)) {
        path.replaceWith(t.numericLiteral(result));
    }
}
\end{lstlisting}

\textbf{Neden exit visitor?} Bottom-up calisir: ic node oence, parent sonra. \code{1 + 2 * 3}'te oence \code{*} degerlendirilir (leaf'lere daha yakin), sonra \code{+}.

\begin{example}[UnaryExpression Katlama]
\begin{lstlisting}[language=JavaScript]
// !0 -> true, !1 -> false (Phase 2)
if (operator === "!" && t.isNumericLiteral(argument)) {
    path.replaceWith(t.booleanLiteral(!argument.value));
}
// void 0 -> undefined
if (operator === "void" && argument.value === 0) {
    path.replaceWith(t.identifier("undefined"));
}
\end{lstlisting}
\end{example}

\subsection{Sinirlar}

Karadul'un constant folder'i sadece \textbf{literal-to-literal} katlar:

\begin{lstlisting}[language=JavaScript]
const A = 10;
const B = A + 5;  // Bu KATLANMAZ: A Identifier, NumericLiteral degil
\end{lstlisting}

Interprocedural constant propagation icin Kildall'in worklist algoritmasi gerekir --- $\bigO{N \cdot D}$ burada $D$ = lattice yueksekligi. Deobfuscation baglaminda literal-to-literal cogunlukla yeterli.


\section{Dead Code Elimination --- Liveness Analizi}
\label{sec:dce-deep}

\subsection{Dataflow Denklemleri}

CFG uezerinde backward analysis:

\begin{equation}
\text{LIVE\_out}[B] = \bigcup_{S \in \text{succ}(B)} \text{LIVE\_in}[S]
\end{equation}

\begin{equation}
\text{LIVE\_in}[B] = \text{use}(B) \cup \big(\text{LIVE\_out}[B] \setminus \text{def}(B)\big)
\end{equation}

\textbf{Fixed-point iteration:} Tuem bloklari $\text{LIVE\_in}$ ve $\text{LIVE\_out}$ degismeyene kadar itere et. Karmasiklik: $\bigO{N \cdot D}$ burada $N$ = basic block sayisi, $D$ = iterasyon sayisi (pratikte 2--5).

\subsection{Karadul'un Basitlestirilmis Yaklasimi}

Tam liveness analysis yerine \textbf{constant-test elimination}:

\begin{lstlisting}[language=JavaScript]
// Phase 3: IfStatement.exit
if (t.isBooleanLiteral(test) && test.value === true) {
    path.replaceWith(path.node.consequent);  // if(true){X} -> X
}
if (t.isBooleanLiteral(test) && test.value === false) {
    if (path.node.alternate) {
        path.replaceWith(path.node.alternate);  // if(false){X}else{Y} -> Y
    } else {
        path.remove();  // if(false){X} -> (silindi)
    }
}
\end{lstlisting}

\textbf{Phase 2 + Phase 3 sinerjisi:} Phase 2 \code{!0}'yi \code{true}'ya cevirir, Phase 3 \code{if(true)} goerueuep temizler. Exit visitor sayesinde tek AST traverse'da her iki faz da calisir.


\section{String Obfuscation Teknikleri ve Coezeuemleri}
\label{sec:string-obf}

\subsection{6 Yaygin Teknik}

\begin{enumerate}
\item \textbf{Hex Encoding:} \code{"\textbackslash x48\textbackslash x65\textbackslash x6c\textbackslash x6c\textbackslash x6f"} $\to$ \code{"Hello"}

Reverse: \code{/\textbackslash\textbackslash x([0-9a-fA-F]\{2\})/g} ile decode.

\item \textbf{Unicode Escape:} \code{"\textbackslash u0048\textbackslash u0065..."} $\to$ \code{"Hello"}

\item \textbf{CharCode Array:} \code{String.fromCharCode(72, 101, 108, 108, 111)} $\to$ \code{"Hello"}

AST'de CallExpression ile \code{String.fromCharCode} tespiti, arguments hepsi NumericLiteral ise derleme zamaninda hesapla.

\item \textbf{Base64:} \code{atob("SGVsbG8=")} $\to$ \code{"Hello"}

\item \textbf{String Array Rotation (obfuscator.io):}
\begin{lstlisting}[language=JavaScript]
var _0x1234 = ["log", "Hello", "console"];
(function(_0xa, _0xb) {
    while (--_0xb) { _0xa.push(_0xa.shift()); }
})(_0x1234, 0x2);
// Simdi _0x1234[0]="console", [1]="log", [2]="Hello"
\end{lstlisting}

Reverse: Rotation miktarini tespit et, arrayi rotate et, referanslari degistir.

\item \textbf{XOR Cipher:} Key bilgisi gerekir. Statik analiz ile key literal ise coezeueuelebilir.
\end{enumerate}

\subsection{Statik Analiz vs Runtime Emulation}

\begin{table}[ht]
\centering
\begin{tabular}{p{3.5cm}p{4cm}p{4cm}}
\toprule
\textbf{\color{sectioncyan}Yaklasim} & \textbf{\color{sectioncyan}Avantaj} & \textbf{\color{sectioncyan}Dezavantaj} \\
\midrule
AST-based statik & Gueuevenli, eval yok & Dinamik key'leri coezemez \\
eval() emulation & Her seyi cozer & Kod enjeksiyon riski \\
Sandboxed VM & Gueuevenli runtime & Performans maliyeti \\
\bottomrule
\end{tabular}
\end{table}

Karadul dogru tercihi yapmis: synchrony AST-based statik analiz kullaniyor, eval \textbf{kullanmiyor}.


\section{Computed-to-Static --- Soundness Analizi}
\label{sec:computed-deep}

Doenueuseueuem: \code{obj["foo"]} $\to$ \code{obj.foo}

\begin{tanim}[Gueuevenli Doenueuseueuem Kosullari]
\begin{enumerate}
\item Property key bir StringLiteral olmali (degisken/expression degil)
\item Key degeri gecerli JavaScript identifier olmali: \code{/\^{}[a-zA-Z\_\$][a-zA-Z0-9\_\$]*\$/}
\item Key degeri reserved word \textbf{olmamali}
\end{enumerate}
\end{tanim}

\textbf{Gueuevenli olmayan} durumlar:
\begin{itemize}
\item \code{obj["class"]} --- reserved word
\item \code{obj["foo-bar"]} --- tire gecerli identifier degil
\item \code{obj[variable]} --- StringLiteral degil
\item \code{obj[fn()]} --- runtime degeri
\end{itemize}

Karadul \textbf{soundness'i completeness'a tercih eder} --- yanlis doenueuseueuem kodun anlamini degistirir ve geri alinmasi cok zor.


\section{Semantik Degisken Isimlendirme --- 25 Kural Motoru}
\label{sec:semantic-rename-deep}

\subsection{Mimari}

\begin{center}
\begin{tikzpicture}[
  box/.style={rectangle, draw=sectioncyan, fill=boxbg, minimum width=5cm, minimum height=0.7cm, text=white, font=\small},
  arr/.style={-{Stealth[length=2.5mm]}, thick, color=gray!60!white},
]
  \node[box] (input) at (0,0) {Minified identifier: e, t, n};
  \node[box] (usage) at (0,-1.2) {Usage Info Collection (tum reference'lar)};
  \node[box] (rules) at (0,-2.4) {10 Kural Motoru (name, confidence)};
  \node[box] (select) at (0,-3.6) {En Yueueksek Confidence Secimi (esik: 0.3)};
  \node[box] (rename) at (0,-4.8) {Scope-Aware Rename (Babel scope.rename)};

  \draw[arr] (input) -- (usage);
  \draw[arr] (usage) -- (rules);
  \draw[arr] (rules) -- (select);
  \draw[arr] (select) -- (rename);
\end{tikzpicture}
\end{center}

\subsection{Kural Detaylari}

\textbf{Kural 1 --- API Call Context:} $\sim$80 method 10 kategoriye ayrilmis:

\begin{table}[ht]
\centering
\footnotesize
\begin{tabular}{llll}
\toprule
\textbf{\color{sectioncyan}Kategori} & \textbf{\color{sectioncyan}Method Oernekleri} & \textbf{\color{sectioncyan}Sonuc} & \textbf{\color{sectioncyan}Conf.} \\
\midrule
String (24) & split, trim, replace, match & \code{str} veya \code{text} & 0.85 \\
Array (25) & push, pop, map, filter, reduce & \code{items} & 0.85 \\
Map/Set & get+set+has & \code{cache} & 0.80 \\
Promise & then+catch & \code{promise} & 0.80 \\
EventEmitter & on+emit & \code{emitter} & 0.80 \\
Stream & pipe+write & \code{stream} & 0.80 \\
Buffer & readUInt8, writeUInt8 & \code{buffer} & 0.85 \\
RegExp & test+exec & \code{regex} & 0.85 \\
Date & getTime, toISOString & \code{date} & 0.80 \\
\bottomrule
\end{tabular}
\caption{API Call Context kural detaylari.}
\end{table}

\textbf{Kural 3 --- Property Access Patterns} (en yueueksek precision):

\begin{table}[ht]
\centering
\footnotesize
\begin{tabular}{lll}
\toprule
\textbf{\color{sectioncyan}Property Pattern} & \textbf{\color{sectioncyan}Sonuc} & \textbf{\color{sectioncyan}Conf.} \\
\midrule
\code{.message} + \code{.stack} & \code{error} (Error object) & 0.90 \\
\code{.body} + \code{.params} + \code{.query} & \code{request} (HTTP) & 0.90 \\
\code{.innerHTML} + \code{.classList} & \code{element} (DOM) & 0.85 \\
\code{.readFileSync} & \code{fs} (Node.js fs module) & 0.90 \\
\code{.resolve} + \code{.join} & \code{path} (Node.js path module) & 0.85 \\
\bottomrule
\end{tabular}
\caption{Property Access pattern'lari.}
\end{table}

\subsection{Neden LLM Kullanilmiyor?}

\begin{enumerate}
\item \textbf{Hiz:} 50K identifier $\times$ $\sim$100ms/istek = 83 dakika (LLM). Heuristik: $< 2$ saniye.
\item \textbf{Determinizm:} Ayni girdi $\to$ her seferinde ayni cikti.
\item \textbf{Maliyet:} API maliyeti $\sim$\$20/analiz.
\item \textbf{Cevrimdisi:} Internet gerektirmez.
\end{enumerate}

Trade-off: LLM ``processUserAuthentication'' gibi akilli isimler ueretebilir. Ama 25 kural pratikte yeterli ve 1000$\times$ hizli.


\section{Webpack Module Extraction --- Detayli}
\label{sec:webpack-deep}

\subsection{Extraction Algoritmasi}

\file{deobfuscate.mjs:543--581}'de:

\begin{enumerate}
\item \textbf{Module object tespiti:} \code{ObjectExpression} visitor ile bueyueuek key-value object'leri bul
\item \textbf{Heuristik:} En az 2 numeric/string key $+$ en az 2 function value $\to$ webpack module object
\item \textbf{Cikarma:} Her key-value cifti bir modul. \code{generate()} ile kaynak kodunu cikar, ayri dosyaya yaz
\end{enumerate}

\begin{lstlisting}[language=JavaScript]
// Heuristik (satir 548-566):
let numericKeyCount = 0, funcValueCount = 0;
for (const prop of props) {
    if (key.type === "NumericLiteral" || key.type === "StringLiteral")
        numericKeyCount++;
    if (prop.value.type === "FunctionExpression")
        funcValueCount++;
}
if (numericKeyCount >= 2 && funcValueCount >= 2) {
    // Webpack module object tespit edildi!
}
\end{lstlisting}

\subsection{esbuild vs webpack}

\begin{table}[ht]
\centering
\begin{tabular}{lll}
\toprule
\textbf{\color{sectioncyan}Ozellik} & \textbf{\color{sectioncyan}webpack} & \textbf{\color{sectioncyan}esbuild} \\
\midrule
Module storage & Object literal & Array veya Object \\
Loader & \code{\_\_webpack\_require\_\_} & \code{\_\_require} veya inline \\
Tree shaking & Sinirli & Agresif (ESM-native) \\
Module ID & Numeric veya string hash & Path-based string \\
\bottomrule
\end{tabular}
\end{table}


\section{Chunked Processing --- 200MB+ Dosyalar}
\label{sec:chunked-deep}

\subsection{Problem}

200MB JS dosyasi icin Babel parse: $\sim$4--8$\times$ bellek kullanimi $\to$ 800MB--1.6GB heap. Node.js default heap 1.5GB $\to$ bellek tasmasi.

\subsection{Neden 512KB Chunk?}

\file{stream-parse.mjs:33}'te: \code{maxBlockKb = 512}

\begin{enumerate}
\item \textbf{Babel bellek:} 512KB $\times$ 8 = 4MB per chunk. 8GB heap icinde yuezlerce chunk islenebilir.
\item \textbf{JS statement boyutu:} Cogu top-level fonksiyon 1--50KB. 512KB bueyueuek fonksiyonlari bile tek chunk'ta tutar.
\item \textbf{Cache:} Apple Silicon L2 cache $\sim$16MB. 512KB rahatca sigar.
\end{enumerate}

\subsection{Boundary Tespiti}

\textbf{Brace depth tracking:} \code{\{} goerueueluuence depth++, \code{\}} goerueuelueuence depth-{}-. \code{depth === 0} iken yeni top-level statement $\to$ split.

\begin{uyari}[Cross-Chunk Reference Sorunu]
Bir chunk'taki degisken baska chunk'ta tanimli olabilir. Karadul her chunk'i bagimsiz isleyip sonuclari birlestirir --- en basit ve performansli ama semantic bilgi kaybina yol acar.
\end{uyari}

\subsection{Acorn Fallback --- 32MB+ Dosyalar}

Babel parser 32MB+ dosyalarda coekmekti. Coezeueuem: ueuec katmanli hata kurtarma:

\begin{enumerate}
\item \textbf{Babel parse:} Normal deneme
\item \textbf{Pre-beautify + Babel:} Babel basarisiz $\to$ js-beautify ile formatlayip tekrar dene
\item \textbf{Acorn fallback:} Babel yine basarisiz $\to$ Acorn parser devreye girer
\end{enumerate}

Bu fix, LM Studio 32MB renderer bundle'ini coezeueuep 4/4 bundle basarisini sagladi.


% --- DEVELOPER AGENT DETAYLI GENISLETME: BOLUM 5+6 ---
% ============================================================================
% BOLUM 5 & 6: JavaScript Deobfuscation + Bayesian Isim Fuzyonu
% ============================================================================
% Bu dosya main.tex'teki Bolum 5 ve 6'nin GENISLETILMIS versiyonudur.
% main.tex'te % ============================================================================
% BOLUM 5 & 6: JavaScript Deobfuscation + Bayesian Isim Fuzyonu
% ============================================================================
% Bu dosya main.tex'teki Bolum 5 ve 6'nin GENISLETILMIS versiyonudur.
% main.tex'te % ============================================================================
% BOLUM 5 & 6: JavaScript Deobfuscation + Bayesian Isim Fuzyonu
% ============================================================================
% Bu dosya main.tex'teki Bolum 5 ve 6'nin GENISLETILMIS versiyonudur.
% main.tex'te \input{chapters/ch05_06_deobf_bayesian} ile dahil edilebilir.
% ONEMLI: \chapter tanimlanmaz, cunku chapter main.tex'te tanimlidir.
%         Burada \section ve \subsection seviyesinde yazilir.
% ============================================================================


% ############################################################################
%
%                    B O L U M   5 :   D E O B F U S C A T I O N
%
% ############################################################################

% Bu bolum, Karadul'un en yenilikci bilesenlerinden birini --- 9 fazli derin
% JavaScript deobfuscation pipeline'ini --- detayli olarak aciklar.
% Kaynak kodlar:
%   - karadul/deobfuscators/manager.py
%   - karadul/deobfuscators/deep_pipeline.py
%   - karadul/deobfuscators/babel_pipeline.py
%   - karadul/deobfuscators/synchrony_wrapper.py
%   - karadul/deobfuscators/string_decryptor.py
%   - karadul/deobfuscators/opaque_predicate.py
%   - karadul/deobfuscators/cff_deflattener.py
%   - scripts/deobfuscate.mjs
%   - scripts/deep-deobfuscate.mjs
%   - scripts/beautify.mjs
%   - scripts/smart-webpack-unpack.mjs
%   - scripts/cursor-enhanced-rename.mjs


% ===================================================================
\section{Obfuscation: Formal Tanim ve Teorik Temel}
\label{sec:obf-formal}
% ===================================================================

\subsection{Semantik Koruma Teoremi}

Kod obfuscation, bir programin kaynak kodunu, davranisini degistirmeden okunamazlastirma islemidir. Bunu formal olarak tanimlayalim.

\begin{tanim}[Obfuscation Fonksiyonu]
Bir obfuscation fonksiyonu $\mathcal{O}: \mathcal{P} \to \mathcal{P}$ asagidaki kosullari saglayan bir donustumdur:
\begin{enumerate}
\item \textbf{Semantik koruma:} Her girdi $x$ icin, $\mathcal{O}(P)(x) = P(x)$
\item \textbf{Polinom yavaslatma:} $|\mathcal{O}(P)| \leq \text{poly}(|P|)$ ve calisma suresi polinomial olarak artar
\item \textbf{Analiz zorlugu:} Herhangi bir analiz $A$ icin, $A(\mathcal{O}(P))$'den elde edilen bilgi, $A(P')$'den ($P'$ erisim saglanmamis bir program) elde edilenden anlamli olarak fazla degildir
\end{enumerate}
\end{tanim}

Bu tanimi daha kompakt yazarsak:

\begin{equation}
\mathcal{O}: \mathcal{P} \to \mathcal{P}, \quad \forall P \in \mathcal{P}: \quad \llbracket \mathcal{O}(P) \rrbracket = \llbracket P \rrbracket
\label{eq:semantic-preservation}
\end{equation}

burada $\llbracket \cdot \rrbracket$ programin denotational semantigini (yani giris-cikis davranisini) gosterir.

\begin{theorem}[Obfuscation Tersini Alma]
Eger $\mathcal{O}$ yalnizca \textbf{syntactic} donusumler uyguluyorsa (string encoding, isim degistirme, control flow yeniden yapilama), o zaman bir ters-obfuscation fonksiyonu $\mathcal{D}$ oyle vardir ki:
\begin{equation}
\mathcal{D}(\mathcal{O}(P)) \approx P
\label{eq:deobf-inverse}
\end{equation}
burada $\approx$ yapisal esdegerlik (structural equivalence) anlamindadir.
\end{theorem}

\textit{Ispat fikri:} Syntactic obfuscation, AST uzerinde terslenebilir donusumler uygular. Her donusum $t_i$ icin bir ters $t_i^{-1}$ mevcuttur. Obfuscation $\mathcal{O} = t_n \circ \cdots \circ t_1$ ise, $\mathcal{D} = t_1^{-1} \circ \cdots \circ t_n^{-1}$ dir. Ancak pratikte $t_i$'ler bilinmez; dolayisiyla Karadul gibi araclar \textbf{heuristik ters donusumleri} uygulamak zorundadir.

\subsection{Obfuscation Karmasiklik Siniflamasi}

Obfuscation tekniklerini bilgi-kuramsal zorluk derecesine gore siniflandiririz:

\begin{equation}
\text{Zorluk}(\mathcal{O}) = -\sum_{i} p_i \log_2 p_i + \text{cfg\_depth}(\mathcal{O}(P))
\label{eq:obf-complexity}
\end{equation}

burada ilk terim Shannon entropisi, ikinci terim control flow graph derinligidir.

\begin{table}[h]
\centering
\begin{tabular}{llcc}
\toprule
\textbf{\color{sectioncyan}Teknik} & \textbf{\color{sectioncyan}Zorluk} & \textbf{\color{sectioncyan}Tersinirlik} & \textbf{\color{sectioncyan}Karadul Kapsami} \\
\midrule
Minification & Dusuk & $\sim$100\% & Beautify \\
String encoding & Dusuk & $\sim$100\% & String unhex \\
Dead code injection & Orta & $\sim$95\% & Dead code elim. \\
String rotation & Orta & $\sim$90\% & Synchrony \\
Opaque predicates & Orta--Yuksek & $\sim$85\% & Pattern + Z3 \\
Control flow flattening & Yuksek & $\sim$70\% & CFF deflattener \\
Virtualization & Cok yuksek & $\sim$30\% & Kismi destek \\
\bottomrule
\end{tabular}
\caption{Obfuscation teknikleri ve zorluk dereceleri.}
\label{tab:obf-techniques}
\end{table}


% ===================================================================
\section{Yaygin Obfuscation Teknikleri}
\label{sec:obf-techniques}
% ===================================================================

\subsection{Minification}

En basit obfuscation formu. Kaynak koddan gereksiz bilgileri (bosluk, yorum, uzun degisken isimleri) kaldirarak dosya boyutunu kucultme ve okunakliligi azaltma islevidir.

\begin{lstlisting}[language=JavaScript, caption={Minification oncesi ve sonrasi}]
// Orijinal (okunabilir):
function calculateTotalPrice(items, taxRate) {
    let subtotal = 0;
    for (const item of items) {
        subtotal += item.price * item.quantity;
    }
    return subtotal * (1 + taxRate);
}

// Minified:
function a(b,c){let d=0;for(const e of b)d+=e.price*e.quantity;return d*(1+c)}
\end{lstlisting}

Minification bilgi kaybina neden olur: degisken isimleri \code{a}, \code{b}, \code{c} orijinal semantigini tasimaz. Ancak \textbf{program davranisi} aynidir --- Denklem~\ref{eq:semantic-preservation} gecerlidir.

\subsection{String Encoding}

String literal'ler cogunlukla iki yontemle gizlenir:

\textbf{Hex encoding:}
\begin{equation}
\text{encode}_{\text{hex}}(s) = \text{concat}_{i=0}^{|s|-1} \texttt{"\textbackslash x"} \cdot \text{hex}(\text{charCodeAt}(s, i))
\end{equation}

\textbf{Unicode encoding:}
\begin{equation}
\text{encode}_{\text{unicode}}(s) = \text{concat}_{i=0}^{|s|-1} \texttt{"\textbackslash u"} \cdot \text{pad}(\text{hex}(\text{codePointAt}(s, i)), 4)
\end{equation}

\textbf{Base64 encoding:}
\begin{equation}
\text{encode}_{\text{b64}}(s) = \texttt{atob("} \cdot \text{base64}(s) \cdot \texttt{")}
\end{equation}

Ornek: \code{"Hello"} $\to$ \code{"\textbackslash x48\textbackslash x65\textbackslash x6c\textbackslash x6c\textbackslash x6f"} $\to$ \code{atob("SGVsbG8=")}

\subsection{Control Flow Flattening (CFF)}

CFF, programin dogal kontrol akisini bir \texttt{while-switch} dispatcher'a donusturur:

\begin{lstlisting}[language=JavaScript, caption={CFF ornegi}]
// Orijinal:
function process(x) {
    let a = x + 1;    // Block 0
    let b = a * 2;    // Block 1
    if (b > 10) {     // Block 2
        return b;     // Block 3
    }
    return a;          // Block 4
}

// CFF uygulanmis:
function process(x) {
    let state = 0, a, b;
    while (true) {
        switch (state) {
            case 0: a = x + 1; state = 1; break;
            case 1: b = a * 2; state = 2; break;
            case 2: state = b > 10 ? 3 : 4; break;
            case 3: return b;
            case 4: return a;
        }
    }
}
\end{lstlisting}

CFF'nin formal modeli bir sonlu durum makinesidir (FSM):

\begin{equation}
\text{CFF}(P) = (Q, \Sigma, \delta, q_0, F)
\end{equation}

burada $Q$ state kuemesi, $\delta: Q \times \Sigma \to Q$ gecis fonksiyonu, $q_0$ baslangic durumu ve $F \subseteq Q$ bitis durumlaridir.

\begin{center}
\begin{tikzpicture}[
  state/.style={circle, draw=sectioncyan, fill=boxbg, minimum size=1.2cm, text=white, font=\small\bfseries},
  arrow/.style={-{Stealth[length=3mm]}, thick, color=gray!70!white},
  every node/.style={font=\small},
]
  \node[state] (q0) at (0,0) {$q_0$};
  \node[state] (q1) at (3,0) {$q_1$};
  \node[state] (q2) at (6,0) {$q_2$};
  \node[state, fill=green!20!black] (q3) at (9,1) {$q_3$};
  \node[state, fill=green!20!black] (q4) at (9,-1) {$q_4$};

  \draw[arrow] (-1.5,0) -- (q0);
  \draw[arrow] (q0) -- node[above, text=gray!60!white] {\texttt{a=x+1}} (q1);
  \draw[arrow] (q1) -- node[above, text=gray!60!white] {\texttt{b=a*2}} (q2);
  \draw[arrow] (q2) -- node[above, text=gray!60!white, sloped] {\texttt{b>10}} (q3);
  \draw[arrow] (q2) -- node[below, text=gray!60!white, sloped] {\texttt{b<=10}} (q4);

  \node[below=0.6cm, font=\tiny, text=cyan!70!white] at (q3.south) {return b};
  \node[below=0.6cm, font=\tiny, text=cyan!70!white] at (q4.south) {return a};
\end{tikzpicture}
\end{center}

\subsection{Opaque Predicates}

Opaque predicate, her zaman ayni sonucu veren (true veya false) ama statik analizin bunu kolayca cikaramayacagi ifadelerdir.

\begin{tanim}[Opaque Predicate]
Bir ifade $\varphi$ icin:
\begin{itemize}
\item $\varphi^T$: Her zaman true ($\forall \sigma: \llbracket \varphi \rrbracket_\sigma = \text{true}$)
\item $\varphi^F$: Her zaman false ($\forall \sigma: \llbracket \varphi \rrbracket_\sigma = \text{false}$)
\item $\varphi^?$: Bazen true, bazen false (gercek dallanma)
\end{itemize}
burada $\sigma$ program state'ini temsil eder.
\end{tanim}

\file{opaque\_predicate.py} dosyasinda tanimlanan oruntuleri inceleyelim:

\begin{table}[h]
\centering
\small
\begin{tabular}{p{5cm}cp{6cm}}
\toprule
\textbf{\color{sectioncyan}Oruuntu} & \textbf{\color{sectioncyan}Sonuc} & \textbf{\color{sectioncyan}Matematik Ispati} \\
\midrule
\code{x*(x+1) \% 2 == 0} & $\varphi^T$ & Ardisik iki sayinin carpimi daima cifttir \\
\code{(x | 1) != 0} & $\varphi^T$ & OR 1 her zaman $\neq 0$ verir \\
\code{(a \& b) == (a \& b)} & $\varphi^T$ & Her ifade kendisiyle esittir \\
\code{2*(x/2) <= x} & $\varphi^T$ & Tamsayi bolumuenun oezelligi: $\lfloor x/2 \rfloor \cdot 2 \leq x$ \\
\code{x*(x+1) \% 2 != 0} & $\varphi^F$ & Yukaridakinin tueuemleri \\
\code{x != x} (int) & $\varphi^F$ & Tamsayi kendisiyle esit (NaN notu asagida) \\
\bottomrule
\end{tabular}
\caption{Opaque predicate oruntuleri ve ispatlar. \textbf{Uyari:} \code{x == x} IEEE 754 NaN icin \texttt{false} dondurur; \file{opaque\_predicate.py} bu nedenle float baglam icin confidence'i 0.60'a dusurur.}
\label{tab:opaque-predicates}
\end{table}

\subsection{Dead Code Injection}

Opaque predicate'ler ile birlesirildiginde, hicbir zaman calistirilmayan sahte kod bloklari eklenir:

\begin{lstlisting}[language=JavaScript, caption={Dead code injection ornegi}]
if (x * (x + 1) % 2 === 0) {
    // Gercek kod (her zaman calisir)
    processPayment(amount);
} else {
    // Dead code (asla calismaz)
    launchMissiles();  // dikkat dagitma amacli
    deleteDatabase();  // analisti yaniltma
}
\end{lstlisting}

\subsection{String Array Rotation}

obfuscator.io gibi araclar tum string'leri bir diziye tasir, diziyi kaydirma (rotation) islemiyle karistirip indeksle erisirir:

\begin{lstlisting}[language=JavaScript, caption={String rotation ornegi}]
// Orijinal:
console.log("Hello");

// Obfuscated:
var _0x1234 = ["Hello", "log", "World"];
(function(_0x, _0y) {
    var _0z = function(_0w) { while(--_0w) { _0x.push(_0x.shift()); } };
    _0z(++_0y);
})(_0x1234, 0x1);  // Diziyi 1 kaydır
console[_0x1234[0]](_0x1234[1]);  // console.log("World") -- indeksler degisti!
\end{lstlisting}

Karadul'un \texttt{synchrony} biliseni bu patterni dogrudan cozer cunku synchrony, obfuscator.io spesifik bir deobfuscation aracidir.


% ===================================================================
\section{Abstract Syntax Tree (AST) Temelleri}
\label{sec:ast-basics}
% ===================================================================

Tum deobfuscation pipeline'i AST uzerinde calisir. Bu bolum AST'nin ne oldugunu, parse tree ile farkini ve ESTree standardini aciklar.

\subsection{BNF ve Context-Free Grammar (CFG)}

JavaScript dili bir context-free grammar (CFG) ile tanimlenir:

\begin{equation}
G = (V, \Sigma, R, S)
\end{equation}

burada $V$ non-terminal semboller, $\Sigma$ terminal semboller, $R$ uretim kurallari ve $S$ baslangic sembolueduer.

JavaScript'in basitlestirilmis BNF'inde \texttt{IfStatement}:

\begin{lstlisting}[basicstyle=\ttfamily\small\color{white}, frame=none]
IfStatement ::= "if" "(" Expression ")" Statement ("else" Statement)?
Expression  ::= AssignmentExpression | Expression "," AssignmentExpression
Statement   ::= Block | IfStatement | ForStatement | ExpressionStatement | ...
\end{lstlisting}

\subsection{Parse Tree vs Abstract Syntax Tree}

\textbf{Parse tree} (concrete syntax tree), gramerin her uretim kuralini icerir --- parantezler, noktalama isaretleri, terminal semboller hepsi agacta gorunur. \textbf{AST} ise semantik olarak anlamsiz node'lari atar.

\begin{center}
\begin{tikzpicture}[
  node distance=0.8cm and 0.5cm,
  treenode/.style={rectangle, rounded corners=2pt, draw=sectioncyan, fill=boxbg, text=white, font=\scriptsize, minimum height=0.6cm, minimum width=1cm},
  leaf/.style={rectangle, rounded corners=2pt, draw=green!60!white, fill=black!80!white, text=green!60!white, font=\scriptsize\ttfamily, minimum height=0.5cm},
  edge from parent/.style={draw=gray!60!white, -{Stealth[length=2mm]}, thin},
  level 1/.style={sibling distance=3.5cm},
  level 2/.style={sibling distance=1.8cm},
  level 3/.style={sibling distance=1.2cm},
]
  % AST tarafı
  \node[treenode, fill=cyan!20!black] {BinaryExpression (+)}
    child { node[treenode] {Identifier}
      child { node[leaf] {a} }
    }
    child { node[treenode] {NumericLiteral}
      child { node[leaf] {1} }
    };

  \node[above=0.3cm, font=\small\color{sectioncyan}] at (0, 1.2) {AST: \texttt{a + 1}};
\end{tikzpicture}
\end{center}

\subsection{ESTree Spesifikasyonu}

JavaScript AST'si icin fiili standart \textbf{ESTree} spesifikasyonudur (\texttt{https://github.com/estree/estree}). Babel parser, ESTree'nin genisletilmis bir versiyonunu uretir:

\begin{itemize}
\item \textbf{Babel AST:} \code{StringLiteral}, \code{NumericLiteral}, \code{ObjectProperty}
\item \textbf{Acorn/ESTree:} \code{Literal}, \code{Property}
\end{itemize}

Bu fark \file{deobfuscate.mjs}'te asagidaki kontrol ile ele alinir:

\begin{lstlisting}[language=JavaScript, caption={Babel vs Acorn AST tespiti (deobfuscate.mjs)}]
const usedAcorn = ast.body !== undefined && !ast.program;
// Babel AST'de ast.program var, Acorn'da yok
\end{lstlisting}

\subsection{Babel Parser Oezellikleri}

\file{deep-deobfuscate.mjs} Babel parser'i su opsiyonlarla cagrilir:

\begin{lstlisting}[language=JavaScript, caption={Babel parser konfigurasyonu}]
ast = parse(source, {
  sourceType: "unambiguous",
  allowImportExportEverywhere: true,
  allowReturnOutsideFunction: true,
  errorRecovery: true,  // Parse hatalarinda devam et
  plugins: ["jsx", "typescript", "decorators-legacy",
    "classProperties", "dynamicImport",
    "optionalChaining", "nullishCoalescingOperator",
    "topLevelAwait", "importMeta"],
});
\end{lstlisting}

\code{errorRecovery: true} kritik bir secimdir. 30MB+ minified dosyalarda parser kismen basarisiz olabilir; bu mod, hatalari toplar ama parse islemini durdurmaz. Boylece kismi AST uzerinde islem yapilabilir.


% ===================================================================
\section{Standart 4 Adimli Deobfuscation Zinciri}
\label{sec:std-chain-detail}
% ===================================================================

\file{manager.py} dosyasindaki \code{DeobfuscationManager} sinifi, standart 4 adimli zinciri orkestre eder:

\begin{center}
\begin{tikzpicture}[
  pstep/.style={rectangle, draw=sectioncyan, fill=boxbg, minimum width=3.2cm, minimum height=1.4cm, text=white, font=\small\bfseries, rounded corners=3pt, align=center},
  arrow/.style={-{Stealth[length=3mm]}, thick, color=gray!70!white},
  fail/.style={-{Stealth[length=2mm]}, dashed, color=red!60!white, thick},
  node distance=1.2cm,
]
  \node[pstep] (input) at (0,0) {Obfuscated\\JS Dosyasi};
  \node[pstep] (beautify) at (4.5,0) {1. Beautify\\(js-beautify)};
  \node[pstep] (synchrony) at (9,0) {2. Synchrony\\(obfuscator.io)};
  \node[pstep] (babel) at (4.5,-3) {3. Babel AST\\(9 faz transform)};
  \node[pstep] (webpack) at (9,-3) {4. Webpack\\Unpack};
  \node[pstep, fill=green!15!black] (output) at (13.5,-3) {Deobfuscated\\Cikti};

  \draw[arrow] (input) -- (beautify);
  \draw[arrow] (beautify) -- (synchrony);
  \draw[arrow] (synchrony) |- (babel);
  \draw[arrow] (babel) -- (webpack);
  \draw[arrow] (webpack) -- (output);

  % Graceful degradation
  \draw[fail] (beautify.south) -- ++(0,-0.8) node[midway, right, font=\tiny, text=red!60!white] {fail} -| (synchrony.south);
  \draw[fail] (synchrony.south) -- ++(0,-0.4) node[midway, right, font=\tiny, text=red!60!white] {fail} -- (babel.north);

  \node[below=0.2cm, font=\tiny, text=gray!50!white] at (beautify.south east) {\file{beautify.mjs}};
  \node[below=0.2cm, font=\tiny, text=gray!50!white] at (synchrony.south east) {\file{synchrony}};
  \node[below=0.2cm, font=\tiny, text=gray!50!white] at (babel.south east) {\file{deobfuscate.mjs}};
  \node[below=0.2cm, font=\tiny, text=gray!50!white] at (webpack.south east) {\file{smart-webpack-unpack.mjs}};
\end{tikzpicture}
\end{center}

\subsection{Graceful Degradation Formuelue}

Her adim basarisiz olabilir. Bu durumda bir onceki basarili cikti sonraki adima gecilir:

\begin{equation}
\text{output}(i) = \begin{cases}
\text{step}_i(\text{output}(i-1)) & \text{eger } \text{step}_i \text{ basarili} \\
\text{output}(i-1) & \text{eger } \text{step}_i \text{ basarisiz (fallback)}
\end{cases}
\label{eq:graceful-degradation}
\end{equation}

Zincirin en az bir adimi basarili olursa, tum zincir ``basarili'' sayilir:

\begin{equation}
\text{success}(\text{chain}) = \bigvee_{i=1}^{n} \text{success}(\text{step}_i)
\label{eq:chain-success}
\end{equation}

\file{manager.py} bunu su kodla implement eder:

\begin{lstlisting}[language=Python, caption={Graceful degradation -- manager.py, run\_chain()}]
for idx, step_name in enumerate(effective_chain, start=1):
    step_output = deob_dir / f"{idx:02d}_{step_name}{input_file.suffix}"
    try:
        step_success = self._execute_step(
            step_name, current_input, step_output, deob_dir)
    except Exception as exc:
        step_success = False
    if step_success and step_output.exists():
        result.steps_completed.append(step_name)
        current_input = step_output  # Sonraki adimin girdisi
    else:
        result.steps_failed.append(step_name)
        # current_input degismez -> onceki basarili cikti kullanilir
\end{lstlisting}

\subsection{Dosya Isimlendirme Konvansiyonu}

Her adim, ciktisini numarali dosya olarak kaydeder:

\begin{center}
\begin{tabular}{ll}
\code{00\_original.js} & Orijinal obfuscated dosya \\
\code{01\_beautify.js} & Beautify sonrasi \\
\code{02\_synchrony.js} & Synchrony sonrasi \\
\code{03\_babel\_transforms.js} & Babel AST sonrasi \\
\code{04\_webpack\_unpack.js} & Webpack modul ayirma sonrasi \\
\end{tabular}
\end{center}

Bu sayede her adimin ciktisi bagimsiz olarak incelenebilir --- debug ve audit icin kritik.


% ===================================================================
\section{Adim 1: Beautify}
\label{sec:beautify}
% ===================================================================

\file{beautify.mjs}, \texttt{js-beautify} kutuphanesini kullanan basit ama kritik bir onisleme adimidir. Minified kod tek satir halindedir; AST parse bile basa\-ri\-siz olabilir cunku Babel'in ternary+arrow+object pattern'lerinde takilma riski vardir.

\subsection{Konfiguerasyon Parametreleri}

\begin{lstlisting}[language=JavaScript, caption={js-beautify konfigurasyonu -- beautify.mjs}]
js_beautify(source, {
  indent_size: 2,
  indent_char: " ",
  max_preserve_newlines: 2,
  brace_style: "collapse,preserve-inline",
  wrap_line_length: 120,
  break_chained_methods: true,
  end_with_newline: true,
});
\end{lstlisting}

\subsection{Etki Analizi}

Beautify'in etkisini olcmek icin genisleme oranini (expansion ratio) tanimlayalim:

\begin{equation}
R_{\text{expand}} = \frac{L_{\text{out}}}{L_{\text{in}}}
\end{equation}

burada $L_{\text{in}}$ girdi satir sayisi, $L_{\text{out}}$ cikti satir sayisi. Tipik degerler:

\begin{itemize}
\item Agresif minification (tek satir): $R_{\text{expand}} \approx 100\text{--}500$
\item Orta minification: $R_{\text{expand}} \approx 2\text{--}5$
\item Zaten formatlenmis: $R_{\text{expand}} \approx 1$
\end{itemize}

\subsection{Fallback Davranisi}

Beautify basarisiz olursa, girdi dosyasi dogrudan ciktiya kopyalanir:

\begin{lstlisting}[language=Python, caption={Beautify fallback -- manager.py}]
def _step_beautify(self, input_file, output_file):
    beautify_script = self._config.scripts_dir / "beautify.mjs"
    if not beautify_script.exists():
        shutil.copy2(input_file, output_file)  # Fallback
        return True
\end{lstlisting}


% ===================================================================
\section{Adim 2: Synchrony Deobfuscation}
\label{sec:synchrony}
% ===================================================================

\file{synchrony\_wrapper.py}, \texttt{synchrony} aracini saran bir Python wrapper'dir. Synchrony, obfuscator.io spesifik bir deobfuscation aracidir ve string rotation, string array, control flow flattening gibi obfuscator.io'ya oezgue teknikleri cozer.

\subsection{Calismaa Mekanizmasi}

\begin{enumerate}
\item \textbf{Shebang soyma:} Eger dosya \code{\#!/usr/bin/env node} ile basliyorsa, synchrony'nin Acorn parser'i bunu parse edemez. Wrapper gecici dosyada shebang'i soyar, islem sonrasi geri ekler.
\item \textbf{Source type deneme:} Synchrony uc modda denenir: \code{both}, \code{module}, \code{script}. ESM dosyalarda \code{ImportDeclaration} hatasi alinirsa \code{module} moduna gecilir.
\item \textbf{Dogrudan dosya ciktisi:} \code{-o} flag'i ile stdout yerine dogrudan dosyaya yazar. Buyuk dosyalarda memory overhead'i onler.
\end{enumerate}

\begin{lstlisting}[language=Python, caption={Synchrony calisma akisi -- synchrony\_wrapper.py}]
source_types = ["both", "module", "script"]
for source_type in source_types:
    cmd = [sync_path, str(actual_input),
           "-o", str(output_file),
           "--type", source_type]
    result = self._runner.run_command(cmd, timeout=timeout)
    if result.success:
        break
    if "ImportDeclaration" in result.stderr:
        continue  # ESM hatasi, module modunu dene
\end{lstlisting}

\subsection{Boyut Analizi}

Synchrony basarili oldugunda dosya boyutu genellikle \textbf{kucueluer} cunku string rotation ve dead code kaldirilir:

\begin{equation}
R_{\text{sync}} = \frac{|\text{output}|}{|\text{input}|} \in [0.3, 1.0]
\end{equation}

$R_{\text{sync}} < 0.5$ ise agresif obfuscation kaldirilmis demektir.


% ===================================================================
\section{9 Fazli Deep Deobfuscation Pipeline}
\label{sec:deep-deobf-detail}
% ===================================================================

\file{deep-deobfuscate.mjs} dosyasi, 10 asamali (9 transform + 1 parse) derin deobfuscation pipeline'i uygular. Phase 2--8, \textbf{visitor-merge optimizasyonu} ile tek bir AST traverse'da birlestirilmistir --- bu webcrack yaklasimi, 7 ayri traverse yerine 1 traverse kullanarak $\sim$3--5x hiz kazanci saglar.

\begin{center}
\begin{tikzpicture}[
  phase/.style={rectangle, draw=sectioncyan, fill=boxbg, minimum width=4cm, minimum height=0.8cm, text=white, font=\scriptsize\bfseries, rounded corners=2pt},
  merged/.style={rectangle, draw=yellow!70!white, fill=boxbg, minimum width=4cm, minimum height=0.8cm, text=yellow!70!white, font=\scriptsize\bfseries, rounded corners=2pt},
  arrow/.style={-{Stealth[length=2mm]}, thin, color=gray!60!white},
  node distance=0.5cm,
]
  \node[phase] (p1) {Phase 1: Parse};
  \node[merged, below=of p1] (p2) {Phase 2: Constant Folding};
  \node[merged, below=of p2] (p3) {Phase 3: Dead Code Elim.};
  \node[merged, below=of p3] (p4) {Phase 4: Computed $\to$ Static};
  \node[merged, below=of p4] (p5) {Phase 5: Comma Split};
  \node[merged, below=of p5] (p6) {Phase 6: Var Decl. Split};
  \node[merged, below=of p6] (p7) {Phase 7: Ternary $\to$ If};
  \node[merged, below=of p7] (p8) {Phase 8: Arrow $\to$ Function};
  \node[phase, below=of p8] (p9) {Phase 9: Smart Rename};
  \node[phase, below=of p9] (p10) {Phase 10: Semantic Rename};

  \draw[arrow] (p1) -- (p2);
  \draw[arrow] (p2) -- (p3);
  \draw[arrow] (p3) -- (p4);
  \draw[arrow] (p4) -- (p5);
  \draw[arrow] (p5) -- (p6);
  \draw[arrow] (p6) -- (p7);
  \draw[arrow] (p7) -- (p8);
  \draw[arrow] (p8) -- (p9);
  \draw[arrow] (p9) -- (p10);

  % Merge bracket
  \draw[yellow!70!white, thick, decorate, decoration={brace, amplitude=8pt, mirror}]
    ([xshift=2.3cm]p2.north east) -- ([xshift=2.3cm]p8.south east)
    node[midway, right=10pt, text=yellow!70!white, font=\scriptsize, align=left] {Tek AST\\traverse\\(visitor-merge)};

  % Separate bracket
  \draw[sectioncyan, thick, decorate, decoration={brace, amplitude=8pt, mirror}]
    ([xshift=2.3cm]p9.north east) -- ([xshift=2.3cm]p10.south east)
    node[midway, right=10pt, text=sectioncyan, font=\scriptsize, align=left] {Ayri\\traverse\\(scope gerekli)};
\end{tikzpicture}
\end{center}

\subsection{Faz 1: Parse (Tolerant Mode)}
\label{sec:phase1}

Babel parser \code{errorRecovery: true} ile calistirilir. Basarisiz olursa iki fallback mevcuttur:

\begin{enumerate}
\item \textbf{Pre-beautify + retry:} Minified kaynak once \code{js-beautify} ile formatlanir, sonra Babel tekrar denenir. 30MB+ dosyalarda Babel'in ternary/arrow/object pattern'lerinde takilma sorunu bu yontemle cogunlukla cozulur.
\item \textbf{Dosya kopyalama:} Her iki yontem de basarisiz olursa, kaynak dosya oldugu gibi output'a kopyalanir. Boylece sonraki pipeline adimlari (webpack unpack) en azindan orijinal dosya uzerinde devam edebilir.
\end{enumerate}

\begin{equation}
\text{parse}(S) = \begin{cases}
\text{Babel}(S) & \text{basarili} \\
\text{Babel}(\text{beautify}(S)) & \text{pre-beautify ile retry} \\
S & \text{fallback: kaynak kopyalama}
\end{cases}
\end{equation}


\subsection{Faz 2: Constant Folding (Sabit Katlama)}
\label{sec:phase2-cf}

Constant folding, derleme/analiz zamaninda hesaplanabilir ifadeleri literal degerleriyle degistirme islemidir.

\subsubsection{Lattice Teorisi ve Abstract Interpretation}

Constant folding, abstract interpretation cercevesinde bir \textbf{flat lattice} uzerinde calisir:

\begin{equation}
L = \{\bot, c_1, c_2, \ldots, c_n, \top\}
\end{equation}

burada:
\begin{itemize}
\item $\bot$: Degisken henuz kullanilmadi (undefined)
\item $c_i$: Degisken sabit $c_i$ degerinde
\item $\top$: Degisken birden fazla deger alabilir (non-constant)
\end{itemize}

\begin{equation}
\text{meet}(x, y) = \begin{cases}
x & \text{eger } x = y \\
\top & \text{eger } x \neq y \text{ ve her ikisi de } \neq \bot \\
y & \text{eger } x = \bot \\
x & \text{eger } y = \bot
\end{cases}
\end{equation}

\subsubsection{Karadul Implementasyonu}

\file{deep-deobfuscate.mjs} Phase 2'de asagidaki donusumler uygulanir:

\begin{table}[h]
\centering
\small
\begin{tabular}{lll}
\toprule
\textbf{\color{sectioncyan}Node Tipi} & \textbf{\color{sectioncyan}Donusum} & \textbf{\color{sectioncyan}Ornek} \\
\midrule
BinaryExpression (+) & String concat & \code{"He"+"llo"} $\to$ \code{"Hello"} \\
BinaryExpression (+) & Numeric add & \code{1+2} $\to$ \code{3} \\
BinaryExpression ($*$, $/$, $-$) & Aritmetik & \code{6*7} $\to$ \code{42} \\
BinaryExpression ($|$,$\&$,$\hat{}$,$\ll$,$\gg$) & Bitwise & \code{0xFF \& 0x0F} $\to$ \code{15} \\
UnaryExpression (!) & Negation & \code{!0} $\to$ \code{true} \\
UnaryExpression (void) & Void & \code{void 0} $\to$ \code{undefined} \\
\bottomrule
\end{tabular}
\caption{Constant folding donusumm tablosu.}
\label{tab:constant-folding}
\end{table}

Guevenlik kontrolu: Bolme islemlerinde \code{right.value !== 0} kontrol edilir ve $\text{Number.isFinite}(result)$ dogrulanir. Overfllow, NaN ve Infinity sonuclari degistirilmez.

\begin{lstlisting}[language=JavaScript, caption={Constant folding visitor -- BinaryExpression}]
BinaryExpression: {
  exit(path) {
    const { left, right, operator } = path.node;
    // String concatenation
    if (operator === "+" && t.isStringLiteral(left)
        && t.isStringLiteral(right)) {
      path.replaceWith(t.stringLiteral(left.value + right.value));
      cnt2++;
    }
    // Numeric operations
    if (t.isNumericLiteral(left) && t.isNumericLiteral(right)) {
      let result;
      switch (operator) {
        case "+": result = left.value + right.value; break;
        case "-": result = left.value - right.value; break;
        // ... (*, /, %, **, |, &, ^, <<, >>, >>>)
      }
      if (result !== undefined && Number.isFinite(result)) {
        path.replaceWith(t.numericLiteral(result));
        cnt2++;
      }
    }
  }
}
\end{lstlisting}

\textbf{Neden \code{exit} visitor?} Bottom-up traversal sayesinde ic node'lar once islenir: \code{(1+2)*3}'te once \code{1+2=3}, sonra \code{3*3=9}. \code{enter} kullansaydik, ic node'lar henuz fold edilmemis olurdu.


\subsection{Faz 3: Dead Code Elimination}
\label{sec:phase3-dce}

\subsubsection{Liveness Analysis Temeli}

Liveness analysis, her program noktasinda hangi degiskenlerin ``canli'' (daha sonra kullanilacak) oldugunu belirler:

\begin{equation}
\text{live\_in}(B) = \text{gen}(B) \cup (\text{live\_out}(B) \setminus \text{kill}(B))
\label{eq:live-in}
\end{equation}
\begin{equation}
\text{live\_out}(B) = \bigcup_{S \in \text{succ}(B)} \text{live\_in}(S)
\label{eq:live-out}
\end{equation}

burada:
\begin{itemize}
\item $\text{gen}(B)$: $B$ blogunda kullanilan (okunan) degiskenler
\item $\text{kill}(B)$: $B$ blogunda tanimlanan (yazilan) degiskenler
\end{itemize}

\subsubsection{Fixed-Point Iteration}

Liveness analysis, denklem sistemi sabit noktaya ulasana kadar iteratif olarak cozulur:

\begin{equation}
\text{live\_in}^{(k+1)}(B) = \text{gen}(B) \cup \left(\text{live\_out}^{(k)}(B) \setminus \text{kill}(B)\right)
\end{equation}

Karmasiklik: $\bigO{|B| \times |V|}$ burada $|B|$ basic block sayisi, $|V|$ degisken sayisi.

\subsubsection{Karadul Implementasyonu}

Phase 3, \code{IfStatement} ve \code{ConditionalExpression} node'larini isle\-r:

\begin{lstlisting}[language=JavaScript, caption={Dead code elimination -- IfStatement}]
IfStatement: {
  exit(path) {
    const test = path.node.test;
    // if (false) {...} -> kaldir veya else branch'i al
    if (t.isBooleanLiteral(test) && test.value === false) {
      if (path.node.alternate) path.replaceWith(path.node.alternate);
      else path.remove();
      cnt3++;
    }
    // if (true) {X} else {Y} -> X
    if (t.isBooleanLiteral(test) && test.value === true) {
      path.replaceWith(path.node.consequent);
      cnt3++;
    }
    // if (0) -> remove, if (1) -> consequent
    if (t.isNumericLiteral(test)) {
      if (test.value === 0) { /* remove or alternate */ }
      else { path.replaceWith(path.node.consequent); }
      cnt3++;
    }
  }
}
\end{lstlisting}

\textbf{Phase 2 ile etkilesim:} Phase 2 \code{!0 $\to$ true} donusumunu yapar. Phase 3, \code{exit} visitor kullandigi icin, ayni traverse'da Phase 2'nin ciktisini gorebilir. Ornegin \code{if(!0)\{X\}} once \code{if(true)\{X\}} olur (Phase 2), sonra \code{X} olur (Phase 3).


\subsection{Faz 4: Computed-to-Static Property}
\label{sec:phase4-cts}

Bu faz, \code{obj["foo"]} gibi computed property erisimlerini \code{obj.foo} gibi statik erisimlere donusturur.

\begin{equation}
\text{obj}[\texttt{"key"}] \xrightarrow{\text{Faz 4}} \text{obj}.\texttt{key}
\quad \text{eger } \texttt{key} \in \text{ValidIdentifier}
\end{equation}

\begin{tanim}[Valid Identifier]
Bir string \code{key} gecerli JavaScript identifier ise ve reserved word degilse donusum guevenlidir:
\begin{equation}
\text{ValidIdentifier}(\texttt{key}) = \texttt{/\^{}[a-zA-Z\_\$][a-zA-Z0-9\_\$]*\$/}.test(\texttt{key}) \wedge \neg\text{isReserved}(\texttt{key})
\end{equation}
\end{tanim}

AST duenuesuemuee:
\begin{itemize}
\item \code{MemberExpression}: \code{computed: true} $\to$ \code{computed: false}
\item \code{ObjectProperty}: Computed key'i identifier'a donustur
\end{itemize}

\begin{lstlisting}[language=JavaScript, caption={Computed to static -- MemberExpression}]
MemberExpression(path) {
  if (!path.node.computed) return;
  const prop = path.node.property;
  if (!t.isStringLiteral(prop)) return;
  if (/^[a-zA-Z_$][a-zA-Z0-9_$]*$/.test(prop.value)
      && !isReservedWord(prop.value)) {
    path.node.computed = false;
    path.node.property = t.identifier(prop.value);
    cnt4++;
  }
}
\end{lstlisting}


\subsection{Faz 5: Comma Expression Splitting}
\label{sec:phase5-comma}

JavaScript'te virgul operatoru (\code{SequenceExpression}) birden fazla ifadeyi tek bir expression'da birlestirir. Obfuscator'lar bunu okunakliligi azaltmak icin kullanir:

\begin{lstlisting}[language=JavaScript]
// Obfuscated:
return a = 1, b = 2, c = 3, process(a, b, c);
// Deobfuscated (Faz 5 sonrasi):
a = 1;
b = 2;
c = 3;
return process(a, b, c);
\end{lstlisting}

Donusum ucg bag\-lam icin uygulanir:

\begin{enumerate}
\item \textbf{ExpressionStatement:} \code{(a, b, c);} $\to$ \code{a; b; c;}
\item \textbf{ReturnStatement:} \code{return (a, b, c);} $\to$ \code{a; b; return c;}
\item \textbf{ThrowStatement:} \code{throw (a, b, c);} $\to$ \code{a; b; throw c;}
\end{enumerate}

\begin{equation}
\text{split}(\texttt{return}(e_1, e_2, \ldots, e_n)) = \left[\text{expr}(e_1), \ldots, \text{expr}(e_{n-1}), \texttt{return}\; e_n\right]
\end{equation}


\subsection{Faz 6: Variable Declaration Splitting}
\label{sec:phase6-vardecl}

Tek bir \code{var/let/const} icindeki birden fazla tanimlama ayri satirlara bolunur:

\begin{lstlisting}[language=JavaScript]
// Oncesi:
var a = 1, b = 2, c = fn(), d = [];
// Sonrasi:
var a = 1;
var b = 2;
var c = fn();
var d = [];
\end{lstlisting}

\textbf{Guvenlik kontrolleri:} For-loop init (\code{for(var i=0, j=0; ...)}), for-in ve for-of icerisindeki deklarasyonlar bolunmez:

\begin{lstlisting}[language=JavaScript, caption={Variable declaration split -- guvenlik kontrolleri}]
if (path.parent.type === "ForStatement" && path.key === "init") return;
if (path.parent.type === "ForInStatement") return;
if (path.parent.type === "ForOfStatement") return;
\end{lstlisting}


\subsection{Faz 7: Ternary $\to$ If/Else}
\label{sec:phase7-ternary}

Uzun ternary ifadeler if/else bloklarina donusturulur. ``Uzun'' tanimi: test + consequent + alternate toplam karakter sayisi $\geq 80$:

\begin{equation}
|\text{code}(\text{test})| + |\text{code}(\text{cons})| + |\text{code}(\text{alt})| \geq 80
\end{equation}

\begin{lstlisting}[language=JavaScript]
// Oncesi (tek satir, okunamaz):
isAuthenticated ? renderDashboard(user, session) : redirectToLogin(returnUrl);
// Sonrasi:
if (isAuthenticated) {
  renderDashboard(user, session);
} else {
  redirectToLogin(returnUrl);
}
\end{lstlisting}


\subsection{Faz 8: Arrow $\to$ Named Function}
\label{sec:phase8-arrow}

\code{const foo = () => \{...\}} formundaki arrow function'lar named function declaration'a donusturulur:

\begin{lstlisting}[language=JavaScript]
// Oncesi:
const processData = (items) => { return items.map(x => x * 2); };
// Sonrasi:
function processData(items) { return items.map(x => x * 2); }
\end{lstlisting}

\textbf{Semantik guevenlik:} \code{this} kullanan arrow function'lar donusturulmez. Arrow function'larda \code{this} lexical scope'a baglidir; normal function'larda ise call-site'a baglidir. Bu semantik fark bozulmamalidir:

\begin{lstlisting}[language=JavaScript, caption={Arrow-to-function this kontrolu}]
const { code: bodyCode } = generate(arrow.body, { compact: true });
if (!bodyCode.includes("this.") && !bodyCode.includes("this[")) {
    // Guvenli: donustur
    const funcDecl = t.functionDeclaration(
      t.identifier(decl.id.name), arrow.params, arrow.body,
      arrow.generator || false, arrow.async || false);
    path.replaceWith(funcDecl);
}
\end{lstlisting}


\subsection{Faz 9: Smart Variable Renaming}
\label{sec:phase9-rename}

Bu faz, minified identifier'lari (1--2 karakter) anlamli isimlere donusturur. LLM \textbf{kullanmaz} --- tamamen heuristik kurallara dayanir. Babel \code{scope.rename()} API'si ile scope-safe renaming yapilir.

\subsubsection{Uec Kural Kategorisi}

\begin{enumerate}
\item \textbf{Module-based:} \code{require("fs")} $\to$ degisken \code{fs} olarak adlandirilir
\item \textbf{Usage-based:} Degisken uzerinde \code{.push()}, \code{.map()} cagriliyorsa $\to$ \code{items}
\item \textbf{Fallback:} Tek harf degiskenler icin varsayilan isimler (\code{i} $\to$ \code{idx}, \code{e} $\to$ \code{elem})
\end{enumerate}

\begin{table}[h]
\centering
\small
\begin{tabular}{p{3.5cm}p{4cm}p{4.5cm}}
\toprule
\textbf{\color{sectioncyan}Pattern} & \textbf{\color{sectioncyan}Tespit Yontemi} & \textbf{\color{sectioncyan}Yeni Isim} \\
\midrule
\code{require("fs")} & CallExpression analizi & \code{fs} \\
\code{require("path")} & CallExpression analizi & \code{path} \\
\code{require("express")} & CallExpression analizi & \code{express} \\
\code{catch(e)} & CatchClause param & \code{error} \\
\code{x.push()}, \code{x.map()} & Usage (Array API) & \code{items} \\
\code{x.trim()}, \code{x.split()} & Usage (String API) & \code{text} \\
\code{x.message}, \code{x.stack} & Property access & \code{err} \\
\code{x.then()}, \code{x.catch()} & Usage (Promise API) & \code{promise} \\
Express \code{(req, res, next)} & Paramettre pozisyonu & \code{request, response, next} \\
esbuild \code{(exports, module)} & Factory pattern & \code{exports, module} \\
\bottomrule
\end{tabular}
\caption{Phase 9 smart rename kural ornekleri.}
\label{tab:phase9-rules}
\end{table}

\subsubsection{Tek Harf Fallback Tablosu}

\begin{lstlisting}[language=JavaScript, caption={Tek harf fallback mapping}]
const defaults = {
  a: "arg", b: "buf", c: "ctx", d: "data",
  e: "elem", f: "fn", g: "gen", h: "handler",
  i: "idx", j: "jdx", k: "key", l: "len",
  m: "mod", n: "count", o: "opts", p: "param",
  q: "query", r: "res", s: "str", t: "tmp",
  u: "usr", v: "val", w: "writer",
};
\end{lstlisting}


\subsection{Faz 10: Semantic Renaming (Deep Usage-Based)}
\label{sec:phase10-semantic}

Phase 10, Phase 9'un basit pattern matching'ini asarak, Babel scope + binding analizi ile her minified identifier'in kullanim baglamini derinlemesine analiz eder. \textbf{12 kural motoru} ile calisir.

\subsubsection{Kullanim Bilgisi Toplama (\_collectUsage10)}

Her identifier icin asagidaki bilgiler toplanir:

\begin{lstlisting}[language=JavaScript, caption={Usage bilgi yapisi}]
{
  methodCalls: Set,       // x.push(), x.split() gibi
  propertyReads: Set,     // x.length, x.name gibi
  propertyWrites: Set,    // x.count = 5 gibi
  calledAs: boolean,      // x() seklinde cagriliyor mu
  comparedWith: string[], // x === "GET" gibi karsilastirmalar
  typeofChecks: string[], // typeof x === "function" gibi
  usedInForOf: boolean,   // for(const y of x) iteratorde mi
  destructured: boolean,  // const {a, b} = x seklinde mi
  indexAccessed: boolean,  // x[0] seklinde mi
  arithmeticOperand: boolean, // x + 1 gibi aritmetikte mi
  thrownAsError: boolean, // throw x seklinde mi
  usedInNew: boolean,     // new x() seklinde mi
}
\end{lstlisting}

\subsubsection{12 Kural Motoru}

\begin{table}[h]
\centering
\small
\begin{tabular}{clcp{4.5cm}}
\toprule
\textbf{\color{sectioncyan}\#} & \textbf{\color{sectioncyan}Kural} & \textbf{\color{sectioncyan}Conf.} & \textbf{\color{sectioncyan}Tespit Kriteri} \\
\midrule
1 & \code{require\_module} & 0.95 & \code{require("fs")} $\to$ \code{fs} \\
2 & \code{new\_error} & 0.90 & \code{new Error()} pattern \\
3 & \code{error\_props} & 0.90 & \code{.message} + \code{.stack} props \\
4 & \code{string\_api} & 0.85 & $\geq$2 string method \\
5 & \code{array\_api} & 0.85 & $\geq$2 array method \\
6 & \code{promise\_api} & 0.85 & \code{.then()} + \code{.catch()} \\
7 & \code{fs\_api} & 0.90 & \code{readFileSync}, \code{writeFile} \\
8 & \code{http\_req} & 0.80 & $\geq$2 of \code{body,params,query,headers} \\
9 & \code{dom\_api} & 0.85 & $\geq$2 DOM property \\
10 & \code{callback} & 0.80 & \code{calledAs} + \code{typeof === "function"} \\
11 & \code{for\_of\_item} & 0.80 & \code{for(const x of iterable)} \\
12 & \code{multi\_compare} & 0.70 & $\geq$3 string karsilastirmasi \\
\bottomrule
\end{tabular}
\caption{Phase 10 semantic renaming kural motoru.}
\label{tab:phase10-rules}
\end{table}

\subsubsection{Phase 10+: Enhanced Rename (cursor-enhanced-rename.mjs)}

\file{cursor-enhanced-rename.mjs}, Phase 10'un kapsamadigi 15 ek kural ile kalan minified identifier'lari isle\-r. \file{deep\_pipeline.py}'da Adim 2.5 olarak calistirilir:

\begin{table}[h]
\centering
\small
\begin{tabular}{clp{5.5cm}}
\toprule
\textbf{\color{sectioncyan}\#} & \textbf{\color{sectioncyan}Kural} & \textbf{\color{sectioncyan}Aciklama} \\
\midrule
11 & Prototype access & \code{x.prototype} $\to$ \code{cls} \\
12 & Binary ops dominant & Cok fazla \code{+}, \code{-}, \code{<<} $\to$ \code{offset} \\
13 & Purely called & Sadece cagriliyor $\to$ \code{handler} \\
14 & Event listener pos & Callback pozisyonu $\to$ \code{eventArg} \\
15 & Member assign & Atama kaynagi $\to$ \code{ref} \\
16 & Single property & Tek prop baskini $\to$ contextual \\
17 & Conditional return & \code{return x ? a : b} $\to$ \code{result} \\
18 & Switch/if chain & Karsilastirma zinciri $\to$ \code{kind} \\
19 & Namespace access & \code{X.C}, \code{X.Q} $\to$ \code{ns} \\
20 & Loop counter & \code{for} init $\to$ \code{idx} \\
21 & Await expression & \code{await x} $\to$ \code{awaited} \\
22 & Default param ctx & Varsayilan deger $\to$ \code{optParam} \\
23 & Clash dedup & Phase 9 artigi $\to$ re-score \\
24 & Length/size dominant & \code{.length} baskini $\to$ \code{collection} \\
25 & Bitwise dominant & \code{|}, \code{\&}, \code{\^{}} $\to$ \code{flags} \\
\bottomrule
\end{tabular}
\caption{Enhanced rename ek kurallari (Kural 11--25).}
\label{tab:enhanced-rules}
\end{table}


% ===================================================================
\section{Webpack Module Extraction}
\label{sec:webpack-unpack}
% ===================================================================

\file{smart-webpack-unpack.mjs}, webpack ve esbuild bundle'larini modullerine ayiran akilli bir module splitter'dir. 7 farkli bundle formatini destekler.

\subsection{Desteklenen Bundle Formatlari}

\begin{enumerate}
\item \textbf{esbuild CJS wrapper:} \code{var X = U((exports, module) => \{...\})}
\item \textbf{esbuild ESM interop:} \code{var Y = Q1(X)}
\item \textbf{esbuild re-export:} \code{lb(target, \{name: () => ref\})}
\item \textbf{Webpack IIFE object:} \code{(function(modules) \{...\})(\{0: function(e,t,n)\{...\}\})}
\item \textbf{Webpack IIFE array:} \code{[function(e,t,n)\{...\}, ...]}
\item \textbf{Webpack 5 named:} \code{\_\_webpack\_modules\_\_ = \{"./src/file.js": ...\}}
\item \textbf{Top-level IIFE:} \code{(function() \{...\})()} --- kapsayici IIFE'leri ayirma
\end{enumerate}

\subsection{\_\_webpack\_require\_\_ Mekanizmasi}

Webpack bundle'larinda moduller \code{\_\_webpack\_require\_\_(id)} ile birbirini cagrilir. Bu pattern AST'de tespit edilir:

\begin{lstlisting}[language=JavaScript, caption={Webpack module tespiti -- deobfuscate.mjs}]
if (callee.name === "__webpack_require__"
    && node.arguments.length >= 1) {
  const arg = node.arguments[0];
  if (arg.type === "NumericLiteral") moduleId = arg.value;
  else if (arg.type === "StringLiteral") moduleId = arg.value;
  if (moduleId !== null && !_webpackIdSet.has(moduleId)) {
    _webpackIdSet.add(moduleId);
    webpackModules.push(moduleId);
  }
}
\end{lstlisting}

\subsection{esbuild vs Webpack Format Farklari}

\begin{table}[h]
\centering
\small
\begin{tabular}{p{3.5cm}p{5cm}p{5cm}}
\toprule
\textbf{\color{sectioncyan}Ozellik} & \textbf{\color{sectioncyan}Webpack} & \textbf{\color{sectioncyan}esbuild} \\
\midrule
Modul sarmalama & IIFE + nesne/dizi & CJS factory wrapper \\
Modul ID & Sayisal (0, 1, 2...) & Degisken adi \\
Dependency & \code{\_\_webpack\_require\_\_} & Dogrudan degisken referansi \\
Entry point & Ilk cagri arg. & Top-level kod \\
Tree shaking & Plugin bazli & Varsayilan \\
\bottomrule
\end{tabular}
\caption{Webpack vs esbuild bundle farklari.}
\label{tab:webpack-vs-esbuild}
\end{table}


% ===================================================================
\section{Buyuk Dosya Isleme: Chunked Processing}
\label{sec:chunked}
% ===================================================================

200MB+ dosyalar icin standart AST parsing bellege sigmaz. \file{deep\_pipeline.py} otomatik olarak \textbf{chunked processing} moduna gecer.

\subsection{Esik Degeri ve Tetikleme}

\begin{equation}
\text{mode} = \begin{cases}
\text{normal} & \text{eger dosya} < 100\text{ MB} \\
\text{chunked} & \text{eger dosya} \geq 100\text{ MB}
\end{cases}
\end{equation}

\file{deep\_pipeline.py}'da \code{LARGE\_FILE\_THRESHOLD\_MB = 100} olarak tanimli.

\subsection{Stream-Parse Pipeline}

\begin{enumerate}
\item \textbf{stream-parse.mjs:} Dosyayi top-level block'lara ayirir (\code{--max-block-kb 512})
\item \textbf{Her block'u bagimsiz isle:} \code{deep-deobfuscate.mjs} her chunk uzerinde ayri calisir (120s timeout/block)
\item \textbf{Birlestir:} Tum chunk ciktilari tek dosyada birlestirilir
\item \textbf{Webpack unpack:} Birlestirilmis dosya uzerinde modul ayirma
\item \textbf{Ikinci pass:} Her modul uzerinde variable renaming
\end{enumerate}

\begin{equation}
T_{\text{chunked}} \approx T_{\text{split}} + n \cdot T_{\text{deob/chunk}} + T_{\text{merge}} + T_{\text{unpack}}
\end{equation}

burada $n$ chunk sayisi, $T_{\text{deob/chunk}}$ chunk basina deobfuscation suresi.

\subsection{Timeout Hesaplaması}

\file{deep\_pipeline.py} dosya boyutuna gore dinamik timeout hesaplar:

\begin{equation}
T_{\text{timeout}} = \max(300, \min(1800, \lceil S_{\text{MB}} \times 100 \rceil)) \text{ saniye}
\label{eq:timeout}
\end{equation}

burada $S_{\text{MB}}$ dosya boyutu megabayt cinsindendir. Minimum 5 dakika, maksimum 30 dakika.


% ===================================================================
\section{Control Flow Flattening (CFF) Deflattening}
\label{sec:cff-deflattening}
% ===================================================================

\file{cff\_deflattener.py}, CFF obfuscation'ini geri alan bir Python modul-udur. Hem C hem JavaScript uzerinde calisir.

\subsection{Tespit Algoritması}

CFF dispatcher pattern regex ile tespit edilir:

\begin{lstlisting}[language=Python, caption={CFF dispatcher tespiti}]
# JS CFF pattern
_WHILE_SWITCH_JS = re.compile(
    r"(?:while\s*\(\s*(?:!0|true|1)\s*\)|for\s*\(\s*;;\s*\))\s*\{"
    r"\s*switch\s*\(\s*(\w+)\s*\)\s*\{",
    re.IGNORECASE,
)
\end{lstlisting}

\subsection{State Transition Graph Cikarma}

Her \code{case} blogundan state gecisleri cikarilarak bir yoenluee graf (directed graph) olusturulur:

\begin{equation}
G = (V, E) \quad \text{burada } V = \{s_0, s_1, \ldots, s_n\}, \quad E = \{(s_i, s_j) : s_i \xrightarrow{\text{state}} s_j\}
\end{equation}

\begin{center}
\begin{tikzpicture}[
  state/.style={circle, draw=sectioncyan, fill=boxbg, minimum size=0.9cm, text=white, font=\small},
  arrow/.style={-{Stealth[length=2.5mm]}, thick, color=gray!70!white},
  every node/.style={font=\scriptsize},
]
  \node[state] (s0) at (0,0) {0};
  \node[state] (s3) at (2.5,0) {3};
  \node[state] (s1) at (5,0) {1};
  \node[state] (s5) at (7.5,0) {5};
  \node[state] (s2) at (5,-1.5) {2};

  \draw[arrow] (s0) -- (s3);
  \draw[arrow] (s3) -- (s1);
  \draw[arrow] (s1) -- (s5);
  \draw[arrow] (s3) -- (s2);

  \node[above=0.2cm, text=green!60!white] at (s0) {entry};
  \node[above=0.2cm, text=red!60!white] at (s5) {exit};
\end{tikzpicture}
\end{center}

\subsection{Topological Sort ile Yeniden Olusturma}

State transition graph'i uzerinde BFS-based topological sort uygulanarak orijinal kod sirasi elde edilir:

\begin{lstlisting}[language=Python, caption={Topological sort -- cff\_deflattener.py}]
def _topological_sort(self, graph, blocks):
    order, visited = [], set()
    entry = blocks[0].state_value
    queue = [entry]
    while queue:
        state = queue.pop(0)
        if state in visited: continue
        visited.add(state)
        order.append(state)
        for ns in graph.get(state, []):
            if ns not in visited:
                queue.append(ns)
    # Ziyaret edilmemis state'leri sona ekle
    for block in blocks:
        if block.state_value not in visited:
            order.append(block.state_value)
    return order
\end{lstlisting}

\subsection{Karmasiklik Analizi}

\begin{equation}
T_{\text{CFF}} = \bigO{|C| + |V| + |E|}
\end{equation}

burada $|C|$ case sayisi, $|V|$ state sayisi (=$|C|$), $|E|$ gecis sayisi. Pratikte $|E| \leq 2|V|$ (her case en fazla 2 gecis yapar --- true/false branch).


% ===================================================================
\section{Opaque Predicate Removal}
\label{sec:opaque-removal}
% ===================================================================

\file{opaque\_predicate.py}, bilinen opaque predicate pattern'lerini regex ile tespit edip dead code'u isaretler veya kaldirir.

\subsection{Tespit Stratejileri}

\begin{enumerate}
\item \textbf{Pattern matching:} 7 always-true + 4 always-false regex pattern (Tablo~\ref{tab:opaque-predicates})
\item \textbf{Constant folding:} Sabit deger kullanan branch'ler (Phase 3 ile entegre)
\item \textbf{Z3 SMT solver (opsiyonel):} Karmasik opaque predicate'lar icin symbolic verification
\end{enumerate}

\subsection{Z3 Entegrasyonu}

Eger z3-solver kurulu ise, pattern matching ile tespit edilemeyen karmasik predicate'lar Z3 ile dogrulanabilir:

\begin{equation}
\text{is\_opaque}(\varphi) = \begin{cases}
\text{true} & \text{eger } Z3.\text{prove}(\forall x: \varphi(x) = \text{true}) \\
\text{false} & \text{eger } Z3.\text{find\_counterexample}(\varphi) \neq \varnothing
\end{cases}
\end{equation}

\file{opaque\_predicate.py}'da Z3 lazy-load yapilir:

\begin{lstlisting}[language=Python, caption={Z3 lazy-load -- opaque\_predicate.py}]
if use_z3:
    try:
        import z3  # noqa: F401
        self._z3_available = True
    except ImportError:
        logger.info("z3-solver bulunamadi, "
                     "pattern-based tespit kullanilacak")
\end{lstlisting}

\subsection{Eliminasyon}

\begin{itemize}
\item \textbf{Always-true:} \code{if(OPAQUE\_TRUE)\{X\} else \{Y\}} $\to$ \code{X}
\item \textbf{Always-false:} \code{if(OPAQUE\_FALSE)\{X\} else \{Y\}} $\to$ \code{Y} (veya sil)
\end{itemize}


% ===================================================================
\section{String Decryption}
\label{sec:string-decrypt}
% ===================================================================

\file{string\_decryptor.py}, binary'lerdeki sifrelenmis string'leri cozen bir moduld-ur. JavaScript deobfuscation ile dogrudan ilgili olmasa da, Karadul'un genel deobfuscation yeteneginin parcasidir.

\subsection{Desteklenen Sifreleme Yoentemleri}

\begin{enumerate}
\item \textbf{XOR single-byte:} \code{buf[i] \^{}= 0xNN}
\item \textbf{XOR multi-byte:} \code{buf[i] \^{}= key[i \% keylen]}
\item \textbf{XOR rolling:} \code{prev = buf[i] \^{} prev}
\item \textbf{RC4:} S-box initialization + swap pattern
\item \textbf{Base64:} \code{atob("SGVsbG8=")}
\item \textbf{Stack string:} Char-by-char push (Ghidra decompile ciktisinda)
\end{enumerate}

\subsection{Stack String Reconstruction}

Ghidra decompile ciktisinda ardisik local degiskenlere tek karakter atamalari ``stack string'' olarak yeniden birlestirilir:

\begin{lstlisting}[language=Python, caption={Stack string ornegi (Ghidra decompile ciktisi)}]
local_28 = 0x48;  # 'H'
local_27 = 0x65;  # 'e'
local_26 = 0x6c;  # 'l'
local_25 = 0x6c;  # 'l'
local_24 = 0x6f;  # 'o'
# -> "Hello" (confidence: 0.85)
\end{lstlisting}

\subsection{XOR Brute-Force}

Kisa string'ler icin tum 256 tek-byte anahtar denenerek en yuksek ``printable'' skorlu sonuc secilir:

\begin{equation}
\text{score}(k) = \frac{|\{c \in \text{decrypt}(d, k) : c \in \text{printable}\}|}{|\text{decrypt}(d, k)|}
\end{equation}

$\text{score}(k) > 0.7$ olan ilk 5 sonuc dondurulur.


% ===================================================================
\section{Acorn Fallback}
\label{sec:acorn-fallback}
% ===================================================================

Babel parser 32MB+ dosyalarda veya ozel syntax pattern'lerinde basarisiz olabilir. \file{deobfuscate.mjs} bu durumda Acorn parser'a duser.

\subsection{Babel vs Acorn Karsilastirmasi}

\begin{table}[h]
\centering
\small
\begin{tabular}{lcc}
\toprule
\textbf{\color{sectioncyan}Ozellik} & \textbf{\color{sectioncyan}Babel} & \textbf{\color{sectioncyan}Acorn} \\
\midrule
AST formati & Babel (genisletilmis) & ESTree (standart) \\
Error recovery & \checkmark & Kismi \\
TypeScript desteegi & \checkmark (plugin) & $\times$ \\
JSX desteegi & \checkmark (plugin) & $\times$ \\
Performans (30MB) & $\sim$45s & $\sim$8s \\
Memory (30MB) & $\sim$4GB & $\sim$1.5GB \\
\bottomrule
\end{tabular}
\caption{Babel vs Acorn parser karsilastirmasi.}
\label{tab:babel-vs-acorn}
\end{table}

Acorn, \file{deobfuscate.mjs}'te lazy-load edilir ve sadece Babel basarisiz olursa yuklenir:

\begin{lstlisting}[language=JavaScript, caption={Acorn lazy-load}]
let acorn = null, acornWalk = null;
function loadAcorn() {
  try {
    acorn = _require("acorn");
    acornWalk = _require("acorn-walk");
    return true;
  } catch { return false; }
}
\end{lstlisting}


\begin{ozet}
JavaScript deobfuscation, 4 adimli standart zincir (beautify $\to$ synchrony $\to$ babel $\to$ webpack) ve 10 fazli derin transform pipeline'indan olusur. Phase 2--8, visitor-merge optimizasyonu ile tek AST traverse'da birlestirilmistir ($\sim$3--5x hiz kazanci). Graceful degradation her adimda basarisizliga tolerans saglar. Faz 9--10'daki LLM-siz semantik isimlendirme toplam 25+ heuristik kuralla etkileyici sonuclar ueretir. 200MB+ dosyalar icin stream-parse ile chunked processing otomatik devreye girer. CFF deflattening, opaque predicate removal ve string decryption ek moduller olarak pipeline'i guclendirir.
\end{ozet}


\clearpage

% ############################################################################
%
%          B O L U M   6 :   B A Y E S I A N   I S I M   F U Z Y O N U
%
% ############################################################################

% Bu bolum, Karadul'un en yenilikci algoritmasi olan Bayesian log-odds isim
% birlestirme sistemini detayli olarak turetir.
% Kaynak kod: karadul/reconstruction/name_merger.py


% ===================================================================
\section{Problem Tanimi: Coklu Kaynak, Tek Isim}
\label{sec:naming-problem-detail}
% ===================================================================

Tersine muhendislikte, bir sembolun (fonksiyon, degisken) gercek ismi bilinmez. Birden fazla \textbf{bagimsiz} kaynak, farkli isimler onerir. Soru sudur: bu onerilerden hangisi dogru, ve nasil birlestiririz?

\begin{tanim}[Isim Fuzyon Problemi]
Verilen:
\begin{itemize}
\item Bir sembol $s$ (orijinal ismi bilinmiyor, orn. \code{FUN\_00401000})
\item $K$ adet kaynak, her biri bir isim onerisi ve confidence degeri ureten: $(n_i, c_i, \text{source}_i)$ icin $i = 1, \ldots, K$
\end{itemize}

Hedef: $s$ icin en uygun ismi $n^*$ sec ve sonucun dogruluk olasililigini $P^*$ hesapla.
\end{tanim}

\begin{example}
\code{FUN\_00401000} icin:
\begin{itemize}
\item \textbf{signature\_db:} \code{SSL\_connect} (confidence: 0.85, $w=1.0$)
\item \textbf{string\_intel:} \code{ssl\_handshake} (confidence: 0.70, $w=0.8$)
\item \textbf{c\_namer:} \code{connect\_tls} (confidence: 0.60, $w=1.0$)
\item \textbf{llm4decompile:} \code{ssl\_init\_connection} (confidence: 0.55, $w=0.5$)
\end{itemize}
Dort farkli isim! Ama hepsi ``SSL baglantisi'' kavramina isaret ediyor. Bayesian fusion bunu nasil cozer?
\end{example}


% ===================================================================
\section{Bayes Teoremi: Temelden Turetme}
\label{sec:bayes-derivation}
% ===================================================================

\subsection{Kosula Bagli Olasilik}

Iki olay $A$ ve $B$ icin:

\begin{equation}
\prob{A \cap B} = \prob{A \given B} \cdot \prob{B} = \prob{B \given A} \cdot \prob{A}
\label{eq:joint-prob}
\end{equation}

Buradan Bayes teoremi turetilir:

\begin{equation}
\boxed{
\prob{H \given E} = \frac{\prob{E \given H} \cdot \prob{H}}{\prob{E}}
}
\label{eq:bayes-theorem}
\end{equation}

\subsection{Isimlendirme Baglaminda Anlami}

\begin{itemize}
\item $H$: ``Bu sembolun gercek ismi $X$'tir'' hipotezi
\item $E_i$: ``Kaynak $i$, $X$ ismini $c_i$ confidence ile onerdi'' kaniti
\item $\prob{H}$: \textbf{Prior} --- hipotez hakkinda onceki inancimiz (baslangicta $0.5$)
\item $\prob{E_i \given H}$: \textbf{Likelihood} --- eger isim gercekten $X$ ise, kaynak $i$'nin bunu $c_i$ ile onermesinin olasiligi. Bunu $\approx c_i$ olarak modelliyoruz.
\item $\prob{E_i \given \neg H}$: Eger isim $X$ \textbf{degilse}, kaynak $i$'nin yine de $X$ onerme olasiligi $\approx 1 - c_i$.
\item $\prob{H \given E_i}$: \textbf{Posterior} --- kanit sonrasi guncellenminsi inanc
\end{itemize}

\subsection{Ardisik Kanit Birlestirme}

Birden fazla kaniti ardisik olarak birlestirmek icin, her kanitin \textbf{likelihood ratio} (LR) degerini kullaniriz:

\begin{equation}
\text{LR}_i = \frac{\prob{E_i \given H}}{\prob{E_i \given \neg H}} = \frac{c_i}{1 - c_i}
\label{eq:lr}
\end{equation}

Ardisik Bayes guncellemesi:

\begin{equation}
\frac{\prob{H \given E_1, E_2, \ldots, E_n}}{\prob{\neg H \given E_1, E_2, \ldots, E_n}} = \frac{\prob{H}}{\prob{\neg H}} \cdot \prod_{i=1}^{n} \text{LR}_i
\end{equation}

\textbf{Eger kaynaklar bagimsiz ise} (naive Bayes varsayimi).


% ===================================================================
\section{Log-Odds Formu: Neden ve Nasil}
\label{sec:logodds-derivation}
% ===================================================================

\subsection{Odds ve Log-Odds}

Odds (bahis orani):

\begin{equation}
\text{odds}(H) = \frac{\prob{H}}{1 - \prob{H}}
\end{equation}

Log-odds (logit):

\begin{equation}
\text{logit}(H) = \log\frac{\prob{H}}{1 - \prob{H}}
\end{equation}

\subsection{Turetme}

Ardisik Bayes guncellemesinin her iki tarafinin logaritmasini alalim:

\begin{align}
\log\frac{\prob{H \given \mathbf{E}}}{\prob{\neg H \given \mathbf{E}}} &= \log\frac{\prob{H}}{\prob{\neg H}} + \sum_{i=1}^{n} \log \text{LR}_i \\
&= \log\frac{\prob{H}}{1 - \prob{H}} + \sum_{i=1}^{n} \log\frac{c_i}{1 - c_i}
\end{align}

\subsection{Korelasyon Duzeltmesi}

Naive Bayes, kaynakların bagimsiz oldugunu varsayar. Ancak bazilari korelasyonludur (ornegin \code{llm4decompile} tum diger kaynaklari dolayli olarak kullanir). Korelasyonlu kaynaklarin etkisini azaltmak icin agirlik $w_i \in (0, 1]$ eklenir:

\begin{equation}
\boxed{
\text{log-odds}(H) = \underbrace{\log\frac{p_0}{1 - p_0}}_{\text{prior}} + \sum_{i=1}^{n} w_i \cdot \underbrace{\log\frac{c_i}{1 - c_i}}_{\text{log-LR}_i}
}
\label{eq:logodds-weighted}
\end{equation}

burada $p_0 = 0.5$ prior (notr baslangic: $\log(0.5/0.5) = 0$).

\subsection{Neden Log Space?}

\begin{enumerate}
\item \textbf{Numerik stabilite:} Carpim yerine toplam. $\prod_i c_i$ ile $n=10$ kaynak icin floating point underflow olusabilir. Toplam bu sorunu ortadan kaldirir.
\item \textbf{Korelasyon duzeltmesi:} $w_i$ ile dogrudan carpma log space'te lineer agirliklamaya denk gelir. Olasilik uzayinda bu $\text{LR}_i^{w_i}$ demektir --- bilgi-kuramsal olarak ``kaynagin bilgi katkisini $w_i$ oranina indir'' anlamina gelir.
\item \textbf{Sigmoid ile donusum:} Sonuc dogal olarak $[0,1]$ araligina doner.
\end{enumerate}


% ===================================================================
\section{Sigmoid Fonksiyonu: Logit'in Tersi}
\label{sec:sigmoid}
% ===================================================================

Log-odds'u tekrar olasiliga donusturmek icin sigmoid ($\sigma$) fonksiyonunu kullaniriz:

\begin{equation}
\prob{H \given \mathbf{E}} = \sigmoid(\text{log-odds}) = \frac{1}{1 + e^{-\text{log-odds}}}
\label{eq:sigmoid-def}
\end{equation}

\subsection{Sigmoid'in Ozellikleri}

\begin{enumerate}
\item $\sigma(0) = 0.5$ (notr: log-odds = 0 ise emin degiliz)
\item $\lim_{x \to \infty} \sigma(x) = 1$ (cok guclu kanit $\to$ kesinlik)
\item $\lim_{x \to -\infty} \sigma(x) = 0$ (cok guclu karsi kanit $\to$ red)
\item $\sigma'(x) = \sigma(x)(1 - \sigma(x))$ (turevi kendisindendir)
\item $\sigma(-x) = 1 - \sigma(x)$ (simetri)
\end{enumerate}

\begin{center}
\begin{tikzpicture}
\begin{axis}[
  width=10cm, height=5cm,
  domain=-6:6,
  samples=100,
  axis lines=middle,
  xlabel={log-odds},
  ylabel={$\sigma(\text{log-odds})$},
  every axis x label/.style={at=(current axis.right of origin), anchor=west},
  every axis y label/.style={at=(current axis.above origin), anchor=south},
  xtick={-4,-2,0,2,4},
  ytick={0,0.25,0.5,0.75,1.0},
  ymin=-0.05, ymax=1.05,
  grid=major,
  grid style={gray!20!black},
  axis line style={gray!60!white},
  tick label style={color=gray!60!white, font=\scriptsize},
  label style={color=sectioncyan, font=\small},
]
  \addplot[cyan, thick] {1/(1+exp(-x))};
  \addplot[red!60!white, dashed, thick] coordinates {(0,0) (0,0.5)};
  \addplot[red!60!white, dashed, thick] coordinates {(-6,0.5) (0,0.5)};
  \node[color=yellow!70!white, font=\scriptsize] at (axis cs:3, 0.3) {$\sigma(x) = \frac{1}{1+e^{-x}}$};
\end{axis}
\end{tikzpicture}
\end{center}

\subsection{Overflow Korumasi}

\file{name\_merger.py}'da sigmoid icin overflow korumasi uygulanir:

\begin{lstlisting}[language=Python, caption={Sigmoid overflow korumasi -- name\_merger.py}]
if log_odds > 20.0:
    p = 1.0       # exp(-20) ~ 2e-9, 1.0'a cok yakin
elif log_odds < -20.0:
    p = 0.0       # exp(20) ~ 5e8, 0.0'a cok yakin
else:
    p = 1.0 / (1.0 + math.exp(-log_odds))
\end{lstlisting}

Ayrica sonuc $[c_{\min}, c_{\max}] = [0.01, 0.99]$ araligina kliplenir --- asla ``\%100 emin'' veya ``\%0 emin'' demiyoruz.


% ===================================================================
\section{7 Kaynak Detayli Aciklama}
\label{sec:sources}
% ===================================================================

\subsection{binary\_extractor ($w = 1.0$)}

Debug string'ler, RTTI (Run-Time Type Information) ve build path bilgileri derleyici tarafindan binary'ye gomulur. Bu bilgi \textbf{ground truth}'a en yakin kaynaktir:

\begin{itemize}
\item Debug string'ler: \code{\_\_func\_\_} makrosu, assert mesajlari
\item RTTI: C++ sinif isimleri (\code{typeinfo name})
\item Build path: \code{/Users/dev/project/src/ssl\_connect.c}
\end{itemize}

$w = 1.0$ cunku bu bilgi derleyiciden dogrudan gelir --- hicbir tahmin yoktur.

\subsection{c\_namer ($w = 1.0$)}

6 katmanli heuristik isimlendirme motoru:

\begin{enumerate}
\item Fonksiyon imzasi analizi (parametre tipleri, donus tipi)
\item String kullanimi (hata mesajlari, log string'leri)
\item Cagrilan fonksiyonlar (API pattern tespiti)
\item Veri akisi analizi (tanimlama-kullanim zincirleri)
\item Algoritma tespiti (hash, sort, compress pattern'leri)
\item Statistiksel isim onerisi
\end{enumerate}

$w = 1.0$ cunku 6 katman zaten kendi icinde filtrelenmis; cift zayiflatma gereksiz.

\subsection{string\_intel ($w = 0.8$)}

Error mesajlari, protocol handler string'leri ve sabit mesajlardan fonksiyon ismi cikarimi:

\begin{lstlisting}[language=Python]
# Ornek: error mesajindan isim cikarimi
"Failed to connect SSL socket" -> ssl_connect (conf: 0.70)
\end{lstlisting}

$w = 0.8$ cunku error mesajlari, debug string'ler ile kismen ortak kaynaklardan beslenir.

\subsection{signature\_db ($w = 1.0$)}

FLIRT (Fast Library Identification and Recognition Technology) ile 1.5M+ kutuphane fonksiyon imzasi eslestirilir. Bu tamamen \textbf{byte-level} bir eslestirmedir:

\begin{equation}
\text{match}(f) = \argmax_{s \in \text{FLIRT\_DB}} \text{similarity}(\text{bytes}(f), \text{bytes}(s))
\end{equation}

$w = 1.0$ cunku byte pattern tamamen bagimsiz bir bilgi kaynagi.

\subsection{swift\_demangle ($w = 1.0$)}

Swift binary'lerinde mangled sembol isimleri (ornegin \code{\_\$s4main7MyClassC9processSSyS2SF}) demangle edilerek orijinal isim elde edilir:

\begin{equation}
\text{demangle}(\code{\_\$s4main7MyClassC9processSSyS2SF}) = \code{main.MyClass.process(String) -> String}
\end{equation}

$w = 1.0$ cunku demangle \textbf{kesin} bilgidir, tahmin degil.

\subsection{source\_matcher ($w = 0.85$)}

NPM paket fingerprinting: Obfuscated JS kodundaki string ve yapi pattern'leri bilinen NPM paketleri ile eslestirilir:

\begin{itemize}
\item Benzersiz string literal'ler (hata mesajlari, API endpoint'leri)
\item Fonksiyon imza pattern'leri (parametre sayisi + kullanim)
\item Export yapisi (modul cikti formati)
\end{itemize}

$w = 0.85$ cunku heuristik bazli ama yuksek kesinlikli.

\subsection{llm4decompile ($w = 0.5$)}

Makine ogrenmesi modeli ile isim tahmini. Neden $w = 0.5$?

\begin{uyari}[LLM Korelasyon Problemi]
LLM modeli, egitim verisinde tum diger kaynaklarin bilgisini dolayli olarak iciyor:
\begin{itemize}
\item Kutuphane imzalari (signature\_db ile korelasyonlu)
\item Debug string pattern'leri (binary\_extractor ile korelasyonlu)
\item Error mesaj kaliplari (string\_intel ile korelasyonlu)
\end{itemize}
Bu nedenle LLM'in bilgi katkisi ``yarisina'' indirilir: $w = 0.5$.
\end{uyari}

\subsection{byte\_pattern ($w = 1.0$)}

Kutuphane fonksiyonlarinin byte-level fuzzy matching ile tespiti. FLIRT'ten farki: FLIRT kesin imza kullanirken, byte\_pattern fuzzy threshold ile calisir. Yine de tamamen bagimsiz bir byte-level kaynak.


% ===================================================================
\section{Korelasyon-Aware Agirliklar: Bilgi Kurami}
\label{sec:correlation-theory}
% ===================================================================

\subsection{Neden $w_i < 1$?}

Iki kaynak $S_i$ ve $S_j$ arasindaki korelasyonu mutual information ile olceriz:

\begin{equation}
I(S_i; S_j) = \sum_{s_i, s_j} p(s_i, s_j) \log\frac{p(s_i, s_j)}{p(s_i) \cdot p(s_j)}
\label{eq:mutual-info}
\end{equation}

Eger $I(S_i; S_j) > 0$ ise, iki kaynak ortak bilgi tasir. Her ikisini de $w=1.0$ ile saymak bu ortak bilgiyi \textbf{iki kez} sayar ve \textbf{overconfidence} olusturur.

Korelasyon duzeltme prensibimiz:

\begin{equation}
w_i = 1 - \frac{I(S_i; S_{\text{others}})}{H(S_i)}
\label{eq:weight-from-mi}
\end{equation}

burada $H(S_i)$ kaynak $i$'nin entropisi. Pratikte Denklem~\ref{eq:weight-from-mi} yerine ampirik agirliklar kullanilir (Tablo~\ref{tab:weights-detail}).

\subsection{Kaynak Agirliklari Tablosu}

\begin{table}[h]
\centering
\small
\begin{tabular}{llp{6.5cm}}
\toprule
\textbf{\color{sectioncyan}Kaynak} & \textbf{\color{sectioncyan}$w_i$} & \textbf{\color{sectioncyan}Gerekce} \\
\midrule
binary\_extractor & 1.0 & Derleyici ciktisi, tamamen bagimsiz \\
c\_namer & 1.0 & 6 katman kendi icinde filtrelenmis \\
string\_intel & 0.8 & Error msg.~binary\_extractor ile kismi korelasyon \\
signature\_db & 1.0 & FLIRT byte pattern, tamamen bagimsiz \\
swift\_demangle & 1.0 & Deterministik demangle, kesin bilgi \\
source\_matcher & 0.85 & NPM heuristik, bagimsiz ama kesin degil \\
llm4decompile & 0.5 & Tum kaynaklari dolayli kullaniyor \\
byte\_pattern & 1.0 & Byte-level fuzzy match, bagimsiz \\
\midrule
\textit{Bilinmeyen kaynak} & 0.7 & Varsayilan fallback \\
\bottomrule
\end{tabular}
\caption{Kaynak agirliklari, korelasyon analizi ve gerekceler.}
\label{tab:weights-detail}
\end{table}

\subsection{Overconfidence Problemi: Somut Ornek}

Naive Bayes ($w_i = 1.0$ her kaynak icin) ile korelasyon-aware karsilastirma:

\begin{align}
\text{Naive } (w_i = 1): \quad \text{log-odds} &= 0 + 1.0 \cdot \log\frac{0.85}{0.15} + 1.0 \cdot \log\frac{0.70}{0.30} + 1.0 \cdot \log\frac{0.60}{0.40} \\
&= 1.735 + 0.847 + 0.405 = 2.987 \\
P_{\text{naive}} &= \sigma(2.987) = 0.952 \quad (\text{\textbf{\%95.2}})
\end{align}

\begin{align}
\text{Corrected}: \quad \text{log-odds} &= 0 + 1.0 \cdot 1.735 + 0.8 \cdot 0.847 + 0.5 \cdot 0.405 \\
&= 1.735 + 0.678 + 0.203 = 2.616 \\
P_{\text{corr}} &= \sigma(2.616) = 0.932 \quad (\text{\textbf{\%93.2}})
\end{align}

\textbf{Fark:} Naive Bayes \%95.2 derken, korelasyon duzeltmesi \%93.2 der --- \%2 daha muhafazakar ve \textbf{daha dogru}. Cunku llm4decompile'in verdigi bilginin bir kismi zaten diger kaynaklarda mevcuttu.


% ===================================================================
\section{Gruplama Stratejisi: 5 Seviye}
\label{sec:grouping-detail}
% ===================================================================

\file{name\_merger.py}'deki \code{\_merge\_candidates()} fonksiyonu, adaylari 5 seviyeli bir hiyerarsiden gecirir. Her seviye bir oncekinden daha gevşek bir eslestirme kullanir.

\begin{center}
\begin{tikzpicture}[
  level/.style={rectangle, draw=sectioncyan, fill=boxbg, minimum width=5cm, minimum height=0.9cm, text=white, font=\small, rounded corners=2pt, align=center},
  arrow/.style={-{Stealth[length=2.5mm]}, thick, color=gray!60!white},
  fail/.style={-{Stealth[length=2mm]}, dashed, color=red!50!white},
  node distance=0.7cm,
]
  \node[level] (l1) {Seviye 1: Exact Multi-Match};
  \node[level, below=of l1] (l2) {Seviye 2: Partial Match};
  \node[level, below=of l2] (l3) {Seviye 3: Semantic Match};
  \node[level, below=of l3] (l4) {Seviye 4: Voting ($\geq$3)};
  \node[level, below=of l4] (l5) {Seviye 5: Single Source};

  \draw[arrow] (l1) -- node[right, font=\tiny, text=red!50!white] {eslesmezse} (l2);
  \draw[arrow] (l2) -- node[right, font=\tiny, text=red!50!white] {eslesmezse} (l3);
  \draw[arrow] (l3) -- node[right, font=\tiny, text=red!50!white] {eslesmezse} (l4);
  \draw[arrow] (l4) -- node[right, font=\tiny, text=red!50!white] {eslesmezse} (l5);

  \node[right=1.5cm, font=\scriptsize, text=green!60!white] at (l1.east) {En yuksek guven};
  \node[right=1.5cm, font=\scriptsize, text=yellow!70!white] at (l3.east) {Orta guven};
  \node[right=1.5cm, font=\scriptsize, text=red!60!white] at (l5.east) {En dusuk guven};
\end{tikzpicture}
\end{center}

\subsection{Seviye 1: Exact Multi-Match}

Ayni normalize isim $\geq 2$ \textbf{farkli} kaynaktan geldiyse, dogrudan Bayesian merge:

\begin{lstlisting}[language=Python, caption={Exact multi-match -- name\_merger.py}]
for norm_name, group in name_groups.items():
    if len(group) >= 2:
        sources = set(c.source for c in group)
        if len(sources) >= 2:
            conf = bayesian_merge(
                [c.confidence for c in group],
                [c.source for c in group], self._cfg)
            return MergedName(..., merge_method="bayesian_multi")
\end{lstlisting}

\subsection{Seviye 2: Partial Match}

Farkli isimler arasinda prefix/suffix varyanti aranir. Ornegin:
\begin{itemize}
\item \code{cycle\_sizes\_default} vs \code{cycle\_size} $\to$ partial match
\item \code{active\_event\_monitor} vs \code{event\_monitor} $\to$ partial match
\end{itemize}

Overlap skoru:

\begin{equation}
\text{overlap}(A, B) = \begin{cases}
\frac{|A_{\text{short}}|}{|A_{\text{long}}|} & \text{eger } A_{\text{short}} \subset A_{\text{long}} \text{ (substring)} \\[6pt]
J(W_A, W_B) = \frac{|W_A \cap W_B|}{|W_A \cup W_B|} & \text{aksi halde (Jaccard)}
\end{cases}
\label{eq:overlap}
\end{equation}

$\text{overlap} \geq 0.5$ ise partial match kabul edilir. Penalty:

\begin{equation}
\text{penalty} = 0.7 + 0.3 \times \text{overlap}
\end{equation}

Tam substring ($\text{overlap} = 1.0$) $\to$ penalty $= 1.0$ (penalty yok). Kismi overlap ($\text{overlap} = 0.5$) $\to$ penalty $= 0.85$.


\subsection{Seviye 3: Semantic Match}

\file{name\_merger.py}'de 20 adet \code{frozenset} ile kelime gruplari tanimlidir. Iki isim, ayni gruba dusen kelimeleri paylasiyorsa ``semantik olarak benzer'' sayilir.

\begin{table}[h]
\centering
\small
\begin{tabular}{p{3cm}p{9cm}}
\toprule
\textbf{\color{sectioncyan}Grup Adi} & \textbf{\color{sectioncyan}Kelimeler} \\
\midrule
Gonderme & send, transmit, write, emit, dispatch, post \\
Alma & recv, receive, read, get, fetch \\
Olusturma & init, initialize, setup, create, construct, new, alloc \\
Yok etme & destroy, cleanup, teardown, free, delete, release, close, dispose \\
Parse/Decode & parse, decode, deserialize, unmarshal, extract \\
Encode & serialize, encode, marshal, format, stringify \\
Isleyici & handle, process, on, dispatch, callback, handler \\
Dogrulama & check, validate, verify, test, assert, is, has \\
Koleksiyon ekle & add, insert, append, push, enqueue, put \\
Koleksiyon sil & remove, delete, erase, pop, dequeue, take \\
Boyut & size, length, count, num, len, total \\
Tampon & buf, buffer, data, bytes, payload, blob \\
Mesaj & msg, message, packet, frame, request, response \\
Hata & err, error, fault, failure, exception \\
Baglam & ctx, context, state, env, session \\
Konfiguerasyon & cfg, config, configuration, settings, options, prefs \\
Baglanti & conn, connection, socket, link, channel \\
Loglama & log, print, trace, debug, info, warn \\
Kilit alma & lock, acquire, enter, begin, grab \\
Kilit birakma & unlock, release, leave, end, drop \\
\bottomrule
\end{tabular}
\caption{Semantik kelime gruplari (20 frozenset).}
\label{tab:semantic-groups}
\end{table}

\subsubsection{Expand-Then-Intersect Algoritmasi}

\begin{enumerate}
\item Her isim kelimelere ayrilir: \code{ssl\_handshake} $\to$ \{ssl, handshake\}
\item Her kelime, ait oldugu semantic grubun \textbf{tum kelimeleri} ile genisletilir:
      \code{handshake} $\notin$ herhangi bir grup; \code{connect} $\to$ \{conn, connection, socket, link, channel\}
\item Genisletilmis kumelerin kesisimi kontrol edilir
\item Kesisim orani $\geq 0.3$ ise semantik benzerlik var
\end{enumerate}

\begin{lstlisting}[language=Python, caption={Semantic similarity -- name\_merger.py}]
def _is_semantically_similar(self, name1, name2):
    words1 = set(self._normalize(name1).split("_"))
    words2 = set(self._normalize(name2).split("_"))
    # Dogrudan kelime kesisimi
    common = words1 & words2
    if common and len(common) / max(len(words1), len(words2)) >= 0.5:
        return True
    # Semantic group genisletme
    expanded1 = set()
    for w in words1:
        if w in self._word_to_group:
            expanded1 |= self._word_to_group[w]
        expanded1.add(w)
    # ... ayni expanded2 icin ...
    overlap = expanded1 & expanded2
    union = expanded1 | expanded2
    return len(overlap) / len(union) >= 0.3
\end{lstlisting}

Semantik match'te confidence \textbf{0.85x penalty} ile carpilir cunku isimler tam eslesmiyor:

\begin{equation}
c_i^{\text{semantic}} = c_i \times 0.85
\end{equation}


\subsection{Seviye 4: Voting}

$\geq 3$ kaynak ayni normalize ismi veriyorsa, cogunluk kazanir:

\begin{equation}
n^* = \argmax_{n} \sum_{i=1}^{K} \mathbb{1}[\text{normalize}(n_i) = n]
\end{equation}

Eger $\text{count}(n^*) \geq 3$ ise, bu adaylar Bayesian merge'e girer.

\subsection{Seviye 5: Single Source}

Tek kaynak kaldiysa, Bayesian merge gerekmez. Ancak kaynagin $w_i$ deger ile duzeltme yapilir:

\begin{equation}
c_{\text{adjusted}} = \sigma\left(w_i \cdot \log\frac{c_i}{1 - c_i}\right)
\end{equation}

$w_i < 1$ ise bu ``tek kaynaga fazla guvenme'' demektir.


% ===================================================================
\section{Normalize Isim Karsilastirmasi}
\label{sec:normalize}
% ===================================================================

Farkli kaynaklar farkli konvansiyonlarda isim uretir. Karsilastirma oncesi normalizasyon sart:

\begin{lstlisting}[language=Python, caption={Isim normalizasyonu -- name\_merger.py}]
def _normalize(self, name):
    # camelCase -> snake_case
    name = re.sub(r"([a-z])([A-Z])", r"\1_\2", name)
    # Prefix'leri kaldir (m_, s_, g_, p_)
    name = re.sub(r"^[msgp]_", "", name)
    return name.lower().strip("_")
\end{lstlisting}

Ornekler:
\begin{itemize}
\item \code{SSLConnect} $\to$ \code{ssl\_connect}
\item \code{m\_sslConnect} $\to$ \code{ssl\_connect}
\item \code{ssl\_connect} $\to$ \code{ssl\_connect} (degismez)
\item \code{g\_SSLHandler} $\to$ \code{ssl\_handler}
\end{itemize}


% ===================================================================
\section{UNK Threshold: Precision vs Recall}
\label{sec:unk-threshold}
% ===================================================================

\subsection{NSA-Grade Felsefe}

\begin{uyari}[Yanlis Isim Tehlikesi]
Tersine muhendislikte yanlis isim, isimsizden \textbf{cok daha kotu}dur:
\begin{itemize}
\item Yanlis isim $\to$ analist yanlis yolda saatler/gunler harcar
\item UNK (isimsiz) $\to$ analist ``bilmiyorum'' der, dikkatli devam eder
\end{itemize}

\textbf{Sonuc:} Dusuk confidence'li isimleri \textbf{atamayiz}. $c < 0.30$ ise UNK.
\end{uyari}

\subsection{ROC Analizi}

UNK threshold secimi, precision ve recall arasinda trade-off'tur:

\begin{equation}
\text{Precision} = \frac{TP}{TP + FP}, \quad \text{Recall} = \frac{TP}{TP + FN}
\end{equation}

\begin{itemize}
\item \textbf{Threshold $\uparrow$:} Precision $\uparrow$ (daha az yanlis isim), Recall $\downarrow$ (daha cok UNK)
\item \textbf{Threshold $\downarrow$:} Precision $\downarrow$ (daha cok yanlis isim), Recall $\uparrow$ (daha az UNK)
\end{itemize}

Karadul'un secimi: $\theta_{\text{UNK}} = 0.30$. Bu, \textbf{Type I error maliyetinin Type II error maliyetinden cok daha yuksek} olmasi varsayimindan gelir:

\begin{equation}
\text{Cost}_{\text{wrong\_name}} \gg \text{Cost}_{\text{UNK}} \implies \theta_{\text{UNK}} \text{ yuksek tutulur}
\end{equation}


% ===================================================================
\section{Somut Hesaplama Ornekleri}
\label{sec:bayesian-examples}
% ===================================================================

\subsection{Senaryo 1: Uc Kaynak Ayni Isim (Exact Multi-Match)}

\code{FUN\_00401000} icin uc kaynak ayni normalize ismi veriyor:

\begin{center}
\begin{tabular}{llcc}
\toprule
Kaynak & Onerilen Isim & $c_i$ & $w_i$ \\
\midrule
signature\_db & \code{SSL\_connect} & 0.85 & 1.0 \\
string\_intel & \code{SSL\_connect} & 0.70 & 0.8 \\
c\_namer & \code{ssl\_connect} & 0.60 & 1.0 \\
\bottomrule
\end{tabular}
\end{center}

Normalize edilince uc isim de \code{ssl\_connect} --- exact multi-match!

\textbf{Adim 1: Prior log-odds}
\begin{equation}
\ell_0 = \log\frac{0.5}{0.5} = 0
\end{equation}

\textbf{Adim 2: Her kaynağin log-LR katkisi}
\begin{align}
\Delta\ell_1 &= w_1 \cdot \log\frac{c_1}{1-c_1} = 1.0 \cdot \log\frac{0.85}{0.15} = 1.0 \cdot 1.7346 = 1.7346 \\[4pt]
\Delta\ell_2 &= w_2 \cdot \log\frac{c_2}{1-c_2} = 0.8 \cdot \log\frac{0.70}{0.30} = 0.8 \cdot 0.8473 = 0.6779 \\[4pt]
\Delta\ell_3 &= w_3 \cdot \log\frac{c_3}{1-c_3} = 1.0 \cdot \log\frac{0.60}{0.40} = 1.0 \cdot 0.4055 = 0.4055
\end{align}

\textbf{Adim 3: Toplam log-odds}
\begin{equation}
\ell = 0 + 1.7346 + 0.6779 + 0.4055 = 2.8180
\end{equation}

\textbf{Adim 4: Sigmoid}
\begin{equation}
P = \sigma(2.8180) = \frac{1}{1 + e^{-2.8180}} = \frac{1}{1 + 0.05975} = 0.9436
\end{equation}

\begin{basari}[Senaryo 1 Sonuc]
Uc kaynak birlestirildiginde \textbf{\%94.4 guvenle} \code{ssl\_connect} ismi atanir. En iyi tek kaynak \%85'ti --- Bayesian fusion \%9.4 puan ekledi.
\end{basari}


\subsection{Senaryo 2: Iki Kaynak Partial Match}

\code{FUN\_00501234} icin iki kaynak benzer ama farkli isimler veriyor:

\begin{center}
\begin{tabular}{llcc}
\toprule
Kaynak & Onerilen Isim & $c_i$ & $w_i$ \\
\midrule
binary\_extractor & \code{process\_message} & 0.80 & 1.0 \\
c\_namer & \code{handle\_message\_data} & 0.65 & 1.0 \\
\bottomrule
\end{tabular}
\end{center}

Normalize: \code{process\_message} vs \code{handle\_message\_data}

\textbf{Kelime kumesi:} $\{$process, message$\}$ vs $\{$handle, message, data$\}$. Jaccard:

\begin{equation}
J = \frac{|\{\text{message}\}|}{|\{\text{process, message, handle, data}\}|} = \frac{1}{4} = 0.25 < 0.5
\end{equation}

Dogrudan Jaccard $< 0.5$, ama ``handle'' ve ``process'' ayni semantic gruba dusur (Tablo~\ref{tab:semantic-groups}: ``handle, process, on, dispatch, callback, handler''). Semantic expansion sonrasi:

Expanded$_1$ = $\{$process, message, handle, on, dispatch, callback, handler$\}$ \\
Expanded$_2$ = $\{$handle, message, data, process, on, dispatch, callback, handler, bytes, payload, blob, buf, buffer$\}$

\begin{equation}
\frac{|\text{Expanded}_1 \cap \text{Expanded}_2|}{|\text{Expanded}_1 \cup \text{Expanded}_2|} = \frac{7}{13} = 0.538 \geq 0.3
\end{equation}

Semantic match! Penalty $= 0.85$:

\begin{align}
\Delta\ell_1 &= 1.0 \cdot \log\frac{0.80 \times 0.85}{1 - 0.80 \times 0.85} = 1.0 \cdot \log\frac{0.68}{0.32} = 0.7538 \\[4pt]
\Delta\ell_2 &= 1.0 \cdot \log\frac{0.65 \times 0.85}{1 - 0.65 \times 0.85} = 1.0 \cdot \log\frac{0.5525}{0.4475} = 0.2107 \\[4pt]
\ell &= 0 + 0.7538 + 0.2107 = 0.9645 \\[4pt]
P &= \sigma(0.9645) = \frac{1}{1 + e^{-0.9645}} = 0.724
\end{align}

\begin{basari}[Senaryo 2 Sonuc]
Semantic match ile iki kaynak birlestirildiginde \textbf{\%72.4 guvenle} \code{process\_message} (daha kisa isim) atanir. Penalty olmadan \%80 olacakti --- semantik belirsizlik \%7.6 dusurdu.
\end{basari}


\subsection{Senaryo 3: Tek Kaynak, Dusuk Confidence $\to$ UNK}

\code{FUN\_00601000} icin sadece bir kaynak var:

\begin{center}
\begin{tabular}{llcc}
\toprule
Kaynak & Onerilen Isim & $c_i$ & $w_i$ \\
\midrule
llm4decompile & \code{init\_module} & 0.45 & 0.5 \\
\bottomrule
\end{tabular}
\end{center}

Tek kaynak, Single Source (Seviye 5):

\begin{align}
\ell &= w_1 \cdot \log\frac{c_1}{1-c_1} = 0.5 \cdot \log\frac{0.45}{0.55} = 0.5 \cdot (-0.2007) = -0.1003 \\[4pt]
P &= \sigma(-0.1003) = \frac{1}{1 + e^{0.1003}} = 0.4749
\end{align}

$P = 0.475 < \theta_{\text{UNK}} = 0.30$? Hayir, $0.475 > 0.30$.

Ama bekleyin --- burada $P < 0.5$, yani model ``bu ismin dogru olma olasiligi \%50'den az'' diyor. Rasyonel bir analist buna guvenm-ez.

Esik degeri $\theta_{\text{UNK}} = 0.30$ ile $P = 0.475 > 0.30$ oldugu icin \textit{teknik olarak} isim atanir. Ancak pratikte bu dusuk confidence'li isim dikkatle ele alinir.

Simdi daha dusuk confidence ile deneyelim: $c_1 = 0.35$:

\begin{align}
\ell &= 0.5 \cdot \log\frac{0.35}{0.65} = 0.5 \cdot (-0.6190) = -0.3095 \\
P &= \sigma(-0.3095) = 0.4233
\end{align}

$P = 0.423 > 0.30$ --- hala UNK sinirinin uzerinde. Ancak $c_1 = 0.15$:

\begin{align}
\ell &= 0.5 \cdot \log\frac{0.15}{0.85} = 0.5 \cdot (-1.7346) = -0.8673 \\
P &= \sigma(-0.8673) = 0.296
\end{align}

$P = 0.296 < 0.30$ $\to$ \textbf{UNK!} Isim atanmaz.

\begin{uyari}[Senaryo 3 Sonuc]
LLM kaynaginin tek basina confidence $= 0.15$, $w = 0.5$ ile verdigi isim UNK olarak kalir. \textbf{Yanlis isim, isimsizden kotudur.}
\end{uyari}


% ===================================================================
\section{Hesaplama Karmasikligi}
\label{sec:complexity}
% ===================================================================

\subsection{Genel Karmasiklik}

$N$ adet sembol ve her sembol icin ortalama $K$ aday icin:

\begin{equation}
T_{\text{merge}} = \bigO{N \cdot K^2}
\end{equation}

Kare terimi, her aday ciftinin partial/semantic karsilastirmasinden gelir.

\subsection{Alt Islem Karmasikliklari}

\begin{table}[h]
\centering
\small
\begin{tabular}{lll}
\toprule
\textbf{\color{sectioncyan}Islem} & \textbf{\color{sectioncyan}Karmasiklik} & \textbf{\color{sectioncyan}Aciklama} \\
\midrule
Normalizasyon & $\bigO{|n|}$ & Regex + lowercase, string uzunlugu ile lineer \\
Exact match & $\bigO{K}$ & Hash-based gruplama \\
Partial overlap & $\bigO{K^2 \cdot |n|}$ & Tum ciftleri karsilastir, her biri $\bigO{|n|}$ \\
Semantic match & $\bigO{K^2 \cdot |W|}$ & Kelime genisletme + kesisim, $|W|$ = kelime sayisi \\
Bayesian merge & $\bigO{K}$ & Log-odds toplami (lineer) \\
\bottomrule
\end{tabular}
\caption{Alt islem karmasikliklari.}
\label{tab:complexity}
\end{table}

\subsection{Pratikte Performans}

Tipik degerler: $N \approx 5000$ (bir binary'deki fonksiyon sayisi), $K \approx 3$ (ortalama kaynak sayisi):

\begin{equation}
T \approx 5000 \times 3^2 = 45{,}000 \text{ islem}
\end{equation}

Python'da bu $< 0.1$s surer --- bottleneck tamamen decompilation ve statik analiz asamalarindadir.


% ===================================================================
\section{Overconfidence Problemi ve Cozumu}
\label{sec:overconfidence}
% ===================================================================

\subsection{Naive Bayes Hatasi}

Naive Bayes, tum kaynaklarin \textbf{kosula bagli bagimsiz} oldugunu varsayar:

\begin{equation}
\prob{E_1, E_2, \ldots, E_n \given H} = \prod_{i=1}^{n} \prob{E_i \given H}
\end{equation}

Bu varsayim genellikle \textbf{yanlistir}. Eger iki kaynak korelasyonlu ise, ayni bilgiyi iki kez saymis oluruz.

\subsection{Korelasyon ile Duzeltme}

Log-odds formulunde $w_i$ korelasyonu absorbe eder:

\begin{equation}
\text{LR}_i^{\text{effective}} = \text{LR}_i^{w_i}
\end{equation}

$w_i = 1$: Kaynak tamamen bagimsiz $\to$ tam bilgi katkisi. \\
$w_i = 0.5$: Kaynak yari korelasyonlu $\to$ bilgi katkisi ``yarisina'' iner. \\
$w_i = 0$: Kaynak tamamen bagimli $\to$ bilgi katkisi sifir (sayilmaz).

Bilgi-kuramsal olarak bu, log-likelihood ratio'nun eksponent ile olceklenmesidir:

\begin{equation}
\text{effective info gain}_i = w_i \cdot \text{KL}(\prob{\cdot \given E_i} \| \prob{\cdot})
\end{equation}

burada KL divergence bilgi kazancini olcer.

\subsection{Sonuc}

\begin{basari}[Bayesian Fusion Avantajlari]
\begin{enumerate}
\item Coklu kaynaklari sistematik olarak birlestirir --- el ile ``en iyi kaynagi sec'' yerine matematiksel optimal
\item Korelasyonu modelleyerek overconfidence'i onler
\item UNK threshold ile yanlis isim atama riskini minimize eder
\item Log-odds formunda numerik olarak stabil
\item Her adimi tamamen deterministik ve auditlenebilir
\end{enumerate}
\end{basari}


\begin{ozet}
Bayesian isim fuzyonu, 7+ bagimsiz kaynaktan gelen isimleri korelasyon-aware log-odds ile birlestirir. 5 seviyeli gruplama (exact $\to$ partial $\to$ semantic $\to$ voting $\to$ single) her durum icin uygun strateji secer. $w_i < 1$ agirliklari korelasyonlu kaynaklarin etkisini bilgi-kuramsal olarak azaltir. UNK threshold (\%30) NSA-grade felsefe ile yanlis isimlendirmeyi onler. Somut orneklerde: uc kaynak exact match ile \%85'ten \%94.4'e, iki kaynak semantic match ile \%72.4'e yukseltildi; tek dusuk-confidence kaynak ise UNK olarak dogru sekilde reddedildi. Karmasiklik $\bigO{N \cdot K^2}$ ile pratikte milisaniye mertebesinde calisir.
\end{ozet}
 ile dahil edilebilir.
% ONEMLI: \chapter tanimlanmaz, cunku chapter main.tex'te tanimlidir.
%         Burada \section ve \subsection seviyesinde yazilir.
% ============================================================================


% ############################################################################
%
%                    B O L U M   5 :   D E O B F U S C A T I O N
%
% ############################################################################

% Bu bolum, Karadul'un en yenilikci bilesenlerinden birini --- 9 fazli derin
% JavaScript deobfuscation pipeline'ini --- detayli olarak aciklar.
% Kaynak kodlar:
%   - karadul/deobfuscators/manager.py
%   - karadul/deobfuscators/deep_pipeline.py
%   - karadul/deobfuscators/babel_pipeline.py
%   - karadul/deobfuscators/synchrony_wrapper.py
%   - karadul/deobfuscators/string_decryptor.py
%   - karadul/deobfuscators/opaque_predicate.py
%   - karadul/deobfuscators/cff_deflattener.py
%   - scripts/deobfuscate.mjs
%   - scripts/deep-deobfuscate.mjs
%   - scripts/beautify.mjs
%   - scripts/smart-webpack-unpack.mjs
%   - scripts/cursor-enhanced-rename.mjs


% ===================================================================
\section{Obfuscation: Formal Tanim ve Teorik Temel}
\label{sec:obf-formal}
% ===================================================================

\subsection{Semantik Koruma Teoremi}

Kod obfuscation, bir programin kaynak kodunu, davranisini degistirmeden okunamazlastirma islemidir. Bunu formal olarak tanimlayalim.

\begin{tanim}[Obfuscation Fonksiyonu]
Bir obfuscation fonksiyonu $\mathcal{O}: \mathcal{P} \to \mathcal{P}$ asagidaki kosullari saglayan bir donustumdur:
\begin{enumerate}
\item \textbf{Semantik koruma:} Her girdi $x$ icin, $\mathcal{O}(P)(x) = P(x)$
\item \textbf{Polinom yavaslatma:} $|\mathcal{O}(P)| \leq \text{poly}(|P|)$ ve calisma suresi polinomial olarak artar
\item \textbf{Analiz zorlugu:} Herhangi bir analiz $A$ icin, $A(\mathcal{O}(P))$'den elde edilen bilgi, $A(P')$'den ($P'$ erisim saglanmamis bir program) elde edilenden anlamli olarak fazla degildir
\end{enumerate}
\end{tanim}

Bu tanimi daha kompakt yazarsak:

\begin{equation}
\mathcal{O}: \mathcal{P} \to \mathcal{P}, \quad \forall P \in \mathcal{P}: \quad \llbracket \mathcal{O}(P) \rrbracket = \llbracket P \rrbracket
\label{eq:semantic-preservation}
\end{equation}

burada $\llbracket \cdot \rrbracket$ programin denotational semantigini (yani giris-cikis davranisini) gosterir.

\begin{theorem}[Obfuscation Tersini Alma]
Eger $\mathcal{O}$ yalnizca \textbf{syntactic} donusumler uyguluyorsa (string encoding, isim degistirme, control flow yeniden yapilama), o zaman bir ters-obfuscation fonksiyonu $\mathcal{D}$ oyle vardir ki:
\begin{equation}
\mathcal{D}(\mathcal{O}(P)) \approx P
\label{eq:deobf-inverse}
\end{equation}
burada $\approx$ yapisal esdegerlik (structural equivalence) anlamindadir.
\end{theorem}

\textit{Ispat fikri:} Syntactic obfuscation, AST uzerinde terslenebilir donusumler uygular. Her donusum $t_i$ icin bir ters $t_i^{-1}$ mevcuttur. Obfuscation $\mathcal{O} = t_n \circ \cdots \circ t_1$ ise, $\mathcal{D} = t_1^{-1} \circ \cdots \circ t_n^{-1}$ dir. Ancak pratikte $t_i$'ler bilinmez; dolayisiyla Karadul gibi araclar \textbf{heuristik ters donusumleri} uygulamak zorundadir.

\subsection{Obfuscation Karmasiklik Siniflamasi}

Obfuscation tekniklerini bilgi-kuramsal zorluk derecesine gore siniflandiririz:

\begin{equation}
\text{Zorluk}(\mathcal{O}) = -\sum_{i} p_i \log_2 p_i + \text{cfg\_depth}(\mathcal{O}(P))
\label{eq:obf-complexity}
\end{equation}

burada ilk terim Shannon entropisi, ikinci terim control flow graph derinligidir.

\begin{table}[h]
\centering
\begin{tabular}{llcc}
\toprule
\textbf{\color{sectioncyan}Teknik} & \textbf{\color{sectioncyan}Zorluk} & \textbf{\color{sectioncyan}Tersinirlik} & \textbf{\color{sectioncyan}Karadul Kapsami} \\
\midrule
Minification & Dusuk & $\sim$100\% & Beautify \\
String encoding & Dusuk & $\sim$100\% & String unhex \\
Dead code injection & Orta & $\sim$95\% & Dead code elim. \\
String rotation & Orta & $\sim$90\% & Synchrony \\
Opaque predicates & Orta--Yuksek & $\sim$85\% & Pattern + Z3 \\
Control flow flattening & Yuksek & $\sim$70\% & CFF deflattener \\
Virtualization & Cok yuksek & $\sim$30\% & Kismi destek \\
\bottomrule
\end{tabular}
\caption{Obfuscation teknikleri ve zorluk dereceleri.}
\label{tab:obf-techniques}
\end{table}


% ===================================================================
\section{Yaygin Obfuscation Teknikleri}
\label{sec:obf-techniques}
% ===================================================================

\subsection{Minification}

En basit obfuscation formu. Kaynak koddan gereksiz bilgileri (bosluk, yorum, uzun degisken isimleri) kaldirarak dosya boyutunu kucultme ve okunakliligi azaltma islevidir.

\begin{lstlisting}[language=JavaScript, caption={Minification oncesi ve sonrasi}]
// Orijinal (okunabilir):
function calculateTotalPrice(items, taxRate) {
    let subtotal = 0;
    for (const item of items) {
        subtotal += item.price * item.quantity;
    }
    return subtotal * (1 + taxRate);
}

// Minified:
function a(b,c){let d=0;for(const e of b)d+=e.price*e.quantity;return d*(1+c)}
\end{lstlisting}

Minification bilgi kaybina neden olur: degisken isimleri \code{a}, \code{b}, \code{c} orijinal semantigini tasimaz. Ancak \textbf{program davranisi} aynidir --- Denklem~\ref{eq:semantic-preservation} gecerlidir.

\subsection{String Encoding}

String literal'ler cogunlukla iki yontemle gizlenir:

\textbf{Hex encoding:}
\begin{equation}
\text{encode}_{\text{hex}}(s) = \text{concat}_{i=0}^{|s|-1} \texttt{"\textbackslash x"} \cdot \text{hex}(\text{charCodeAt}(s, i))
\end{equation}

\textbf{Unicode encoding:}
\begin{equation}
\text{encode}_{\text{unicode}}(s) = \text{concat}_{i=0}^{|s|-1} \texttt{"\textbackslash u"} \cdot \text{pad}(\text{hex}(\text{codePointAt}(s, i)), 4)
\end{equation}

\textbf{Base64 encoding:}
\begin{equation}
\text{encode}_{\text{b64}}(s) = \texttt{atob("} \cdot \text{base64}(s) \cdot \texttt{")}
\end{equation}

Ornek: \code{"Hello"} $\to$ \code{"\textbackslash x48\textbackslash x65\textbackslash x6c\textbackslash x6c\textbackslash x6f"} $\to$ \code{atob("SGVsbG8=")}

\subsection{Control Flow Flattening (CFF)}

CFF, programin dogal kontrol akisini bir \texttt{while-switch} dispatcher'a donusturur:

\begin{lstlisting}[language=JavaScript, caption={CFF ornegi}]
// Orijinal:
function process(x) {
    let a = x + 1;    // Block 0
    let b = a * 2;    // Block 1
    if (b > 10) {     // Block 2
        return b;     // Block 3
    }
    return a;          // Block 4
}

// CFF uygulanmis:
function process(x) {
    let state = 0, a, b;
    while (true) {
        switch (state) {
            case 0: a = x + 1; state = 1; break;
            case 1: b = a * 2; state = 2; break;
            case 2: state = b > 10 ? 3 : 4; break;
            case 3: return b;
            case 4: return a;
        }
    }
}
\end{lstlisting}

CFF'nin formal modeli bir sonlu durum makinesidir (FSM):

\begin{equation}
\text{CFF}(P) = (Q, \Sigma, \delta, q_0, F)
\end{equation}

burada $Q$ state kuemesi, $\delta: Q \times \Sigma \to Q$ gecis fonksiyonu, $q_0$ baslangic durumu ve $F \subseteq Q$ bitis durumlaridir.

\begin{center}
\begin{tikzpicture}[
  state/.style={circle, draw=sectioncyan, fill=boxbg, minimum size=1.2cm, text=white, font=\small\bfseries},
  arrow/.style={-{Stealth[length=3mm]}, thick, color=gray!70!white},
  every node/.style={font=\small},
]
  \node[state] (q0) at (0,0) {$q_0$};
  \node[state] (q1) at (3,0) {$q_1$};
  \node[state] (q2) at (6,0) {$q_2$};
  \node[state, fill=green!20!black] (q3) at (9,1) {$q_3$};
  \node[state, fill=green!20!black] (q4) at (9,-1) {$q_4$};

  \draw[arrow] (-1.5,0) -- (q0);
  \draw[arrow] (q0) -- node[above, text=gray!60!white] {\texttt{a=x+1}} (q1);
  \draw[arrow] (q1) -- node[above, text=gray!60!white] {\texttt{b=a*2}} (q2);
  \draw[arrow] (q2) -- node[above, text=gray!60!white, sloped] {\texttt{b>10}} (q3);
  \draw[arrow] (q2) -- node[below, text=gray!60!white, sloped] {\texttt{b<=10}} (q4);

  \node[below=0.6cm, font=\tiny, text=cyan!70!white] at (q3.south) {return b};
  \node[below=0.6cm, font=\tiny, text=cyan!70!white] at (q4.south) {return a};
\end{tikzpicture}
\end{center}

\subsection{Opaque Predicates}

Opaque predicate, her zaman ayni sonucu veren (true veya false) ama statik analizin bunu kolayca cikaramayacagi ifadelerdir.

\begin{tanim}[Opaque Predicate]
Bir ifade $\varphi$ icin:
\begin{itemize}
\item $\varphi^T$: Her zaman true ($\forall \sigma: \llbracket \varphi \rrbracket_\sigma = \text{true}$)
\item $\varphi^F$: Her zaman false ($\forall \sigma: \llbracket \varphi \rrbracket_\sigma = \text{false}$)
\item $\varphi^?$: Bazen true, bazen false (gercek dallanma)
\end{itemize}
burada $\sigma$ program state'ini temsil eder.
\end{tanim}

\file{opaque\_predicate.py} dosyasinda tanimlanan oruntuleri inceleyelim:

\begin{table}[h]
\centering
\small
\begin{tabular}{p{5cm}cp{6cm}}
\toprule
\textbf{\color{sectioncyan}Oruuntu} & \textbf{\color{sectioncyan}Sonuc} & \textbf{\color{sectioncyan}Matematik Ispati} \\
\midrule
\code{x*(x+1) \% 2 == 0} & $\varphi^T$ & Ardisik iki sayinin carpimi daima cifttir \\
\code{(x | 1) != 0} & $\varphi^T$ & OR 1 her zaman $\neq 0$ verir \\
\code{(a \& b) == (a \& b)} & $\varphi^T$ & Her ifade kendisiyle esittir \\
\code{2*(x/2) <= x} & $\varphi^T$ & Tamsayi bolumuenun oezelligi: $\lfloor x/2 \rfloor \cdot 2 \leq x$ \\
\code{x*(x+1) \% 2 != 0} & $\varphi^F$ & Yukaridakinin tueuemleri \\
\code{x != x} (int) & $\varphi^F$ & Tamsayi kendisiyle esit (NaN notu asagida) \\
\bottomrule
\end{tabular}
\caption{Opaque predicate oruntuleri ve ispatlar. \textbf{Uyari:} \code{x == x} IEEE 754 NaN icin \texttt{false} dondurur; \file{opaque\_predicate.py} bu nedenle float baglam icin confidence'i 0.60'a dusurur.}
\label{tab:opaque-predicates}
\end{table}

\subsection{Dead Code Injection}

Opaque predicate'ler ile birlesirildiginde, hicbir zaman calistirilmayan sahte kod bloklari eklenir:

\begin{lstlisting}[language=JavaScript, caption={Dead code injection ornegi}]
if (x * (x + 1) % 2 === 0) {
    // Gercek kod (her zaman calisir)
    processPayment(amount);
} else {
    // Dead code (asla calismaz)
    launchMissiles();  // dikkat dagitma amacli
    deleteDatabase();  // analisti yaniltma
}
\end{lstlisting}

\subsection{String Array Rotation}

obfuscator.io gibi araclar tum string'leri bir diziye tasir, diziyi kaydirma (rotation) islemiyle karistirip indeksle erisirir:

\begin{lstlisting}[language=JavaScript, caption={String rotation ornegi}]
// Orijinal:
console.log("Hello");

// Obfuscated:
var _0x1234 = ["Hello", "log", "World"];
(function(_0x, _0y) {
    var _0z = function(_0w) { while(--_0w) { _0x.push(_0x.shift()); } };
    _0z(++_0y);
})(_0x1234, 0x1);  // Diziyi 1 kaydır
console[_0x1234[0]](_0x1234[1]);  // console.log("World") -- indeksler degisti!
\end{lstlisting}

Karadul'un \texttt{synchrony} biliseni bu patterni dogrudan cozer cunku synchrony, obfuscator.io spesifik bir deobfuscation aracidir.


% ===================================================================
\section{Abstract Syntax Tree (AST) Temelleri}
\label{sec:ast-basics}
% ===================================================================

Tum deobfuscation pipeline'i AST uzerinde calisir. Bu bolum AST'nin ne oldugunu, parse tree ile farkini ve ESTree standardini aciklar.

\subsection{BNF ve Context-Free Grammar (CFG)}

JavaScript dili bir context-free grammar (CFG) ile tanimlenir:

\begin{equation}
G = (V, \Sigma, R, S)
\end{equation}

burada $V$ non-terminal semboller, $\Sigma$ terminal semboller, $R$ uretim kurallari ve $S$ baslangic sembolueduer.

JavaScript'in basitlestirilmis BNF'inde \texttt{IfStatement}:

\begin{lstlisting}[basicstyle=\ttfamily\small\color{white}, frame=none]
IfStatement ::= "if" "(" Expression ")" Statement ("else" Statement)?
Expression  ::= AssignmentExpression | Expression "," AssignmentExpression
Statement   ::= Block | IfStatement | ForStatement | ExpressionStatement | ...
\end{lstlisting}

\subsection{Parse Tree vs Abstract Syntax Tree}

\textbf{Parse tree} (concrete syntax tree), gramerin her uretim kuralini icerir --- parantezler, noktalama isaretleri, terminal semboller hepsi agacta gorunur. \textbf{AST} ise semantik olarak anlamsiz node'lari atar.

\begin{center}
\begin{tikzpicture}[
  node distance=0.8cm and 0.5cm,
  treenode/.style={rectangle, rounded corners=2pt, draw=sectioncyan, fill=boxbg, text=white, font=\scriptsize, minimum height=0.6cm, minimum width=1cm},
  leaf/.style={rectangle, rounded corners=2pt, draw=green!60!white, fill=black!80!white, text=green!60!white, font=\scriptsize\ttfamily, minimum height=0.5cm},
  edge from parent/.style={draw=gray!60!white, -{Stealth[length=2mm]}, thin},
  level 1/.style={sibling distance=3.5cm},
  level 2/.style={sibling distance=1.8cm},
  level 3/.style={sibling distance=1.2cm},
]
  % AST tarafı
  \node[treenode, fill=cyan!20!black] {BinaryExpression (+)}
    child { node[treenode] {Identifier}
      child { node[leaf] {a} }
    }
    child { node[treenode] {NumericLiteral}
      child { node[leaf] {1} }
    };

  \node[above=0.3cm, font=\small\color{sectioncyan}] at (0, 1.2) {AST: \texttt{a + 1}};
\end{tikzpicture}
\end{center}

\subsection{ESTree Spesifikasyonu}

JavaScript AST'si icin fiili standart \textbf{ESTree} spesifikasyonudur (\texttt{https://github.com/estree/estree}). Babel parser, ESTree'nin genisletilmis bir versiyonunu uretir:

\begin{itemize}
\item \textbf{Babel AST:} \code{StringLiteral}, \code{NumericLiteral}, \code{ObjectProperty}
\item \textbf{Acorn/ESTree:} \code{Literal}, \code{Property}
\end{itemize}

Bu fark \file{deobfuscate.mjs}'te asagidaki kontrol ile ele alinir:

\begin{lstlisting}[language=JavaScript, caption={Babel vs Acorn AST tespiti (deobfuscate.mjs)}]
const usedAcorn = ast.body !== undefined && !ast.program;
// Babel AST'de ast.program var, Acorn'da yok
\end{lstlisting}

\subsection{Babel Parser Oezellikleri}

\file{deep-deobfuscate.mjs} Babel parser'i su opsiyonlarla cagrilir:

\begin{lstlisting}[language=JavaScript, caption={Babel parser konfigurasyonu}]
ast = parse(source, {
  sourceType: "unambiguous",
  allowImportExportEverywhere: true,
  allowReturnOutsideFunction: true,
  errorRecovery: true,  // Parse hatalarinda devam et
  plugins: ["jsx", "typescript", "decorators-legacy",
    "classProperties", "dynamicImport",
    "optionalChaining", "nullishCoalescingOperator",
    "topLevelAwait", "importMeta"],
});
\end{lstlisting}

\code{errorRecovery: true} kritik bir secimdir. 30MB+ minified dosyalarda parser kismen basarisiz olabilir; bu mod, hatalari toplar ama parse islemini durdurmaz. Boylece kismi AST uzerinde islem yapilabilir.


% ===================================================================
\section{Standart 4 Adimli Deobfuscation Zinciri}
\label{sec:std-chain-detail}
% ===================================================================

\file{manager.py} dosyasindaki \code{DeobfuscationManager} sinifi, standart 4 adimli zinciri orkestre eder:

\begin{center}
\begin{tikzpicture}[
  pstep/.style={rectangle, draw=sectioncyan, fill=boxbg, minimum width=3.2cm, minimum height=1.4cm, text=white, font=\small\bfseries, rounded corners=3pt, align=center},
  arrow/.style={-{Stealth[length=3mm]}, thick, color=gray!70!white},
  fail/.style={-{Stealth[length=2mm]}, dashed, color=red!60!white, thick},
  node distance=1.2cm,
]
  \node[pstep] (input) at (0,0) {Obfuscated\\JS Dosyasi};
  \node[pstep] (beautify) at (4.5,0) {1. Beautify\\(js-beautify)};
  \node[pstep] (synchrony) at (9,0) {2. Synchrony\\(obfuscator.io)};
  \node[pstep] (babel) at (4.5,-3) {3. Babel AST\\(9 faz transform)};
  \node[pstep] (webpack) at (9,-3) {4. Webpack\\Unpack};
  \node[pstep, fill=green!15!black] (output) at (13.5,-3) {Deobfuscated\\Cikti};

  \draw[arrow] (input) -- (beautify);
  \draw[arrow] (beautify) -- (synchrony);
  \draw[arrow] (synchrony) |- (babel);
  \draw[arrow] (babel) -- (webpack);
  \draw[arrow] (webpack) -- (output);

  % Graceful degradation
  \draw[fail] (beautify.south) -- ++(0,-0.8) node[midway, right, font=\tiny, text=red!60!white] {fail} -| (synchrony.south);
  \draw[fail] (synchrony.south) -- ++(0,-0.4) node[midway, right, font=\tiny, text=red!60!white] {fail} -- (babel.north);

  \node[below=0.2cm, font=\tiny, text=gray!50!white] at (beautify.south east) {\file{beautify.mjs}};
  \node[below=0.2cm, font=\tiny, text=gray!50!white] at (synchrony.south east) {\file{synchrony}};
  \node[below=0.2cm, font=\tiny, text=gray!50!white] at (babel.south east) {\file{deobfuscate.mjs}};
  \node[below=0.2cm, font=\tiny, text=gray!50!white] at (webpack.south east) {\file{smart-webpack-unpack.mjs}};
\end{tikzpicture}
\end{center}

\subsection{Graceful Degradation Formuelue}

Her adim basarisiz olabilir. Bu durumda bir onceki basarili cikti sonraki adima gecilir:

\begin{equation}
\text{output}(i) = \begin{cases}
\text{step}_i(\text{output}(i-1)) & \text{eger } \text{step}_i \text{ basarili} \\
\text{output}(i-1) & \text{eger } \text{step}_i \text{ basarisiz (fallback)}
\end{cases}
\label{eq:graceful-degradation}
\end{equation}

Zincirin en az bir adimi basarili olursa, tum zincir ``basarili'' sayilir:

\begin{equation}
\text{success}(\text{chain}) = \bigvee_{i=1}^{n} \text{success}(\text{step}_i)
\label{eq:chain-success}
\end{equation}

\file{manager.py} bunu su kodla implement eder:

\begin{lstlisting}[language=Python, caption={Graceful degradation -- manager.py, run\_chain()}]
for idx, step_name in enumerate(effective_chain, start=1):
    step_output = deob_dir / f"{idx:02d}_{step_name}{input_file.suffix}"
    try:
        step_success = self._execute_step(
            step_name, current_input, step_output, deob_dir)
    except Exception as exc:
        step_success = False
    if step_success and step_output.exists():
        result.steps_completed.append(step_name)
        current_input = step_output  # Sonraki adimin girdisi
    else:
        result.steps_failed.append(step_name)
        # current_input degismez -> onceki basarili cikti kullanilir
\end{lstlisting}

\subsection{Dosya Isimlendirme Konvansiyonu}

Her adim, ciktisini numarali dosya olarak kaydeder:

\begin{center}
\begin{tabular}{ll}
\code{00\_original.js} & Orijinal obfuscated dosya \\
\code{01\_beautify.js} & Beautify sonrasi \\
\code{02\_synchrony.js} & Synchrony sonrasi \\
\code{03\_babel\_transforms.js} & Babel AST sonrasi \\
\code{04\_webpack\_unpack.js} & Webpack modul ayirma sonrasi \\
\end{tabular}
\end{center}

Bu sayede her adimin ciktisi bagimsiz olarak incelenebilir --- debug ve audit icin kritik.


% ===================================================================
\section{Adim 1: Beautify}
\label{sec:beautify}
% ===================================================================

\file{beautify.mjs}, \texttt{js-beautify} kutuphanesini kullanan basit ama kritik bir onisleme adimidir. Minified kod tek satir halindedir; AST parse bile basa\-ri\-siz olabilir cunku Babel'in ternary+arrow+object pattern'lerinde takilma riski vardir.

\subsection{Konfiguerasyon Parametreleri}

\begin{lstlisting}[language=JavaScript, caption={js-beautify konfigurasyonu -- beautify.mjs}]
js_beautify(source, {
  indent_size: 2,
  indent_char: " ",
  max_preserve_newlines: 2,
  brace_style: "collapse,preserve-inline",
  wrap_line_length: 120,
  break_chained_methods: true,
  end_with_newline: true,
});
\end{lstlisting}

\subsection{Etki Analizi}

Beautify'in etkisini olcmek icin genisleme oranini (expansion ratio) tanimlayalim:

\begin{equation}
R_{\text{expand}} = \frac{L_{\text{out}}}{L_{\text{in}}}
\end{equation}

burada $L_{\text{in}}$ girdi satir sayisi, $L_{\text{out}}$ cikti satir sayisi. Tipik degerler:

\begin{itemize}
\item Agresif minification (tek satir): $R_{\text{expand}} \approx 100\text{--}500$
\item Orta minification: $R_{\text{expand}} \approx 2\text{--}5$
\item Zaten formatlenmis: $R_{\text{expand}} \approx 1$
\end{itemize}

\subsection{Fallback Davranisi}

Beautify basarisiz olursa, girdi dosyasi dogrudan ciktiya kopyalanir:

\begin{lstlisting}[language=Python, caption={Beautify fallback -- manager.py}]
def _step_beautify(self, input_file, output_file):
    beautify_script = self._config.scripts_dir / "beautify.mjs"
    if not beautify_script.exists():
        shutil.copy2(input_file, output_file)  # Fallback
        return True
\end{lstlisting}


% ===================================================================
\section{Adim 2: Synchrony Deobfuscation}
\label{sec:synchrony}
% ===================================================================

\file{synchrony\_wrapper.py}, \texttt{synchrony} aracini saran bir Python wrapper'dir. Synchrony, obfuscator.io spesifik bir deobfuscation aracidir ve string rotation, string array, control flow flattening gibi obfuscator.io'ya oezgue teknikleri cozer.

\subsection{Calismaa Mekanizmasi}

\begin{enumerate}
\item \textbf{Shebang soyma:} Eger dosya \code{\#!/usr/bin/env node} ile basliyorsa, synchrony'nin Acorn parser'i bunu parse edemez. Wrapper gecici dosyada shebang'i soyar, islem sonrasi geri ekler.
\item \textbf{Source type deneme:} Synchrony uc modda denenir: \code{both}, \code{module}, \code{script}. ESM dosyalarda \code{ImportDeclaration} hatasi alinirsa \code{module} moduna gecilir.
\item \textbf{Dogrudan dosya ciktisi:} \code{-o} flag'i ile stdout yerine dogrudan dosyaya yazar. Buyuk dosyalarda memory overhead'i onler.
\end{enumerate}

\begin{lstlisting}[language=Python, caption={Synchrony calisma akisi -- synchrony\_wrapper.py}]
source_types = ["both", "module", "script"]
for source_type in source_types:
    cmd = [sync_path, str(actual_input),
           "-o", str(output_file),
           "--type", source_type]
    result = self._runner.run_command(cmd, timeout=timeout)
    if result.success:
        break
    if "ImportDeclaration" in result.stderr:
        continue  # ESM hatasi, module modunu dene
\end{lstlisting}

\subsection{Boyut Analizi}

Synchrony basarili oldugunda dosya boyutu genellikle \textbf{kucueluer} cunku string rotation ve dead code kaldirilir:

\begin{equation}
R_{\text{sync}} = \frac{|\text{output}|}{|\text{input}|} \in [0.3, 1.0]
\end{equation}

$R_{\text{sync}} < 0.5$ ise agresif obfuscation kaldirilmis demektir.


% ===================================================================
\section{9 Fazli Deep Deobfuscation Pipeline}
\label{sec:deep-deobf-detail}
% ===================================================================

\file{deep-deobfuscate.mjs} dosyasi, 10 asamali (9 transform + 1 parse) derin deobfuscation pipeline'i uygular. Phase 2--8, \textbf{visitor-merge optimizasyonu} ile tek bir AST traverse'da birlestirilmistir --- bu webcrack yaklasimi, 7 ayri traverse yerine 1 traverse kullanarak $\sim$3--5x hiz kazanci saglar.

\begin{center}
\begin{tikzpicture}[
  phase/.style={rectangle, draw=sectioncyan, fill=boxbg, minimum width=4cm, minimum height=0.8cm, text=white, font=\scriptsize\bfseries, rounded corners=2pt},
  merged/.style={rectangle, draw=yellow!70!white, fill=boxbg, minimum width=4cm, minimum height=0.8cm, text=yellow!70!white, font=\scriptsize\bfseries, rounded corners=2pt},
  arrow/.style={-{Stealth[length=2mm]}, thin, color=gray!60!white},
  node distance=0.5cm,
]
  \node[phase] (p1) {Phase 1: Parse};
  \node[merged, below=of p1] (p2) {Phase 2: Constant Folding};
  \node[merged, below=of p2] (p3) {Phase 3: Dead Code Elim.};
  \node[merged, below=of p3] (p4) {Phase 4: Computed $\to$ Static};
  \node[merged, below=of p4] (p5) {Phase 5: Comma Split};
  \node[merged, below=of p5] (p6) {Phase 6: Var Decl. Split};
  \node[merged, below=of p6] (p7) {Phase 7: Ternary $\to$ If};
  \node[merged, below=of p7] (p8) {Phase 8: Arrow $\to$ Function};
  \node[phase, below=of p8] (p9) {Phase 9: Smart Rename};
  \node[phase, below=of p9] (p10) {Phase 10: Semantic Rename};

  \draw[arrow] (p1) -- (p2);
  \draw[arrow] (p2) -- (p3);
  \draw[arrow] (p3) -- (p4);
  \draw[arrow] (p4) -- (p5);
  \draw[arrow] (p5) -- (p6);
  \draw[arrow] (p6) -- (p7);
  \draw[arrow] (p7) -- (p8);
  \draw[arrow] (p8) -- (p9);
  \draw[arrow] (p9) -- (p10);

  % Merge bracket
  \draw[yellow!70!white, thick, decorate, decoration={brace, amplitude=8pt, mirror}]
    ([xshift=2.3cm]p2.north east) -- ([xshift=2.3cm]p8.south east)
    node[midway, right=10pt, text=yellow!70!white, font=\scriptsize, align=left] {Tek AST\\traverse\\(visitor-merge)};

  % Separate bracket
  \draw[sectioncyan, thick, decorate, decoration={brace, amplitude=8pt, mirror}]
    ([xshift=2.3cm]p9.north east) -- ([xshift=2.3cm]p10.south east)
    node[midway, right=10pt, text=sectioncyan, font=\scriptsize, align=left] {Ayri\\traverse\\(scope gerekli)};
\end{tikzpicture}
\end{center}

\subsection{Faz 1: Parse (Tolerant Mode)}
\label{sec:phase1}

Babel parser \code{errorRecovery: true} ile calistirilir. Basarisiz olursa iki fallback mevcuttur:

\begin{enumerate}
\item \textbf{Pre-beautify + retry:} Minified kaynak once \code{js-beautify} ile formatlanir, sonra Babel tekrar denenir. 30MB+ dosyalarda Babel'in ternary/arrow/object pattern'lerinde takilma sorunu bu yontemle cogunlukla cozulur.
\item \textbf{Dosya kopyalama:} Her iki yontem de basarisiz olursa, kaynak dosya oldugu gibi output'a kopyalanir. Boylece sonraki pipeline adimlari (webpack unpack) en azindan orijinal dosya uzerinde devam edebilir.
\end{enumerate}

\begin{equation}
\text{parse}(S) = \begin{cases}
\text{Babel}(S) & \text{basarili} \\
\text{Babel}(\text{beautify}(S)) & \text{pre-beautify ile retry} \\
S & \text{fallback: kaynak kopyalama}
\end{cases}
\end{equation}


\subsection{Faz 2: Constant Folding (Sabit Katlama)}
\label{sec:phase2-cf}

Constant folding, derleme/analiz zamaninda hesaplanabilir ifadeleri literal degerleriyle degistirme islemidir.

\subsubsection{Lattice Teorisi ve Abstract Interpretation}

Constant folding, abstract interpretation cercevesinde bir \textbf{flat lattice} uzerinde calisir:

\begin{equation}
L = \{\bot, c_1, c_2, \ldots, c_n, \top\}
\end{equation}

burada:
\begin{itemize}
\item $\bot$: Degisken henuz kullanilmadi (undefined)
\item $c_i$: Degisken sabit $c_i$ degerinde
\item $\top$: Degisken birden fazla deger alabilir (non-constant)
\end{itemize}

\begin{equation}
\text{meet}(x, y) = \begin{cases}
x & \text{eger } x = y \\
\top & \text{eger } x \neq y \text{ ve her ikisi de } \neq \bot \\
y & \text{eger } x = \bot \\
x & \text{eger } y = \bot
\end{cases}
\end{equation}

\subsubsection{Karadul Implementasyonu}

\file{deep-deobfuscate.mjs} Phase 2'de asagidaki donusumler uygulanir:

\begin{table}[h]
\centering
\small
\begin{tabular}{lll}
\toprule
\textbf{\color{sectioncyan}Node Tipi} & \textbf{\color{sectioncyan}Donusum} & \textbf{\color{sectioncyan}Ornek} \\
\midrule
BinaryExpression (+) & String concat & \code{"He"+"llo"} $\to$ \code{"Hello"} \\
BinaryExpression (+) & Numeric add & \code{1+2} $\to$ \code{3} \\
BinaryExpression ($*$, $/$, $-$) & Aritmetik & \code{6*7} $\to$ \code{42} \\
BinaryExpression ($|$,$\&$,$\hat{}$,$\ll$,$\gg$) & Bitwise & \code{0xFF \& 0x0F} $\to$ \code{15} \\
UnaryExpression (!) & Negation & \code{!0} $\to$ \code{true} \\
UnaryExpression (void) & Void & \code{void 0} $\to$ \code{undefined} \\
\bottomrule
\end{tabular}
\caption{Constant folding donusumm tablosu.}
\label{tab:constant-folding}
\end{table}

Guevenlik kontrolu: Bolme islemlerinde \code{right.value !== 0} kontrol edilir ve $\text{Number.isFinite}(result)$ dogrulanir. Overfllow, NaN ve Infinity sonuclari degistirilmez.

\begin{lstlisting}[language=JavaScript, caption={Constant folding visitor -- BinaryExpression}]
BinaryExpression: {
  exit(path) {
    const { left, right, operator } = path.node;
    // String concatenation
    if (operator === "+" && t.isStringLiteral(left)
        && t.isStringLiteral(right)) {
      path.replaceWith(t.stringLiteral(left.value + right.value));
      cnt2++;
    }
    // Numeric operations
    if (t.isNumericLiteral(left) && t.isNumericLiteral(right)) {
      let result;
      switch (operator) {
        case "+": result = left.value + right.value; break;
        case "-": result = left.value - right.value; break;
        // ... (*, /, %, **, |, &, ^, <<, >>, >>>)
      }
      if (result !== undefined && Number.isFinite(result)) {
        path.replaceWith(t.numericLiteral(result));
        cnt2++;
      }
    }
  }
}
\end{lstlisting}

\textbf{Neden \code{exit} visitor?} Bottom-up traversal sayesinde ic node'lar once islenir: \code{(1+2)*3}'te once \code{1+2=3}, sonra \code{3*3=9}. \code{enter} kullansaydik, ic node'lar henuz fold edilmemis olurdu.


\subsection{Faz 3: Dead Code Elimination}
\label{sec:phase3-dce}

\subsubsection{Liveness Analysis Temeli}

Liveness analysis, her program noktasinda hangi degiskenlerin ``canli'' (daha sonra kullanilacak) oldugunu belirler:

\begin{equation}
\text{live\_in}(B) = \text{gen}(B) \cup (\text{live\_out}(B) \setminus \text{kill}(B))
\label{eq:live-in}
\end{equation}
\begin{equation}
\text{live\_out}(B) = \bigcup_{S \in \text{succ}(B)} \text{live\_in}(S)
\label{eq:live-out}
\end{equation}

burada:
\begin{itemize}
\item $\text{gen}(B)$: $B$ blogunda kullanilan (okunan) degiskenler
\item $\text{kill}(B)$: $B$ blogunda tanimlanan (yazilan) degiskenler
\end{itemize}

\subsubsection{Fixed-Point Iteration}

Liveness analysis, denklem sistemi sabit noktaya ulasana kadar iteratif olarak cozulur:

\begin{equation}
\text{live\_in}^{(k+1)}(B) = \text{gen}(B) \cup \left(\text{live\_out}^{(k)}(B) \setminus \text{kill}(B)\right)
\end{equation}

Karmasiklik: $\bigO{|B| \times |V|}$ burada $|B|$ basic block sayisi, $|V|$ degisken sayisi.

\subsubsection{Karadul Implementasyonu}

Phase 3, \code{IfStatement} ve \code{ConditionalExpression} node'larini isle\-r:

\begin{lstlisting}[language=JavaScript, caption={Dead code elimination -- IfStatement}]
IfStatement: {
  exit(path) {
    const test = path.node.test;
    // if (false) {...} -> kaldir veya else branch'i al
    if (t.isBooleanLiteral(test) && test.value === false) {
      if (path.node.alternate) path.replaceWith(path.node.alternate);
      else path.remove();
      cnt3++;
    }
    // if (true) {X} else {Y} -> X
    if (t.isBooleanLiteral(test) && test.value === true) {
      path.replaceWith(path.node.consequent);
      cnt3++;
    }
    // if (0) -> remove, if (1) -> consequent
    if (t.isNumericLiteral(test)) {
      if (test.value === 0) { /* remove or alternate */ }
      else { path.replaceWith(path.node.consequent); }
      cnt3++;
    }
  }
}
\end{lstlisting}

\textbf{Phase 2 ile etkilesim:} Phase 2 \code{!0 $\to$ true} donusumunu yapar. Phase 3, \code{exit} visitor kullandigi icin, ayni traverse'da Phase 2'nin ciktisini gorebilir. Ornegin \code{if(!0)\{X\}} once \code{if(true)\{X\}} olur (Phase 2), sonra \code{X} olur (Phase 3).


\subsection{Faz 4: Computed-to-Static Property}
\label{sec:phase4-cts}

Bu faz, \code{obj["foo"]} gibi computed property erisimlerini \code{obj.foo} gibi statik erisimlere donusturur.

\begin{equation}
\text{obj}[\texttt{"key"}] \xrightarrow{\text{Faz 4}} \text{obj}.\texttt{key}
\quad \text{eger } \texttt{key} \in \text{ValidIdentifier}
\end{equation}

\begin{tanim}[Valid Identifier]
Bir string \code{key} gecerli JavaScript identifier ise ve reserved word degilse donusum guevenlidir:
\begin{equation}
\text{ValidIdentifier}(\texttt{key}) = \texttt{/\^{}[a-zA-Z\_\$][a-zA-Z0-9\_\$]*\$/}.test(\texttt{key}) \wedge \neg\text{isReserved}(\texttt{key})
\end{equation}
\end{tanim}

AST duenuesuemuee:
\begin{itemize}
\item \code{MemberExpression}: \code{computed: true} $\to$ \code{computed: false}
\item \code{ObjectProperty}: Computed key'i identifier'a donustur
\end{itemize}

\begin{lstlisting}[language=JavaScript, caption={Computed to static -- MemberExpression}]
MemberExpression(path) {
  if (!path.node.computed) return;
  const prop = path.node.property;
  if (!t.isStringLiteral(prop)) return;
  if (/^[a-zA-Z_$][a-zA-Z0-9_$]*$/.test(prop.value)
      && !isReservedWord(prop.value)) {
    path.node.computed = false;
    path.node.property = t.identifier(prop.value);
    cnt4++;
  }
}
\end{lstlisting}


\subsection{Faz 5: Comma Expression Splitting}
\label{sec:phase5-comma}

JavaScript'te virgul operatoru (\code{SequenceExpression}) birden fazla ifadeyi tek bir expression'da birlestirir. Obfuscator'lar bunu okunakliligi azaltmak icin kullanir:

\begin{lstlisting}[language=JavaScript]
// Obfuscated:
return a = 1, b = 2, c = 3, process(a, b, c);
// Deobfuscated (Faz 5 sonrasi):
a = 1;
b = 2;
c = 3;
return process(a, b, c);
\end{lstlisting}

Donusum ucg bag\-lam icin uygulanir:

\begin{enumerate}
\item \textbf{ExpressionStatement:} \code{(a, b, c);} $\to$ \code{a; b; c;}
\item \textbf{ReturnStatement:} \code{return (a, b, c);} $\to$ \code{a; b; return c;}
\item \textbf{ThrowStatement:} \code{throw (a, b, c);} $\to$ \code{a; b; throw c;}
\end{enumerate}

\begin{equation}
\text{split}(\texttt{return}(e_1, e_2, \ldots, e_n)) = \left[\text{expr}(e_1), \ldots, \text{expr}(e_{n-1}), \texttt{return}\; e_n\right]
\end{equation}


\subsection{Faz 6: Variable Declaration Splitting}
\label{sec:phase6-vardecl}

Tek bir \code{var/let/const} icindeki birden fazla tanimlama ayri satirlara bolunur:

\begin{lstlisting}[language=JavaScript]
// Oncesi:
var a = 1, b = 2, c = fn(), d = [];
// Sonrasi:
var a = 1;
var b = 2;
var c = fn();
var d = [];
\end{lstlisting}

\textbf{Guvenlik kontrolleri:} For-loop init (\code{for(var i=0, j=0; ...)}), for-in ve for-of icerisindeki deklarasyonlar bolunmez:

\begin{lstlisting}[language=JavaScript, caption={Variable declaration split -- guvenlik kontrolleri}]
if (path.parent.type === "ForStatement" && path.key === "init") return;
if (path.parent.type === "ForInStatement") return;
if (path.parent.type === "ForOfStatement") return;
\end{lstlisting}


\subsection{Faz 7: Ternary $\to$ If/Else}
\label{sec:phase7-ternary}

Uzun ternary ifadeler if/else bloklarina donusturulur. ``Uzun'' tanimi: test + consequent + alternate toplam karakter sayisi $\geq 80$:

\begin{equation}
|\text{code}(\text{test})| + |\text{code}(\text{cons})| + |\text{code}(\text{alt})| \geq 80
\end{equation}

\begin{lstlisting}[language=JavaScript]
// Oncesi (tek satir, okunamaz):
isAuthenticated ? renderDashboard(user, session) : redirectToLogin(returnUrl);
// Sonrasi:
if (isAuthenticated) {
  renderDashboard(user, session);
} else {
  redirectToLogin(returnUrl);
}
\end{lstlisting}


\subsection{Faz 8: Arrow $\to$ Named Function}
\label{sec:phase8-arrow}

\code{const foo = () => \{...\}} formundaki arrow function'lar named function declaration'a donusturulur:

\begin{lstlisting}[language=JavaScript]
// Oncesi:
const processData = (items) => { return items.map(x => x * 2); };
// Sonrasi:
function processData(items) { return items.map(x => x * 2); }
\end{lstlisting}

\textbf{Semantik guevenlik:} \code{this} kullanan arrow function'lar donusturulmez. Arrow function'larda \code{this} lexical scope'a baglidir; normal function'larda ise call-site'a baglidir. Bu semantik fark bozulmamalidir:

\begin{lstlisting}[language=JavaScript, caption={Arrow-to-function this kontrolu}]
const { code: bodyCode } = generate(arrow.body, { compact: true });
if (!bodyCode.includes("this.") && !bodyCode.includes("this[")) {
    // Guvenli: donustur
    const funcDecl = t.functionDeclaration(
      t.identifier(decl.id.name), arrow.params, arrow.body,
      arrow.generator || false, arrow.async || false);
    path.replaceWith(funcDecl);
}
\end{lstlisting}


\subsection{Faz 9: Smart Variable Renaming}
\label{sec:phase9-rename}

Bu faz, minified identifier'lari (1--2 karakter) anlamli isimlere donusturur. LLM \textbf{kullanmaz} --- tamamen heuristik kurallara dayanir. Babel \code{scope.rename()} API'si ile scope-safe renaming yapilir.

\subsubsection{Uec Kural Kategorisi}

\begin{enumerate}
\item \textbf{Module-based:} \code{require("fs")} $\to$ degisken \code{fs} olarak adlandirilir
\item \textbf{Usage-based:} Degisken uzerinde \code{.push()}, \code{.map()} cagriliyorsa $\to$ \code{items}
\item \textbf{Fallback:} Tek harf degiskenler icin varsayilan isimler (\code{i} $\to$ \code{idx}, \code{e} $\to$ \code{elem})
\end{enumerate}

\begin{table}[h]
\centering
\small
\begin{tabular}{p{3.5cm}p{4cm}p{4.5cm}}
\toprule
\textbf{\color{sectioncyan}Pattern} & \textbf{\color{sectioncyan}Tespit Yontemi} & \textbf{\color{sectioncyan}Yeni Isim} \\
\midrule
\code{require("fs")} & CallExpression analizi & \code{fs} \\
\code{require("path")} & CallExpression analizi & \code{path} \\
\code{require("express")} & CallExpression analizi & \code{express} \\
\code{catch(e)} & CatchClause param & \code{error} \\
\code{x.push()}, \code{x.map()} & Usage (Array API) & \code{items} \\
\code{x.trim()}, \code{x.split()} & Usage (String API) & \code{text} \\
\code{x.message}, \code{x.stack} & Property access & \code{err} \\
\code{x.then()}, \code{x.catch()} & Usage (Promise API) & \code{promise} \\
Express \code{(req, res, next)} & Paramettre pozisyonu & \code{request, response, next} \\
esbuild \code{(exports, module)} & Factory pattern & \code{exports, module} \\
\bottomrule
\end{tabular}
\caption{Phase 9 smart rename kural ornekleri.}
\label{tab:phase9-rules}
\end{table}

\subsubsection{Tek Harf Fallback Tablosu}

\begin{lstlisting}[language=JavaScript, caption={Tek harf fallback mapping}]
const defaults = {
  a: "arg", b: "buf", c: "ctx", d: "data",
  e: "elem", f: "fn", g: "gen", h: "handler",
  i: "idx", j: "jdx", k: "key", l: "len",
  m: "mod", n: "count", o: "opts", p: "param",
  q: "query", r: "res", s: "str", t: "tmp",
  u: "usr", v: "val", w: "writer",
};
\end{lstlisting}


\subsection{Faz 10: Semantic Renaming (Deep Usage-Based)}
\label{sec:phase10-semantic}

Phase 10, Phase 9'un basit pattern matching'ini asarak, Babel scope + binding analizi ile her minified identifier'in kullanim baglamini derinlemesine analiz eder. \textbf{12 kural motoru} ile calisir.

\subsubsection{Kullanim Bilgisi Toplama (\_collectUsage10)}

Her identifier icin asagidaki bilgiler toplanir:

\begin{lstlisting}[language=JavaScript, caption={Usage bilgi yapisi}]
{
  methodCalls: Set,       // x.push(), x.split() gibi
  propertyReads: Set,     // x.length, x.name gibi
  propertyWrites: Set,    // x.count = 5 gibi
  calledAs: boolean,      // x() seklinde cagriliyor mu
  comparedWith: string[], // x === "GET" gibi karsilastirmalar
  typeofChecks: string[], // typeof x === "function" gibi
  usedInForOf: boolean,   // for(const y of x) iteratorde mi
  destructured: boolean,  // const {a, b} = x seklinde mi
  indexAccessed: boolean,  // x[0] seklinde mi
  arithmeticOperand: boolean, // x + 1 gibi aritmetikte mi
  thrownAsError: boolean, // throw x seklinde mi
  usedInNew: boolean,     // new x() seklinde mi
}
\end{lstlisting}

\subsubsection{12 Kural Motoru}

\begin{table}[h]
\centering
\small
\begin{tabular}{clcp{4.5cm}}
\toprule
\textbf{\color{sectioncyan}\#} & \textbf{\color{sectioncyan}Kural} & \textbf{\color{sectioncyan}Conf.} & \textbf{\color{sectioncyan}Tespit Kriteri} \\
\midrule
1 & \code{require\_module} & 0.95 & \code{require("fs")} $\to$ \code{fs} \\
2 & \code{new\_error} & 0.90 & \code{new Error()} pattern \\
3 & \code{error\_props} & 0.90 & \code{.message} + \code{.stack} props \\
4 & \code{string\_api} & 0.85 & $\geq$2 string method \\
5 & \code{array\_api} & 0.85 & $\geq$2 array method \\
6 & \code{promise\_api} & 0.85 & \code{.then()} + \code{.catch()} \\
7 & \code{fs\_api} & 0.90 & \code{readFileSync}, \code{writeFile} \\
8 & \code{http\_req} & 0.80 & $\geq$2 of \code{body,params,query,headers} \\
9 & \code{dom\_api} & 0.85 & $\geq$2 DOM property \\
10 & \code{callback} & 0.80 & \code{calledAs} + \code{typeof === "function"} \\
11 & \code{for\_of\_item} & 0.80 & \code{for(const x of iterable)} \\
12 & \code{multi\_compare} & 0.70 & $\geq$3 string karsilastirmasi \\
\bottomrule
\end{tabular}
\caption{Phase 10 semantic renaming kural motoru.}
\label{tab:phase10-rules}
\end{table}

\subsubsection{Phase 10+: Enhanced Rename (cursor-enhanced-rename.mjs)}

\file{cursor-enhanced-rename.mjs}, Phase 10'un kapsamadigi 15 ek kural ile kalan minified identifier'lari isle\-r. \file{deep\_pipeline.py}'da Adim 2.5 olarak calistirilir:

\begin{table}[h]
\centering
\small
\begin{tabular}{clp{5.5cm}}
\toprule
\textbf{\color{sectioncyan}\#} & \textbf{\color{sectioncyan}Kural} & \textbf{\color{sectioncyan}Aciklama} \\
\midrule
11 & Prototype access & \code{x.prototype} $\to$ \code{cls} \\
12 & Binary ops dominant & Cok fazla \code{+}, \code{-}, \code{<<} $\to$ \code{offset} \\
13 & Purely called & Sadece cagriliyor $\to$ \code{handler} \\
14 & Event listener pos & Callback pozisyonu $\to$ \code{eventArg} \\
15 & Member assign & Atama kaynagi $\to$ \code{ref} \\
16 & Single property & Tek prop baskini $\to$ contextual \\
17 & Conditional return & \code{return x ? a : b} $\to$ \code{result} \\
18 & Switch/if chain & Karsilastirma zinciri $\to$ \code{kind} \\
19 & Namespace access & \code{X.C}, \code{X.Q} $\to$ \code{ns} \\
20 & Loop counter & \code{for} init $\to$ \code{idx} \\
21 & Await expression & \code{await x} $\to$ \code{awaited} \\
22 & Default param ctx & Varsayilan deger $\to$ \code{optParam} \\
23 & Clash dedup & Phase 9 artigi $\to$ re-score \\
24 & Length/size dominant & \code{.length} baskini $\to$ \code{collection} \\
25 & Bitwise dominant & \code{|}, \code{\&}, \code{\^{}} $\to$ \code{flags} \\
\bottomrule
\end{tabular}
\caption{Enhanced rename ek kurallari (Kural 11--25).}
\label{tab:enhanced-rules}
\end{table}


% ===================================================================
\section{Webpack Module Extraction}
\label{sec:webpack-unpack}
% ===================================================================

\file{smart-webpack-unpack.mjs}, webpack ve esbuild bundle'larini modullerine ayiran akilli bir module splitter'dir. 7 farkli bundle formatini destekler.

\subsection{Desteklenen Bundle Formatlari}

\begin{enumerate}
\item \textbf{esbuild CJS wrapper:} \code{var X = U((exports, module) => \{...\})}
\item \textbf{esbuild ESM interop:} \code{var Y = Q1(X)}
\item \textbf{esbuild re-export:} \code{lb(target, \{name: () => ref\})}
\item \textbf{Webpack IIFE object:} \code{(function(modules) \{...\})(\{0: function(e,t,n)\{...\}\})}
\item \textbf{Webpack IIFE array:} \code{[function(e,t,n)\{...\}, ...]}
\item \textbf{Webpack 5 named:} \code{\_\_webpack\_modules\_\_ = \{"./src/file.js": ...\}}
\item \textbf{Top-level IIFE:} \code{(function() \{...\})()} --- kapsayici IIFE'leri ayirma
\end{enumerate}

\subsection{\_\_webpack\_require\_\_ Mekanizmasi}

Webpack bundle'larinda moduller \code{\_\_webpack\_require\_\_(id)} ile birbirini cagrilir. Bu pattern AST'de tespit edilir:

\begin{lstlisting}[language=JavaScript, caption={Webpack module tespiti -- deobfuscate.mjs}]
if (callee.name === "__webpack_require__"
    && node.arguments.length >= 1) {
  const arg = node.arguments[0];
  if (arg.type === "NumericLiteral") moduleId = arg.value;
  else if (arg.type === "StringLiteral") moduleId = arg.value;
  if (moduleId !== null && !_webpackIdSet.has(moduleId)) {
    _webpackIdSet.add(moduleId);
    webpackModules.push(moduleId);
  }
}
\end{lstlisting}

\subsection{esbuild vs Webpack Format Farklari}

\begin{table}[h]
\centering
\small
\begin{tabular}{p{3.5cm}p{5cm}p{5cm}}
\toprule
\textbf{\color{sectioncyan}Ozellik} & \textbf{\color{sectioncyan}Webpack} & \textbf{\color{sectioncyan}esbuild} \\
\midrule
Modul sarmalama & IIFE + nesne/dizi & CJS factory wrapper \\
Modul ID & Sayisal (0, 1, 2...) & Degisken adi \\
Dependency & \code{\_\_webpack\_require\_\_} & Dogrudan degisken referansi \\
Entry point & Ilk cagri arg. & Top-level kod \\
Tree shaking & Plugin bazli & Varsayilan \\
\bottomrule
\end{tabular}
\caption{Webpack vs esbuild bundle farklari.}
\label{tab:webpack-vs-esbuild}
\end{table}


% ===================================================================
\section{Buyuk Dosya Isleme: Chunked Processing}
\label{sec:chunked}
% ===================================================================

200MB+ dosyalar icin standart AST parsing bellege sigmaz. \file{deep\_pipeline.py} otomatik olarak \textbf{chunked processing} moduna gecer.

\subsection{Esik Degeri ve Tetikleme}

\begin{equation}
\text{mode} = \begin{cases}
\text{normal} & \text{eger dosya} < 100\text{ MB} \\
\text{chunked} & \text{eger dosya} \geq 100\text{ MB}
\end{cases}
\end{equation}

\file{deep\_pipeline.py}'da \code{LARGE\_FILE\_THRESHOLD\_MB = 100} olarak tanimli.

\subsection{Stream-Parse Pipeline}

\begin{enumerate}
\item \textbf{stream-parse.mjs:} Dosyayi top-level block'lara ayirir (\code{--max-block-kb 512})
\item \textbf{Her block'u bagimsiz isle:} \code{deep-deobfuscate.mjs} her chunk uzerinde ayri calisir (120s timeout/block)
\item \textbf{Birlestir:} Tum chunk ciktilari tek dosyada birlestirilir
\item \textbf{Webpack unpack:} Birlestirilmis dosya uzerinde modul ayirma
\item \textbf{Ikinci pass:} Her modul uzerinde variable renaming
\end{enumerate}

\begin{equation}
T_{\text{chunked}} \approx T_{\text{split}} + n \cdot T_{\text{deob/chunk}} + T_{\text{merge}} + T_{\text{unpack}}
\end{equation}

burada $n$ chunk sayisi, $T_{\text{deob/chunk}}$ chunk basina deobfuscation suresi.

\subsection{Timeout Hesaplaması}

\file{deep\_pipeline.py} dosya boyutuna gore dinamik timeout hesaplar:

\begin{equation}
T_{\text{timeout}} = \max(300, \min(1800, \lceil S_{\text{MB}} \times 100 \rceil)) \text{ saniye}
\label{eq:timeout}
\end{equation}

burada $S_{\text{MB}}$ dosya boyutu megabayt cinsindendir. Minimum 5 dakika, maksimum 30 dakika.


% ===================================================================
\section{Control Flow Flattening (CFF) Deflattening}
\label{sec:cff-deflattening}
% ===================================================================

\file{cff\_deflattener.py}, CFF obfuscation'ini geri alan bir Python modul-udur. Hem C hem JavaScript uzerinde calisir.

\subsection{Tespit Algoritması}

CFF dispatcher pattern regex ile tespit edilir:

\begin{lstlisting}[language=Python, caption={CFF dispatcher tespiti}]
# JS CFF pattern
_WHILE_SWITCH_JS = re.compile(
    r"(?:while\s*\(\s*(?:!0|true|1)\s*\)|for\s*\(\s*;;\s*\))\s*\{"
    r"\s*switch\s*\(\s*(\w+)\s*\)\s*\{",
    re.IGNORECASE,
)
\end{lstlisting}

\subsection{State Transition Graph Cikarma}

Her \code{case} blogundan state gecisleri cikarilarak bir yoenluee graf (directed graph) olusturulur:

\begin{equation}
G = (V, E) \quad \text{burada } V = \{s_0, s_1, \ldots, s_n\}, \quad E = \{(s_i, s_j) : s_i \xrightarrow{\text{state}} s_j\}
\end{equation}

\begin{center}
\begin{tikzpicture}[
  state/.style={circle, draw=sectioncyan, fill=boxbg, minimum size=0.9cm, text=white, font=\small},
  arrow/.style={-{Stealth[length=2.5mm]}, thick, color=gray!70!white},
  every node/.style={font=\scriptsize},
]
  \node[state] (s0) at (0,0) {0};
  \node[state] (s3) at (2.5,0) {3};
  \node[state] (s1) at (5,0) {1};
  \node[state] (s5) at (7.5,0) {5};
  \node[state] (s2) at (5,-1.5) {2};

  \draw[arrow] (s0) -- (s3);
  \draw[arrow] (s3) -- (s1);
  \draw[arrow] (s1) -- (s5);
  \draw[arrow] (s3) -- (s2);

  \node[above=0.2cm, text=green!60!white] at (s0) {entry};
  \node[above=0.2cm, text=red!60!white] at (s5) {exit};
\end{tikzpicture}
\end{center}

\subsection{Topological Sort ile Yeniden Olusturma}

State transition graph'i uzerinde BFS-based topological sort uygulanarak orijinal kod sirasi elde edilir:

\begin{lstlisting}[language=Python, caption={Topological sort -- cff\_deflattener.py}]
def _topological_sort(self, graph, blocks):
    order, visited = [], set()
    entry = blocks[0].state_value
    queue = [entry]
    while queue:
        state = queue.pop(0)
        if state in visited: continue
        visited.add(state)
        order.append(state)
        for ns in graph.get(state, []):
            if ns not in visited:
                queue.append(ns)
    # Ziyaret edilmemis state'leri sona ekle
    for block in blocks:
        if block.state_value not in visited:
            order.append(block.state_value)
    return order
\end{lstlisting}

\subsection{Karmasiklik Analizi}

\begin{equation}
T_{\text{CFF}} = \bigO{|C| + |V| + |E|}
\end{equation}

burada $|C|$ case sayisi, $|V|$ state sayisi (=$|C|$), $|E|$ gecis sayisi. Pratikte $|E| \leq 2|V|$ (her case en fazla 2 gecis yapar --- true/false branch).


% ===================================================================
\section{Opaque Predicate Removal}
\label{sec:opaque-removal}
% ===================================================================

\file{opaque\_predicate.py}, bilinen opaque predicate pattern'lerini regex ile tespit edip dead code'u isaretler veya kaldirir.

\subsection{Tespit Stratejileri}

\begin{enumerate}
\item \textbf{Pattern matching:} 7 always-true + 4 always-false regex pattern (Tablo~\ref{tab:opaque-predicates})
\item \textbf{Constant folding:} Sabit deger kullanan branch'ler (Phase 3 ile entegre)
\item \textbf{Z3 SMT solver (opsiyonel):} Karmasik opaque predicate'lar icin symbolic verification
\end{enumerate}

\subsection{Z3 Entegrasyonu}

Eger z3-solver kurulu ise, pattern matching ile tespit edilemeyen karmasik predicate'lar Z3 ile dogrulanabilir:

\begin{equation}
\text{is\_opaque}(\varphi) = \begin{cases}
\text{true} & \text{eger } Z3.\text{prove}(\forall x: \varphi(x) = \text{true}) \\
\text{false} & \text{eger } Z3.\text{find\_counterexample}(\varphi) \neq \varnothing
\end{cases}
\end{equation}

\file{opaque\_predicate.py}'da Z3 lazy-load yapilir:

\begin{lstlisting}[language=Python, caption={Z3 lazy-load -- opaque\_predicate.py}]
if use_z3:
    try:
        import z3  # noqa: F401
        self._z3_available = True
    except ImportError:
        logger.info("z3-solver bulunamadi, "
                     "pattern-based tespit kullanilacak")
\end{lstlisting}

\subsection{Eliminasyon}

\begin{itemize}
\item \textbf{Always-true:} \code{if(OPAQUE\_TRUE)\{X\} else \{Y\}} $\to$ \code{X}
\item \textbf{Always-false:} \code{if(OPAQUE\_FALSE)\{X\} else \{Y\}} $\to$ \code{Y} (veya sil)
\end{itemize}


% ===================================================================
\section{String Decryption}
\label{sec:string-decrypt}
% ===================================================================

\file{string\_decryptor.py}, binary'lerdeki sifrelenmis string'leri cozen bir moduld-ur. JavaScript deobfuscation ile dogrudan ilgili olmasa da, Karadul'un genel deobfuscation yeteneginin parcasidir.

\subsection{Desteklenen Sifreleme Yoentemleri}

\begin{enumerate}
\item \textbf{XOR single-byte:} \code{buf[i] \^{}= 0xNN}
\item \textbf{XOR multi-byte:} \code{buf[i] \^{}= key[i \% keylen]}
\item \textbf{XOR rolling:} \code{prev = buf[i] \^{} prev}
\item \textbf{RC4:} S-box initialization + swap pattern
\item \textbf{Base64:} \code{atob("SGVsbG8=")}
\item \textbf{Stack string:} Char-by-char push (Ghidra decompile ciktisinda)
\end{enumerate}

\subsection{Stack String Reconstruction}

Ghidra decompile ciktisinda ardisik local degiskenlere tek karakter atamalari ``stack string'' olarak yeniden birlestirilir:

\begin{lstlisting}[language=Python, caption={Stack string ornegi (Ghidra decompile ciktisi)}]
local_28 = 0x48;  # 'H'
local_27 = 0x65;  # 'e'
local_26 = 0x6c;  # 'l'
local_25 = 0x6c;  # 'l'
local_24 = 0x6f;  # 'o'
# -> "Hello" (confidence: 0.85)
\end{lstlisting}

\subsection{XOR Brute-Force}

Kisa string'ler icin tum 256 tek-byte anahtar denenerek en yuksek ``printable'' skorlu sonuc secilir:

\begin{equation}
\text{score}(k) = \frac{|\{c \in \text{decrypt}(d, k) : c \in \text{printable}\}|}{|\text{decrypt}(d, k)|}
\end{equation}

$\text{score}(k) > 0.7$ olan ilk 5 sonuc dondurulur.


% ===================================================================
\section{Acorn Fallback}
\label{sec:acorn-fallback}
% ===================================================================

Babel parser 32MB+ dosyalarda veya ozel syntax pattern'lerinde basarisiz olabilir. \file{deobfuscate.mjs} bu durumda Acorn parser'a duser.

\subsection{Babel vs Acorn Karsilastirmasi}

\begin{table}[h]
\centering
\small
\begin{tabular}{lcc}
\toprule
\textbf{\color{sectioncyan}Ozellik} & \textbf{\color{sectioncyan}Babel} & \textbf{\color{sectioncyan}Acorn} \\
\midrule
AST formati & Babel (genisletilmis) & ESTree (standart) \\
Error recovery & \checkmark & Kismi \\
TypeScript desteegi & \checkmark (plugin) & $\times$ \\
JSX desteegi & \checkmark (plugin) & $\times$ \\
Performans (30MB) & $\sim$45s & $\sim$8s \\
Memory (30MB) & $\sim$4GB & $\sim$1.5GB \\
\bottomrule
\end{tabular}
\caption{Babel vs Acorn parser karsilastirmasi.}
\label{tab:babel-vs-acorn}
\end{table}

Acorn, \file{deobfuscate.mjs}'te lazy-load edilir ve sadece Babel basarisiz olursa yuklenir:

\begin{lstlisting}[language=JavaScript, caption={Acorn lazy-load}]
let acorn = null, acornWalk = null;
function loadAcorn() {
  try {
    acorn = _require("acorn");
    acornWalk = _require("acorn-walk");
    return true;
  } catch { return false; }
}
\end{lstlisting}


\begin{ozet}
JavaScript deobfuscation, 4 adimli standart zincir (beautify $\to$ synchrony $\to$ babel $\to$ webpack) ve 10 fazli derin transform pipeline'indan olusur. Phase 2--8, visitor-merge optimizasyonu ile tek AST traverse'da birlestirilmistir ($\sim$3--5x hiz kazanci). Graceful degradation her adimda basarisizliga tolerans saglar. Faz 9--10'daki LLM-siz semantik isimlendirme toplam 25+ heuristik kuralla etkileyici sonuclar ueretir. 200MB+ dosyalar icin stream-parse ile chunked processing otomatik devreye girer. CFF deflattening, opaque predicate removal ve string decryption ek moduller olarak pipeline'i guclendirir.
\end{ozet}


\clearpage

% ############################################################################
%
%          B O L U M   6 :   B A Y E S I A N   I S I M   F U Z Y O N U
%
% ############################################################################

% Bu bolum, Karadul'un en yenilikci algoritmasi olan Bayesian log-odds isim
% birlestirme sistemini detayli olarak turetir.
% Kaynak kod: karadul/reconstruction/name_merger.py


% ===================================================================
\section{Problem Tanimi: Coklu Kaynak, Tek Isim}
\label{sec:naming-problem-detail}
% ===================================================================

Tersine muhendislikte, bir sembolun (fonksiyon, degisken) gercek ismi bilinmez. Birden fazla \textbf{bagimsiz} kaynak, farkli isimler onerir. Soru sudur: bu onerilerden hangisi dogru, ve nasil birlestiririz?

\begin{tanim}[Isim Fuzyon Problemi]
Verilen:
\begin{itemize}
\item Bir sembol $s$ (orijinal ismi bilinmiyor, orn. \code{FUN\_00401000})
\item $K$ adet kaynak, her biri bir isim onerisi ve confidence degeri ureten: $(n_i, c_i, \text{source}_i)$ icin $i = 1, \ldots, K$
\end{itemize}

Hedef: $s$ icin en uygun ismi $n^*$ sec ve sonucun dogruluk olasililigini $P^*$ hesapla.
\end{tanim}

\begin{example}
\code{FUN\_00401000} icin:
\begin{itemize}
\item \textbf{signature\_db:} \code{SSL\_connect} (confidence: 0.85, $w=1.0$)
\item \textbf{string\_intel:} \code{ssl\_handshake} (confidence: 0.70, $w=0.8$)
\item \textbf{c\_namer:} \code{connect\_tls} (confidence: 0.60, $w=1.0$)
\item \textbf{llm4decompile:} \code{ssl\_init\_connection} (confidence: 0.55, $w=0.5$)
\end{itemize}
Dort farkli isim! Ama hepsi ``SSL baglantisi'' kavramina isaret ediyor. Bayesian fusion bunu nasil cozer?
\end{example}


% ===================================================================
\section{Bayes Teoremi: Temelden Turetme}
\label{sec:bayes-derivation}
% ===================================================================

\subsection{Kosula Bagli Olasilik}

Iki olay $A$ ve $B$ icin:

\begin{equation}
\prob{A \cap B} = \prob{A \given B} \cdot \prob{B} = \prob{B \given A} \cdot \prob{A}
\label{eq:joint-prob}
\end{equation}

Buradan Bayes teoremi turetilir:

\begin{equation}
\boxed{
\prob{H \given E} = \frac{\prob{E \given H} \cdot \prob{H}}{\prob{E}}
}
\label{eq:bayes-theorem}
\end{equation}

\subsection{Isimlendirme Baglaminda Anlami}

\begin{itemize}
\item $H$: ``Bu sembolun gercek ismi $X$'tir'' hipotezi
\item $E_i$: ``Kaynak $i$, $X$ ismini $c_i$ confidence ile onerdi'' kaniti
\item $\prob{H}$: \textbf{Prior} --- hipotez hakkinda onceki inancimiz (baslangicta $0.5$)
\item $\prob{E_i \given H}$: \textbf{Likelihood} --- eger isim gercekten $X$ ise, kaynak $i$'nin bunu $c_i$ ile onermesinin olasiligi. Bunu $\approx c_i$ olarak modelliyoruz.
\item $\prob{E_i \given \neg H}$: Eger isim $X$ \textbf{degilse}, kaynak $i$'nin yine de $X$ onerme olasiligi $\approx 1 - c_i$.
\item $\prob{H \given E_i}$: \textbf{Posterior} --- kanit sonrasi guncellenminsi inanc
\end{itemize}

\subsection{Ardisik Kanit Birlestirme}

Birden fazla kaniti ardisik olarak birlestirmek icin, her kanitin \textbf{likelihood ratio} (LR) degerini kullaniriz:

\begin{equation}
\text{LR}_i = \frac{\prob{E_i \given H}}{\prob{E_i \given \neg H}} = \frac{c_i}{1 - c_i}
\label{eq:lr}
\end{equation}

Ardisik Bayes guncellemesi:

\begin{equation}
\frac{\prob{H \given E_1, E_2, \ldots, E_n}}{\prob{\neg H \given E_1, E_2, \ldots, E_n}} = \frac{\prob{H}}{\prob{\neg H}} \cdot \prod_{i=1}^{n} \text{LR}_i
\end{equation}

\textbf{Eger kaynaklar bagimsiz ise} (naive Bayes varsayimi).


% ===================================================================
\section{Log-Odds Formu: Neden ve Nasil}
\label{sec:logodds-derivation}
% ===================================================================

\subsection{Odds ve Log-Odds}

Odds (bahis orani):

\begin{equation}
\text{odds}(H) = \frac{\prob{H}}{1 - \prob{H}}
\end{equation}

Log-odds (logit):

\begin{equation}
\text{logit}(H) = \log\frac{\prob{H}}{1 - \prob{H}}
\end{equation}

\subsection{Turetme}

Ardisik Bayes guncellemesinin her iki tarafinin logaritmasini alalim:

\begin{align}
\log\frac{\prob{H \given \mathbf{E}}}{\prob{\neg H \given \mathbf{E}}} &= \log\frac{\prob{H}}{\prob{\neg H}} + \sum_{i=1}^{n} \log \text{LR}_i \\
&= \log\frac{\prob{H}}{1 - \prob{H}} + \sum_{i=1}^{n} \log\frac{c_i}{1 - c_i}
\end{align}

\subsection{Korelasyon Duzeltmesi}

Naive Bayes, kaynakların bagimsiz oldugunu varsayar. Ancak bazilari korelasyonludur (ornegin \code{llm4decompile} tum diger kaynaklari dolayli olarak kullanir). Korelasyonlu kaynaklarin etkisini azaltmak icin agirlik $w_i \in (0, 1]$ eklenir:

\begin{equation}
\boxed{
\text{log-odds}(H) = \underbrace{\log\frac{p_0}{1 - p_0}}_{\text{prior}} + \sum_{i=1}^{n} w_i \cdot \underbrace{\log\frac{c_i}{1 - c_i}}_{\text{log-LR}_i}
}
\label{eq:logodds-weighted}
\end{equation}

burada $p_0 = 0.5$ prior (notr baslangic: $\log(0.5/0.5) = 0$).

\subsection{Neden Log Space?}

\begin{enumerate}
\item \textbf{Numerik stabilite:} Carpim yerine toplam. $\prod_i c_i$ ile $n=10$ kaynak icin floating point underflow olusabilir. Toplam bu sorunu ortadan kaldirir.
\item \textbf{Korelasyon duzeltmesi:} $w_i$ ile dogrudan carpma log space'te lineer agirliklamaya denk gelir. Olasilik uzayinda bu $\text{LR}_i^{w_i}$ demektir --- bilgi-kuramsal olarak ``kaynagin bilgi katkisini $w_i$ oranina indir'' anlamina gelir.
\item \textbf{Sigmoid ile donusum:} Sonuc dogal olarak $[0,1]$ araligina doner.
\end{enumerate}


% ===================================================================
\section{Sigmoid Fonksiyonu: Logit'in Tersi}
\label{sec:sigmoid}
% ===================================================================

Log-odds'u tekrar olasiliga donusturmek icin sigmoid ($\sigma$) fonksiyonunu kullaniriz:

\begin{equation}
\prob{H \given \mathbf{E}} = \sigmoid(\text{log-odds}) = \frac{1}{1 + e^{-\text{log-odds}}}
\label{eq:sigmoid-def}
\end{equation}

\subsection{Sigmoid'in Ozellikleri}

\begin{enumerate}
\item $\sigma(0) = 0.5$ (notr: log-odds = 0 ise emin degiliz)
\item $\lim_{x \to \infty} \sigma(x) = 1$ (cok guclu kanit $\to$ kesinlik)
\item $\lim_{x \to -\infty} \sigma(x) = 0$ (cok guclu karsi kanit $\to$ red)
\item $\sigma'(x) = \sigma(x)(1 - \sigma(x))$ (turevi kendisindendir)
\item $\sigma(-x) = 1 - \sigma(x)$ (simetri)
\end{enumerate}

\begin{center}
\begin{tikzpicture}
\begin{axis}[
  width=10cm, height=5cm,
  domain=-6:6,
  samples=100,
  axis lines=middle,
  xlabel={log-odds},
  ylabel={$\sigma(\text{log-odds})$},
  every axis x label/.style={at=(current axis.right of origin), anchor=west},
  every axis y label/.style={at=(current axis.above origin), anchor=south},
  xtick={-4,-2,0,2,4},
  ytick={0,0.25,0.5,0.75,1.0},
  ymin=-0.05, ymax=1.05,
  grid=major,
  grid style={gray!20!black},
  axis line style={gray!60!white},
  tick label style={color=gray!60!white, font=\scriptsize},
  label style={color=sectioncyan, font=\small},
]
  \addplot[cyan, thick] {1/(1+exp(-x))};
  \addplot[red!60!white, dashed, thick] coordinates {(0,0) (0,0.5)};
  \addplot[red!60!white, dashed, thick] coordinates {(-6,0.5) (0,0.5)};
  \node[color=yellow!70!white, font=\scriptsize] at (axis cs:3, 0.3) {$\sigma(x) = \frac{1}{1+e^{-x}}$};
\end{axis}
\end{tikzpicture}
\end{center}

\subsection{Overflow Korumasi}

\file{name\_merger.py}'da sigmoid icin overflow korumasi uygulanir:

\begin{lstlisting}[language=Python, caption={Sigmoid overflow korumasi -- name\_merger.py}]
if log_odds > 20.0:
    p = 1.0       # exp(-20) ~ 2e-9, 1.0'a cok yakin
elif log_odds < -20.0:
    p = 0.0       # exp(20) ~ 5e8, 0.0'a cok yakin
else:
    p = 1.0 / (1.0 + math.exp(-log_odds))
\end{lstlisting}

Ayrica sonuc $[c_{\min}, c_{\max}] = [0.01, 0.99]$ araligina kliplenir --- asla ``\%100 emin'' veya ``\%0 emin'' demiyoruz.


% ===================================================================
\section{7 Kaynak Detayli Aciklama}
\label{sec:sources}
% ===================================================================

\subsection{binary\_extractor ($w = 1.0$)}

Debug string'ler, RTTI (Run-Time Type Information) ve build path bilgileri derleyici tarafindan binary'ye gomulur. Bu bilgi \textbf{ground truth}'a en yakin kaynaktir:

\begin{itemize}
\item Debug string'ler: \code{\_\_func\_\_} makrosu, assert mesajlari
\item RTTI: C++ sinif isimleri (\code{typeinfo name})
\item Build path: \code{/Users/dev/project/src/ssl\_connect.c}
\end{itemize}

$w = 1.0$ cunku bu bilgi derleyiciden dogrudan gelir --- hicbir tahmin yoktur.

\subsection{c\_namer ($w = 1.0$)}

6 katmanli heuristik isimlendirme motoru:

\begin{enumerate}
\item Fonksiyon imzasi analizi (parametre tipleri, donus tipi)
\item String kullanimi (hata mesajlari, log string'leri)
\item Cagrilan fonksiyonlar (API pattern tespiti)
\item Veri akisi analizi (tanimlama-kullanim zincirleri)
\item Algoritma tespiti (hash, sort, compress pattern'leri)
\item Statistiksel isim onerisi
\end{enumerate}

$w = 1.0$ cunku 6 katman zaten kendi icinde filtrelenmis; cift zayiflatma gereksiz.

\subsection{string\_intel ($w = 0.8$)}

Error mesajlari, protocol handler string'leri ve sabit mesajlardan fonksiyon ismi cikarimi:

\begin{lstlisting}[language=Python]
# Ornek: error mesajindan isim cikarimi
"Failed to connect SSL socket" -> ssl_connect (conf: 0.70)
\end{lstlisting}

$w = 0.8$ cunku error mesajlari, debug string'ler ile kismen ortak kaynaklardan beslenir.

\subsection{signature\_db ($w = 1.0$)}

FLIRT (Fast Library Identification and Recognition Technology) ile 1.5M+ kutuphane fonksiyon imzasi eslestirilir. Bu tamamen \textbf{byte-level} bir eslestirmedir:

\begin{equation}
\text{match}(f) = \argmax_{s \in \text{FLIRT\_DB}} \text{similarity}(\text{bytes}(f), \text{bytes}(s))
\end{equation}

$w = 1.0$ cunku byte pattern tamamen bagimsiz bir bilgi kaynagi.

\subsection{swift\_demangle ($w = 1.0$)}

Swift binary'lerinde mangled sembol isimleri (ornegin \code{\_\$s4main7MyClassC9processSSyS2SF}) demangle edilerek orijinal isim elde edilir:

\begin{equation}
\text{demangle}(\code{\_\$s4main7MyClassC9processSSyS2SF}) = \code{main.MyClass.process(String) -> String}
\end{equation}

$w = 1.0$ cunku demangle \textbf{kesin} bilgidir, tahmin degil.

\subsection{source\_matcher ($w = 0.85$)}

NPM paket fingerprinting: Obfuscated JS kodundaki string ve yapi pattern'leri bilinen NPM paketleri ile eslestirilir:

\begin{itemize}
\item Benzersiz string literal'ler (hata mesajlari, API endpoint'leri)
\item Fonksiyon imza pattern'leri (parametre sayisi + kullanim)
\item Export yapisi (modul cikti formati)
\end{itemize}

$w = 0.85$ cunku heuristik bazli ama yuksek kesinlikli.

\subsection{llm4decompile ($w = 0.5$)}

Makine ogrenmesi modeli ile isim tahmini. Neden $w = 0.5$?

\begin{uyari}[LLM Korelasyon Problemi]
LLM modeli, egitim verisinde tum diger kaynaklarin bilgisini dolayli olarak iciyor:
\begin{itemize}
\item Kutuphane imzalari (signature\_db ile korelasyonlu)
\item Debug string pattern'leri (binary\_extractor ile korelasyonlu)
\item Error mesaj kaliplari (string\_intel ile korelasyonlu)
\end{itemize}
Bu nedenle LLM'in bilgi katkisi ``yarisina'' indirilir: $w = 0.5$.
\end{uyari}

\subsection{byte\_pattern ($w = 1.0$)}

Kutuphane fonksiyonlarinin byte-level fuzzy matching ile tespiti. FLIRT'ten farki: FLIRT kesin imza kullanirken, byte\_pattern fuzzy threshold ile calisir. Yine de tamamen bagimsiz bir byte-level kaynak.


% ===================================================================
\section{Korelasyon-Aware Agirliklar: Bilgi Kurami}
\label{sec:correlation-theory}
% ===================================================================

\subsection{Neden $w_i < 1$?}

Iki kaynak $S_i$ ve $S_j$ arasindaki korelasyonu mutual information ile olceriz:

\begin{equation}
I(S_i; S_j) = \sum_{s_i, s_j} p(s_i, s_j) \log\frac{p(s_i, s_j)}{p(s_i) \cdot p(s_j)}
\label{eq:mutual-info}
\end{equation}

Eger $I(S_i; S_j) > 0$ ise, iki kaynak ortak bilgi tasir. Her ikisini de $w=1.0$ ile saymak bu ortak bilgiyi \textbf{iki kez} sayar ve \textbf{overconfidence} olusturur.

Korelasyon duzeltme prensibimiz:

\begin{equation}
w_i = 1 - \frac{I(S_i; S_{\text{others}})}{H(S_i)}
\label{eq:weight-from-mi}
\end{equation}

burada $H(S_i)$ kaynak $i$'nin entropisi. Pratikte Denklem~\ref{eq:weight-from-mi} yerine ampirik agirliklar kullanilir (Tablo~\ref{tab:weights-detail}).

\subsection{Kaynak Agirliklari Tablosu}

\begin{table}[h]
\centering
\small
\begin{tabular}{llp{6.5cm}}
\toprule
\textbf{\color{sectioncyan}Kaynak} & \textbf{\color{sectioncyan}$w_i$} & \textbf{\color{sectioncyan}Gerekce} \\
\midrule
binary\_extractor & 1.0 & Derleyici ciktisi, tamamen bagimsiz \\
c\_namer & 1.0 & 6 katman kendi icinde filtrelenmis \\
string\_intel & 0.8 & Error msg.~binary\_extractor ile kismi korelasyon \\
signature\_db & 1.0 & FLIRT byte pattern, tamamen bagimsiz \\
swift\_demangle & 1.0 & Deterministik demangle, kesin bilgi \\
source\_matcher & 0.85 & NPM heuristik, bagimsiz ama kesin degil \\
llm4decompile & 0.5 & Tum kaynaklari dolayli kullaniyor \\
byte\_pattern & 1.0 & Byte-level fuzzy match, bagimsiz \\
\midrule
\textit{Bilinmeyen kaynak} & 0.7 & Varsayilan fallback \\
\bottomrule
\end{tabular}
\caption{Kaynak agirliklari, korelasyon analizi ve gerekceler.}
\label{tab:weights-detail}
\end{table}

\subsection{Overconfidence Problemi: Somut Ornek}

Naive Bayes ($w_i = 1.0$ her kaynak icin) ile korelasyon-aware karsilastirma:

\begin{align}
\text{Naive } (w_i = 1): \quad \text{log-odds} &= 0 + 1.0 \cdot \log\frac{0.85}{0.15} + 1.0 \cdot \log\frac{0.70}{0.30} + 1.0 \cdot \log\frac{0.60}{0.40} \\
&= 1.735 + 0.847 + 0.405 = 2.987 \\
P_{\text{naive}} &= \sigma(2.987) = 0.952 \quad (\text{\textbf{\%95.2}})
\end{align}

\begin{align}
\text{Corrected}: \quad \text{log-odds} &= 0 + 1.0 \cdot 1.735 + 0.8 \cdot 0.847 + 0.5 \cdot 0.405 \\
&= 1.735 + 0.678 + 0.203 = 2.616 \\
P_{\text{corr}} &= \sigma(2.616) = 0.932 \quad (\text{\textbf{\%93.2}})
\end{align}

\textbf{Fark:} Naive Bayes \%95.2 derken, korelasyon duzeltmesi \%93.2 der --- \%2 daha muhafazakar ve \textbf{daha dogru}. Cunku llm4decompile'in verdigi bilginin bir kismi zaten diger kaynaklarda mevcuttu.


% ===================================================================
\section{Gruplama Stratejisi: 5 Seviye}
\label{sec:grouping-detail}
% ===================================================================

\file{name\_merger.py}'deki \code{\_merge\_candidates()} fonksiyonu, adaylari 5 seviyeli bir hiyerarsiden gecirir. Her seviye bir oncekinden daha gevşek bir eslestirme kullanir.

\begin{center}
\begin{tikzpicture}[
  level/.style={rectangle, draw=sectioncyan, fill=boxbg, minimum width=5cm, minimum height=0.9cm, text=white, font=\small, rounded corners=2pt, align=center},
  arrow/.style={-{Stealth[length=2.5mm]}, thick, color=gray!60!white},
  fail/.style={-{Stealth[length=2mm]}, dashed, color=red!50!white},
  node distance=0.7cm,
]
  \node[level] (l1) {Seviye 1: Exact Multi-Match};
  \node[level, below=of l1] (l2) {Seviye 2: Partial Match};
  \node[level, below=of l2] (l3) {Seviye 3: Semantic Match};
  \node[level, below=of l3] (l4) {Seviye 4: Voting ($\geq$3)};
  \node[level, below=of l4] (l5) {Seviye 5: Single Source};

  \draw[arrow] (l1) -- node[right, font=\tiny, text=red!50!white] {eslesmezse} (l2);
  \draw[arrow] (l2) -- node[right, font=\tiny, text=red!50!white] {eslesmezse} (l3);
  \draw[arrow] (l3) -- node[right, font=\tiny, text=red!50!white] {eslesmezse} (l4);
  \draw[arrow] (l4) -- node[right, font=\tiny, text=red!50!white] {eslesmezse} (l5);

  \node[right=1.5cm, font=\scriptsize, text=green!60!white] at (l1.east) {En yuksek guven};
  \node[right=1.5cm, font=\scriptsize, text=yellow!70!white] at (l3.east) {Orta guven};
  \node[right=1.5cm, font=\scriptsize, text=red!60!white] at (l5.east) {En dusuk guven};
\end{tikzpicture}
\end{center}

\subsection{Seviye 1: Exact Multi-Match}

Ayni normalize isim $\geq 2$ \textbf{farkli} kaynaktan geldiyse, dogrudan Bayesian merge:

\begin{lstlisting}[language=Python, caption={Exact multi-match -- name\_merger.py}]
for norm_name, group in name_groups.items():
    if len(group) >= 2:
        sources = set(c.source for c in group)
        if len(sources) >= 2:
            conf = bayesian_merge(
                [c.confidence for c in group],
                [c.source for c in group], self._cfg)
            return MergedName(..., merge_method="bayesian_multi")
\end{lstlisting}

\subsection{Seviye 2: Partial Match}

Farkli isimler arasinda prefix/suffix varyanti aranir. Ornegin:
\begin{itemize}
\item \code{cycle\_sizes\_default} vs \code{cycle\_size} $\to$ partial match
\item \code{active\_event\_monitor} vs \code{event\_monitor} $\to$ partial match
\end{itemize}

Overlap skoru:

\begin{equation}
\text{overlap}(A, B) = \begin{cases}
\frac{|A_{\text{short}}|}{|A_{\text{long}}|} & \text{eger } A_{\text{short}} \subset A_{\text{long}} \text{ (substring)} \\[6pt]
J(W_A, W_B) = \frac{|W_A \cap W_B|}{|W_A \cup W_B|} & \text{aksi halde (Jaccard)}
\end{cases}
\label{eq:overlap}
\end{equation}

$\text{overlap} \geq 0.5$ ise partial match kabul edilir. Penalty:

\begin{equation}
\text{penalty} = 0.7 + 0.3 \times \text{overlap}
\end{equation}

Tam substring ($\text{overlap} = 1.0$) $\to$ penalty $= 1.0$ (penalty yok). Kismi overlap ($\text{overlap} = 0.5$) $\to$ penalty $= 0.85$.


\subsection{Seviye 3: Semantic Match}

\file{name\_merger.py}'de 20 adet \code{frozenset} ile kelime gruplari tanimlidir. Iki isim, ayni gruba dusen kelimeleri paylasiyorsa ``semantik olarak benzer'' sayilir.

\begin{table}[h]
\centering
\small
\begin{tabular}{p{3cm}p{9cm}}
\toprule
\textbf{\color{sectioncyan}Grup Adi} & \textbf{\color{sectioncyan}Kelimeler} \\
\midrule
Gonderme & send, transmit, write, emit, dispatch, post \\
Alma & recv, receive, read, get, fetch \\
Olusturma & init, initialize, setup, create, construct, new, alloc \\
Yok etme & destroy, cleanup, teardown, free, delete, release, close, dispose \\
Parse/Decode & parse, decode, deserialize, unmarshal, extract \\
Encode & serialize, encode, marshal, format, stringify \\
Isleyici & handle, process, on, dispatch, callback, handler \\
Dogrulama & check, validate, verify, test, assert, is, has \\
Koleksiyon ekle & add, insert, append, push, enqueue, put \\
Koleksiyon sil & remove, delete, erase, pop, dequeue, take \\
Boyut & size, length, count, num, len, total \\
Tampon & buf, buffer, data, bytes, payload, blob \\
Mesaj & msg, message, packet, frame, request, response \\
Hata & err, error, fault, failure, exception \\
Baglam & ctx, context, state, env, session \\
Konfiguerasyon & cfg, config, configuration, settings, options, prefs \\
Baglanti & conn, connection, socket, link, channel \\
Loglama & log, print, trace, debug, info, warn \\
Kilit alma & lock, acquire, enter, begin, grab \\
Kilit birakma & unlock, release, leave, end, drop \\
\bottomrule
\end{tabular}
\caption{Semantik kelime gruplari (20 frozenset).}
\label{tab:semantic-groups}
\end{table}

\subsubsection{Expand-Then-Intersect Algoritmasi}

\begin{enumerate}
\item Her isim kelimelere ayrilir: \code{ssl\_handshake} $\to$ \{ssl, handshake\}
\item Her kelime, ait oldugu semantic grubun \textbf{tum kelimeleri} ile genisletilir:
      \code{handshake} $\notin$ herhangi bir grup; \code{connect} $\to$ \{conn, connection, socket, link, channel\}
\item Genisletilmis kumelerin kesisimi kontrol edilir
\item Kesisim orani $\geq 0.3$ ise semantik benzerlik var
\end{enumerate}

\begin{lstlisting}[language=Python, caption={Semantic similarity -- name\_merger.py}]
def _is_semantically_similar(self, name1, name2):
    words1 = set(self._normalize(name1).split("_"))
    words2 = set(self._normalize(name2).split("_"))
    # Dogrudan kelime kesisimi
    common = words1 & words2
    if common and len(common) / max(len(words1), len(words2)) >= 0.5:
        return True
    # Semantic group genisletme
    expanded1 = set()
    for w in words1:
        if w in self._word_to_group:
            expanded1 |= self._word_to_group[w]
        expanded1.add(w)
    # ... ayni expanded2 icin ...
    overlap = expanded1 & expanded2
    union = expanded1 | expanded2
    return len(overlap) / len(union) >= 0.3
\end{lstlisting}

Semantik match'te confidence \textbf{0.85x penalty} ile carpilir cunku isimler tam eslesmiyor:

\begin{equation}
c_i^{\text{semantic}} = c_i \times 0.85
\end{equation}


\subsection{Seviye 4: Voting}

$\geq 3$ kaynak ayni normalize ismi veriyorsa, cogunluk kazanir:

\begin{equation}
n^* = \argmax_{n} \sum_{i=1}^{K} \mathbb{1}[\text{normalize}(n_i) = n]
\end{equation}

Eger $\text{count}(n^*) \geq 3$ ise, bu adaylar Bayesian merge'e girer.

\subsection{Seviye 5: Single Source}

Tek kaynak kaldiysa, Bayesian merge gerekmez. Ancak kaynagin $w_i$ deger ile duzeltme yapilir:

\begin{equation}
c_{\text{adjusted}} = \sigma\left(w_i \cdot \log\frac{c_i}{1 - c_i}\right)
\end{equation}

$w_i < 1$ ise bu ``tek kaynaga fazla guvenme'' demektir.


% ===================================================================
\section{Normalize Isim Karsilastirmasi}
\label{sec:normalize}
% ===================================================================

Farkli kaynaklar farkli konvansiyonlarda isim uretir. Karsilastirma oncesi normalizasyon sart:

\begin{lstlisting}[language=Python, caption={Isim normalizasyonu -- name\_merger.py}]
def _normalize(self, name):
    # camelCase -> snake_case
    name = re.sub(r"([a-z])([A-Z])", r"\1_\2", name)
    # Prefix'leri kaldir (m_, s_, g_, p_)
    name = re.sub(r"^[msgp]_", "", name)
    return name.lower().strip("_")
\end{lstlisting}

Ornekler:
\begin{itemize}
\item \code{SSLConnect} $\to$ \code{ssl\_connect}
\item \code{m\_sslConnect} $\to$ \code{ssl\_connect}
\item \code{ssl\_connect} $\to$ \code{ssl\_connect} (degismez)
\item \code{g\_SSLHandler} $\to$ \code{ssl\_handler}
\end{itemize}


% ===================================================================
\section{UNK Threshold: Precision vs Recall}
\label{sec:unk-threshold}
% ===================================================================

\subsection{NSA-Grade Felsefe}

\begin{uyari}[Yanlis Isim Tehlikesi]
Tersine muhendislikte yanlis isim, isimsizden \textbf{cok daha kotu}dur:
\begin{itemize}
\item Yanlis isim $\to$ analist yanlis yolda saatler/gunler harcar
\item UNK (isimsiz) $\to$ analist ``bilmiyorum'' der, dikkatli devam eder
\end{itemize}

\textbf{Sonuc:} Dusuk confidence'li isimleri \textbf{atamayiz}. $c < 0.30$ ise UNK.
\end{uyari}

\subsection{ROC Analizi}

UNK threshold secimi, precision ve recall arasinda trade-off'tur:

\begin{equation}
\text{Precision} = \frac{TP}{TP + FP}, \quad \text{Recall} = \frac{TP}{TP + FN}
\end{equation}

\begin{itemize}
\item \textbf{Threshold $\uparrow$:} Precision $\uparrow$ (daha az yanlis isim), Recall $\downarrow$ (daha cok UNK)
\item \textbf{Threshold $\downarrow$:} Precision $\downarrow$ (daha cok yanlis isim), Recall $\uparrow$ (daha az UNK)
\end{itemize}

Karadul'un secimi: $\theta_{\text{UNK}} = 0.30$. Bu, \textbf{Type I error maliyetinin Type II error maliyetinden cok daha yuksek} olmasi varsayimindan gelir:

\begin{equation}
\text{Cost}_{\text{wrong\_name}} \gg \text{Cost}_{\text{UNK}} \implies \theta_{\text{UNK}} \text{ yuksek tutulur}
\end{equation}


% ===================================================================
\section{Somut Hesaplama Ornekleri}
\label{sec:bayesian-examples}
% ===================================================================

\subsection{Senaryo 1: Uc Kaynak Ayni Isim (Exact Multi-Match)}

\code{FUN\_00401000} icin uc kaynak ayni normalize ismi veriyor:

\begin{center}
\begin{tabular}{llcc}
\toprule
Kaynak & Onerilen Isim & $c_i$ & $w_i$ \\
\midrule
signature\_db & \code{SSL\_connect} & 0.85 & 1.0 \\
string\_intel & \code{SSL\_connect} & 0.70 & 0.8 \\
c\_namer & \code{ssl\_connect} & 0.60 & 1.0 \\
\bottomrule
\end{tabular}
\end{center}

Normalize edilince uc isim de \code{ssl\_connect} --- exact multi-match!

\textbf{Adim 1: Prior log-odds}
\begin{equation}
\ell_0 = \log\frac{0.5}{0.5} = 0
\end{equation}

\textbf{Adim 2: Her kaynağin log-LR katkisi}
\begin{align}
\Delta\ell_1 &= w_1 \cdot \log\frac{c_1}{1-c_1} = 1.0 \cdot \log\frac{0.85}{0.15} = 1.0 \cdot 1.7346 = 1.7346 \\[4pt]
\Delta\ell_2 &= w_2 \cdot \log\frac{c_2}{1-c_2} = 0.8 \cdot \log\frac{0.70}{0.30} = 0.8 \cdot 0.8473 = 0.6779 \\[4pt]
\Delta\ell_3 &= w_3 \cdot \log\frac{c_3}{1-c_3} = 1.0 \cdot \log\frac{0.60}{0.40} = 1.0 \cdot 0.4055 = 0.4055
\end{align}

\textbf{Adim 3: Toplam log-odds}
\begin{equation}
\ell = 0 + 1.7346 + 0.6779 + 0.4055 = 2.8180
\end{equation}

\textbf{Adim 4: Sigmoid}
\begin{equation}
P = \sigma(2.8180) = \frac{1}{1 + e^{-2.8180}} = \frac{1}{1 + 0.05975} = 0.9436
\end{equation}

\begin{basari}[Senaryo 1 Sonuc]
Uc kaynak birlestirildiginde \textbf{\%94.4 guvenle} \code{ssl\_connect} ismi atanir. En iyi tek kaynak \%85'ti --- Bayesian fusion \%9.4 puan ekledi.
\end{basari}


\subsection{Senaryo 2: Iki Kaynak Partial Match}

\code{FUN\_00501234} icin iki kaynak benzer ama farkli isimler veriyor:

\begin{center}
\begin{tabular}{llcc}
\toprule
Kaynak & Onerilen Isim & $c_i$ & $w_i$ \\
\midrule
binary\_extractor & \code{process\_message} & 0.80 & 1.0 \\
c\_namer & \code{handle\_message\_data} & 0.65 & 1.0 \\
\bottomrule
\end{tabular}
\end{center}

Normalize: \code{process\_message} vs \code{handle\_message\_data}

\textbf{Kelime kumesi:} $\{$process, message$\}$ vs $\{$handle, message, data$\}$. Jaccard:

\begin{equation}
J = \frac{|\{\text{message}\}|}{|\{\text{process, message, handle, data}\}|} = \frac{1}{4} = 0.25 < 0.5
\end{equation}

Dogrudan Jaccard $< 0.5$, ama ``handle'' ve ``process'' ayni semantic gruba dusur (Tablo~\ref{tab:semantic-groups}: ``handle, process, on, dispatch, callback, handler''). Semantic expansion sonrasi:

Expanded$_1$ = $\{$process, message, handle, on, dispatch, callback, handler$\}$ \\
Expanded$_2$ = $\{$handle, message, data, process, on, dispatch, callback, handler, bytes, payload, blob, buf, buffer$\}$

\begin{equation}
\frac{|\text{Expanded}_1 \cap \text{Expanded}_2|}{|\text{Expanded}_1 \cup \text{Expanded}_2|} = \frac{7}{13} = 0.538 \geq 0.3
\end{equation}

Semantic match! Penalty $= 0.85$:

\begin{align}
\Delta\ell_1 &= 1.0 \cdot \log\frac{0.80 \times 0.85}{1 - 0.80 \times 0.85} = 1.0 \cdot \log\frac{0.68}{0.32} = 0.7538 \\[4pt]
\Delta\ell_2 &= 1.0 \cdot \log\frac{0.65 \times 0.85}{1 - 0.65 \times 0.85} = 1.0 \cdot \log\frac{0.5525}{0.4475} = 0.2107 \\[4pt]
\ell &= 0 + 0.7538 + 0.2107 = 0.9645 \\[4pt]
P &= \sigma(0.9645) = \frac{1}{1 + e^{-0.9645}} = 0.724
\end{align}

\begin{basari}[Senaryo 2 Sonuc]
Semantic match ile iki kaynak birlestirildiginde \textbf{\%72.4 guvenle} \code{process\_message} (daha kisa isim) atanir. Penalty olmadan \%80 olacakti --- semantik belirsizlik \%7.6 dusurdu.
\end{basari}


\subsection{Senaryo 3: Tek Kaynak, Dusuk Confidence $\to$ UNK}

\code{FUN\_00601000} icin sadece bir kaynak var:

\begin{center}
\begin{tabular}{llcc}
\toprule
Kaynak & Onerilen Isim & $c_i$ & $w_i$ \\
\midrule
llm4decompile & \code{init\_module} & 0.45 & 0.5 \\
\bottomrule
\end{tabular}
\end{center}

Tek kaynak, Single Source (Seviye 5):

\begin{align}
\ell &= w_1 \cdot \log\frac{c_1}{1-c_1} = 0.5 \cdot \log\frac{0.45}{0.55} = 0.5 \cdot (-0.2007) = -0.1003 \\[4pt]
P &= \sigma(-0.1003) = \frac{1}{1 + e^{0.1003}} = 0.4749
\end{align}

$P = 0.475 < \theta_{\text{UNK}} = 0.30$? Hayir, $0.475 > 0.30$.

Ama bekleyin --- burada $P < 0.5$, yani model ``bu ismin dogru olma olasiligi \%50'den az'' diyor. Rasyonel bir analist buna guvenm-ez.

Esik degeri $\theta_{\text{UNK}} = 0.30$ ile $P = 0.475 > 0.30$ oldugu icin \textit{teknik olarak} isim atanir. Ancak pratikte bu dusuk confidence'li isim dikkatle ele alinir.

Simdi daha dusuk confidence ile deneyelim: $c_1 = 0.35$:

\begin{align}
\ell &= 0.5 \cdot \log\frac{0.35}{0.65} = 0.5 \cdot (-0.6190) = -0.3095 \\
P &= \sigma(-0.3095) = 0.4233
\end{align}

$P = 0.423 > 0.30$ --- hala UNK sinirinin uzerinde. Ancak $c_1 = 0.15$:

\begin{align}
\ell &= 0.5 \cdot \log\frac{0.15}{0.85} = 0.5 \cdot (-1.7346) = -0.8673 \\
P &= \sigma(-0.8673) = 0.296
\end{align}

$P = 0.296 < 0.30$ $\to$ \textbf{UNK!} Isim atanmaz.

\begin{uyari}[Senaryo 3 Sonuc]
LLM kaynaginin tek basina confidence $= 0.15$, $w = 0.5$ ile verdigi isim UNK olarak kalir. \textbf{Yanlis isim, isimsizden kotudur.}
\end{uyari}


% ===================================================================
\section{Hesaplama Karmasikligi}
\label{sec:complexity}
% ===================================================================

\subsection{Genel Karmasiklik}

$N$ adet sembol ve her sembol icin ortalama $K$ aday icin:

\begin{equation}
T_{\text{merge}} = \bigO{N \cdot K^2}
\end{equation}

Kare terimi, her aday ciftinin partial/semantic karsilastirmasinden gelir.

\subsection{Alt Islem Karmasikliklari}

\begin{table}[h]
\centering
\small
\begin{tabular}{lll}
\toprule
\textbf{\color{sectioncyan}Islem} & \textbf{\color{sectioncyan}Karmasiklik} & \textbf{\color{sectioncyan}Aciklama} \\
\midrule
Normalizasyon & $\bigO{|n|}$ & Regex + lowercase, string uzunlugu ile lineer \\
Exact match & $\bigO{K}$ & Hash-based gruplama \\
Partial overlap & $\bigO{K^2 \cdot |n|}$ & Tum ciftleri karsilastir, her biri $\bigO{|n|}$ \\
Semantic match & $\bigO{K^2 \cdot |W|}$ & Kelime genisletme + kesisim, $|W|$ = kelime sayisi \\
Bayesian merge & $\bigO{K}$ & Log-odds toplami (lineer) \\
\bottomrule
\end{tabular}
\caption{Alt islem karmasikliklari.}
\label{tab:complexity}
\end{table}

\subsection{Pratikte Performans}

Tipik degerler: $N \approx 5000$ (bir binary'deki fonksiyon sayisi), $K \approx 3$ (ortalama kaynak sayisi):

\begin{equation}
T \approx 5000 \times 3^2 = 45{,}000 \text{ islem}
\end{equation}

Python'da bu $< 0.1$s surer --- bottleneck tamamen decompilation ve statik analiz asamalarindadir.


% ===================================================================
\section{Overconfidence Problemi ve Cozumu}
\label{sec:overconfidence}
% ===================================================================

\subsection{Naive Bayes Hatasi}

Naive Bayes, tum kaynaklarin \textbf{kosula bagli bagimsiz} oldugunu varsayar:

\begin{equation}
\prob{E_1, E_2, \ldots, E_n \given H} = \prod_{i=1}^{n} \prob{E_i \given H}
\end{equation}

Bu varsayim genellikle \textbf{yanlistir}. Eger iki kaynak korelasyonlu ise, ayni bilgiyi iki kez saymis oluruz.

\subsection{Korelasyon ile Duzeltme}

Log-odds formulunde $w_i$ korelasyonu absorbe eder:

\begin{equation}
\text{LR}_i^{\text{effective}} = \text{LR}_i^{w_i}
\end{equation}

$w_i = 1$: Kaynak tamamen bagimsiz $\to$ tam bilgi katkisi. \\
$w_i = 0.5$: Kaynak yari korelasyonlu $\to$ bilgi katkisi ``yarisina'' iner. \\
$w_i = 0$: Kaynak tamamen bagimli $\to$ bilgi katkisi sifir (sayilmaz).

Bilgi-kuramsal olarak bu, log-likelihood ratio'nun eksponent ile olceklenmesidir:

\begin{equation}
\text{effective info gain}_i = w_i \cdot \text{KL}(\prob{\cdot \given E_i} \| \prob{\cdot})
\end{equation}

burada KL divergence bilgi kazancini olcer.

\subsection{Sonuc}

\begin{basari}[Bayesian Fusion Avantajlari]
\begin{enumerate}
\item Coklu kaynaklari sistematik olarak birlestirir --- el ile ``en iyi kaynagi sec'' yerine matematiksel optimal
\item Korelasyonu modelleyerek overconfidence'i onler
\item UNK threshold ile yanlis isim atama riskini minimize eder
\item Log-odds formunda numerik olarak stabil
\item Her adimi tamamen deterministik ve auditlenebilir
\end{enumerate}
\end{basari}


\begin{ozet}
Bayesian isim fuzyonu, 7+ bagimsiz kaynaktan gelen isimleri korelasyon-aware log-odds ile birlestirir. 5 seviyeli gruplama (exact $\to$ partial $\to$ semantic $\to$ voting $\to$ single) her durum icin uygun strateji secer. $w_i < 1$ agirliklari korelasyonlu kaynaklarin etkisini bilgi-kuramsal olarak azaltir. UNK threshold (\%30) NSA-grade felsefe ile yanlis isimlendirmeyi onler. Somut orneklerde: uc kaynak exact match ile \%85'ten \%94.4'e, iki kaynak semantic match ile \%72.4'e yukseltildi; tek dusuk-confidence kaynak ise UNK olarak dogru sekilde reddedildi. Karmasiklik $\bigO{N \cdot K^2}$ ile pratikte milisaniye mertebesinde calisir.
\end{ozet}
 ile dahil edilebilir.
% ONEMLI: \chapter tanimlanmaz, cunku chapter main.tex'te tanimlidir.
%         Burada \section ve \subsection seviyesinde yazilir.
% ============================================================================


% ############################################################################
%
%                    B O L U M   5 :   D E O B F U S C A T I O N
%
% ############################################################################

% Bu bolum, Karadul'un en yenilikci bilesenlerinden birini --- 9 fazli derin
% JavaScript deobfuscation pipeline'ini --- detayli olarak aciklar.
% Kaynak kodlar:
%   - karadul/deobfuscators/manager.py
%   - karadul/deobfuscators/deep_pipeline.py
%   - karadul/deobfuscators/babel_pipeline.py
%   - karadul/deobfuscators/synchrony_wrapper.py
%   - karadul/deobfuscators/string_decryptor.py
%   - karadul/deobfuscators/opaque_predicate.py
%   - karadul/deobfuscators/cff_deflattener.py
%   - scripts/deobfuscate.mjs
%   - scripts/deep-deobfuscate.mjs
%   - scripts/beautify.mjs
%   - scripts/smart-webpack-unpack.mjs
%   - scripts/cursor-enhanced-rename.mjs


% ===================================================================
\section{Obfuscation: Formal Tanim ve Teorik Temel}
\label{sec:obf-formal}
% ===================================================================

\subsection{Semantik Koruma Teoremi}

Kod obfuscation, bir programin kaynak kodunu, davranisini degistirmeden okunamazlastirma islemidir. Bunu formal olarak tanimlayalim.

\begin{tanim}[Obfuscation Fonksiyonu]
Bir obfuscation fonksiyonu $\mathcal{O}: \mathcal{P} \to \mathcal{P}$ asagidaki kosullari saglayan bir donustumdur:
\begin{enumerate}
\item \textbf{Semantik koruma:} Her girdi $x$ icin, $\mathcal{O}(P)(x) = P(x)$
\item \textbf{Polinom yavaslatma:} $|\mathcal{O}(P)| \leq \text{poly}(|P|)$ ve calisma suresi polinomial olarak artar
\item \textbf{Analiz zorlugu:} Herhangi bir analiz $A$ icin, $A(\mathcal{O}(P))$'den elde edilen bilgi, $A(P')$'den ($P'$ erisim saglanmamis bir program) elde edilenden anlamli olarak fazla degildir
\end{enumerate}
\end{tanim}

Bu tanimi daha kompakt yazarsak:

\begin{equation}
\mathcal{O}: \mathcal{P} \to \mathcal{P}, \quad \forall P \in \mathcal{P}: \quad \llbracket \mathcal{O}(P) \rrbracket = \llbracket P \rrbracket
\label{eq:semantic-preservation}
\end{equation}

burada $\llbracket \cdot \rrbracket$ programin denotational semantigini (yani giris-cikis davranisini) gosterir.

\begin{theorem}[Obfuscation Tersini Alma]
Eger $\mathcal{O}$ yalnizca \textbf{syntactic} donusumler uyguluyorsa (string encoding, isim degistirme, control flow yeniden yapilama), o zaman bir ters-obfuscation fonksiyonu $\mathcal{D}$ oyle vardir ki:
\begin{equation}
\mathcal{D}(\mathcal{O}(P)) \approx P
\label{eq:deobf-inverse}
\end{equation}
burada $\approx$ yapisal esdegerlik (structural equivalence) anlamindadir.
\end{theorem}

\textit{Ispat fikri:} Syntactic obfuscation, AST uzerinde terslenebilir donusumler uygular. Her donusum $t_i$ icin bir ters $t_i^{-1}$ mevcuttur. Obfuscation $\mathcal{O} = t_n \circ \cdots \circ t_1$ ise, $\mathcal{D} = t_1^{-1} \circ \cdots \circ t_n^{-1}$ dir. Ancak pratikte $t_i$'ler bilinmez; dolayisiyla Karadul gibi araclar \textbf{heuristik ters donusumleri} uygulamak zorundadir.

\subsection{Obfuscation Karmasiklik Siniflamasi}

Obfuscation tekniklerini bilgi-kuramsal zorluk derecesine gore siniflandiririz:

\begin{equation}
\text{Zorluk}(\mathcal{O}) = -\sum_{i} p_i \log_2 p_i + \text{cfg\_depth}(\mathcal{O}(P))
\label{eq:obf-complexity}
\end{equation}

burada ilk terim Shannon entropisi, ikinci terim control flow graph derinligidir.

\begin{table}[h]
\centering
\begin{tabular}{llcc}
\toprule
\textbf{\color{sectioncyan}Teknik} & \textbf{\color{sectioncyan}Zorluk} & \textbf{\color{sectioncyan}Tersinirlik} & \textbf{\color{sectioncyan}Karadul Kapsami} \\
\midrule
Minification & Dusuk & $\sim$100\% & Beautify \\
String encoding & Dusuk & $\sim$100\% & String unhex \\
Dead code injection & Orta & $\sim$95\% & Dead code elim. \\
String rotation & Orta & $\sim$90\% & Synchrony \\
Opaque predicates & Orta--Yuksek & $\sim$85\% & Pattern + Z3 \\
Control flow flattening & Yuksek & $\sim$70\% & CFF deflattener \\
Virtualization & Cok yuksek & $\sim$30\% & Kismi destek \\
\bottomrule
\end{tabular}
\caption{Obfuscation teknikleri ve zorluk dereceleri.}
\label{tab:obf-techniques}
\end{table}


% ===================================================================
\section{Yaygin Obfuscation Teknikleri}
\label{sec:obf-techniques}
% ===================================================================

\subsection{Minification}

En basit obfuscation formu. Kaynak koddan gereksiz bilgileri (bosluk, yorum, uzun degisken isimleri) kaldirarak dosya boyutunu kucultme ve okunakliligi azaltma islevidir.

\begin{lstlisting}[language=JavaScript, caption={Minification oncesi ve sonrasi}]
// Orijinal (okunabilir):
function calculateTotalPrice(items, taxRate) {
    let subtotal = 0;
    for (const item of items) {
        subtotal += item.price * item.quantity;
    }
    return subtotal * (1 + taxRate);
}

// Minified:
function a(b,c){let d=0;for(const e of b)d+=e.price*e.quantity;return d*(1+c)}
\end{lstlisting}

Minification bilgi kaybina neden olur: degisken isimleri \code{a}, \code{b}, \code{c} orijinal semantigini tasimaz. Ancak \textbf{program davranisi} aynidir --- Denklem~\ref{eq:semantic-preservation} gecerlidir.

\subsection{String Encoding}

String literal'ler cogunlukla iki yontemle gizlenir:

\textbf{Hex encoding:}
\begin{equation}
\text{encode}_{\text{hex}}(s) = \text{concat}_{i=0}^{|s|-1} \texttt{"\textbackslash x"} \cdot \text{hex}(\text{charCodeAt}(s, i))
\end{equation}

\textbf{Unicode encoding:}
\begin{equation}
\text{encode}_{\text{unicode}}(s) = \text{concat}_{i=0}^{|s|-1} \texttt{"\textbackslash u"} \cdot \text{pad}(\text{hex}(\text{codePointAt}(s, i)), 4)
\end{equation}

\textbf{Base64 encoding:}
\begin{equation}
\text{encode}_{\text{b64}}(s) = \texttt{atob("} \cdot \text{base64}(s) \cdot \texttt{")}
\end{equation}

Ornek: \code{"Hello"} $\to$ \code{"\textbackslash x48\textbackslash x65\textbackslash x6c\textbackslash x6c\textbackslash x6f"} $\to$ \code{atob("SGVsbG8=")}

\subsection{Control Flow Flattening (CFF)}

CFF, programin dogal kontrol akisini bir \texttt{while-switch} dispatcher'a donusturur:

\begin{lstlisting}[language=JavaScript, caption={CFF ornegi}]
// Orijinal:
function process(x) {
    let a = x + 1;    // Block 0
    let b = a * 2;    // Block 1
    if (b > 10) {     // Block 2
        return b;     // Block 3
    }
    return a;          // Block 4
}

// CFF uygulanmis:
function process(x) {
    let state = 0, a, b;
    while (true) {
        switch (state) {
            case 0: a = x + 1; state = 1; break;
            case 1: b = a * 2; state = 2; break;
            case 2: state = b > 10 ? 3 : 4; break;
            case 3: return b;
            case 4: return a;
        }
    }
}
\end{lstlisting}

CFF'nin formal modeli bir sonlu durum makinesidir (FSM):

\begin{equation}
\text{CFF}(P) = (Q, \Sigma, \delta, q_0, F)
\end{equation}

burada $Q$ state kuemesi, $\delta: Q \times \Sigma \to Q$ gecis fonksiyonu, $q_0$ baslangic durumu ve $F \subseteq Q$ bitis durumlaridir.

\begin{center}
\begin{tikzpicture}[
  state/.style={circle, draw=sectioncyan, fill=boxbg, minimum size=1.2cm, text=white, font=\small\bfseries},
  arrow/.style={-{Stealth[length=3mm]}, thick, color=gray!70!white},
  every node/.style={font=\small},
]
  \node[state] (q0) at (0,0) {$q_0$};
  \node[state] (q1) at (3,0) {$q_1$};
  \node[state] (q2) at (6,0) {$q_2$};
  \node[state, fill=green!20!black] (q3) at (9,1) {$q_3$};
  \node[state, fill=green!20!black] (q4) at (9,-1) {$q_4$};

  \draw[arrow] (-1.5,0) -- (q0);
  \draw[arrow] (q0) -- node[above, text=gray!60!white] {\texttt{a=x+1}} (q1);
  \draw[arrow] (q1) -- node[above, text=gray!60!white] {\texttt{b=a*2}} (q2);
  \draw[arrow] (q2) -- node[above, text=gray!60!white, sloped] {\texttt{b>10}} (q3);
  \draw[arrow] (q2) -- node[below, text=gray!60!white, sloped] {\texttt{b<=10}} (q4);

  \node[below=0.6cm, font=\tiny, text=cyan!70!white] at (q3.south) {return b};
  \node[below=0.6cm, font=\tiny, text=cyan!70!white] at (q4.south) {return a};
\end{tikzpicture}
\end{center}

\subsection{Opaque Predicates}

Opaque predicate, her zaman ayni sonucu veren (true veya false) ama statik analizin bunu kolayca cikaramayacagi ifadelerdir.

\begin{tanim}[Opaque Predicate]
Bir ifade $\varphi$ icin:
\begin{itemize}
\item $\varphi^T$: Her zaman true ($\forall \sigma: \llbracket \varphi \rrbracket_\sigma = \text{true}$)
\item $\varphi^F$: Her zaman false ($\forall \sigma: \llbracket \varphi \rrbracket_\sigma = \text{false}$)
\item $\varphi^?$: Bazen true, bazen false (gercek dallanma)
\end{itemize}
burada $\sigma$ program state'ini temsil eder.
\end{tanim}

\file{opaque\_predicate.py} dosyasinda tanimlanan oruntuleri inceleyelim:

\begin{table}[h]
\centering
\small
\begin{tabular}{p{5cm}cp{6cm}}
\toprule
\textbf{\color{sectioncyan}Oruuntu} & \textbf{\color{sectioncyan}Sonuc} & \textbf{\color{sectioncyan}Matematik Ispati} \\
\midrule
\code{x*(x+1) \% 2 == 0} & $\varphi^T$ & Ardisik iki sayinin carpimi daima cifttir \\
\code{(x | 1) != 0} & $\varphi^T$ & OR 1 her zaman $\neq 0$ verir \\
\code{(a \& b) == (a \& b)} & $\varphi^T$ & Her ifade kendisiyle esittir \\
\code{2*(x/2) <= x} & $\varphi^T$ & Tamsayi bolumuenun oezelligi: $\lfloor x/2 \rfloor \cdot 2 \leq x$ \\
\code{x*(x+1) \% 2 != 0} & $\varphi^F$ & Yukaridakinin tueuemleri \\
\code{x != x} (int) & $\varphi^F$ & Tamsayi kendisiyle esit (NaN notu asagida) \\
\bottomrule
\end{tabular}
\caption{Opaque predicate oruntuleri ve ispatlar. \textbf{Uyari:} \code{x == x} IEEE 754 NaN icin \texttt{false} dondurur; \file{opaque\_predicate.py} bu nedenle float baglam icin confidence'i 0.60'a dusurur.}
\label{tab:opaque-predicates}
\end{table}

\subsection{Dead Code Injection}

Opaque predicate'ler ile birlesirildiginde, hicbir zaman calistirilmayan sahte kod bloklari eklenir:

\begin{lstlisting}[language=JavaScript, caption={Dead code injection ornegi}]
if (x * (x + 1) % 2 === 0) {
    // Gercek kod (her zaman calisir)
    processPayment(amount);
} else {
    // Dead code (asla calismaz)
    launchMissiles();  // dikkat dagitma amacli
    deleteDatabase();  // analisti yaniltma
}
\end{lstlisting}

\subsection{String Array Rotation}

obfuscator.io gibi araclar tum string'leri bir diziye tasir, diziyi kaydirma (rotation) islemiyle karistirip indeksle erisirir:

\begin{lstlisting}[language=JavaScript, caption={String rotation ornegi}]
// Orijinal:
console.log("Hello");

// Obfuscated:
var _0x1234 = ["Hello", "log", "World"];
(function(_0x, _0y) {
    var _0z = function(_0w) { while(--_0w) { _0x.push(_0x.shift()); } };
    _0z(++_0y);
})(_0x1234, 0x1);  // Diziyi 1 kaydır
console[_0x1234[0]](_0x1234[1]);  // console.log("World") -- indeksler degisti!
\end{lstlisting}

Karadul'un \texttt{synchrony} biliseni bu patterni dogrudan cozer cunku synchrony, obfuscator.io spesifik bir deobfuscation aracidir.


% ===================================================================
\section{Abstract Syntax Tree (AST) Temelleri}
\label{sec:ast-basics}
% ===================================================================

Tum deobfuscation pipeline'i AST uzerinde calisir. Bu bolum AST'nin ne oldugunu, parse tree ile farkini ve ESTree standardini aciklar.

\subsection{BNF ve Context-Free Grammar (CFG)}

JavaScript dili bir context-free grammar (CFG) ile tanimlenir:

\begin{equation}
G = (V, \Sigma, R, S)
\end{equation}

burada $V$ non-terminal semboller, $\Sigma$ terminal semboller, $R$ uretim kurallari ve $S$ baslangic sembolueduer.

JavaScript'in basitlestirilmis BNF'inde \texttt{IfStatement}:

\begin{lstlisting}[basicstyle=\ttfamily\small\color{white}, frame=none]
IfStatement ::= "if" "(" Expression ")" Statement ("else" Statement)?
Expression  ::= AssignmentExpression | Expression "," AssignmentExpression
Statement   ::= Block | IfStatement | ForStatement | ExpressionStatement | ...
\end{lstlisting}

\subsection{Parse Tree vs Abstract Syntax Tree}

\textbf{Parse tree} (concrete syntax tree), gramerin her uretim kuralini icerir --- parantezler, noktalama isaretleri, terminal semboller hepsi agacta gorunur. \textbf{AST} ise semantik olarak anlamsiz node'lari atar.

\begin{center}
\begin{tikzpicture}[
  node distance=0.8cm and 0.5cm,
  treenode/.style={rectangle, rounded corners=2pt, draw=sectioncyan, fill=boxbg, text=white, font=\scriptsize, minimum height=0.6cm, minimum width=1cm},
  leaf/.style={rectangle, rounded corners=2pt, draw=green!60!white, fill=black!80!white, text=green!60!white, font=\scriptsize\ttfamily, minimum height=0.5cm},
  edge from parent/.style={draw=gray!60!white, -{Stealth[length=2mm]}, thin},
  level 1/.style={sibling distance=3.5cm},
  level 2/.style={sibling distance=1.8cm},
  level 3/.style={sibling distance=1.2cm},
]
  % AST tarafı
  \node[treenode, fill=cyan!20!black] {BinaryExpression (+)}
    child { node[treenode] {Identifier}
      child { node[leaf] {a} }
    }
    child { node[treenode] {NumericLiteral}
      child { node[leaf] {1} }
    };

  \node[above=0.3cm, font=\small\color{sectioncyan}] at (0, 1.2) {AST: \texttt{a + 1}};
\end{tikzpicture}
\end{center}

\subsection{ESTree Spesifikasyonu}

JavaScript AST'si icin fiili standart \textbf{ESTree} spesifikasyonudur (\texttt{https://github.com/estree/estree}). Babel parser, ESTree'nin genisletilmis bir versiyonunu uretir:

\begin{itemize}
\item \textbf{Babel AST:} \code{StringLiteral}, \code{NumericLiteral}, \code{ObjectProperty}
\item \textbf{Acorn/ESTree:} \code{Literal}, \code{Property}
\end{itemize}

Bu fark \file{deobfuscate.mjs}'te asagidaki kontrol ile ele alinir:

\begin{lstlisting}[language=JavaScript, caption={Babel vs Acorn AST tespiti (deobfuscate.mjs)}]
const usedAcorn = ast.body !== undefined && !ast.program;
// Babel AST'de ast.program var, Acorn'da yok
\end{lstlisting}

\subsection{Babel Parser Oezellikleri}

\file{deep-deobfuscate.mjs} Babel parser'i su opsiyonlarla cagrilir:

\begin{lstlisting}[language=JavaScript, caption={Babel parser konfigurasyonu}]
ast = parse(source, {
  sourceType: "unambiguous",
  allowImportExportEverywhere: true,
  allowReturnOutsideFunction: true,
  errorRecovery: true,  // Parse hatalarinda devam et
  plugins: ["jsx", "typescript", "decorators-legacy",
    "classProperties", "dynamicImport",
    "optionalChaining", "nullishCoalescingOperator",
    "topLevelAwait", "importMeta"],
});
\end{lstlisting}

\code{errorRecovery: true} kritik bir secimdir. 30MB+ minified dosyalarda parser kismen basarisiz olabilir; bu mod, hatalari toplar ama parse islemini durdurmaz. Boylece kismi AST uzerinde islem yapilabilir.


% ===================================================================
\section{Standart 4 Adimli Deobfuscation Zinciri}
\label{sec:std-chain-detail}
% ===================================================================

\file{manager.py} dosyasindaki \code{DeobfuscationManager} sinifi, standart 4 adimli zinciri orkestre eder:

\begin{center}
\begin{tikzpicture}[
  pstep/.style={rectangle, draw=sectioncyan, fill=boxbg, minimum width=3.2cm, minimum height=1.4cm, text=white, font=\small\bfseries, rounded corners=3pt, align=center},
  arrow/.style={-{Stealth[length=3mm]}, thick, color=gray!70!white},
  fail/.style={-{Stealth[length=2mm]}, dashed, color=red!60!white, thick},
  node distance=1.2cm,
]
  \node[pstep] (input) at (0,0) {Obfuscated\\JS Dosyasi};
  \node[pstep] (beautify) at (4.5,0) {1. Beautify\\(js-beautify)};
  \node[pstep] (synchrony) at (9,0) {2. Synchrony\\(obfuscator.io)};
  \node[pstep] (babel) at (4.5,-3) {3. Babel AST\\(9 faz transform)};
  \node[pstep] (webpack) at (9,-3) {4. Webpack\\Unpack};
  \node[pstep, fill=green!15!black] (output) at (13.5,-3) {Deobfuscated\\Cikti};

  \draw[arrow] (input) -- (beautify);
  \draw[arrow] (beautify) -- (synchrony);
  \draw[arrow] (synchrony) |- (babel);
  \draw[arrow] (babel) -- (webpack);
  \draw[arrow] (webpack) -- (output);

  % Graceful degradation
  \draw[fail] (beautify.south) -- ++(0,-0.8) node[midway, right, font=\tiny, text=red!60!white] {fail} -| (synchrony.south);
  \draw[fail] (synchrony.south) -- ++(0,-0.4) node[midway, right, font=\tiny, text=red!60!white] {fail} -- (babel.north);

  \node[below=0.2cm, font=\tiny, text=gray!50!white] at (beautify.south east) {\file{beautify.mjs}};
  \node[below=0.2cm, font=\tiny, text=gray!50!white] at (synchrony.south east) {\file{synchrony}};
  \node[below=0.2cm, font=\tiny, text=gray!50!white] at (babel.south east) {\file{deobfuscate.mjs}};
  \node[below=0.2cm, font=\tiny, text=gray!50!white] at (webpack.south east) {\file{smart-webpack-unpack.mjs}};
\end{tikzpicture}
\end{center}

\subsection{Graceful Degradation Formuelue}

Her adim basarisiz olabilir. Bu durumda bir onceki basarili cikti sonraki adima gecilir:

\begin{equation}
\text{output}(i) = \begin{cases}
\text{step}_i(\text{output}(i-1)) & \text{eger } \text{step}_i \text{ basarili} \\
\text{output}(i-1) & \text{eger } \text{step}_i \text{ basarisiz (fallback)}
\end{cases}
\label{eq:graceful-degradation}
\end{equation}

Zincirin en az bir adimi basarili olursa, tum zincir ``basarili'' sayilir:

\begin{equation}
\text{success}(\text{chain}) = \bigvee_{i=1}^{n} \text{success}(\text{step}_i)
\label{eq:chain-success}
\end{equation}

\file{manager.py} bunu su kodla implement eder:

\begin{lstlisting}[language=Python, caption={Graceful degradation -- manager.py, run\_chain()}]
for idx, step_name in enumerate(effective_chain, start=1):
    step_output = deob_dir / f"{idx:02d}_{step_name}{input_file.suffix}"
    try:
        step_success = self._execute_step(
            step_name, current_input, step_output, deob_dir)
    except Exception as exc:
        step_success = False
    if step_success and step_output.exists():
        result.steps_completed.append(step_name)
        current_input = step_output  # Sonraki adimin girdisi
    else:
        result.steps_failed.append(step_name)
        # current_input degismez -> onceki basarili cikti kullanilir
\end{lstlisting}

\subsection{Dosya Isimlendirme Konvansiyonu}

Her adim, ciktisini numarali dosya olarak kaydeder:

\begin{center}
\begin{tabular}{ll}
\code{00\_original.js} & Orijinal obfuscated dosya \\
\code{01\_beautify.js} & Beautify sonrasi \\
\code{02\_synchrony.js} & Synchrony sonrasi \\
\code{03\_babel\_transforms.js} & Babel AST sonrasi \\
\code{04\_webpack\_unpack.js} & Webpack modul ayirma sonrasi \\
\end{tabular}
\end{center}

Bu sayede her adimin ciktisi bagimsiz olarak incelenebilir --- debug ve audit icin kritik.


% ===================================================================
\section{Adim 1: Beautify}
\label{sec:beautify}
% ===================================================================

\file{beautify.mjs}, \texttt{js-beautify} kutuphanesini kullanan basit ama kritik bir onisleme adimidir. Minified kod tek satir halindedir; AST parse bile basa\-ri\-siz olabilir cunku Babel'in ternary+arrow+object pattern'lerinde takilma riski vardir.

\subsection{Konfiguerasyon Parametreleri}

\begin{lstlisting}[language=JavaScript, caption={js-beautify konfigurasyonu -- beautify.mjs}]
js_beautify(source, {
  indent_size: 2,
  indent_char: " ",
  max_preserve_newlines: 2,
  brace_style: "collapse,preserve-inline",
  wrap_line_length: 120,
  break_chained_methods: true,
  end_with_newline: true,
});
\end{lstlisting}

\subsection{Etki Analizi}

Beautify'in etkisini olcmek icin genisleme oranini (expansion ratio) tanimlayalim:

\begin{equation}
R_{\text{expand}} = \frac{L_{\text{out}}}{L_{\text{in}}}
\end{equation}

burada $L_{\text{in}}$ girdi satir sayisi, $L_{\text{out}}$ cikti satir sayisi. Tipik degerler:

\begin{itemize}
\item Agresif minification (tek satir): $R_{\text{expand}} \approx 100\text{--}500$
\item Orta minification: $R_{\text{expand}} \approx 2\text{--}5$
\item Zaten formatlenmis: $R_{\text{expand}} \approx 1$
\end{itemize}

\subsection{Fallback Davranisi}

Beautify basarisiz olursa, girdi dosyasi dogrudan ciktiya kopyalanir:

\begin{lstlisting}[language=Python, caption={Beautify fallback -- manager.py}]
def _step_beautify(self, input_file, output_file):
    beautify_script = self._config.scripts_dir / "beautify.mjs"
    if not beautify_script.exists():
        shutil.copy2(input_file, output_file)  # Fallback
        return True
\end{lstlisting}


% ===================================================================
\section{Adim 2: Synchrony Deobfuscation}
\label{sec:synchrony}
% ===================================================================

\file{synchrony\_wrapper.py}, \texttt{synchrony} aracini saran bir Python wrapper'dir. Synchrony, obfuscator.io spesifik bir deobfuscation aracidir ve string rotation, string array, control flow flattening gibi obfuscator.io'ya oezgue teknikleri cozer.

\subsection{Calismaa Mekanizmasi}

\begin{enumerate}
\item \textbf{Shebang soyma:} Eger dosya \code{\#!/usr/bin/env node} ile basliyorsa, synchrony'nin Acorn parser'i bunu parse edemez. Wrapper gecici dosyada shebang'i soyar, islem sonrasi geri ekler.
\item \textbf{Source type deneme:} Synchrony uc modda denenir: \code{both}, \code{module}, \code{script}. ESM dosyalarda \code{ImportDeclaration} hatasi alinirsa \code{module} moduna gecilir.
\item \textbf{Dogrudan dosya ciktisi:} \code{-o} flag'i ile stdout yerine dogrudan dosyaya yazar. Buyuk dosyalarda memory overhead'i onler.
\end{enumerate}

\begin{lstlisting}[language=Python, caption={Synchrony calisma akisi -- synchrony\_wrapper.py}]
source_types = ["both", "module", "script"]
for source_type in source_types:
    cmd = [sync_path, str(actual_input),
           "-o", str(output_file),
           "--type", source_type]
    result = self._runner.run_command(cmd, timeout=timeout)
    if result.success:
        break
    if "ImportDeclaration" in result.stderr:
        continue  # ESM hatasi, module modunu dene
\end{lstlisting}

\subsection{Boyut Analizi}

Synchrony basarili oldugunda dosya boyutu genellikle \textbf{kucueluer} cunku string rotation ve dead code kaldirilir:

\begin{equation}
R_{\text{sync}} = \frac{|\text{output}|}{|\text{input}|} \in [0.3, 1.0]
\end{equation}

$R_{\text{sync}} < 0.5$ ise agresif obfuscation kaldirilmis demektir.


% ===================================================================
\section{9 Fazli Deep Deobfuscation Pipeline}
\label{sec:deep-deobf-detail}
% ===================================================================

\file{deep-deobfuscate.mjs} dosyasi, 10 asamali (9 transform + 1 parse) derin deobfuscation pipeline'i uygular. Phase 2--8, \textbf{visitor-merge optimizasyonu} ile tek bir AST traverse'da birlestirilmistir --- bu webcrack yaklasimi, 7 ayri traverse yerine 1 traverse kullanarak $\sim$3--5x hiz kazanci saglar.

\begin{center}
\begin{tikzpicture}[
  phase/.style={rectangle, draw=sectioncyan, fill=boxbg, minimum width=4cm, minimum height=0.8cm, text=white, font=\scriptsize\bfseries, rounded corners=2pt},
  merged/.style={rectangle, draw=yellow!70!white, fill=boxbg, minimum width=4cm, minimum height=0.8cm, text=yellow!70!white, font=\scriptsize\bfseries, rounded corners=2pt},
  arrow/.style={-{Stealth[length=2mm]}, thin, color=gray!60!white},
  node distance=0.5cm,
]
  \node[phase] (p1) {Phase 1: Parse};
  \node[merged, below=of p1] (p2) {Phase 2: Constant Folding};
  \node[merged, below=of p2] (p3) {Phase 3: Dead Code Elim.};
  \node[merged, below=of p3] (p4) {Phase 4: Computed $\to$ Static};
  \node[merged, below=of p4] (p5) {Phase 5: Comma Split};
  \node[merged, below=of p5] (p6) {Phase 6: Var Decl. Split};
  \node[merged, below=of p6] (p7) {Phase 7: Ternary $\to$ If};
  \node[merged, below=of p7] (p8) {Phase 8: Arrow $\to$ Function};
  \node[phase, below=of p8] (p9) {Phase 9: Smart Rename};
  \node[phase, below=of p9] (p10) {Phase 10: Semantic Rename};

  \draw[arrow] (p1) -- (p2);
  \draw[arrow] (p2) -- (p3);
  \draw[arrow] (p3) -- (p4);
  \draw[arrow] (p4) -- (p5);
  \draw[arrow] (p5) -- (p6);
  \draw[arrow] (p6) -- (p7);
  \draw[arrow] (p7) -- (p8);
  \draw[arrow] (p8) -- (p9);
  \draw[arrow] (p9) -- (p10);

  % Merge bracket
  \draw[yellow!70!white, thick, decorate, decoration={brace, amplitude=8pt, mirror}]
    ([xshift=2.3cm]p2.north east) -- ([xshift=2.3cm]p8.south east)
    node[midway, right=10pt, text=yellow!70!white, font=\scriptsize, align=left] {Tek AST\\traverse\\(visitor-merge)};

  % Separate bracket
  \draw[sectioncyan, thick, decorate, decoration={brace, amplitude=8pt, mirror}]
    ([xshift=2.3cm]p9.north east) -- ([xshift=2.3cm]p10.south east)
    node[midway, right=10pt, text=sectioncyan, font=\scriptsize, align=left] {Ayri\\traverse\\(scope gerekli)};
\end{tikzpicture}
\end{center}

\subsection{Faz 1: Parse (Tolerant Mode)}
\label{sec:phase1}

Babel parser \code{errorRecovery: true} ile calistirilir. Basarisiz olursa iki fallback mevcuttur:

\begin{enumerate}
\item \textbf{Pre-beautify + retry:} Minified kaynak once \code{js-beautify} ile formatlanir, sonra Babel tekrar denenir. 30MB+ dosyalarda Babel'in ternary/arrow/object pattern'lerinde takilma sorunu bu yontemle cogunlukla cozulur.
\item \textbf{Dosya kopyalama:} Her iki yontem de basarisiz olursa, kaynak dosya oldugu gibi output'a kopyalanir. Boylece sonraki pipeline adimlari (webpack unpack) en azindan orijinal dosya uzerinde devam edebilir.
\end{enumerate}

\begin{equation}
\text{parse}(S) = \begin{cases}
\text{Babel}(S) & \text{basarili} \\
\text{Babel}(\text{beautify}(S)) & \text{pre-beautify ile retry} \\
S & \text{fallback: kaynak kopyalama}
\end{cases}
\end{equation}


\subsection{Faz 2: Constant Folding (Sabit Katlama)}
\label{sec:phase2-cf}

Constant folding, derleme/analiz zamaninda hesaplanabilir ifadeleri literal degerleriyle degistirme islemidir.

\subsubsection{Lattice Teorisi ve Abstract Interpretation}

Constant folding, abstract interpretation cercevesinde bir \textbf{flat lattice} uzerinde calisir:

\begin{equation}
L = \{\bot, c_1, c_2, \ldots, c_n, \top\}
\end{equation}

burada:
\begin{itemize}
\item $\bot$: Degisken henuz kullanilmadi (undefined)
\item $c_i$: Degisken sabit $c_i$ degerinde
\item $\top$: Degisken birden fazla deger alabilir (non-constant)
\end{itemize}

\begin{equation}
\text{meet}(x, y) = \begin{cases}
x & \text{eger } x = y \\
\top & \text{eger } x \neq y \text{ ve her ikisi de } \neq \bot \\
y & \text{eger } x = \bot \\
x & \text{eger } y = \bot
\end{cases}
\end{equation}

\subsubsection{Karadul Implementasyonu}

\file{deep-deobfuscate.mjs} Phase 2'de asagidaki donusumler uygulanir:

\begin{table}[h]
\centering
\small
\begin{tabular}{lll}
\toprule
\textbf{\color{sectioncyan}Node Tipi} & \textbf{\color{sectioncyan}Donusum} & \textbf{\color{sectioncyan}Ornek} \\
\midrule
BinaryExpression (+) & String concat & \code{"He"+"llo"} $\to$ \code{"Hello"} \\
BinaryExpression (+) & Numeric add & \code{1+2} $\to$ \code{3} \\
BinaryExpression ($*$, $/$, $-$) & Aritmetik & \code{6*7} $\to$ \code{42} \\
BinaryExpression ($|$,$\&$,$\hat{}$,$\ll$,$\gg$) & Bitwise & \code{0xFF \& 0x0F} $\to$ \code{15} \\
UnaryExpression (!) & Negation & \code{!0} $\to$ \code{true} \\
UnaryExpression (void) & Void & \code{void 0} $\to$ \code{undefined} \\
\bottomrule
\end{tabular}
\caption{Constant folding donusumm tablosu.}
\label{tab:constant-folding}
\end{table}

Guevenlik kontrolu: Bolme islemlerinde \code{right.value !== 0} kontrol edilir ve $\text{Number.isFinite}(result)$ dogrulanir. Overfllow, NaN ve Infinity sonuclari degistirilmez.

\begin{lstlisting}[language=JavaScript, caption={Constant folding visitor -- BinaryExpression}]
BinaryExpression: {
  exit(path) {
    const { left, right, operator } = path.node;
    // String concatenation
    if (operator === "+" && t.isStringLiteral(left)
        && t.isStringLiteral(right)) {
      path.replaceWith(t.stringLiteral(left.value + right.value));
      cnt2++;
    }
    // Numeric operations
    if (t.isNumericLiteral(left) && t.isNumericLiteral(right)) {
      let result;
      switch (operator) {
        case "+": result = left.value + right.value; break;
        case "-": result = left.value - right.value; break;
        // ... (*, /, %, **, |, &, ^, <<, >>, >>>)
      }
      if (result !== undefined && Number.isFinite(result)) {
        path.replaceWith(t.numericLiteral(result));
        cnt2++;
      }
    }
  }
}
\end{lstlisting}

\textbf{Neden \code{exit} visitor?} Bottom-up traversal sayesinde ic node'lar once islenir: \code{(1+2)*3}'te once \code{1+2=3}, sonra \code{3*3=9}. \code{enter} kullansaydik, ic node'lar henuz fold edilmemis olurdu.


\subsection{Faz 3: Dead Code Elimination}
\label{sec:phase3-dce}

\subsubsection{Liveness Analysis Temeli}

Liveness analysis, her program noktasinda hangi degiskenlerin ``canli'' (daha sonra kullanilacak) oldugunu belirler:

\begin{equation}
\text{live\_in}(B) = \text{gen}(B) \cup (\text{live\_out}(B) \setminus \text{kill}(B))
\label{eq:live-in}
\end{equation}
\begin{equation}
\text{live\_out}(B) = \bigcup_{S \in \text{succ}(B)} \text{live\_in}(S)
\label{eq:live-out}
\end{equation}

burada:
\begin{itemize}
\item $\text{gen}(B)$: $B$ blogunda kullanilan (okunan) degiskenler
\item $\text{kill}(B)$: $B$ blogunda tanimlanan (yazilan) degiskenler
\end{itemize}

\subsubsection{Fixed-Point Iteration}

Liveness analysis, denklem sistemi sabit noktaya ulasana kadar iteratif olarak cozulur:

\begin{equation}
\text{live\_in}^{(k+1)}(B) = \text{gen}(B) \cup \left(\text{live\_out}^{(k)}(B) \setminus \text{kill}(B)\right)
\end{equation}

Karmasiklik: $\bigO{|B| \times |V|}$ burada $|B|$ basic block sayisi, $|V|$ degisken sayisi.

\subsubsection{Karadul Implementasyonu}

Phase 3, \code{IfStatement} ve \code{ConditionalExpression} node'larini isle\-r:

\begin{lstlisting}[language=JavaScript, caption={Dead code elimination -- IfStatement}]
IfStatement: {
  exit(path) {
    const test = path.node.test;
    // if (false) {...} -> kaldir veya else branch'i al
    if (t.isBooleanLiteral(test) && test.value === false) {
      if (path.node.alternate) path.replaceWith(path.node.alternate);
      else path.remove();
      cnt3++;
    }
    // if (true) {X} else {Y} -> X
    if (t.isBooleanLiteral(test) && test.value === true) {
      path.replaceWith(path.node.consequent);
      cnt3++;
    }
    // if (0) -> remove, if (1) -> consequent
    if (t.isNumericLiteral(test)) {
      if (test.value === 0) { /* remove or alternate */ }
      else { path.replaceWith(path.node.consequent); }
      cnt3++;
    }
  }
}
\end{lstlisting}

\textbf{Phase 2 ile etkilesim:} Phase 2 \code{!0 $\to$ true} donusumunu yapar. Phase 3, \code{exit} visitor kullandigi icin, ayni traverse'da Phase 2'nin ciktisini gorebilir. Ornegin \code{if(!0)\{X\}} once \code{if(true)\{X\}} olur (Phase 2), sonra \code{X} olur (Phase 3).


\subsection{Faz 4: Computed-to-Static Property}
\label{sec:phase4-cts}

Bu faz, \code{obj["foo"]} gibi computed property erisimlerini \code{obj.foo} gibi statik erisimlere donusturur.

\begin{equation}
\text{obj}[\texttt{"key"}] \xrightarrow{\text{Faz 4}} \text{obj}.\texttt{key}
\quad \text{eger } \texttt{key} \in \text{ValidIdentifier}
\end{equation}

\begin{tanim}[Valid Identifier]
Bir string \code{key} gecerli JavaScript identifier ise ve reserved word degilse donusum guevenlidir:
\begin{equation}
\text{ValidIdentifier}(\texttt{key}) = \texttt{/\^{}[a-zA-Z\_\$][a-zA-Z0-9\_\$]*\$/}.test(\texttt{key}) \wedge \neg\text{isReserved}(\texttt{key})
\end{equation}
\end{tanim}

AST duenuesuemuee:
\begin{itemize}
\item \code{MemberExpression}: \code{computed: true} $\to$ \code{computed: false}
\item \code{ObjectProperty}: Computed key'i identifier'a donustur
\end{itemize}

\begin{lstlisting}[language=JavaScript, caption={Computed to static -- MemberExpression}]
MemberExpression(path) {
  if (!path.node.computed) return;
  const prop = path.node.property;
  if (!t.isStringLiteral(prop)) return;
  if (/^[a-zA-Z_$][a-zA-Z0-9_$]*$/.test(prop.value)
      && !isReservedWord(prop.value)) {
    path.node.computed = false;
    path.node.property = t.identifier(prop.value);
    cnt4++;
  }
}
\end{lstlisting}


\subsection{Faz 5: Comma Expression Splitting}
\label{sec:phase5-comma}

JavaScript'te virgul operatoru (\code{SequenceExpression}) birden fazla ifadeyi tek bir expression'da birlestirir. Obfuscator'lar bunu okunakliligi azaltmak icin kullanir:

\begin{lstlisting}[language=JavaScript]
// Obfuscated:
return a = 1, b = 2, c = 3, process(a, b, c);
// Deobfuscated (Faz 5 sonrasi):
a = 1;
b = 2;
c = 3;
return process(a, b, c);
\end{lstlisting}

Donusum ucg bag\-lam icin uygulanir:

\begin{enumerate}
\item \textbf{ExpressionStatement:} \code{(a, b, c);} $\to$ \code{a; b; c;}
\item \textbf{ReturnStatement:} \code{return (a, b, c);} $\to$ \code{a; b; return c;}
\item \textbf{ThrowStatement:} \code{throw (a, b, c);} $\to$ \code{a; b; throw c;}
\end{enumerate}

\begin{equation}
\text{split}(\texttt{return}(e_1, e_2, \ldots, e_n)) = \left[\text{expr}(e_1), \ldots, \text{expr}(e_{n-1}), \texttt{return}\; e_n\right]
\end{equation}


\subsection{Faz 6: Variable Declaration Splitting}
\label{sec:phase6-vardecl}

Tek bir \code{var/let/const} icindeki birden fazla tanimlama ayri satirlara bolunur:

\begin{lstlisting}[language=JavaScript]
// Oncesi:
var a = 1, b = 2, c = fn(), d = [];
// Sonrasi:
var a = 1;
var b = 2;
var c = fn();
var d = [];
\end{lstlisting}

\textbf{Guvenlik kontrolleri:} For-loop init (\code{for(var i=0, j=0; ...)}), for-in ve for-of icerisindeki deklarasyonlar bolunmez:

\begin{lstlisting}[language=JavaScript, caption={Variable declaration split -- guvenlik kontrolleri}]
if (path.parent.type === "ForStatement" && path.key === "init") return;
if (path.parent.type === "ForInStatement") return;
if (path.parent.type === "ForOfStatement") return;
\end{lstlisting}


\subsection{Faz 7: Ternary $\to$ If/Else}
\label{sec:phase7-ternary}

Uzun ternary ifadeler if/else bloklarina donusturulur. ``Uzun'' tanimi: test + consequent + alternate toplam karakter sayisi $\geq 80$:

\begin{equation}
|\text{code}(\text{test})| + |\text{code}(\text{cons})| + |\text{code}(\text{alt})| \geq 80
\end{equation}

\begin{lstlisting}[language=JavaScript]
// Oncesi (tek satir, okunamaz):
isAuthenticated ? renderDashboard(user, session) : redirectToLogin(returnUrl);
// Sonrasi:
if (isAuthenticated) {
  renderDashboard(user, session);
} else {
  redirectToLogin(returnUrl);
}
\end{lstlisting}


\subsection{Faz 8: Arrow $\to$ Named Function}
\label{sec:phase8-arrow}

\code{const foo = () => \{...\}} formundaki arrow function'lar named function declaration'a donusturulur:

\begin{lstlisting}[language=JavaScript]
// Oncesi:
const processData = (items) => { return items.map(x => x * 2); };
// Sonrasi:
function processData(items) { return items.map(x => x * 2); }
\end{lstlisting}

\textbf{Semantik guevenlik:} \code{this} kullanan arrow function'lar donusturulmez. Arrow function'larda \code{this} lexical scope'a baglidir; normal function'larda ise call-site'a baglidir. Bu semantik fark bozulmamalidir:

\begin{lstlisting}[language=JavaScript, caption={Arrow-to-function this kontrolu}]
const { code: bodyCode } = generate(arrow.body, { compact: true });
if (!bodyCode.includes("this.") && !bodyCode.includes("this[")) {
    // Guvenli: donustur
    const funcDecl = t.functionDeclaration(
      t.identifier(decl.id.name), arrow.params, arrow.body,
      arrow.generator || false, arrow.async || false);
    path.replaceWith(funcDecl);
}
\end{lstlisting}


\subsection{Faz 9: Smart Variable Renaming}
\label{sec:phase9-rename}

Bu faz, minified identifier'lari (1--2 karakter) anlamli isimlere donusturur. LLM \textbf{kullanmaz} --- tamamen heuristik kurallara dayanir. Babel \code{scope.rename()} API'si ile scope-safe renaming yapilir.

\subsubsection{Uec Kural Kategorisi}

\begin{enumerate}
\item \textbf{Module-based:} \code{require("fs")} $\to$ degisken \code{fs} olarak adlandirilir
\item \textbf{Usage-based:} Degisken uzerinde \code{.push()}, \code{.map()} cagriliyorsa $\to$ \code{items}
\item \textbf{Fallback:} Tek harf degiskenler icin varsayilan isimler (\code{i} $\to$ \code{idx}, \code{e} $\to$ \code{elem})
\end{enumerate}

\begin{table}[h]
\centering
\small
\begin{tabular}{p{3.5cm}p{4cm}p{4.5cm}}
\toprule
\textbf{\color{sectioncyan}Pattern} & \textbf{\color{sectioncyan}Tespit Yontemi} & \textbf{\color{sectioncyan}Yeni Isim} \\
\midrule
\code{require("fs")} & CallExpression analizi & \code{fs} \\
\code{require("path")} & CallExpression analizi & \code{path} \\
\code{require("express")} & CallExpression analizi & \code{express} \\
\code{catch(e)} & CatchClause param & \code{error} \\
\code{x.push()}, \code{x.map()} & Usage (Array API) & \code{items} \\
\code{x.trim()}, \code{x.split()} & Usage (String API) & \code{text} \\
\code{x.message}, \code{x.stack} & Property access & \code{err} \\
\code{x.then()}, \code{x.catch()} & Usage (Promise API) & \code{promise} \\
Express \code{(req, res, next)} & Paramettre pozisyonu & \code{request, response, next} \\
esbuild \code{(exports, module)} & Factory pattern & \code{exports, module} \\
\bottomrule
\end{tabular}
\caption{Phase 9 smart rename kural ornekleri.}
\label{tab:phase9-rules}
\end{table}

\subsubsection{Tek Harf Fallback Tablosu}

\begin{lstlisting}[language=JavaScript, caption={Tek harf fallback mapping}]
const defaults = {
  a: "arg", b: "buf", c: "ctx", d: "data",
  e: "elem", f: "fn", g: "gen", h: "handler",
  i: "idx", j: "jdx", k: "key", l: "len",
  m: "mod", n: "count", o: "opts", p: "param",
  q: "query", r: "res", s: "str", t: "tmp",
  u: "usr", v: "val", w: "writer",
};
\end{lstlisting}


\subsection{Faz 10: Semantic Renaming (Deep Usage-Based)}
\label{sec:phase10-semantic}

Phase 10, Phase 9'un basit pattern matching'ini asarak, Babel scope + binding analizi ile her minified identifier'in kullanim baglamini derinlemesine analiz eder. \textbf{12 kural motoru} ile calisir.

\subsubsection{Kullanim Bilgisi Toplama (\_collectUsage10)}

Her identifier icin asagidaki bilgiler toplanir:

\begin{lstlisting}[language=JavaScript, caption={Usage bilgi yapisi}]
{
  methodCalls: Set,       // x.push(), x.split() gibi
  propertyReads: Set,     // x.length, x.name gibi
  propertyWrites: Set,    // x.count = 5 gibi
  calledAs: boolean,      // x() seklinde cagriliyor mu
  comparedWith: string[], // x === "GET" gibi karsilastirmalar
  typeofChecks: string[], // typeof x === "function" gibi
  usedInForOf: boolean,   // for(const y of x) iteratorde mi
  destructured: boolean,  // const {a, b} = x seklinde mi
  indexAccessed: boolean,  // x[0] seklinde mi
  arithmeticOperand: boolean, // x + 1 gibi aritmetikte mi
  thrownAsError: boolean, // throw x seklinde mi
  usedInNew: boolean,     // new x() seklinde mi
}
\end{lstlisting}

\subsubsection{12 Kural Motoru}

\begin{table}[h]
\centering
\small
\begin{tabular}{clcp{4.5cm}}
\toprule
\textbf{\color{sectioncyan}\#} & \textbf{\color{sectioncyan}Kural} & \textbf{\color{sectioncyan}Conf.} & \textbf{\color{sectioncyan}Tespit Kriteri} \\
\midrule
1 & \code{require\_module} & 0.95 & \code{require("fs")} $\to$ \code{fs} \\
2 & \code{new\_error} & 0.90 & \code{new Error()} pattern \\
3 & \code{error\_props} & 0.90 & \code{.message} + \code{.stack} props \\
4 & \code{string\_api} & 0.85 & $\geq$2 string method \\
5 & \code{array\_api} & 0.85 & $\geq$2 array method \\
6 & \code{promise\_api} & 0.85 & \code{.then()} + \code{.catch()} \\
7 & \code{fs\_api} & 0.90 & \code{readFileSync}, \code{writeFile} \\
8 & \code{http\_req} & 0.80 & $\geq$2 of \code{body,params,query,headers} \\
9 & \code{dom\_api} & 0.85 & $\geq$2 DOM property \\
10 & \code{callback} & 0.80 & \code{calledAs} + \code{typeof === "function"} \\
11 & \code{for\_of\_item} & 0.80 & \code{for(const x of iterable)} \\
12 & \code{multi\_compare} & 0.70 & $\geq$3 string karsilastirmasi \\
\bottomrule
\end{tabular}
\caption{Phase 10 semantic renaming kural motoru.}
\label{tab:phase10-rules}
\end{table}

\subsubsection{Phase 10+: Enhanced Rename (cursor-enhanced-rename.mjs)}

\file{cursor-enhanced-rename.mjs}, Phase 10'un kapsamadigi 15 ek kural ile kalan minified identifier'lari isle\-r. \file{deep\_pipeline.py}'da Adim 2.5 olarak calistirilir:

\begin{table}[h]
\centering
\small
\begin{tabular}{clp{5.5cm}}
\toprule
\textbf{\color{sectioncyan}\#} & \textbf{\color{sectioncyan}Kural} & \textbf{\color{sectioncyan}Aciklama} \\
\midrule
11 & Prototype access & \code{x.prototype} $\to$ \code{cls} \\
12 & Binary ops dominant & Cok fazla \code{+}, \code{-}, \code{<<} $\to$ \code{offset} \\
13 & Purely called & Sadece cagriliyor $\to$ \code{handler} \\
14 & Event listener pos & Callback pozisyonu $\to$ \code{eventArg} \\
15 & Member assign & Atama kaynagi $\to$ \code{ref} \\
16 & Single property & Tek prop baskini $\to$ contextual \\
17 & Conditional return & \code{return x ? a : b} $\to$ \code{result} \\
18 & Switch/if chain & Karsilastirma zinciri $\to$ \code{kind} \\
19 & Namespace access & \code{X.C}, \code{X.Q} $\to$ \code{ns} \\
20 & Loop counter & \code{for} init $\to$ \code{idx} \\
21 & Await expression & \code{await x} $\to$ \code{awaited} \\
22 & Default param ctx & Varsayilan deger $\to$ \code{optParam} \\
23 & Clash dedup & Phase 9 artigi $\to$ re-score \\
24 & Length/size dominant & \code{.length} baskini $\to$ \code{collection} \\
25 & Bitwise dominant & \code{|}, \code{\&}, \code{\^{}} $\to$ \code{flags} \\
\bottomrule
\end{tabular}
\caption{Enhanced rename ek kurallari (Kural 11--25).}
\label{tab:enhanced-rules}
\end{table}


% ===================================================================
\section{Webpack Module Extraction}
\label{sec:webpack-unpack}
% ===================================================================

\file{smart-webpack-unpack.mjs}, webpack ve esbuild bundle'larini modullerine ayiran akilli bir module splitter'dir. 7 farkli bundle formatini destekler.

\subsection{Desteklenen Bundle Formatlari}

\begin{enumerate}
\item \textbf{esbuild CJS wrapper:} \code{var X = U((exports, module) => \{...\})}
\item \textbf{esbuild ESM interop:} \code{var Y = Q1(X)}
\item \textbf{esbuild re-export:} \code{lb(target, \{name: () => ref\})}
\item \textbf{Webpack IIFE object:} \code{(function(modules) \{...\})(\{0: function(e,t,n)\{...\}\})}
\item \textbf{Webpack IIFE array:} \code{[function(e,t,n)\{...\}, ...]}
\item \textbf{Webpack 5 named:} \code{\_\_webpack\_modules\_\_ = \{"./src/file.js": ...\}}
\item \textbf{Top-level IIFE:} \code{(function() \{...\})()} --- kapsayici IIFE'leri ayirma
\end{enumerate}

\subsection{\_\_webpack\_require\_\_ Mekanizmasi}

Webpack bundle'larinda moduller \code{\_\_webpack\_require\_\_(id)} ile birbirini cagrilir. Bu pattern AST'de tespit edilir:

\begin{lstlisting}[language=JavaScript, caption={Webpack module tespiti -- deobfuscate.mjs}]
if (callee.name === "__webpack_require__"
    && node.arguments.length >= 1) {
  const arg = node.arguments[0];
  if (arg.type === "NumericLiteral") moduleId = arg.value;
  else if (arg.type === "StringLiteral") moduleId = arg.value;
  if (moduleId !== null && !_webpackIdSet.has(moduleId)) {
    _webpackIdSet.add(moduleId);
    webpackModules.push(moduleId);
  }
}
\end{lstlisting}

\subsection{esbuild vs Webpack Format Farklari}

\begin{table}[h]
\centering
\small
\begin{tabular}{p{3.5cm}p{5cm}p{5cm}}
\toprule
\textbf{\color{sectioncyan}Ozellik} & \textbf{\color{sectioncyan}Webpack} & \textbf{\color{sectioncyan}esbuild} \\
\midrule
Modul sarmalama & IIFE + nesne/dizi & CJS factory wrapper \\
Modul ID & Sayisal (0, 1, 2...) & Degisken adi \\
Dependency & \code{\_\_webpack\_require\_\_} & Dogrudan degisken referansi \\
Entry point & Ilk cagri arg. & Top-level kod \\
Tree shaking & Plugin bazli & Varsayilan \\
\bottomrule
\end{tabular}
\caption{Webpack vs esbuild bundle farklari.}
\label{tab:webpack-vs-esbuild}
\end{table}


% ===================================================================
\section{Buyuk Dosya Isleme: Chunked Processing}
\label{sec:chunked}
% ===================================================================

200MB+ dosyalar icin standart AST parsing bellege sigmaz. \file{deep\_pipeline.py} otomatik olarak \textbf{chunked processing} moduna gecer.

\subsection{Esik Degeri ve Tetikleme}

\begin{equation}
\text{mode} = \begin{cases}
\text{normal} & \text{eger dosya} < 100\text{ MB} \\
\text{chunked} & \text{eger dosya} \geq 100\text{ MB}
\end{cases}
\end{equation}

\file{deep\_pipeline.py}'da \code{LARGE\_FILE\_THRESHOLD\_MB = 100} olarak tanimli.

\subsection{Stream-Parse Pipeline}

\begin{enumerate}
\item \textbf{stream-parse.mjs:} Dosyayi top-level block'lara ayirir (\code{--max-block-kb 512})
\item \textbf{Her block'u bagimsiz isle:} \code{deep-deobfuscate.mjs} her chunk uzerinde ayri calisir (120s timeout/block)
\item \textbf{Birlestir:} Tum chunk ciktilari tek dosyada birlestirilir
\item \textbf{Webpack unpack:} Birlestirilmis dosya uzerinde modul ayirma
\item \textbf{Ikinci pass:} Her modul uzerinde variable renaming
\end{enumerate}

\begin{equation}
T_{\text{chunked}} \approx T_{\text{split}} + n \cdot T_{\text{deob/chunk}} + T_{\text{merge}} + T_{\text{unpack}}
\end{equation}

burada $n$ chunk sayisi, $T_{\text{deob/chunk}}$ chunk basina deobfuscation suresi.

\subsection{Timeout Hesaplaması}

\file{deep\_pipeline.py} dosya boyutuna gore dinamik timeout hesaplar:

\begin{equation}
T_{\text{timeout}} = \max(300, \min(1800, \lceil S_{\text{MB}} \times 100 \rceil)) \text{ saniye}
\label{eq:timeout}
\end{equation}

burada $S_{\text{MB}}$ dosya boyutu megabayt cinsindendir. Minimum 5 dakika, maksimum 30 dakika.


% ===================================================================
\section{Control Flow Flattening (CFF) Deflattening}
\label{sec:cff-deflattening}
% ===================================================================

\file{cff\_deflattener.py}, CFF obfuscation'ini geri alan bir Python modul-udur. Hem C hem JavaScript uzerinde calisir.

\subsection{Tespit Algoritması}

CFF dispatcher pattern regex ile tespit edilir:

\begin{lstlisting}[language=Python, caption={CFF dispatcher tespiti}]
# JS CFF pattern
_WHILE_SWITCH_JS = re.compile(
    r"(?:while\s*\(\s*(?:!0|true|1)\s*\)|for\s*\(\s*;;\s*\))\s*\{"
    r"\s*switch\s*\(\s*(\w+)\s*\)\s*\{",
    re.IGNORECASE,
)
\end{lstlisting}

\subsection{State Transition Graph Cikarma}

Her \code{case} blogundan state gecisleri cikarilarak bir yoenluee graf (directed graph) olusturulur:

\begin{equation}
G = (V, E) \quad \text{burada } V = \{s_0, s_1, \ldots, s_n\}, \quad E = \{(s_i, s_j) : s_i \xrightarrow{\text{state}} s_j\}
\end{equation}

\begin{center}
\begin{tikzpicture}[
  state/.style={circle, draw=sectioncyan, fill=boxbg, minimum size=0.9cm, text=white, font=\small},
  arrow/.style={-{Stealth[length=2.5mm]}, thick, color=gray!70!white},
  every node/.style={font=\scriptsize},
]
  \node[state] (s0) at (0,0) {0};
  \node[state] (s3) at (2.5,0) {3};
  \node[state] (s1) at (5,0) {1};
  \node[state] (s5) at (7.5,0) {5};
  \node[state] (s2) at (5,-1.5) {2};

  \draw[arrow] (s0) -- (s3);
  \draw[arrow] (s3) -- (s1);
  \draw[arrow] (s1) -- (s5);
  \draw[arrow] (s3) -- (s2);

  \node[above=0.2cm, text=green!60!white] at (s0) {entry};
  \node[above=0.2cm, text=red!60!white] at (s5) {exit};
\end{tikzpicture}
\end{center}

\subsection{Topological Sort ile Yeniden Olusturma}

State transition graph'i uzerinde BFS-based topological sort uygulanarak orijinal kod sirasi elde edilir:

\begin{lstlisting}[language=Python, caption={Topological sort -- cff\_deflattener.py}]
def _topological_sort(self, graph, blocks):
    order, visited = [], set()
    entry = blocks[0].state_value
    queue = [entry]
    while queue:
        state = queue.pop(0)
        if state in visited: continue
        visited.add(state)
        order.append(state)
        for ns in graph.get(state, []):
            if ns not in visited:
                queue.append(ns)
    # Ziyaret edilmemis state'leri sona ekle
    for block in blocks:
        if block.state_value not in visited:
            order.append(block.state_value)
    return order
\end{lstlisting}

\subsection{Karmasiklik Analizi}

\begin{equation}
T_{\text{CFF}} = \bigO{|C| + |V| + |E|}
\end{equation}

burada $|C|$ case sayisi, $|V|$ state sayisi (=$|C|$), $|E|$ gecis sayisi. Pratikte $|E| \leq 2|V|$ (her case en fazla 2 gecis yapar --- true/false branch).


% ===================================================================
\section{Opaque Predicate Removal}
\label{sec:opaque-removal}
% ===================================================================

\file{opaque\_predicate.py}, bilinen opaque predicate pattern'lerini regex ile tespit edip dead code'u isaretler veya kaldirir.

\subsection{Tespit Stratejileri}

\begin{enumerate}
\item \textbf{Pattern matching:} 7 always-true + 4 always-false regex pattern (Tablo~\ref{tab:opaque-predicates})
\item \textbf{Constant folding:} Sabit deger kullanan branch'ler (Phase 3 ile entegre)
\item \textbf{Z3 SMT solver (opsiyonel):} Karmasik opaque predicate'lar icin symbolic verification
\end{enumerate}

\subsection{Z3 Entegrasyonu}

Eger z3-solver kurulu ise, pattern matching ile tespit edilemeyen karmasik predicate'lar Z3 ile dogrulanabilir:

\begin{equation}
\text{is\_opaque}(\varphi) = \begin{cases}
\text{true} & \text{eger } Z3.\text{prove}(\forall x: \varphi(x) = \text{true}) \\
\text{false} & \text{eger } Z3.\text{find\_counterexample}(\varphi) \neq \varnothing
\end{cases}
\end{equation}

\file{opaque\_predicate.py}'da Z3 lazy-load yapilir:

\begin{lstlisting}[language=Python, caption={Z3 lazy-load -- opaque\_predicate.py}]
if use_z3:
    try:
        import z3  # noqa: F401
        self._z3_available = True
    except ImportError:
        logger.info("z3-solver bulunamadi, "
                     "pattern-based tespit kullanilacak")
\end{lstlisting}

\subsection{Eliminasyon}

\begin{itemize}
\item \textbf{Always-true:} \code{if(OPAQUE\_TRUE)\{X\} else \{Y\}} $\to$ \code{X}
\item \textbf{Always-false:} \code{if(OPAQUE\_FALSE)\{X\} else \{Y\}} $\to$ \code{Y} (veya sil)
\end{itemize}


% ===================================================================
\section{String Decryption}
\label{sec:string-decrypt}
% ===================================================================

\file{string\_decryptor.py}, binary'lerdeki sifrelenmis string'leri cozen bir moduld-ur. JavaScript deobfuscation ile dogrudan ilgili olmasa da, Karadul'un genel deobfuscation yeteneginin parcasidir.

\subsection{Desteklenen Sifreleme Yoentemleri}

\begin{enumerate}
\item \textbf{XOR single-byte:} \code{buf[i] \^{}= 0xNN}
\item \textbf{XOR multi-byte:} \code{buf[i] \^{}= key[i \% keylen]}
\item \textbf{XOR rolling:} \code{prev = buf[i] \^{} prev}
\item \textbf{RC4:} S-box initialization + swap pattern
\item \textbf{Base64:} \code{atob("SGVsbG8=")}
\item \textbf{Stack string:} Char-by-char push (Ghidra decompile ciktisinda)
\end{enumerate}

\subsection{Stack String Reconstruction}

Ghidra decompile ciktisinda ardisik local degiskenlere tek karakter atamalari ``stack string'' olarak yeniden birlestirilir:

\begin{lstlisting}[language=Python, caption={Stack string ornegi (Ghidra decompile ciktisi)}]
local_28 = 0x48;  # 'H'
local_27 = 0x65;  # 'e'
local_26 = 0x6c;  # 'l'
local_25 = 0x6c;  # 'l'
local_24 = 0x6f;  # 'o'
# -> "Hello" (confidence: 0.85)
\end{lstlisting}

\subsection{XOR Brute-Force}

Kisa string'ler icin tum 256 tek-byte anahtar denenerek en yuksek ``printable'' skorlu sonuc secilir:

\begin{equation}
\text{score}(k) = \frac{|\{c \in \text{decrypt}(d, k) : c \in \text{printable}\}|}{|\text{decrypt}(d, k)|}
\end{equation}

$\text{score}(k) > 0.7$ olan ilk 5 sonuc dondurulur.


% ===================================================================
\section{Acorn Fallback}
\label{sec:acorn-fallback}
% ===================================================================

Babel parser 32MB+ dosyalarda veya ozel syntax pattern'lerinde basarisiz olabilir. \file{deobfuscate.mjs} bu durumda Acorn parser'a duser.

\subsection{Babel vs Acorn Karsilastirmasi}

\begin{table}[h]
\centering
\small
\begin{tabular}{lcc}
\toprule
\textbf{\color{sectioncyan}Ozellik} & \textbf{\color{sectioncyan}Babel} & \textbf{\color{sectioncyan}Acorn} \\
\midrule
AST formati & Babel (genisletilmis) & ESTree (standart) \\
Error recovery & \checkmark & Kismi \\
TypeScript desteegi & \checkmark (plugin) & $\times$ \\
JSX desteegi & \checkmark (plugin) & $\times$ \\
Performans (30MB) & $\sim$45s & $\sim$8s \\
Memory (30MB) & $\sim$4GB & $\sim$1.5GB \\
\bottomrule
\end{tabular}
\caption{Babel vs Acorn parser karsilastirmasi.}
\label{tab:babel-vs-acorn}
\end{table}

Acorn, \file{deobfuscate.mjs}'te lazy-load edilir ve sadece Babel basarisiz olursa yuklenir:

\begin{lstlisting}[language=JavaScript, caption={Acorn lazy-load}]
let acorn = null, acornWalk = null;
function loadAcorn() {
  try {
    acorn = _require("acorn");
    acornWalk = _require("acorn-walk");
    return true;
  } catch { return false; }
}
\end{lstlisting}


\begin{ozet}
JavaScript deobfuscation, 4 adimli standart zincir (beautify $\to$ synchrony $\to$ babel $\to$ webpack) ve 10 fazli derin transform pipeline'indan olusur. Phase 2--8, visitor-merge optimizasyonu ile tek AST traverse'da birlestirilmistir ($\sim$3--5x hiz kazanci). Graceful degradation her adimda basarisizliga tolerans saglar. Faz 9--10'daki LLM-siz semantik isimlendirme toplam 25+ heuristik kuralla etkileyici sonuclar ueretir. 200MB+ dosyalar icin stream-parse ile chunked processing otomatik devreye girer. CFF deflattening, opaque predicate removal ve string decryption ek moduller olarak pipeline'i guclendirir.
\end{ozet}


\clearpage

% ############################################################################
%
%          B O L U M   6 :   B A Y E S I A N   I S I M   F U Z Y O N U
%
% ############################################################################

% Bu bolum, Karadul'un en yenilikci algoritmasi olan Bayesian log-odds isim
% birlestirme sistemini detayli olarak turetir.
% Kaynak kod: karadul/reconstruction/name_merger.py


% ===================================================================
\section{Problem Tanimi: Coklu Kaynak, Tek Isim}
\label{sec:naming-problem-detail}
% ===================================================================

Tersine muhendislikte, bir sembolun (fonksiyon, degisken) gercek ismi bilinmez. Birden fazla \textbf{bagimsiz} kaynak, farkli isimler onerir. Soru sudur: bu onerilerden hangisi dogru, ve nasil birlestiririz?

\begin{tanim}[Isim Fuzyon Problemi]
Verilen:
\begin{itemize}
\item Bir sembol $s$ (orijinal ismi bilinmiyor, orn. \code{FUN\_00401000})
\item $K$ adet kaynak, her biri bir isim onerisi ve confidence degeri ureten: $(n_i, c_i, \text{source}_i)$ icin $i = 1, \ldots, K$
\end{itemize}

Hedef: $s$ icin en uygun ismi $n^*$ sec ve sonucun dogruluk olasililigini $P^*$ hesapla.
\end{tanim}

\begin{example}
\code{FUN\_00401000} icin:
\begin{itemize}
\item \textbf{signature\_db:} \code{SSL\_connect} (confidence: 0.85, $w=1.0$)
\item \textbf{string\_intel:} \code{ssl\_handshake} (confidence: 0.70, $w=0.8$)
\item \textbf{c\_namer:} \code{connect\_tls} (confidence: 0.60, $w=1.0$)
\item \textbf{llm4decompile:} \code{ssl\_init\_connection} (confidence: 0.55, $w=0.5$)
\end{itemize}
Dort farkli isim! Ama hepsi ``SSL baglantisi'' kavramina isaret ediyor. Bayesian fusion bunu nasil cozer?
\end{example}


% ===================================================================
\section{Bayes Teoremi: Temelden Turetme}
\label{sec:bayes-derivation}
% ===================================================================

\subsection{Kosula Bagli Olasilik}

Iki olay $A$ ve $B$ icin:

\begin{equation}
\prob{A \cap B} = \prob{A \given B} \cdot \prob{B} = \prob{B \given A} \cdot \prob{A}
\label{eq:joint-prob}
\end{equation}

Buradan Bayes teoremi turetilir:

\begin{equation}
\boxed{
\prob{H \given E} = \frac{\prob{E \given H} \cdot \prob{H}}{\prob{E}}
}
\label{eq:bayes-theorem}
\end{equation}

\subsection{Isimlendirme Baglaminda Anlami}

\begin{itemize}
\item $H$: ``Bu sembolun gercek ismi $X$'tir'' hipotezi
\item $E_i$: ``Kaynak $i$, $X$ ismini $c_i$ confidence ile onerdi'' kaniti
\item $\prob{H}$: \textbf{Prior} --- hipotez hakkinda onceki inancimiz (baslangicta $0.5$)
\item $\prob{E_i \given H}$: \textbf{Likelihood} --- eger isim gercekten $X$ ise, kaynak $i$'nin bunu $c_i$ ile onermesinin olasiligi. Bunu $\approx c_i$ olarak modelliyoruz.
\item $\prob{E_i \given \neg H}$: Eger isim $X$ \textbf{degilse}, kaynak $i$'nin yine de $X$ onerme olasiligi $\approx 1 - c_i$.
\item $\prob{H \given E_i}$: \textbf{Posterior} --- kanit sonrasi guncellenminsi inanc
\end{itemize}

\subsection{Ardisik Kanit Birlestirme}

Birden fazla kaniti ardisik olarak birlestirmek icin, her kanitin \textbf{likelihood ratio} (LR) degerini kullaniriz:

\begin{equation}
\text{LR}_i = \frac{\prob{E_i \given H}}{\prob{E_i \given \neg H}} = \frac{c_i}{1 - c_i}
\label{eq:lr}
\end{equation}

Ardisik Bayes guncellemesi:

\begin{equation}
\frac{\prob{H \given E_1, E_2, \ldots, E_n}}{\prob{\neg H \given E_1, E_2, \ldots, E_n}} = \frac{\prob{H}}{\prob{\neg H}} \cdot \prod_{i=1}^{n} \text{LR}_i
\end{equation}

\textbf{Eger kaynaklar bagimsiz ise} (naive Bayes varsayimi).


% ===================================================================
\section{Log-Odds Formu: Neden ve Nasil}
\label{sec:logodds-derivation}
% ===================================================================

\subsection{Odds ve Log-Odds}

Odds (bahis orani):

\begin{equation}
\text{odds}(H) = \frac{\prob{H}}{1 - \prob{H}}
\end{equation}

Log-odds (logit):

\begin{equation}
\text{logit}(H) = \log\frac{\prob{H}}{1 - \prob{H}}
\end{equation}

\subsection{Turetme}

Ardisik Bayes guncellemesinin her iki tarafinin logaritmasini alalim:

\begin{align}
\log\frac{\prob{H \given \mathbf{E}}}{\prob{\neg H \given \mathbf{E}}} &= \log\frac{\prob{H}}{\prob{\neg H}} + \sum_{i=1}^{n} \log \text{LR}_i \\
&= \log\frac{\prob{H}}{1 - \prob{H}} + \sum_{i=1}^{n} \log\frac{c_i}{1 - c_i}
\end{align}

\subsection{Korelasyon Duzeltmesi}

Naive Bayes, kaynakların bagimsiz oldugunu varsayar. Ancak bazilari korelasyonludur (ornegin \code{llm4decompile} tum diger kaynaklari dolayli olarak kullanir). Korelasyonlu kaynaklarin etkisini azaltmak icin agirlik $w_i \in (0, 1]$ eklenir:

\begin{equation}
\boxed{
\text{log-odds}(H) = \underbrace{\log\frac{p_0}{1 - p_0}}_{\text{prior}} + \sum_{i=1}^{n} w_i \cdot \underbrace{\log\frac{c_i}{1 - c_i}}_{\text{log-LR}_i}
}
\label{eq:logodds-weighted}
\end{equation}

burada $p_0 = 0.5$ prior (notr baslangic: $\log(0.5/0.5) = 0$).

\subsection{Neden Log Space?}

\begin{enumerate}
\item \textbf{Numerik stabilite:} Carpim yerine toplam. $\prod_i c_i$ ile $n=10$ kaynak icin floating point underflow olusabilir. Toplam bu sorunu ortadan kaldirir.
\item \textbf{Korelasyon duzeltmesi:} $w_i$ ile dogrudan carpma log space'te lineer agirliklamaya denk gelir. Olasilik uzayinda bu $\text{LR}_i^{w_i}$ demektir --- bilgi-kuramsal olarak ``kaynagin bilgi katkisini $w_i$ oranina indir'' anlamina gelir.
\item \textbf{Sigmoid ile donusum:} Sonuc dogal olarak $[0,1]$ araligina doner.
\end{enumerate}


% ===================================================================
\section{Sigmoid Fonksiyonu: Logit'in Tersi}
\label{sec:sigmoid}
% ===================================================================

Log-odds'u tekrar olasiliga donusturmek icin sigmoid ($\sigma$) fonksiyonunu kullaniriz:

\begin{equation}
\prob{H \given \mathbf{E}} = \sigmoid(\text{log-odds}) = \frac{1}{1 + e^{-\text{log-odds}}}
\label{eq:sigmoid-def}
\end{equation}

\subsection{Sigmoid'in Ozellikleri}

\begin{enumerate}
\item $\sigma(0) = 0.5$ (notr: log-odds = 0 ise emin degiliz)
\item $\lim_{x \to \infty} \sigma(x) = 1$ (cok guclu kanit $\to$ kesinlik)
\item $\lim_{x \to -\infty} \sigma(x) = 0$ (cok guclu karsi kanit $\to$ red)
\item $\sigma'(x) = \sigma(x)(1 - \sigma(x))$ (turevi kendisindendir)
\item $\sigma(-x) = 1 - \sigma(x)$ (simetri)
\end{enumerate}

\begin{center}
\begin{tikzpicture}
\begin{axis}[
  width=10cm, height=5cm,
  domain=-6:6,
  samples=100,
  axis lines=middle,
  xlabel={log-odds},
  ylabel={$\sigma(\text{log-odds})$},
  every axis x label/.style={at=(current axis.right of origin), anchor=west},
  every axis y label/.style={at=(current axis.above origin), anchor=south},
  xtick={-4,-2,0,2,4},
  ytick={0,0.25,0.5,0.75,1.0},
  ymin=-0.05, ymax=1.05,
  grid=major,
  grid style={gray!20!black},
  axis line style={gray!60!white},
  tick label style={color=gray!60!white, font=\scriptsize},
  label style={color=sectioncyan, font=\small},
]
  \addplot[cyan, thick] {1/(1+exp(-x))};
  \addplot[red!60!white, dashed, thick] coordinates {(0,0) (0,0.5)};
  \addplot[red!60!white, dashed, thick] coordinates {(-6,0.5) (0,0.5)};
  \node[color=yellow!70!white, font=\scriptsize] at (axis cs:3, 0.3) {$\sigma(x) = \frac{1}{1+e^{-x}}$};
\end{axis}
\end{tikzpicture}
\end{center}

\subsection{Overflow Korumasi}

\file{name\_merger.py}'da sigmoid icin overflow korumasi uygulanir:

\begin{lstlisting}[language=Python, caption={Sigmoid overflow korumasi -- name\_merger.py}]
if log_odds > 20.0:
    p = 1.0       # exp(-20) ~ 2e-9, 1.0'a cok yakin
elif log_odds < -20.0:
    p = 0.0       # exp(20) ~ 5e8, 0.0'a cok yakin
else:
    p = 1.0 / (1.0 + math.exp(-log_odds))
\end{lstlisting}

Ayrica sonuc $[c_{\min}, c_{\max}] = [0.01, 0.99]$ araligina kliplenir --- asla ``\%100 emin'' veya ``\%0 emin'' demiyoruz.


% ===================================================================
\section{7 Kaynak Detayli Aciklama}
\label{sec:sources}
% ===================================================================

\subsection{binary\_extractor ($w = 1.0$)}

Debug string'ler, RTTI (Run-Time Type Information) ve build path bilgileri derleyici tarafindan binary'ye gomulur. Bu bilgi \textbf{ground truth}'a en yakin kaynaktir:

\begin{itemize}
\item Debug string'ler: \code{\_\_func\_\_} makrosu, assert mesajlari
\item RTTI: C++ sinif isimleri (\code{typeinfo name})
\item Build path: \code{/Users/dev/project/src/ssl\_connect.c}
\end{itemize}

$w = 1.0$ cunku bu bilgi derleyiciden dogrudan gelir --- hicbir tahmin yoktur.

\subsection{c\_namer ($w = 1.0$)}

6 katmanli heuristik isimlendirme motoru:

\begin{enumerate}
\item Fonksiyon imzasi analizi (parametre tipleri, donus tipi)
\item String kullanimi (hata mesajlari, log string'leri)
\item Cagrilan fonksiyonlar (API pattern tespiti)
\item Veri akisi analizi (tanimlama-kullanim zincirleri)
\item Algoritma tespiti (hash, sort, compress pattern'leri)
\item Statistiksel isim onerisi
\end{enumerate}

$w = 1.0$ cunku 6 katman zaten kendi icinde filtrelenmis; cift zayiflatma gereksiz.

\subsection{string\_intel ($w = 0.8$)}

Error mesajlari, protocol handler string'leri ve sabit mesajlardan fonksiyon ismi cikarimi:

\begin{lstlisting}[language=Python]
# Ornek: error mesajindan isim cikarimi
"Failed to connect SSL socket" -> ssl_connect (conf: 0.70)
\end{lstlisting}

$w = 0.8$ cunku error mesajlari, debug string'ler ile kismen ortak kaynaklardan beslenir.

\subsection{signature\_db ($w = 1.0$)}

FLIRT (Fast Library Identification and Recognition Technology) ile 1.5M+ kutuphane fonksiyon imzasi eslestirilir. Bu tamamen \textbf{byte-level} bir eslestirmedir:

\begin{equation}
\text{match}(f) = \argmax_{s \in \text{FLIRT\_DB}} \text{similarity}(\text{bytes}(f), \text{bytes}(s))
\end{equation}

$w = 1.0$ cunku byte pattern tamamen bagimsiz bir bilgi kaynagi.

\subsection{swift\_demangle ($w = 1.0$)}

Swift binary'lerinde mangled sembol isimleri (ornegin \code{\_\$s4main7MyClassC9processSSyS2SF}) demangle edilerek orijinal isim elde edilir:

\begin{equation}
\text{demangle}(\code{\_\$s4main7MyClassC9processSSyS2SF}) = \code{main.MyClass.process(String) -> String}
\end{equation}

$w = 1.0$ cunku demangle \textbf{kesin} bilgidir, tahmin degil.

\subsection{source\_matcher ($w = 0.85$)}

NPM paket fingerprinting: Obfuscated JS kodundaki string ve yapi pattern'leri bilinen NPM paketleri ile eslestirilir:

\begin{itemize}
\item Benzersiz string literal'ler (hata mesajlari, API endpoint'leri)
\item Fonksiyon imza pattern'leri (parametre sayisi + kullanim)
\item Export yapisi (modul cikti formati)
\end{itemize}

$w = 0.85$ cunku heuristik bazli ama yuksek kesinlikli.

\subsection{llm4decompile ($w = 0.5$)}

Makine ogrenmesi modeli ile isim tahmini. Neden $w = 0.5$?

\begin{uyari}[LLM Korelasyon Problemi]
LLM modeli, egitim verisinde tum diger kaynaklarin bilgisini dolayli olarak iciyor:
\begin{itemize}
\item Kutuphane imzalari (signature\_db ile korelasyonlu)
\item Debug string pattern'leri (binary\_extractor ile korelasyonlu)
\item Error mesaj kaliplari (string\_intel ile korelasyonlu)
\end{itemize}
Bu nedenle LLM'in bilgi katkisi ``yarisina'' indirilir: $w = 0.5$.
\end{uyari}

\subsection{byte\_pattern ($w = 1.0$)}

Kutuphane fonksiyonlarinin byte-level fuzzy matching ile tespiti. FLIRT'ten farki: FLIRT kesin imza kullanirken, byte\_pattern fuzzy threshold ile calisir. Yine de tamamen bagimsiz bir byte-level kaynak.


% ===================================================================
\section{Korelasyon-Aware Agirliklar: Bilgi Kurami}
\label{sec:correlation-theory}
% ===================================================================

\subsection{Neden $w_i < 1$?}

Iki kaynak $S_i$ ve $S_j$ arasindaki korelasyonu mutual information ile olceriz:

\begin{equation}
I(S_i; S_j) = \sum_{s_i, s_j} p(s_i, s_j) \log\frac{p(s_i, s_j)}{p(s_i) \cdot p(s_j)}
\label{eq:mutual-info}
\end{equation}

Eger $I(S_i; S_j) > 0$ ise, iki kaynak ortak bilgi tasir. Her ikisini de $w=1.0$ ile saymak bu ortak bilgiyi \textbf{iki kez} sayar ve \textbf{overconfidence} olusturur.

Korelasyon duzeltme prensibimiz:

\begin{equation}
w_i = 1 - \frac{I(S_i; S_{\text{others}})}{H(S_i)}
\label{eq:weight-from-mi}
\end{equation}

burada $H(S_i)$ kaynak $i$'nin entropisi. Pratikte Denklem~\ref{eq:weight-from-mi} yerine ampirik agirliklar kullanilir (Tablo~\ref{tab:weights-detail}).

\subsection{Kaynak Agirliklari Tablosu}

\begin{table}[h]
\centering
\small
\begin{tabular}{llp{6.5cm}}
\toprule
\textbf{\color{sectioncyan}Kaynak} & \textbf{\color{sectioncyan}$w_i$} & \textbf{\color{sectioncyan}Gerekce} \\
\midrule
binary\_extractor & 1.0 & Derleyici ciktisi, tamamen bagimsiz \\
c\_namer & 1.0 & 6 katman kendi icinde filtrelenmis \\
string\_intel & 0.8 & Error msg.~binary\_extractor ile kismi korelasyon \\
signature\_db & 1.0 & FLIRT byte pattern, tamamen bagimsiz \\
swift\_demangle & 1.0 & Deterministik demangle, kesin bilgi \\
source\_matcher & 0.85 & NPM heuristik, bagimsiz ama kesin degil \\
llm4decompile & 0.5 & Tum kaynaklari dolayli kullaniyor \\
byte\_pattern & 1.0 & Byte-level fuzzy match, bagimsiz \\
\midrule
\textit{Bilinmeyen kaynak} & 0.7 & Varsayilan fallback \\
\bottomrule
\end{tabular}
\caption{Kaynak agirliklari, korelasyon analizi ve gerekceler.}
\label{tab:weights-detail}
\end{table}

\subsection{Overconfidence Problemi: Somut Ornek}

Naive Bayes ($w_i = 1.0$ her kaynak icin) ile korelasyon-aware karsilastirma:

\begin{align}
\text{Naive } (w_i = 1): \quad \text{log-odds} &= 0 + 1.0 \cdot \log\frac{0.85}{0.15} + 1.0 \cdot \log\frac{0.70}{0.30} + 1.0 \cdot \log\frac{0.60}{0.40} \\
&= 1.735 + 0.847 + 0.405 = 2.987 \\
P_{\text{naive}} &= \sigma(2.987) = 0.952 \quad (\text{\textbf{\%95.2}})
\end{align}

\begin{align}
\text{Corrected}: \quad \text{log-odds} &= 0 + 1.0 \cdot 1.735 + 0.8 \cdot 0.847 + 0.5 \cdot 0.405 \\
&= 1.735 + 0.678 + 0.203 = 2.616 \\
P_{\text{corr}} &= \sigma(2.616) = 0.932 \quad (\text{\textbf{\%93.2}})
\end{align}

\textbf{Fark:} Naive Bayes \%95.2 derken, korelasyon duzeltmesi \%93.2 der --- \%2 daha muhafazakar ve \textbf{daha dogru}. Cunku llm4decompile'in verdigi bilginin bir kismi zaten diger kaynaklarda mevcuttu.


% ===================================================================
\section{Gruplama Stratejisi: 5 Seviye}
\label{sec:grouping-detail}
% ===================================================================

\file{name\_merger.py}'deki \code{\_merge\_candidates()} fonksiyonu, adaylari 5 seviyeli bir hiyerarsiden gecirir. Her seviye bir oncekinden daha gevşek bir eslestirme kullanir.

\begin{center}
\begin{tikzpicture}[
  level/.style={rectangle, draw=sectioncyan, fill=boxbg, minimum width=5cm, minimum height=0.9cm, text=white, font=\small, rounded corners=2pt, align=center},
  arrow/.style={-{Stealth[length=2.5mm]}, thick, color=gray!60!white},
  fail/.style={-{Stealth[length=2mm]}, dashed, color=red!50!white},
  node distance=0.7cm,
]
  \node[level] (l1) {Seviye 1: Exact Multi-Match};
  \node[level, below=of l1] (l2) {Seviye 2: Partial Match};
  \node[level, below=of l2] (l3) {Seviye 3: Semantic Match};
  \node[level, below=of l3] (l4) {Seviye 4: Voting ($\geq$3)};
  \node[level, below=of l4] (l5) {Seviye 5: Single Source};

  \draw[arrow] (l1) -- node[right, font=\tiny, text=red!50!white] {eslesmezse} (l2);
  \draw[arrow] (l2) -- node[right, font=\tiny, text=red!50!white] {eslesmezse} (l3);
  \draw[arrow] (l3) -- node[right, font=\tiny, text=red!50!white] {eslesmezse} (l4);
  \draw[arrow] (l4) -- node[right, font=\tiny, text=red!50!white] {eslesmezse} (l5);

  \node[right=1.5cm, font=\scriptsize, text=green!60!white] at (l1.east) {En yuksek guven};
  \node[right=1.5cm, font=\scriptsize, text=yellow!70!white] at (l3.east) {Orta guven};
  \node[right=1.5cm, font=\scriptsize, text=red!60!white] at (l5.east) {En dusuk guven};
\end{tikzpicture}
\end{center}

\subsection{Seviye 1: Exact Multi-Match}

Ayni normalize isim $\geq 2$ \textbf{farkli} kaynaktan geldiyse, dogrudan Bayesian merge:

\begin{lstlisting}[language=Python, caption={Exact multi-match -- name\_merger.py}]
for norm_name, group in name_groups.items():
    if len(group) >= 2:
        sources = set(c.source for c in group)
        if len(sources) >= 2:
            conf = bayesian_merge(
                [c.confidence for c in group],
                [c.source for c in group], self._cfg)
            return MergedName(..., merge_method="bayesian_multi")
\end{lstlisting}

\subsection{Seviye 2: Partial Match}

Farkli isimler arasinda prefix/suffix varyanti aranir. Ornegin:
\begin{itemize}
\item \code{cycle\_sizes\_default} vs \code{cycle\_size} $\to$ partial match
\item \code{active\_event\_monitor} vs \code{event\_monitor} $\to$ partial match
\end{itemize}

Overlap skoru:

\begin{equation}
\text{overlap}(A, B) = \begin{cases}
\frac{|A_{\text{short}}|}{|A_{\text{long}}|} & \text{eger } A_{\text{short}} \subset A_{\text{long}} \text{ (substring)} \\[6pt]
J(W_A, W_B) = \frac{|W_A \cap W_B|}{|W_A \cup W_B|} & \text{aksi halde (Jaccard)}
\end{cases}
\label{eq:overlap}
\end{equation}

$\text{overlap} \geq 0.5$ ise partial match kabul edilir. Penalty:

\begin{equation}
\text{penalty} = 0.7 + 0.3 \times \text{overlap}
\end{equation}

Tam substring ($\text{overlap} = 1.0$) $\to$ penalty $= 1.0$ (penalty yok). Kismi overlap ($\text{overlap} = 0.5$) $\to$ penalty $= 0.85$.


\subsection{Seviye 3: Semantic Match}

\file{name\_merger.py}'de 20 adet \code{frozenset} ile kelime gruplari tanimlidir. Iki isim, ayni gruba dusen kelimeleri paylasiyorsa ``semantik olarak benzer'' sayilir.

\begin{table}[h]
\centering
\small
\begin{tabular}{p{3cm}p{9cm}}
\toprule
\textbf{\color{sectioncyan}Grup Adi} & \textbf{\color{sectioncyan}Kelimeler} \\
\midrule
Gonderme & send, transmit, write, emit, dispatch, post \\
Alma & recv, receive, read, get, fetch \\
Olusturma & init, initialize, setup, create, construct, new, alloc \\
Yok etme & destroy, cleanup, teardown, free, delete, release, close, dispose \\
Parse/Decode & parse, decode, deserialize, unmarshal, extract \\
Encode & serialize, encode, marshal, format, stringify \\
Isleyici & handle, process, on, dispatch, callback, handler \\
Dogrulama & check, validate, verify, test, assert, is, has \\
Koleksiyon ekle & add, insert, append, push, enqueue, put \\
Koleksiyon sil & remove, delete, erase, pop, dequeue, take \\
Boyut & size, length, count, num, len, total \\
Tampon & buf, buffer, data, bytes, payload, blob \\
Mesaj & msg, message, packet, frame, request, response \\
Hata & err, error, fault, failure, exception \\
Baglam & ctx, context, state, env, session \\
Konfiguerasyon & cfg, config, configuration, settings, options, prefs \\
Baglanti & conn, connection, socket, link, channel \\
Loglama & log, print, trace, debug, info, warn \\
Kilit alma & lock, acquire, enter, begin, grab \\
Kilit birakma & unlock, release, leave, end, drop \\
\bottomrule
\end{tabular}
\caption{Semantik kelime gruplari (20 frozenset).}
\label{tab:semantic-groups}
\end{table}

\subsubsection{Expand-Then-Intersect Algoritmasi}

\begin{enumerate}
\item Her isim kelimelere ayrilir: \code{ssl\_handshake} $\to$ \{ssl, handshake\}
\item Her kelime, ait oldugu semantic grubun \textbf{tum kelimeleri} ile genisletilir:
      \code{handshake} $\notin$ herhangi bir grup; \code{connect} $\to$ \{conn, connection, socket, link, channel\}
\item Genisletilmis kumelerin kesisimi kontrol edilir
\item Kesisim orani $\geq 0.3$ ise semantik benzerlik var
\end{enumerate}

\begin{lstlisting}[language=Python, caption={Semantic similarity -- name\_merger.py}]
def _is_semantically_similar(self, name1, name2):
    words1 = set(self._normalize(name1).split("_"))
    words2 = set(self._normalize(name2).split("_"))
    # Dogrudan kelime kesisimi
    common = words1 & words2
    if common and len(common) / max(len(words1), len(words2)) >= 0.5:
        return True
    # Semantic group genisletme
    expanded1 = set()
    for w in words1:
        if w in self._word_to_group:
            expanded1 |= self._word_to_group[w]
        expanded1.add(w)
    # ... ayni expanded2 icin ...
    overlap = expanded1 & expanded2
    union = expanded1 | expanded2
    return len(overlap) / len(union) >= 0.3
\end{lstlisting}

Semantik match'te confidence \textbf{0.85x penalty} ile carpilir cunku isimler tam eslesmiyor:

\begin{equation}
c_i^{\text{semantic}} = c_i \times 0.85
\end{equation}


\subsection{Seviye 4: Voting}

$\geq 3$ kaynak ayni normalize ismi veriyorsa, cogunluk kazanir:

\begin{equation}
n^* = \argmax_{n} \sum_{i=1}^{K} \mathbb{1}[\text{normalize}(n_i) = n]
\end{equation}

Eger $\text{count}(n^*) \geq 3$ ise, bu adaylar Bayesian merge'e girer.

\subsection{Seviye 5: Single Source}

Tek kaynak kaldiysa, Bayesian merge gerekmez. Ancak kaynagin $w_i$ deger ile duzeltme yapilir:

\begin{equation}
c_{\text{adjusted}} = \sigma\left(w_i \cdot \log\frac{c_i}{1 - c_i}\right)
\end{equation}

$w_i < 1$ ise bu ``tek kaynaga fazla guvenme'' demektir.


% ===================================================================
\section{Normalize Isim Karsilastirmasi}
\label{sec:normalize}
% ===================================================================

Farkli kaynaklar farkli konvansiyonlarda isim uretir. Karsilastirma oncesi normalizasyon sart:

\begin{lstlisting}[language=Python, caption={Isim normalizasyonu -- name\_merger.py}]
def _normalize(self, name):
    # camelCase -> snake_case
    name = re.sub(r"([a-z])([A-Z])", r"\1_\2", name)
    # Prefix'leri kaldir (m_, s_, g_, p_)
    name = re.sub(r"^[msgp]_", "", name)
    return name.lower().strip("_")
\end{lstlisting}

Ornekler:
\begin{itemize}
\item \code{SSLConnect} $\to$ \code{ssl\_connect}
\item \code{m\_sslConnect} $\to$ \code{ssl\_connect}
\item \code{ssl\_connect} $\to$ \code{ssl\_connect} (degismez)
\item \code{g\_SSLHandler} $\to$ \code{ssl\_handler}
\end{itemize}


% ===================================================================
\section{UNK Threshold: Precision vs Recall}
\label{sec:unk-threshold}
% ===================================================================

\subsection{NSA-Grade Felsefe}

\begin{uyari}[Yanlis Isim Tehlikesi]
Tersine muhendislikte yanlis isim, isimsizden \textbf{cok daha kotu}dur:
\begin{itemize}
\item Yanlis isim $\to$ analist yanlis yolda saatler/gunler harcar
\item UNK (isimsiz) $\to$ analist ``bilmiyorum'' der, dikkatli devam eder
\end{itemize}

\textbf{Sonuc:} Dusuk confidence'li isimleri \textbf{atamayiz}. $c < 0.30$ ise UNK.
\end{uyari}

\subsection{ROC Analizi}

UNK threshold secimi, precision ve recall arasinda trade-off'tur:

\begin{equation}
\text{Precision} = \frac{TP}{TP + FP}, \quad \text{Recall} = \frac{TP}{TP + FN}
\end{equation}

\begin{itemize}
\item \textbf{Threshold $\uparrow$:} Precision $\uparrow$ (daha az yanlis isim), Recall $\downarrow$ (daha cok UNK)
\item \textbf{Threshold $\downarrow$:} Precision $\downarrow$ (daha cok yanlis isim), Recall $\uparrow$ (daha az UNK)
\end{itemize}

Karadul'un secimi: $\theta_{\text{UNK}} = 0.30$. Bu, \textbf{Type I error maliyetinin Type II error maliyetinden cok daha yuksek} olmasi varsayimindan gelir:

\begin{equation}
\text{Cost}_{\text{wrong\_name}} \gg \text{Cost}_{\text{UNK}} \implies \theta_{\text{UNK}} \text{ yuksek tutulur}
\end{equation}


% ===================================================================
\section{Somut Hesaplama Ornekleri}
\label{sec:bayesian-examples}
% ===================================================================

\subsection{Senaryo 1: Uc Kaynak Ayni Isim (Exact Multi-Match)}

\code{FUN\_00401000} icin uc kaynak ayni normalize ismi veriyor:

\begin{center}
\begin{tabular}{llcc}
\toprule
Kaynak & Onerilen Isim & $c_i$ & $w_i$ \\
\midrule
signature\_db & \code{SSL\_connect} & 0.85 & 1.0 \\
string\_intel & \code{SSL\_connect} & 0.70 & 0.8 \\
c\_namer & \code{ssl\_connect} & 0.60 & 1.0 \\
\bottomrule
\end{tabular}
\end{center}

Normalize edilince uc isim de \code{ssl\_connect} --- exact multi-match!

\textbf{Adim 1: Prior log-odds}
\begin{equation}
\ell_0 = \log\frac{0.5}{0.5} = 0
\end{equation}

\textbf{Adim 2: Her kaynağin log-LR katkisi}
\begin{align}
\Delta\ell_1 &= w_1 \cdot \log\frac{c_1}{1-c_1} = 1.0 \cdot \log\frac{0.85}{0.15} = 1.0 \cdot 1.7346 = 1.7346 \\[4pt]
\Delta\ell_2 &= w_2 \cdot \log\frac{c_2}{1-c_2} = 0.8 \cdot \log\frac{0.70}{0.30} = 0.8 \cdot 0.8473 = 0.6779 \\[4pt]
\Delta\ell_3 &= w_3 \cdot \log\frac{c_3}{1-c_3} = 1.0 \cdot \log\frac{0.60}{0.40} = 1.0 \cdot 0.4055 = 0.4055
\end{align}

\textbf{Adim 3: Toplam log-odds}
\begin{equation}
\ell = 0 + 1.7346 + 0.6779 + 0.4055 = 2.8180
\end{equation}

\textbf{Adim 4: Sigmoid}
\begin{equation}
P = \sigma(2.8180) = \frac{1}{1 + e^{-2.8180}} = \frac{1}{1 + 0.05975} = 0.9436
\end{equation}

\begin{basari}[Senaryo 1 Sonuc]
Uc kaynak birlestirildiginde \textbf{\%94.4 guvenle} \code{ssl\_connect} ismi atanir. En iyi tek kaynak \%85'ti --- Bayesian fusion \%9.4 puan ekledi.
\end{basari}


\subsection{Senaryo 2: Iki Kaynak Partial Match}

\code{FUN\_00501234} icin iki kaynak benzer ama farkli isimler veriyor:

\begin{center}
\begin{tabular}{llcc}
\toprule
Kaynak & Onerilen Isim & $c_i$ & $w_i$ \\
\midrule
binary\_extractor & \code{process\_message} & 0.80 & 1.0 \\
c\_namer & \code{handle\_message\_data} & 0.65 & 1.0 \\
\bottomrule
\end{tabular}
\end{center}

Normalize: \code{process\_message} vs \code{handle\_message\_data}

\textbf{Kelime kumesi:} $\{$process, message$\}$ vs $\{$handle, message, data$\}$. Jaccard:

\begin{equation}
J = \frac{|\{\text{message}\}|}{|\{\text{process, message, handle, data}\}|} = \frac{1}{4} = 0.25 < 0.5
\end{equation}

Dogrudan Jaccard $< 0.5$, ama ``handle'' ve ``process'' ayni semantic gruba dusur (Tablo~\ref{tab:semantic-groups}: ``handle, process, on, dispatch, callback, handler''). Semantic expansion sonrasi:

Expanded$_1$ = $\{$process, message, handle, on, dispatch, callback, handler$\}$ \\
Expanded$_2$ = $\{$handle, message, data, process, on, dispatch, callback, handler, bytes, payload, blob, buf, buffer$\}$

\begin{equation}
\frac{|\text{Expanded}_1 \cap \text{Expanded}_2|}{|\text{Expanded}_1 \cup \text{Expanded}_2|} = \frac{7}{13} = 0.538 \geq 0.3
\end{equation}

Semantic match! Penalty $= 0.85$:

\begin{align}
\Delta\ell_1 &= 1.0 \cdot \log\frac{0.80 \times 0.85}{1 - 0.80 \times 0.85} = 1.0 \cdot \log\frac{0.68}{0.32} = 0.7538 \\[4pt]
\Delta\ell_2 &= 1.0 \cdot \log\frac{0.65 \times 0.85}{1 - 0.65 \times 0.85} = 1.0 \cdot \log\frac{0.5525}{0.4475} = 0.2107 \\[4pt]
\ell &= 0 + 0.7538 + 0.2107 = 0.9645 \\[4pt]
P &= \sigma(0.9645) = \frac{1}{1 + e^{-0.9645}} = 0.724
\end{align}

\begin{basari}[Senaryo 2 Sonuc]
Semantic match ile iki kaynak birlestirildiginde \textbf{\%72.4 guvenle} \code{process\_message} (daha kisa isim) atanir. Penalty olmadan \%80 olacakti --- semantik belirsizlik \%7.6 dusurdu.
\end{basari}


\subsection{Senaryo 3: Tek Kaynak, Dusuk Confidence $\to$ UNK}

\code{FUN\_00601000} icin sadece bir kaynak var:

\begin{center}
\begin{tabular}{llcc}
\toprule
Kaynak & Onerilen Isim & $c_i$ & $w_i$ \\
\midrule
llm4decompile & \code{init\_module} & 0.45 & 0.5 \\
\bottomrule
\end{tabular}
\end{center}

Tek kaynak, Single Source (Seviye 5):

\begin{align}
\ell &= w_1 \cdot \log\frac{c_1}{1-c_1} = 0.5 \cdot \log\frac{0.45}{0.55} = 0.5 \cdot (-0.2007) = -0.1003 \\[4pt]
P &= \sigma(-0.1003) = \frac{1}{1 + e^{0.1003}} = 0.4749
\end{align}

$P = 0.475 < \theta_{\text{UNK}} = 0.30$? Hayir, $0.475 > 0.30$.

Ama bekleyin --- burada $P < 0.5$, yani model ``bu ismin dogru olma olasiligi \%50'den az'' diyor. Rasyonel bir analist buna guvenm-ez.

Esik degeri $\theta_{\text{UNK}} = 0.30$ ile $P = 0.475 > 0.30$ oldugu icin \textit{teknik olarak} isim atanir. Ancak pratikte bu dusuk confidence'li isim dikkatle ele alinir.

Simdi daha dusuk confidence ile deneyelim: $c_1 = 0.35$:

\begin{align}
\ell &= 0.5 \cdot \log\frac{0.35}{0.65} = 0.5 \cdot (-0.6190) = -0.3095 \\
P &= \sigma(-0.3095) = 0.4233
\end{align}

$P = 0.423 > 0.30$ --- hala UNK sinirinin uzerinde. Ancak $c_1 = 0.15$:

\begin{align}
\ell &= 0.5 \cdot \log\frac{0.15}{0.85} = 0.5 \cdot (-1.7346) = -0.8673 \\
P &= \sigma(-0.8673) = 0.296
\end{align}

$P = 0.296 < 0.30$ $\to$ \textbf{UNK!} Isim atanmaz.

\begin{uyari}[Senaryo 3 Sonuc]
LLM kaynaginin tek basina confidence $= 0.15$, $w = 0.5$ ile verdigi isim UNK olarak kalir. \textbf{Yanlis isim, isimsizden kotudur.}
\end{uyari}


% ===================================================================
\section{Hesaplama Karmasikligi}
\label{sec:complexity}
% ===================================================================

\subsection{Genel Karmasiklik}

$N$ adet sembol ve her sembol icin ortalama $K$ aday icin:

\begin{equation}
T_{\text{merge}} = \bigO{N \cdot K^2}
\end{equation}

Kare terimi, her aday ciftinin partial/semantic karsilastirmasinden gelir.

\subsection{Alt Islem Karmasikliklari}

\begin{table}[h]
\centering
\small
\begin{tabular}{lll}
\toprule
\textbf{\color{sectioncyan}Islem} & \textbf{\color{sectioncyan}Karmasiklik} & \textbf{\color{sectioncyan}Aciklama} \\
\midrule
Normalizasyon & $\bigO{|n|}$ & Regex + lowercase, string uzunlugu ile lineer \\
Exact match & $\bigO{K}$ & Hash-based gruplama \\
Partial overlap & $\bigO{K^2 \cdot |n|}$ & Tum ciftleri karsilastir, her biri $\bigO{|n|}$ \\
Semantic match & $\bigO{K^2 \cdot |W|}$ & Kelime genisletme + kesisim, $|W|$ = kelime sayisi \\
Bayesian merge & $\bigO{K}$ & Log-odds toplami (lineer) \\
\bottomrule
\end{tabular}
\caption{Alt islem karmasikliklari.}
\label{tab:complexity}
\end{table}

\subsection{Pratikte Performans}

Tipik degerler: $N \approx 5000$ (bir binary'deki fonksiyon sayisi), $K \approx 3$ (ortalama kaynak sayisi):

\begin{equation}
T \approx 5000 \times 3^2 = 45{,}000 \text{ islem}
\end{equation}

Python'da bu $< 0.1$s surer --- bottleneck tamamen decompilation ve statik analiz asamalarindadir.


% ===================================================================
\section{Overconfidence Problemi ve Cozumu}
\label{sec:overconfidence}
% ===================================================================

\subsection{Naive Bayes Hatasi}

Naive Bayes, tum kaynaklarin \textbf{kosula bagli bagimsiz} oldugunu varsayar:

\begin{equation}
\prob{E_1, E_2, \ldots, E_n \given H} = \prod_{i=1}^{n} \prob{E_i \given H}
\end{equation}

Bu varsayim genellikle \textbf{yanlistir}. Eger iki kaynak korelasyonlu ise, ayni bilgiyi iki kez saymis oluruz.

\subsection{Korelasyon ile Duzeltme}

Log-odds formulunde $w_i$ korelasyonu absorbe eder:

\begin{equation}
\text{LR}_i^{\text{effective}} = \text{LR}_i^{w_i}
\end{equation}

$w_i = 1$: Kaynak tamamen bagimsiz $\to$ tam bilgi katkisi. \\
$w_i = 0.5$: Kaynak yari korelasyonlu $\to$ bilgi katkisi ``yarisina'' iner. \\
$w_i = 0$: Kaynak tamamen bagimli $\to$ bilgi katkisi sifir (sayilmaz).

Bilgi-kuramsal olarak bu, log-likelihood ratio'nun eksponent ile olceklenmesidir:

\begin{equation}
\text{effective info gain}_i = w_i \cdot \text{KL}(\prob{\cdot \given E_i} \| \prob{\cdot})
\end{equation}

burada KL divergence bilgi kazancini olcer.

\subsection{Sonuc}

\begin{basari}[Bayesian Fusion Avantajlari]
\begin{enumerate}
\item Coklu kaynaklari sistematik olarak birlestirir --- el ile ``en iyi kaynagi sec'' yerine matematiksel optimal
\item Korelasyonu modelleyerek overconfidence'i onler
\item UNK threshold ile yanlis isim atama riskini minimize eder
\item Log-odds formunda numerik olarak stabil
\item Her adimi tamamen deterministik ve auditlenebilir
\end{enumerate}
\end{basari}


\begin{ozet}
Bayesian isim fuzyonu, 7+ bagimsiz kaynaktan gelen isimleri korelasyon-aware log-odds ile birlestirir. 5 seviyeli gruplama (exact $\to$ partial $\to$ semantic $\to$ voting $\to$ single) her durum icin uygun strateji secer. $w_i < 1$ agirliklari korelasyonlu kaynaklarin etkisini bilgi-kuramsal olarak azaltir. UNK threshold (\%30) NSA-grade felsefe ile yanlis isimlendirmeyi onler. Somut orneklerde: uc kaynak exact match ile \%85'ten \%94.4'e, iki kaynak semantic match ile \%72.4'e yukseltildi; tek dusuk-confidence kaynak ise UNK olarak dogru sekilde reddedildi. Karmasiklik $\bigO{N \cdot K^2}$ ile pratikte milisaniye mertebesinde calisir.
\end{ozet}



% ============================================================================
% BOLUM 6: BAYESIAN ISIM FUZYONU
% ============================================================================
\chapter{Bayesian Isim Fuzyonu}
\label{ch:bayesian}

Bu boluem, Karadul'un en yenilikci algoritmasi olan Bayesian log-odds isim birlestirme sistemini detayli olarak tueret\-ir.

\section{Problem Tanimi}
\label{sec:naming-problem}

Tersine muehendislikte, bir fonksiyonun gercek ismi bilinmez. Birden fazla kaynak farkli isim oenerir:

\begin{example}
\code{FUN\_00401000} icin:
\begin{itemize}
\item \textbf{signature\_db:} \code{SSL\_connect} (confidence: 0.85)
\item \textbf{string\_intel:} \code{ssl\_handshake} (confidence: 0.70)
\item \textbf{c\_namer:} \code{connect\_tls} (confidence: 0.60)
\end{itemize}
Hangisini sececegiz? Hepsini nasil birlestirecegiz?
\end{example}

\section{Bayesian Inference Temelleri}
\label{sec:bayes-temelleri}

\subsection{Bayes Teoremi}

\begin{equation}
\prob{H \given E} = \frac{\prob{E \given H} \cdot \prob{H}}{\prob{E}}
\label{eq:bayes}
\end{equation}

burada:
\begin{itemize}
\item $H$: Hipotez (``bu fonksiyonun ismi $X$'tir'')
\item $E$: Kanit (kaynaktan gelen isim onerisi)
\item $\prob{H}$: Prior (oenceki bilgi --- baslangiccta $0.5$)
\item $\prob{E \given H}$: Likelihood (kaynag\-in guveenilirli\-gi)
\item $\prob{H \given E}$: Posterior (guencellenmis inanc)
\end{itemize}

\subsection{Log-Odds Formu}

Birden fazla kaniti birlestirmek icin log-odds formunu kullaniriz:

\begin{equation}
\boxed{
\text{log-odds}(H) = \log\frac{\prob{H}}{1 - \prob{H}} + \sum_{i=1}^{n} w_i \cdot \log\frac{c_i}{1 - c_i}
}
\label{eq:logodds}
\end{equation}

burada:
\begin{itemize}
\item $c_i$: $i$-inci kaynag\-in confidence degeri ($[0.01, 0.99]$ icinde)
\item $w_i$: $i$-inci kaynag\-in agirligi (korelasyon duzeltmesi)
\item Prior $= 0.5$ $\Rightarrow$ $\log(0.5/0.5) = 0$ (noetral baslangic)
\end{itemize}

Final olasilik:

\begin{equation}
\prob{H \given \text{tuem kanitlar}} = \sigmoid(\text{log-odds}) = \frac{1}{1 + e^{-\text{log-odds}}}
\label{eq:sigmoid}
\end{equation}

\subsection{Neden Log-Odds?}

\begin{enumerate}
\item \textbf{Numerik stabilite:} Carpim yerine toplam --- floating point underflow onlenir.
\item \textbf{Korelasyon duzeltmesi:} $w_i$ ile dogrudan agirlik verme muemkuen --- naive Bayes'te bu zor.
\item \textbf{Sigmoid:} Sonucu dogal olarak $[0, 1]$ araligina eristirir.
\end{enumerate}

\section{Korelasyon-Aware Agirliklar}
\label{sec:correlation}

\subsection{Neden $w_i < 1$?}

Iki kaynak korelasyonlu ise, ayni bilgiyi iki kez saymis oluruz. Bu \textbf{overconfidence} (asiri guveen) yaratir.

\begin{tanim}[Korelasyon Duzeltmesi]
$w_i$ agirligi, kaynag\-in bagimsizlik derecesini yansitir:
\begin{equation}
w_i = \begin{cases}
1.0 & \text{tamamen bagimsiz (bilgi kaybetsiz)} \\
0.5 & \text{yueksek korelasyon (bilgi yarisina indirilir)} \\
0.0 & \text{tam bagimli (bilgi tamamen tekrar)}
\end{cases}
\end{equation}
\end{tanim}

\subsection{Kaynak Agirliklari Tablosu}

\begin{table}[h]
\centering
\begin{tabular}{llp{7cm}}
\toprule
\textbf{\color{sectioncyan}Kaynak} & \textbf{\color{sectioncyan}$w_i$} & \textbf{\color{sectioncyan}Neden} \\
\midrule
binary\_extractor & 1.0 & Debug string/RTTI: derleyiciden dogrudan \\
c\_namer & 1.0 & 6 katman zaten filtrelenmis \\
string\_intel & 0.8 & Error msg.\ benzer kaynaklarla kismen ortak \\
signature\_db & 1.0 & FLIRT: tamamen bagimsiz byte pattern \\
swift\_demangle & 1.0 & Mangled name: kesin bilgi \\
source\_matcher & 0.85 & NPM kaynak: bagimsiz ama heuristik \\
llm4decompile & 0.5 & Tuem kaynaklari dolayli olarak kullaniyor \\
byte\_pattern & 1.0 & Byte-level fuzzy match: bagimsiz \\
\bottomrule
\end{tabular}
\caption{Kaynak agirliklari ve gerekcceleri.}
\label{tab:weights}
\end{table}

\section{Somut Hesaplama Oernegi}
\label{sec:bayesian-example}

Uec kaynak ayni fonksiyon icin isim oeneriyor:

\begin{align}
\text{signature\_db:} \quad & c_1 = 0.85, \quad w_1 = 1.0 \\
\text{string\_intel:} \quad & c_2 = 0.70, \quad w_2 = 0.8 \\
\text{llm4decompile:} \quad & c_3 = 0.60, \quad w_3 = 0.5
\end{align}

Adim adim hesaplama:

\begin{align}
\text{log-odds} &= \underbrace{0}_{\text{prior}} + 1.0 \cdot \log\frac{0.85}{0.15} + 0.8 \cdot \log\frac{0.70}{0.30} + 0.5 \cdot \log\frac{0.60}{0.40} \\
&= 0 + 1.0 \cdot 1.735 + 0.8 \cdot 0.847 + 0.5 \cdot 0.405 \\
&= 1.735 + 0.678 + 0.203 \\
&= 2.616
\end{align}

\begin{equation}
\prob{dogru} = \sigmoid(2.616) = \frac{1}{1 + e^{-2.616}} = \frac{1}{1 + 0.073} = 0.932
\end{equation}

\begin{basari}[Sonuc]
Uec kaynak birlestirildiginde \textbf{\%93.2 guvenle} isim atanir. Tek basina en iyisi \%85'ti --- Bayesian fusion \%8.2 puan ekledi.
\end{basari}

\subsection{Ikinci Oernek: Tam Bagimsiz Kaynaklar}

Tuem $w_i = 1.0$ olsaydi (naive Bayes):

\begin{align}
\text{log-odds} &= 0 + 1.0 \cdot \log\frac{0.85}{0.15} + 1.0 \cdot \log\frac{0.70}{0.30} + 1.0 \cdot \log\frac{0.60}{0.40} \\
&= 1.735 + 0.847 + 0.405 = 2.987
\end{align}

\begin{equation}
\prob{dogru}_{\text{naive}} = \sigmoid(2.987) = 0.952
\end{equation}

Karsilastirma: korelasyon duzeltmesi \textbf{\%95.2'yi \%93.2'ye duesuerdue} --- 2 puanlik ``overconfidence'' oenlendi. Bu kueceuk goruenebilir ama binlerce sembol uezerinde yanlis isim oranini onemli oelcuede azaltir.

\subsection{Uecuencue Oernek: Tek Kaynak, Duesueek Confidence}

Yalnizca \code{string\_intel} kaynagi, $c = 0.45$, $w = 0.8$:

\begin{align}
\text{log-odds} &= 0 + 0.8 \cdot \log\frac{0.45}{0.55} = 0.8 \cdot (-0.200) = -0.160 \\
\prob{dogru} &= \sigmoid(-0.160) = 0.460
\end{align}

$0.460 > 0.30$ (UNK threshold) ama $< 0.70$ (pipeline threshold). \textbf{Sonuc:} name\_merger seviyesinde tutulur ama son ciktida UNK olarak gosterilir --- ikili threshold sistemi devreye girer.

\section{Gruplama Stratejisi}
\label{sec:grouping}

Adaylar sirasiyla su stratejilerle gruplanir:

\begin{enumerate}
\item \textbf{Exact multi-match:} Ayni isim $\geq 2$ farkli kaynaktan $\to$ dogrudan Bayesian merge
\item \textbf{Partial match:} Prefix/suffix varyanti (Jaccard $\geq 0.5$) $\to$ penalty ile merge
\item \textbf{Semantic match:} Benzer anlamli isimler (20 kelime grubu) $\to$ 0.85x penalty
\item \textbf{Voting:} $\geq 3$ kaynak ayni $\to$ Bayesian merge
\item \textbf{Single source:} Tek aday $\to$ $w$ ile duzeltilmis confidence
\end{enumerate}

\subsection{Partial Overlap Hesabi}

Iki normalize isim arasindaki benzerlik:

\begin{equation}
\text{overlap}(A, B) = \begin{cases}
|A| / |B| & \text{eger } A \subset B \text{ (substring)} \\
J(W_A, W_B) & \text{aksi halde (Jaccard)}
\end{cases}
\end{equation}

burada $J(A, B) = |A \cap B| / |A \cup B|$ Jaccard benzerlik katsayisi.

\subsection{UNK Threshold}

\begin{uyari}[NSA-Grade Felsefe]
``Yanlis isim, isimsizden koetueduer.'' Bu nedenle confidence $< 0.30$ ise isim \textbf{atanmaz} (UNK olarak kalir).

Type I error (yanlis isim) maliyeti $\gg$ Type II error (UNK bira\-kma) maliyeti.
\end{uyari}

\subsection{Ikili Threshold Sistemi}

Karadul'da iki farkli threshold var:

\begin{enumerate}
\item \textbf{0.30 (name\_merger seviyesi):} ``Bu ismin dogru olma olasiligi en az \%30 olmali, yoksa UNK birak.'' Duesueek goruenuyor ama...
\item \textbf{0.70 (pipeline seviyesi):} ``Son ciktida sadece \%70+ confidence isimleri goester.''
\end{enumerate}

0.30--0.70 arasindaki isimler \textit{``olasi ama emin degilim''} kategorisinde tutulur --- interaktif modda analiste sunulabilir, otomatik modda goesterilmez.

\subsection{ROC Perspektifi}

Farkli threshold'larin tahmini performansi:

\begin{table}[ht]
\centering
\begin{tabular}{lll}
\toprule
\textbf{\color{sectioncyan}Threshold} & \textbf{\color{sectioncyan}Precision} & \textbf{\color{sectioncyan}Recall} \\
\midrule
0.20 & $\sim$\%60 & $\sim$\%95 \\
0.30 & $\sim$\%75 & $\sim$\%85 \\
0.50 & $\sim$\%90 & $\sim$\%60 \\
0.70 & $\sim$\%97 & $\sim$\%35 \\
\bottomrule
\end{tabular}
\caption{Threshold vs precision-recall tradeoff (tahmini).}
\end{table}

\subsection{Hesaplama Karmasikligi}

$N$ sembol, $M$ kaynak, $K$ aday/sembol:

\begin{itemize}
\item Bayesian merge (her sembol icin): $\bigO{M \cdot K}$
\item Gruplama (her sembol icin): $\bigO{K^2}$ (partial match icin tuem ciftler)
\item \textbf{Toplam:} $\bigO{N \cdot (M \cdot K + K^2)}$
\end{itemize}

100.000 sembol, ortalama 1.5 aday: $100{,}000 \times (8 \times 1.5 + 1.5^2) = 1{,}425{,}000$ islem $\to$ $< 1$ saniye.

\begin{ozet}
Bayesian isim fuzyonu, 7 bagimsiz kaynaktan gelen isimleri korelasyon-aware log-odds ile birlestirir. $w_i < 1$ agirliklari korelasyonlu kaynaklarin etkisini azaltir. UNK threshold (\%30) yanlis isimlendirmeyi ooenler. Somut oeernekte: uec kaynak birlestirildiginde \%85'ten \%93.2'ye iyilesme goezlendi.
\end{ozet}


% ============================================================================
% BOLUM 7: KAYNAK KOD REKONSTRUKSIYONU
% ============================================================================
\chapter{Kaynak Kod Rekonstruksiyonu}
\label{ch:reconstruct}

Bu boluem, isimlendirilmis sembollerin anlamli kaynak koda doenuesteruelmesi suerecini aciklar.

\section{6 Katmanli C Naming}
\label{sec:c-namer}

\file{c\_namer.py} dosyasinda implement edilen 6 katmanli deterministic naming sistemi:

\begin{center}
\begin{tikzpicture}[
  layer/.style={rectangle, draw=sectioncyan, fill=boxbg, minimum width=10cm, minimum height=0.8cm, text=white, font=\small},
  arrow/.style={-{Stealth[length=2mm]}, thick, color=gray!50!white},
]
  \node[layer] (l1) at (0,0) {Katman 1: Symbol/Export Table Isimleri};
  \node[layer] (l2) at (0,-1.2) {Katman 2: String Context (komsu string'ler)};
  \node[layer] (l3) at (0,-2.4) {Katman 3: API Call Pattern Matching};
  \node[layer] (l4) at (0,-3.6) {Katman 4: Call Graph Structure};
  \node[layer] (l5) at (0,-4.8) {Katman 5: Dataflow Analysis};
  \node[layer] (l6) at (0,-6.0) {Katman 6: Type-Based Inference};

  \draw[arrow] (l1) -- (l2);
  \draw[arrow] (l2) -- (l3);
  \draw[arrow] (l3) -- (l4);
  \draw[arrow] (l4) -- (l5);
  \draw[arrow] (l5) -- (l6);

  \node[right=0.3cm, font=\tiny, text=green!50!white] at (l1.east) {En yuksek guven};
  \node[right=0.3cm, font=\tiny, text=yellow!50!white] at (l3.east) {Orta guven};
  \node[right=0.3cm, font=\tiny, text=red!50!white] at (l6.east) {En dusuk guven};
\end{tikzpicture}
\end{center}

Her katman, oenceki katmanin bulamadigi semboller icin ek bilgi saglar. Cascade classifier mantigi: kolay olanlar oence, zor olanlar sona.

\section{FLIRT Signature Matching}
\label{sec:flirt}

IDA Pro'nun FLIRT (Fast Library Identification and Recognition Technology) teknolojisi:

\begin{equation}
\text{match}(f) = \begin{cases}
\text{library\_name} & \text{eger } \text{bytes}(f)[0:n] = \text{pattern}[0:n] \text{ ve CRC dogrulandi} \\
\text{UNK} & \text{aksi halde}
\end{cases}
\end{equation}

Karadul'un signature veritabani 1.5M+ imza, 4.200+ kuetuephane icerir.

\section{NPM Source Matching}
\label{sec:npm-match}

Minified JS fonksiyonlarini orijinal npm paketleri ile eslestirir:

\begin{enumerate}
\item Fonksiyon AST structural hash hesapla
\item npm registry'de benzer paketleri ara
\item Kaynak kodu indir (unpkg veya npm pack)
\item AST-level structural matching yap
\item Benzerlik $\geq 0.65$ ise eslestir
\end{enumerate}

Benzerlik olcuetue (Jaccard uzerinde AST node turleri):

\begin{equation}
\text{sim}(f_1, f_2) = \frac{|\text{nodes}(f_1) \cap \text{nodes}(f_2)|}{|\text{nodes}(f_1) \cup \text{nodes}(f_2)|}
\end{equation}

\section{Algoritma Tespiti}
\label{sec:algo-id}

Binary icindeki kriptografik, hash ve sikistirma algoritmalarini tespit etme:

\begin{itemize}
\item \textbf{AES:} S-box byte pattern (0x63, 0x7c, 0x77, 0x7b, ...)
\item \textbf{SHA-256:} Initial hash values ($H_0 = \texttt{0x6a09e667}$, ...)
\item \textbf{CRC32:} Polynomial tablo degerleri
\item \textbf{zlib:} Adler-32 ve deflate magic
\end{itemize}

\begin{ozet}
Kaynak kod rekonstruksiyonu, 6 katmanli C naming, FLIRT signature matching, NPM source matching ve algoritma tespiti bilesenlerin\-den olusur. Her bilesen ayri bir naming kaynagi olarak Bayesian fuzyona beslenir.
\end{ozet}

% --- DEVELOPER AGENT DETAYLI GENISLETME: BOLUM 7+8+9 ---
% ============================================================================
% BOLUM 7: KAYNAK KOD REKONSTRUKSIYONU
% BOLUM 8: GUVENLIK ANALIZ ALGORITMALARI
% BOLUM 9: PERFORMANS VE OLCEKLENDIRME
% ============================================================================
% Dosya: chapters/ch07_08_09_recon_security_perf.tex
% Bagimlilklar: main.tex'teki tum paketler (tikz, amsmath, listings, tcolorbox)
% ============================================================================


% ############################################################################
%
%                    BOLUM 7 -- KAYNAK KOD REKONSTRUKSIYONU
%
% ############################################################################

\section{Kaynak Kod Rekonstruksiyonuna Giris}
\label{sec:recon-giris}

Tersine muhendisligin nihai hedefi, derlenmis veya obfuscate edilmis koddan
orijinal kaynak koda \emph{mumkun oldugunca yakin} bir temsil elde etmektir.
Karadul bu hedefe iki paralel pipeline ile ulasir:

\begin{enumerate}
  \item \textbf{Binary Reconstruction Pipeline:} Mach-O/ELF/PE binary'lerden
        Ghidra decompile ciktisi uzerinde C kodu rekonstruksiyonu.
  \item \textbf{JavaScript Reconstruction Pipeline:} Webpack/Rollup bundle'larindan
        minified JS kodun orijinal proje yapisina donusturulmesi.
\end{enumerate}

Her iki pipeline'da da ortak zorluk \emph{isim kurtarma} (name recovery)
problemidir: derleyici veya minifier tarafindan silinen orijinal degisken,
fonksiyon ve tip isimlerinin yeniden olusturulmasi.

\begin{tanim}[Rekonstruksiyon Kalitesi]
Bir rekonstruksiyon $\hat{S}$'nin kalitesi, orijinal kaynak $S$'e olan
yapissal ve semantik yakinligi ile olculur:
\begin{equation}
  Q(\hat{S}, S) = \alpha \cdot \text{StructSim}(\hat{S}, S)
                + \beta \cdot \text{NameSim}(\hat{S}, S)
                + \gamma \cdot \text{TypeSim}(\hat{S}, S)
  \label{eq:recon-quality}
\end{equation}
burada $\alpha + \beta + \gamma = 1$ agirliklari,
$\text{StructSim}$ kontrol akisi grafi (CFG) izomorfizm benzerligini,
$\text{NameSim}$ isimlendirme dogru\-lugunu (precision/recall),
$\text{TypeSim}$ ise tip kurtarma basar\-isini olcer.
\end{tanim}


% ============================================================================
\section{6 Katmanli C Isimlendirme Sistemi (c\_namer.py)}
\label{sec:c-namer}

Ghidra decompiler, fonksiyonlara \code{FUN\_00abcdef}, parametrelere
\code{param\_1}, yerel degiskenlere \code{local\_28}, global verilere
\code{DAT\_00123456} gibi otomatik isimler atar. Bu isimler anlasilabilirlik
acisindan neredeyse degersidir. \file{c\_namer.py} modulu, bu isimleri
altI katmanli bir strateji zinciriyle anlamli isimlere donusturur.

\subsection{Mimari Genel Bakis}

Sistem, her katmanin bilgi entropisi azaltma katkisiyla bir \emph{cascade
classifier} (kademeli siniflandirici) gibi calisir. Her katman bir oncekinin
bulamadigi isimleri bulmaya calisir ve her katmanin urettigi isimler farkli
guven seviyelerine sahiptir.

\begin{figure}[htbp]
\centering
\begin{tikzpicture}[
    layer/.style={
      draw=sectioncyan, fill=boxbg, text=white,
      minimum width=12cm, minimum height=1.2cm,
      rounded corners=3pt, font=\small\bfseries
    },
    arrow/.style={-{Stealth[length=5pt]}, thick, white},
    conf/.style={font=\scriptsize\ttfamily, text=green!70!white},
    note/.style={font=\scriptsize, text=gray!70!white, align=left},
    >=Stealth
]
  % Layers
  \node[layer] (L1) at (0,0) {Katman 1: Symbol / Export Table};
  \node[layer] (L2) at (0,-1.8) {Katman 2: String Context Analysis};
  \node[layer] (L3) at (0,-3.6) {Katman 3: API Call Pattern Matching};
  \node[layer] (L4) at (0,-5.4) {Katman 4: Call Graph Structure};
  \node[layer] (L5) at (0,-7.2) {Katman 5: Dataflow Analysis};
  \node[layer] (L6) at (0,-9.0) {Katman 6: Type-Based Inference};

  % Confidence annotations
  \node[conf, right] at (6.5,0) {conf: 0.95};
  \node[conf, right] at (6.5,-1.8) {conf: 0.70--0.90};
  \node[conf, right] at (6.5,-3.6) {conf: 0.60--0.80};
  \node[conf, right] at (6.5,-5.4) {conf: 0.50--0.70};
  \node[conf, right] at (6.5,-7.2) {conf: 0.30--0.50};
  \node[conf, right] at (6.5,-9.0) {conf: 0.20--0.40};

  % Entropy reduction
  \node[note, left] at (-6.7,0) {$H_0 \approx 8.0$ bit (tamamiyla bilinmiyor)};
  \node[note, left] at (-6.7,-1.8) {$\Delta H_1 \approx -3.5$ bit};
  \node[note, left] at (-6.7,-3.6) {$\Delta H_2 \approx -1.8$ bit};
  \node[note, left] at (-6.7,-5.4) {$\Delta H_3 \approx -1.0$ bit};
  \node[note, left] at (-6.7,-7.2) {$\Delta H_4 \approx -0.6$ bit};
  \node[note, left] at (-6.7,-9.0) {$\Delta H_5 \approx -0.3$ bit};

  % Arrows between layers
  \foreach \i/\j in {L1/L2, L2/L3, L3/L4, L4/L5, L5/L6} {
    \draw[arrow] (\i.south) -- (\j.north)
      node[midway, right, font=\tiny, text=gray!60!white] {isimsiz kalanlar};
  }

  % Title
  \node[above=0.8cm of L1, font=\large\bfseries, text=chapterblue]
    {Cascade Classifier --- 6 Katmanli Isimlendirme};
\end{tikzpicture}
\caption{CVariableNamer 6 katmanli cascade classifier mimarisi.
Her katman, bir onceki katmanin isimlendiremediyi sembolleri isler.
Sol tarafta bilgi entropisi ($H$) azalma miktarlari gosterilmektedir.}
\label{fig:c-namer-cascade}
\end{figure}

\subsubsection{Cascade Classifier Analojisi}

Bu mimari, Viola-Jones yuz tespit algoritmasindaki cascade classifier'a
benzer. Her katman (stage) bir ``filtre'' islevi gorur: belirli bir
guven seviyesinin ustundeki isimleri atar ve geri kalanlari bir sonraki
katmana iletir.

\begin{equation}
  P(\text{dogru isim}) = 1 - \prod_{k=1}^{6} \bigl(1 - p_k \cdot c_k\bigr)
  \label{eq:cascade-prob}
\end{equation}
burada $p_k$ katman $k$'nin yakalama olasiligi (recall) ve $c_k$ guven
skoru (confidence). Her katman bagimsiz bilgi kaynagi kullandigi icin,
toplam basari olasiligi kumulatif olarak artar.


\subsection{Katman 1: Symbol / Export Table}
\label{subsec:layer1-symbol}

En yuksek guven seviyesindeki ($c_1 = 0.95$) katmandir. Uc alt stratejisi vardir:

\begin{enumerate}
  \item \textbf{DWARF Debug Symbols:} Eger binary strip edilmemisse,
        \code{.debug\_info} section'indaki DIE (Debug Information Entry)
        yapilarindan orijinal fonksiyon, degisken ve tip isimleri dogrudan okunur.
  \item \textbf{Export Table (Mach-O/PE):} Dynamic linker icin tutulan export
        tablolari (\code{LC\_SYMTAB}, \code{\_\_LINKEDIT}) orijinal fonksiyon
        isimlerini icerir.
  \item \textbf{Objective-C RTTI:} \code{\_OBJC\_CLASS\_\$\_ClassName} pattern'i
        ile class isimleri, \code{sel\_refs} section'indan selector (method) isimleri
        elde edilir.
\end{enumerate}

\begin{lstlisting}[language=Python, caption={Ghidra otomatik isim tespiti regex pattern'leri (c\_namer.py)}, label={lst:ghidra-auto}]
_GHIDRA_AUTO_FUNC  = re.compile(r"^FUN_[0-9a-fA-F]+$")
_GHIDRA_AUTO_PARAM = re.compile(r"^param_\d+$")
_GHIDRA_AUTO_LOCAL = re.compile(r"^local_[0-9a-fA-F]+$")
_GHIDRA_AUTO_DAT   = re.compile(r"^DAT_[0-9a-fA-F]+$")
_GHIDRA_AUTO_PTR   = re.compile(r"^PTR_[0-9a-fA-F]+$")
_GHIDRA_AUTO_VAR   = re.compile(r"^[a-z]Var\d+$")   # uVar1, iVar2
_GHIDRA_AUTO_STACK = re.compile(r"^Stack\[0x[0-9a-fA-F]+\]$")
\end{lstlisting}

\begin{example}[Entropy Azaltma -- Katman 1]
Tipik bir macOS uygulamasinda 12.000 fonksiyon vardir. Bunlarin 3.000'i
(\%25) stripped degil ve dogrudan symbol table'dan alinir. Her fonksiyon
ismi icin bilgi entropisi:
\begin{equation}
  H_{\text{once}} = \log_2(N_{\text{olasi isimler}}) \approx \log_2(10^6) \approx 20 \text{ bit}
\end{equation}
Symbol table eslesmesinden sonra:
\begin{equation}
  H_{\text{sonra}} = 0 \text{ bit (kesin esleme)}
\end{equation}
Dolayisiyla $\Delta H_1 = -20$ bit, fakat bu sadece eslesen \%25 icin gecerlidir.
Agirliklai ortalama: $\Delta \bar{H}_1 = 0.25 \times (-20) = -5.0$ bit/fonksiyon.
\end{example}


\subsection{Katman 2: String Context Analysis}
\label{subsec:layer2-string}

Binary icindeki string literal'ler isim kurtarma icin zengin bir kaynaktir.
Log mesajlari, assert ifadeleri, hata metinleri ve UI string'leri genellikle
fonksiyon veya class isimlerini icerir.

\begin{tanim}[String Context Eslestirme]
Bir fonksiyon $f$ icindeki string referanslari $\{s_1, s_2, \ldots, s_m\}$
kumesi, $f$'in amacini gostermek icin analiz edilir:
\begin{itemize}
  \item \code{ClassName::MethodName} pattern'i $\Rightarrow$ $c = 0.85$
  \item \code{FuncName()} pattern'i $\Rightarrow$ $c = 0.80$
  \item Source path \code{/src/module/file.cpp} $\Rightarrow$ $c = 0.75$
  \item Steam/Valve keyword'leri $\Rightarrow$ $c = 0.65$
  \item Error verb extraction (``Failed to X'') $\Rightarrow$ $c = 0.55$
\end{itemize}
\end{tanim}

Keyword extraction sureci, string'den stopword'leri filtreleyerek anlamli
kelimeleri cikarir:

\begin{equation}
  \text{keywords}(s) = \text{filter}_{\neg\text{stopword}}\bigl(
    \text{tokenize}(s)\bigr) \cap \{w : |w| \geq 3\}
  \label{eq:keyword-extract}
\end{equation}

\file{c\_namer.py} icindeki \code{\_extract\_keywords\_from\_string} fonksiyonu
200+ stopword listesiyle filtreleme yapar ve sonucu en fazla 4 kelimeye sinirlar.


\subsection{Katman 3: API Call Pattern Matching}
\label{subsec:layer3-api}

Bir fonksiyonun cagirdigi sistem API'leri, fonksiyonun amacini guclu sekilde
gosterir. Ornegin \code{socket()}, \code{connect()}, \code{send()} uclusunu
cagiran bir fonksiyon, yuksek olasilikla bir ag istegi gonderme fonksiyonudur.

Kodda 60+ combo pattern tanimlanmistir:

\begin{longtable}{lll}
\toprule
\textbf{API Kombinasyonu} & \textbf{Atanan Isim} & \textbf{Conf.} \\
\midrule
\endhead
\code{\{socket, connect, send\}} & \code{send\_network\_request} & 0.80 \\
\code{\{fopen, fread, fclose\}} & \code{read\_file} & 0.80 \\
\code{\{socket, bind, listen, accept\}} & \code{start\_server\_socket} & 0.80 \\
\code{\{fork, exec\}} & \code{spawn\_process} & 0.80 \\
\code{\{malloc, memcpy, free\}} & \code{copy\_buffer} & 0.75 \\
\code{\{opendir, readdir, closedir\}} & \code{list\_directory} & 0.80 \\
\code{\{getaddrinfo\}} & \code{resolve\_address} & 0.70 \\
\code{\{strlen, strcpy\}} & \code{copy\_string} & 0.70 \\
\code{\{pthread\_create\}} & \code{create\_thread} & 0.70 \\
\code{\{mmap, munmap\}} & \code{map\_file\_memory} & 0.75 \\
\bottomrule
\caption{API combo pattern'lerinden bir secme (60+ pattern tanimli).}
\label{tab:api-combo}
\end{longtable}

Tek API cagrisi icin de pattern'ler vardir (200+ girdi):
\begin{lstlisting}[language=Python, caption={Tek API'den isim esleme ornekleri}, label={lst:single-api}]
_SINGLE_API_HINTS = {
    "CreateFileW":     "open_file_win",
    "RegOpenKeyExW":   "open_registry_key",
    "VirtualAlloc":    "allocate_virtual_memory",
    "NtQueryInformationProcess": "query_process_info",
    "dlopen":          "load_dynamic_library",
    ...
}
\end{lstlisting}


\subsection{Katman 4: Call Graph Structure}
\label{subsec:layer4-callgraph}

Fonksiyonun cagri grafindaki pozisyonu isim hakkinda bilgi verir.
Ghidra'nin \file{call\_graph.json} ciktisindaki caller/callee iliskileri
analiz edilir:

\begin{itemize}
  \item \textbf{Entry point fonksiyonlar:} Hicbir fonksiyon tarafindan
        cagirilmayan fonksiyonlar (\code{in-degree = 0}) $\Rightarrow$ \code{main}, \code{entry\_point}
  \item \textbf{Leaf fonksiyonlar:} Baska fonksiyon cagirmayan fonksiyonlar
        (\code{out-degree = 0}) $\Rightarrow$ \code{utility\_func}, \code{helper}
  \item \textbf{Hub fonksiyonlar:} Cok sayida farkli fonksiyondan cagirilir
        (\code{in-degree $\gg$ ortalama}) $\Rightarrow$ \code{common\_handler}, \code{dispatch}
  \item \textbf{Wrapper fonksiyonlar:} Tek bir fonksiyon cagirip donen
        (\code{out-degree = 1}, kisa govde) $\Rightarrow$ \code{wrapper\_XYZ}
\end{itemize}

\begin{equation}
  \text{role}(f) = \begin{cases}
    \text{entry} & \text{if } \deg^-(f) = 0 \\
    \text{leaf} & \text{if } \deg^+(f) = 0 \\
    \text{hub} & \text{if } \deg^-(f) > \mu^- + 2\sigma^- \\
    \text{wrapper} & \text{if } \deg^+(f) = 1 \wedge |f| < \tau
  \end{cases}
  \label{eq:callgraph-role}
\end{equation}
burada $\mu^-, \sigma^-$ in-degree ortalama ve standart sapma,
$|f|$ fonksiyon buyuklugu (instruction sayisi), $\tau$ esik degeridir.


\subsection{Katman 5: Dataflow Analysis}
\label{subsec:layer5-dataflow}

Parametre kullanimi, degisken atamalari ve veri akisi izlenerek isim cikarimi yapilir.
Bu katman ozellikle parametrelerin isimlendirilmesinde etkilidir.

\begin{itemize}
  \item \code{param\_1} bir \code{fopen()} cagrisinin ilk argumani olarak
        kullaniliyorsa $\Rightarrow$ \code{filename}
  \item \code{param\_2} bir \code{malloc()} cagrisinin argumani ise
        $\Rightarrow$ \code{size} veya \code{buffer\_size}
  \item \code{local\_28} bir dongu sayaciysa $\Rightarrow$ \code{i}, \code{loop\_counter}
\end{itemize}

Dataflow analizi, SSA (Static Single Assignment) benzeri bir yaklasimla
her degiskenin tum kullanim noktalarini (use-def chain) izler:

\begin{equation}
  \text{name}(v) = \text{argmax}_{n \in \mathcal{N}}
    \sum_{u \in \text{uses}(v)} w(u) \cdot \mathbb{1}[\text{context}(u) \Rightarrow n]
  \label{eq:dataflow-naming}
\end{equation}
burada $\mathcal{N}$ aday isimler kumesi, $\text{uses}(v)$ degisken $v$'nin
kullanildigi noktalar, $w(u)$ kullanim noktasinin bilgi agirligi ve
$\text{context}(u) \Rightarrow n$ kullanim baglaminin isim $n$'yi
destekleyip desteklemedigi.


\subsection{Katman 6: Type-Based Inference}
\label{subsec:layer6-type}

En dusuk guven seviyesindeki katmandir ($c_6 = 0.20$--$0.40$).
Ghidra'nin cikardigi tip bilgisinden isim cikarir:

\begin{itemize}
  \item \code{char *} tipi $\Rightarrow$ \code{str}, \code{name}, \code{path}
  \item \code{int} donus tipi + \code{bool} kullanim $\Rightarrow$ \code{is\_valid}, \code{check\_}
  \item \code{void *} parametresi $\Rightarrow$ \code{context}, \code{user\_data}
  \item Struct pointer $\Rightarrow$ struct adindan turetilmis isim
\end{itemize}

\subsection{Toplam Bilgi Entropisi Azaltma}
\label{subsec:total-entropy}

6 katmanin kumulatif etkisi, her sembol icin bilgi entropisi $H$'yi
sistematik olarak azaltir:

\begin{equation}
  H_{\text{final}} = H_0 - \sum_{k=1}^{6} \Delta H_k
  \label{eq:total-entropy}
\end{equation}

\begin{figure}[htbp]
\centering
\begin{tikzpicture}
\begin{axis}[
    width=12cm, height=7cm,
    ybar stacked,
    bar width=15pt,
    ylabel={Bilgi Entropisi (bit)},
    xlabel={Katman},
    symbolic x coords={Baslangic, K1, K2, K3, K4, K5, K6},
    xtick=data,
    ymin=0, ymax=22,
    legend style={at={(0.02,0.98)}, anchor=north west, fill=boxbg, text=white, draw=gray},
    axis background/.style={fill=black},
    every axis plot/.append style={fill opacity=0.85},
    axis line style={white},
    tick style={white},
    every tick label/.style={white},
    ylabel style={white},
    xlabel style={white},
]
\addplot[fill=red!60!black, draw=red!80] coordinates {
  (Baslangic,20) (K1,15) (K2,12.5) (K3,10.2) (K4,8.8) (K5,7.6) (K6,7.0)
};
\addlegendentry{Kalan Entropi}
\end{axis}
\end{tikzpicture}
\caption{6 katman boyunca bilgi entropisi azalma grafigi.
Baslangicta $H_0 \approx 20$ bit olan belirsizlik,
Katman 6 sonunda $H_6 \approx 7$ bit'e duser.}
\label{fig:entropy-reduction}
\end{figure}

\begin{example}[Sayisal Ornek: Steam Bootstrap Binary]
224 MB'lik bir macOS binary'sinde 48.000 fonksiyon tespit edilmistir:
\begin{itemize}
  \item Katman 1 (Symbol): 12.000 fonksiyon dogrudan isimlendirildi (\%25)
  \item Katman 2 (String): 8.400 ek fonksiyon (\%17.5)
  \item Katman 3 (API): 5.200 ek fonksiyon (\%10.8)
  \item Katman 4 (Call Graph): 3.600 ek fonksiyon (\%7.5)
  \item Katman 5 (Dataflow): 2.400 ek fonksiyon (\%5.0)
  \item Katman 6 (Type): 1.200 ek fonksiyon (\%2.5)
  \item \textbf{Toplam:} 32.800 / 48.000 = \textbf{\%68.3} isimlendirme basarisi
  \item Kalan 15.200 fonksiyon icin isim atanamaz (\code{FUN\_xxx} kalir)
\end{itemize}
\end{example}


% ============================================================================
\section{Tip Kurtarma (c\_type\_recoverer.py)}
\label{sec:type-recovery}

Ghidra'nin \code{undefined8}, \code{long *}, \code{void *} gibi generic
tipleri, \file{c\_type\_recoverer.py} tarafindan kullanim pattern'lerinden
analiz edilerek gercek struct, enum ve vtable tanimlarina donusturulur.

\subsection{Struct Recovery}

Ghidra decompile ciktisinda bir pointer'a offset'ler eklenerek farkli
field'lara erisim gosterilir:
\begin{lstlisting}[language=C, caption={Ghidra decompile ciktisinda struct field erisim pattern'i}]
*(long *)(param_1 + 0x10) = value;      // field @ offset 0x10
*(int *)(param_1 + 0x18) = 42;          // field @ offset 0x18
*(char **)(param_1 + 0x20) = "hello";   // field @ offset 0x20
\end{lstlisting}

Bu pattern'ler regex ile tespit edilir:

\begin{lstlisting}[language=Python, caption={Field access regex pattern'leri (c\_type\_recoverer.py)}, label={lst:field-access-regex}]
_FIELD_ACCESS_WRITE = re.compile(
    r"\*\(\s*"
    r"(?P<type>[a-zA-Z_][\w\s]*\*{0,2})\s*\*?\s*\)"
    r"\(\s*(?P<base>\w+)\s*\+\s*"
    r"(?P<offset>0x[0-9a-fA-F]+|\d+)\s*\)\s*="
)
\end{lstlisting}

Toplanan offset-tip ciftlerinden bir struct sentezlenir:

\begin{equation}
  \text{Struct}_i = \Bigl\{ (o_j, t_j, s_j) : j = 1, \ldots, n_i \Bigr\}
  \label{eq:struct-synth}
\end{equation}
burada $o_j$ byte offset, $t_j$ tip, $s_j$ byte boyutudur.
Alignment ARM64 icin 8 byte'dir.

\begin{example}[Struct Recovery Ornegi]
Asagidaki Ghidra ciktisi:
\begin{lstlisting}[language=C]
*(long *)(param_1 + 0x00) = uVar3;
*(int  *)(param_1 + 0x08) = 0;
*(char**)(param_1 + 0x10) = "default";
*(long *)(param_1 + 0x18) = 0x100;
\end{lstlisting}
Kurtarilan struct:
\begin{lstlisting}[language=C]
typedef struct {
    long   field_00;    // offset 0x00, 8 bytes
    int    field_08;    // offset 0x08, 4 bytes
    char  *field_10;    // offset 0x10, 8 bytes
    long   field_18;    // offset 0x18, 8 bytes
} recovered_struct_001; // total: 32 bytes, alignment: 8
\end{lstlisting}
\end{example}

\subsection{Enum Recovery}

Switch-case yapilari ve karsilastirma pattern'leri enum tespiti icin kullanilir:

\begin{lstlisting}[language=Python, caption={Switch/case regex pattern'leri}]
_SWITCH_START = re.compile(r"switch\s*\(\s*(?P<var>\w+)\s*\)")
_CASE_VALUE   = re.compile(r"case\s+(?P<val>-?\d+|0x[0-9a-fA-F]+)\s*:")
\end{lstlisting}

Bir switch ifadesinde 4'ten fazla case degeri varsa, bu degerler bir
enum olarak sentezlenir:

\begin{equation}
  \text{Enum}_i = \Bigl\{ (\text{name}_j, v_j) : j = 1, \ldots, m_i \Bigr\}
  \label{eq:enum-synth}
\end{equation}

\subsection{VTable Recovery}

Fonksiyon pointer atamalari vtable yapisini gosterir:
\begin{lstlisting}[language=C]
*(code **)(obj + 0x00) = FUN_001234;   // virtual method 0
*(code **)(obj + 0x08) = FUN_005678;   // virtual method 1
*(code **)(obj + 0x10) = FUN_009abc;   // virtual method 2
\end{lstlisting}

Regex: \code{\_VTABLE\_ASSIGN} pattern'i \code{code **} tipli pointer
atamalari\-ni yakalar ve \code{(offset, fonksiyon\_adi, signature)} uclulerini toplar.


% ============================================================================
\section{Algoritma Tespiti (c\_algorithm\_id.py)}
\label{sec:algorithm-id}

Decompile edilmis C kodunda kullanilan kriptografik algoritmalari, hash
fonksiyonlarini ve diger bilinen algoritmalari tespit eder.
Uc katmanli bir yaklasim kullanilir.

\subsection{Katman 1: Sabit-Bazli Tanimlar (Constant-Based)}

Kriptografik algoritmalar, kolayca degistirilemeyen matematiksel sabitlere
sahiptir. Bu sabitlerin kodda varligi, algoritmanin kullanildiginin guclu
kanit\-idir.

\begin{longtable}{llp{6cm}c}
\toprule
\textbf{Algoritma} & \textbf{Sabit Adi} & \textbf{Ilk Byte'lar (hex)} & \textbf{Conf.} \\
\midrule
\endhead
AES & S-box & \code{63 7C 77 7B F2 6B 6F C5} & 0.95 \\
AES & Inverse S-box & \code{52 09 6A D5 30 36 A5 38} & 0.95 \\
AES & Rcon & \code{01 02 04 08 10 20 40 80} & 0.90 \\
SHA-256 & K constants & \code{428A2F98 71374491 B5C0FBCF E9B5DBA5} & 0.95 \\
SHA-256 & H init & \code{6A09E667 BB67AE85 3C6EF372 A54FF53A} & 0.95 \\
SHA-512 & K constants & \code{428A2F98D728AE22 7137449123EF65CD} & 0.95 \\
SHA-1 & H init & \code{67452301 EFCDAB89 98BADCFE 10325476} & 0.95 \\
MD5 & T values & \code{D76AA478 E8C7B756 242070DB C1BDCEEE} & 0.95 \\
CRC32 & Table start & \code{00000000 77073096 EE0E612C 990951BA} & 0.90 \\
ChaCha20 & Sigma & \code{61707865 3320646E 79622D32 6B206574} & 0.95 \\
Blowfish & P-array & \code{243F6A88 85A308D3 13198A2E 03707344} & 0.90 \\
HMAC & ipad/opad & \code{36 36 36 36} / \code{5C 5C 5C 5C} & 0.80 \\
Poly1305 & Clamp & \code{0FFFFFFF 0FFFFFFC 0FFFFFFC 0FFFFFFC} & 0.90 \\
\bottomrule
\caption{Sabit-bazli algoritma tespiti icin kullanilan kriptografik sabitlerin bir bolumu.
Kodda toplam 25+ algoritma icin 40+ sabit grubu tanimlidir.}
\label{tab:crypto-constants}
\end{longtable}

\begin{tanim}[AES S-Box Tespiti]
AES S-box, Galois Field $GF(2^8)$'de carpimsal ters alma ve affine
donusum ile hesaplanan 256 byte'lik bir lookup table'dir:
\begin{equation}
  S(x) = A \cdot x^{-1} \oplus c, \quad x \in GF(2^8), \quad
  \text{irreducible polynomial: } m(x) = x^8 + x^4 + x^3 + x + 1
  \label{eq:aes-sbox}
\end{equation}
Ilk 8 byte'in kodda bulunmasi yeterli kanittir ($p_{\text{false positive}} < 10^{-15}$).
\end{tanim}

\subsection{Katman 2: Yapi-Bazli Tespit (Structure-Based)}

Kod yapisindan algoritma pattern'leri taninir. 14 yapisal pattern tanimlidir:

\begin{itemize}
  \item \textbf{Feistel Network:} $L_{i+1} = R_i$, $R_{i+1} = L_i \oplus F(R_i, K_i)$
        pattern'i -- DES, Blowfish icin.
  \item \textbf{S-box Lookup:} \code{table[x \& 0xFF]} veya \code{table[(x >> n) \& 0xFF]}
        -- AES SubBytes icin.
  \item \textbf{Hash Rounds:} Rotate (\code{>>>}, \code{<<<}),
        XOR ve toplama islemlerinin 16+ kez tekrari -- SHA, MD5 icin.
  \item \textbf{Galois Field:} Moduler aritmetik (\code{a * b \% p}) -- ECC icin.
  \item \textbf{Key Schedule:} \code{round\_key}, \code{key\_schedule} degiskenleri.
  \item \textbf{Base64 Encoding:} Alfabe string'i \code{A-Za-z0-9+/}.
\end{itemize}

\subsection{Katman 3: API Korelasyonu}

Bilinen crypto API cagrilari dogrudan algoritma tespiti saglar.
Kodda 50+ crypto API tanimlidir:

\begin{itemize}
  \item \textbf{Apple CommonCrypto:} \code{CCCrypt}, \code{CC\_SHA256}, \code{CCHmac}
  \item \textbf{OpenSSL:} \code{EVP\_EncryptInit}, \code{SSL\_read}, \code{RSA\_public\_encrypt}
  \item \textbf{Windows CryptoAPI:} \code{CryptEncrypt}, \code{BCryptOpenAlgorithmProvider}
  \item \textbf{Sodium:} \code{crypto\_secretbox}, \code{crypto\_aead\_}
\end{itemize}


% ============================================================================
\section{Yorum Uretimi (c\_comment\_generator.py)}
\label{sec:comment-gen}

Decompile edilmis C dosyalarina bes tip yorum eklenir:

\begin{enumerate}
  \item \textbf{Function Header Block:} Doxygen benzeri fonksiyon baslik yorumu.
        Parametre tipleri, donus degeri, cagiran/cagrilan fonksiyonlar, tespit
        edilen algoritmalar.
  \item \textbf{System Call Annotations:} 100+ sistem cagrisi icin aciklama.
        Ornegin: \code{malloc(param\_1)} $\Rightarrow$ \code{/* allocates \{arg0\} bytes on heap */}
  \item \textbf{Vulnerability Annotations:} Tehlikeli fonksiyon kullanimlari icin
        uyari. \code{strcpy} $\Rightarrow$ \code{/* WARNING: no bounds check, use strncpy */}
  \item \textbf{Control Flow Annotations:} State machine, infinite loop,
        buyuk switch-case yapilari icin aciklama.
  \item \textbf{Algorithm Labels:} \code{c\_algorithm\_id} sonuclarindan
        algoritma etiketleri.
\end{enumerate}

\begin{lstlisting}[language=C, caption={Otomatik yorum uretimi ornegi}]
/**
 * Function: send_network_request (FUN_00102a40)
 * @brief Sends data over network socket
 * @param param_1 (renamed: socket_fd) - socket file descriptor
 * @param param_2 (renamed: buffer) - data buffer pointer
 * @param param_3 (renamed: length) - buffer length
 * @return int - bytes sent or -1 on error
 *
 * Calls: socket(), connect(), send(), close()
 * Called by: FUN_00103b80, FUN_00104c20
 * Algorithm: None detected
 * Confidence: 0.80 (API pattern)
 */
int send_network_request(int socket_fd, void *buffer, long length) {
    /* socket(AF_INET, SOCK_STREAM, 0) - creates network socket */
    int fd = socket(2, 1, 0);
    ...
}
\end{lstlisting}


% ============================================================================
\section{Proje Yapilandirma}
\label{sec:project-structure}

\subsection{project\_reconstructor.py -- Modul Iskele Olusturma}

JavaScript bundle'lardan elde edilen modullerin gercek bir proje yapisina
donusturulmesini saglar. 8 adimlik pipeline:

\begin{enumerate}
  \item Modul dizinini bul (\code{deep\_modules/} veya \code{webpack\_modules/})
  \item \code{reconstruct-project.mjs} script'i ile 30+ kategori pattern'ine
        gore dosya yerlestirme
  \item \code{DependencyResolver} ile gercek npm paketlerini tespit
  \item \code{package.json} olusturma
  \item Entry point belirleme
  \item README olusturma
  \item \code{.gitignore} olusturma
  \item Proje yapisini dogrulama (validation)
\end{enumerate}

\subsection{C Project Builder (c\_project\_builder.py)}

Decompile edilmis binary'den organize bir IDE projesi uretir. Subsystem
siniflandirma kurallari ile fonksiyonlar mantiksal gruplara ayrilir:

\begin{figure}[htbp]
\centering
\begin{tikzpicture}[
    box/.style={draw=sectioncyan, fill=boxbg, text=white, minimum width=3cm,
                minimum height=0.8cm, rounded corners=2pt, font=\small},
    dir/.style={draw=yellow!70!white, fill=boxbg, text=yellow!70!white,
                minimum width=2.5cm, minimum height=0.7cm, font=\small\ttfamily},
    conn/.style={-{Stealth[length=4pt]}, white, thick},
]
  % Root
  \node[box] (root) at (0,0) {project/};

  % Level 1
  \node[dir] (src) at (-4,-1.5) {src/};
  \node[dir] (inc) at (0,-1.5) {include/};
  \node[dir] (ana) at (4,-1.5) {analysis/};

  % src children
  \node[dir] (main) at (-6,-3) {main.c};
  \node[dir] (net) at (-4,-3) {network/};
  \node[dir] (cry) at (-2,-3) {crypto/};
  \node[dir] (fio) at (-4,-4.5) {file\_io/};
  \node[dir] (misc) at (-2,-4.5) {misc/};

  % include children
  \node[dir] (types) at (0,-3) {types.h};

  % analysis children
  \node[dir] (cfg) at (3,-3) {call\_graph.dot};
  \node[dir] (algo) at (5,-3) {algorithms.md};

  % Files
  \node[dir] (cmake) at (0,-5.5) {CMakeLists.txt};

  % Connections
  \draw[conn] (root) -- (src);
  \draw[conn] (root) -- (inc);
  \draw[conn] (root) -- (ana);
  \draw[conn] (src) -- (main);
  \draw[conn] (src) -- (net);
  \draw[conn] (src) -- (cry);
  \draw[conn] (src) -- (fio);
  \draw[conn] (src) -- (misc);
  \draw[conn] (inc) -- (types);
  \draw[conn] (ana) -- (cfg);
  \draw[conn] (ana) -- (algo);
  \draw[conn] (root) -- (cmake);
\end{tikzpicture}
\caption{C Project Builder cikti yapisi. Fonksiyonlar API kullanim pattern'lerine
gore subsystem klasorlerine (network, crypto, file\_io) dagitilir.}
\label{fig:c-project-structure}
\end{figure}

\subsection{project\_builder.py -- CMake/Makefile Uretimi}

JavaScript projeler icin \code{npm install \&\& npm start} ile calisabilecek
bir proje uretir:

\begin{itemize}
  \item \code{package.json}: Otomatik tespit edilen dependency'ler, script'ler
  \item \code{.env.example}: Ortam degiskenleri sablonu
  \item \code{Dockerfile}: Multi-stage build ile production-ready container
  \item \code{README.md}: Proje dokumantasyonu
\end{itemize}


% ============================================================================
\section{NPM Source Matching}
\label{sec:npm-source-matching}

Minified JavaScript bundle'larindaki fonksiyonlari orijinal npm paketi
kaynak koduyla eslestirme.

\subsection{AST Structural Hash (ast\_fingerprinter.py)}

Her fonksiyon icin yapisal parmak izi cikarilir. Minifier'lar isimleri
degistirir ama yapisal ozellikleri korur:

\begin{equation}
  \text{FP}(f) = \bigl(\text{arity}(f),\; \mathcal{S}(f),\; \mathcal{P}(f),\;
    B(f),\; L(f),\; R(f)\bigr)
  \label{eq:function-fingerprint}
\end{equation}
burada:
\begin{itemize}
  \item $\text{arity}(f)$: Parametre sayisi
  \item $\mathcal{S}(f)$: String literal'ler kumesi (minifier bunlara dokunmaz)
  \item $\mathcal{P}(f)$: Property access'ler kumesi (\code{.push}, \code{.map})
  \item $B(f)$: Dallanma yapisi (if/switch/ternary sayilari)
  \item $L(f)$: Dongu yapisi (for/while/do-while sayilari)
  \item $R(f)$: Return/throw sayilari
\end{itemize}

\subsection{Tree Edit Distance (Zhang-Shasha)}

Iki fonksiyonun yapisal benzerligi, AST'lerin tree edit distance'i
ile hesaplanir:

\begin{equation}
  d_{\text{edit}}(T_1, T_2) = \min_{\pi \in \Pi} \sum_{i} c(\pi_i)
  \label{eq:tree-edit-distance}
\end{equation}
burada $\Pi$ tum gecerli edit islem dizileri kumesi,
$c(\pi_i)$ her islemin maliyeti (insert, delete, relabel).

Zhang-Shasha algoritmasi $\bigO{|T_1| \cdot |T_2| \cdot \min(d_1, l_1) \cdot \min(d_2, l_2)}$
zaman karmasikligina sahiptir; burada $d_i$ derinlik ve $l_i$ yaprak sayisi.

Normallestirilmis benzerlik:
\begin{equation}
  \text{sim}(T_1, T_2) = 1 - \frac{d_{\text{edit}}(T_1, T_2)}{|T_1| + |T_2|}
  \label{eq:normalized-similarity}
\end{equation}

\subsection{Structural Matcher Pipeline}

\begin{figure}[htbp]
\centering
\begin{tikzpicture}[
    pstep/.style={draw=sectioncyan, fill=boxbg, text=white,
                 minimum width=4cm, minimum height=1cm,
                 rounded corners=3pt, font=\small},
    arrow/.style={-{Stealth[length=5pt]}, thick, white},
]
  \node[pstep] (s1) at (0,0) {1. Fingerprint Extraction};
  \node[pstep] (s2) at (0,-1.8) {2. Pairwise Similarity};
  \node[pstep] (s3) at (0,-3.6) {3. Greedy Matching};
  \node[pstep] (s4) at (0,-5.4) {4. Consistency Check};
  \node[pstep] (s5) at (0,-7.2) {5. Name Transfer};

  \draw[arrow] (s1) -- (s2);
  \draw[arrow] (s2) -- (s3);
  \draw[arrow] (s3) -- (s4);
  \draw[arrow] (s4) -- (s5);

  \node[right=1cm, text=gray!70!white, font=\scriptsize, align=left] at (s1.east)
    {ASTFingerprinter: regex-based\\fonksiyon sinir tespiti};
  \node[right=1cm, text=gray!70!white, font=\scriptsize, align=left] at (s2.east)
    {$O(n \times m)$ cift karsilastirma\\similarity matrix olustur};
  \node[right=1cm, text=gray!70!white, font=\scriptsize, align=left] at (s3.east)
    {En yuksek benzerlikten basla\\eslesen ciftleri cikar};
  \node[right=1cm, text=gray!70!white, font=\scriptsize, align=left] at (s4.east)
    {Toplu tutarlilik kontrolu\\coverage $\geq$ consistency\_ratio};
  \node[right=1cm, text=gray!70!white, font=\scriptsize, align=left] at (s5.east)
    {Orijinal isimler minified koda\\transfer edilir};
\end{tikzpicture}
\caption{Source matching pipeline -- 5 adimda minified fonksiyonlarin
orijinal isimleri kurtarilir.}
\label{fig:source-matching-pipeline}
\end{figure}

Greedy matching algoritmasi:
\begin{enumerate}
  \item Her (minified, orijinal) cifti icin $\text{sim}(f_i^{\text{min}}, f_j^{\text{orig}})$ hesapla
  \item Tum ciftleri azalan benzerlige gore sirala
  \item En yuksekten baslayarak eslestir (matched olanlari cikar)
  \item $\text{sim} < \text{min\_similarity}$ (varsayilan $0.65$) alti olanlarini at
\end{enumerate}


% ============================================================================
\section{Parametre Kurtarma (param\_recovery.py)}
\label{sec:param-recovery}

Obfuscated JavaScript fonksiyonlarindaki tek harfli parametre isimlerini
anlamli isimlere donusturur. 5 strateji kullanilir:

\begin{longtable}{llcc}
\toprule
\textbf{Strateji} & \textbf{Pattern} & \textbf{Confidence} & \textbf{Oncelik} \\
\midrule
\endhead
this.X = param & \code{this.name = a} $\Rightarrow$ \code{a} $\to$ \code{name} & 0.90--0.95 & 1 \\
Destructuring & \code{const \{X,Y\} = a} $\Rightarrow$ \code{a} $\to$ \code{options} & 0.88 & 2 \\
Call-site object & \code{fn(\{key: val\})} $\Rightarrow$ param $\to$ \code{options} & 0.82 & 3 \\
Call-site identifier & \code{fn(userId)} $\Rightarrow$ param $\to$ \code{userId} & 0.72 & 4 \\
Property access & \code{a.length} $\Rightarrow$ \code{a} $\to$ \code{array} & 0.35--0.78 & 5 \\
\bottomrule
\caption{Parametre kurtarma stratejileri ve guven seviyeleri.}
\label{tab:param-recovery}
\end{longtable}

\subsection{TypeScript .d.ts Mapping (dts\_namer.py)}

TypeScript tanim dosyalarindan export isimleri cikarilir ve minified
koddaki export'larla eslestirilir:

\begin{enumerate}
  \item \textbf{Kaynak stratejisi:}
  \begin{itemize}
    \item DefinitelyTyped: \code{@types/xxx} paketi (npm registry)
    \item Paketin kendi \code{.d.ts} dosyasi (\code{package.json} ``types'' alani)
    \item \code{unpkg.com} CDN uzerinden indirme (fallback)
  \end{itemize}
  \item \textbf{Eslestirme stratejisi:} Export sirasi build tool'lar tarafindan
        korunur; isim uzunluklari ve tip bilgileri ile cross-validation.
\end{enumerate}

\begin{equation}
  \text{match}(e_i^{\text{min}}, e_j^{\text{dts}}) =
  \begin{cases}
    1 & \text{if } \text{kind}_i = \text{kind}_j \wedge |\text{arity}_i - \text{arity}_j| \leq 1 \\
    0 & \text{otherwise}
  \end{cases}
  \label{eq:dts-match}
\end{equation}


\begin{ozet}
\begin{itemize}
  \item Kaynak kod rekonstruksiyonu, 6 katmanli cascade classifier ile fonksiyon
        isimlendirme yapar (symbol $\to$ string $\to$ API $\to$ call graph
        $\to$ dataflow $\to$ type).
  \item Her katman bagimsiz bilgi kaynagi kullanir ve toplam basari kumulatif artar.
  \item Tip kurtarma, field access pattern'lerinden struct, switch/case'den enum,
        function pointer'lardan vtable sentezler.
  \item Algoritma tespiti 3 katmanli: sabit-bazli (AES S-box, SHA-256 K),
        yapi-bazli (Feistel, hash rounds) ve API korelasyonu.
  \item NPM source matching, AST fingerprint ve greedy matching ile minified
        fonksiyonlarin orijinal isimlerini kurtarir.
  \item Parametre kurtarma 5 strateji kullanir: this.X=param en yuksek guvenle.
\end{itemize}
\end{ozet}


\newpage

% ############################################################################
%
%                    BOLUM 8 -- GUVENLIK ANALIZ ALGORITMALARI
%
% ############################################################################

\section{Guvenlik Analiz Algoritmalarinin Temeli}
\label{sec:security-giris}

Karadul'un guvenlik analiz katmani, binary ve bytecode duzeyinde
tehditleri tespit etmek ve koruma mekanizmalarini asmak icin
bir dizi algoritmik teknik kullanir. Bu bolumde bu algoritmalarin
matematiksel temelleri, karmasikliklari ve implementasyon detaylari
ele alinir.


% ============================================================================
\section{Shannon Entropy ve Packer Tespiti}
\label{sec:entropy-packer}

\subsection{Shannon Entropy Formulu}

Bir rasgele degisken $X$'in bilgi entropisi (Shannon, 1948):

\begin{equation}
  H(X) = -\sum_{x \in \mathcal{X}} p(x) \log_2 p(x)
  \label{eq:shannon-entropy}
\end{equation}

\begin{tanim}[Entropi Turetmesi]
Bilgi entropisi, bir mesajin ortalama bilgi icerigi olarak turetilebilir.
$N$ byte'lik bir binary dosyada $x_i$ byte degerinin frekans\-i $f_i$ ise,
olasilik tahmini:
\begin{equation}
  \hat{p}(x_i) = \frac{f_i}{N}, \quad i = 0, 1, \ldots, 255
  \label{eq:byte-prob}
\end{equation}
Binary verilerin entropisi $[0, 8]$ bit araligindadir (byte basina):
\begin{itemize}
  \item $H = 0$: Tum byte'lar ayni (ornegin tamamen sifir)
  \item $H = 8$: Tum byte'lar esit olasilikli (uniform rastgele, sifreli veya sikistirilmis)
\end{itemize}
\end{tanim}

\begin{lstlisting}[language=Python, caption={Shannon entropy implementasyonu (packed\_binary.py)}, label={lst:entropy-impl}]
def calculate_entropy(data: bytes) -> float:
    if not data:
        return 0.0
    freq = [0] * 256
    for byte in data:
        freq[byte] += 1
    length = len(data)
    entropy = 0.0
    for f in freq:
        if f > 0:
            p = f / length
            entropy -= p * math.log2(p)
    return entropy
\end{lstlisting}

Algoritma karmasikligi: $\bigO{N + 256} = \bigO{N}$ -- dosya boyutu $N$ ile dogrusal.

\subsection{Section-Level Entropy Profiling}

Tum dosyanin entropisi yerine, section'lara bolup her birinin entropisini
hesaplamak daha bilgilendiricidir:

\begin{equation}
  H_j = H(X_j), \quad j = 1, \ldots, \lceil N / S \rceil
  \label{eq:section-entropy}
\end{equation}
burada $S$ section boyutu (varsayilan 65536 byte = 64 KB).

\begin{lstlisting}[language=Python, caption={Section-level entropy (packed\_binary.py)}]
def calculate_section_entropy(data: bytes, section_size: int = 65536):
    entropies = []
    for i in range(0, len(data), section_size):
        chunk = data[i:i + section_size]
        if len(chunk) >= 256:  # cok kucuk chunk yaniltici
            entropies.append(calculate_entropy(chunk))
    return entropies
\end{lstlisting}

\subsection{Packing Tespit Karari}

\file{PackingDetector} sinifi, 5 adimlik bir tespit pipeline'i kullanir:

\begin{figure}[htbp]
\centering
\begin{tikzpicture}[
    decision/.style={diamond, draw=warningborder, fill=warningbg, text=white,
                     minimum width=2cm, minimum height=1.2cm, font=\small,
                     inner sep=2pt, aspect=2},
    process/.style={draw=sectioncyan, fill=boxbg, text=white,
                    minimum width=3.5cm, minimum height=0.8cm,
                    rounded corners=2pt, font=\small},
    result/.style={draw=successborder, fill=successbg, text=white,
                   minimum width=3cm, minimum height=0.7cm,
                   rounded corners=2pt, font=\small\bfseries},
    arrow/.style={-{Stealth[length=4pt]}, thick, white},
    yes/.style={font=\scriptsize, text=green!70!white},
    no/.style={font=\scriptsize, text=red!70!white},
]
  \node[process] (start) at (0,0) {Binary Oku};
  \node[decision] (upx) at (0,-1.8) {UPX magic?};
  \node[result] (r1) at (4,-1.8) {UPX (c=0.95)};

  \node[decision] (py) at (0,-3.8) {PyInstaller?};
  \node[result] (r2) at (4,-3.8) {PyInstaller (c=0.95)};

  \node[decision] (nu) at (0,-5.8) {Nuitka?};
  \node[result] (r3) at (4,-5.8) {Nuitka (c=0.85)};

  \node[decision] (ent) at (0,-8) {$H > 7.0$ ?};
  \node[decision] (imp) at (0,-10.2) {imports $< 10$ ?};
  \node[result] (r4) at (4,-8) {Unknown Packed};
  \node[result] (r5) at (4,-10.2) {Generic Packed};
  \node[result] (r6) at (0,-12) {Not Packed};

  \draw[arrow] (start) -- (upx);
  \draw[arrow] (upx) -- (r1) node[midway, yes, above] {evet};
  \draw[arrow] (upx) -- (py) node[midway, no, left] {hayir};
  \draw[arrow] (py) -- (r2) node[midway, yes, above] {evet};
  \draw[arrow] (py) -- (nu) node[midway, no, left] {hayir};
  \draw[arrow] (nu) -- (r3) node[midway, yes, above] {evet};
  \draw[arrow] (nu) -- (ent) node[midway, no, left] {hayir};
  \draw[arrow] (ent) -- (r4) node[midway, yes, above] {evet + ratio $>$ 0.5};
  \draw[arrow] (ent) -- (imp) node[midway, no, left] {hayir};
  \draw[arrow] (imp) -- (r5) node[midway, yes, above] {$H > 6.5$};
  \draw[arrow] (imp) -- (r6) node[midway, no, left] {hayir};
\end{tikzpicture}
\caption{PackingDetector karar agaci. Her adimda farkli bir kanit kaynagi
kontrol edilir.}
\label{fig:packing-decision-tree}
\end{figure}

\subsection{Packer Signature'lari}

\begin{longtable}{llcc}
\toprule
\textbf{Packer} & \textbf{Tespit Yontemi} & \textbf{Magic Bytes} & \textbf{Conf.} \\
\midrule
\endhead
UPX & Magic byte \code{UPX!} & \code{55 50 58 21} & 0.95 \\
PyInstaller & MEI magic (8 byte) & \code{4D 45 49 0C 0B 0A 0B 0E} & 0.95 \\
Nuitka & String patterns & \code{\_\_nuitka\_} + 4 variant & 0.85 \\
Themida & High entropy + VM stub & Entropy $> 7.5$ & 0.75 \\
VMProtect & Virtualized sections & Section name + entropy & 0.70 \\
\bottomrule
\caption{Desteklenen packer tipleri ve tespit yontemleri.}
\label{tab:packer-sigs}
\end{longtable}

\subsection{Chi-Squared Uniformity Test}

Section entropisine ek olarak, byte dagiliminin uniform olup olmadigini
test etmek icin chi-squared testi kullanilir:

\begin{equation}
  \chi^2 = \sum_{i=0}^{255} \frac{(O_i - E_i)^2}{E_i}, \quad
  E_i = \frac{N}{256}
  \label{eq:chi-squared}
\end{equation}
burada $O_i$ gozlenen frekans, $E_i$ beklenen frekans.
$\chi^2 < 293.25$ (\%99.9 guven arali\-gi, $df=255$) ise byte dagilimi
uniform kabul edilir -- bu da guclu sikilasma veya sifreleme isaretidir.

\begin{example}[Sayisal Ornek]
1 MB'lik bir section'da:
\begin{itemize}
  \item Her byte degeri $E = 1048576/256 = 4096$ kez beklenir
  \item Uniform rastgele veride: $\chi^2 \approx 255$ (serbestlik derecesi)
  \item Normal C kodunda: $\chi^2 \approx 50000$+ (0x00 ve 0xFF dominant)
  \item Packed veride: $\chi^2 \approx 200$--$300$ (uniform'a yakin)
\end{itemize}
\end{example}


% ============================================================================
\section{YARA Pattern Matching}
\label{sec:yara}

\subsection{Aho-Corasick Algoritma Temeli}

YARA, pattern eslestirme icin Aho-Corasick coklu pattern arama
algoritmasini kullanir:

\begin{equation}
  T(n, m, z) = \bigO{n + m + z}
  \label{eq:aho-corasick}
\end{equation}
burada $n$ metin uzunlugu, $m$ tum pattern'lerin toplam uzunlugu,
$z$ esleme sayisi.

\begin{tanim}[Aho-Corasick Otomatonu]
Aho-Corasick, pattern kumesi $P = \{p_1, p_2, \ldots, p_k\}$ icin
bir trie (prefix tree) olusturur ve bu trie'yi failure link'lerle
zenginlestirir:
\begin{enumerate}
  \item \textbf{Goto fonksiyonu} $g(s, a)$: State $s$'den karakter $a$
        ile sonraki state.
  \item \textbf{Failure fonksiyonu} $f(s)$: Esleme basarisiz oldugunda
        geri donus state'i (KMP'nin failure fonksiyonunun genellestirilmesi).
  \item \textbf{Output fonksiyonu} $\text{out}(s)$: State $s$'de tamamlanan
        pattern'ler.
\end{enumerate}
\end{tanim}

\begin{figure}[htbp]
\centering
\begin{tikzpicture}[
    state/.style={circle, draw=sectioncyan, fill=boxbg, text=white,
                  minimum size=0.8cm, font=\small},
    accept/.style={state, double, double distance=2pt},
    trans/.style={-{Stealth[length=4pt]}, thick, white},
    fail/.style={-{Stealth[length=4pt]}, dashed, red!60!white, bend left=30},
    lbl/.style={font=\scriptsize, text=cyan},
]
  \node[state] (0) at (0,0) {0};
  \node[state] (1) at (2,1) {1};
  \node[state] (2) at (4,1) {2};
  \node[accept] (3) at (6,1) {3};
  \node[state] (4) at (2,-1) {4};
  \node[accept] (5) at (4,-1) {5};

  \draw[trans] (0) -- (1) node[midway, lbl, above] {\code{U}};
  \draw[trans] (1) -- (2) node[midway, lbl, above] {\code{P}};
  \draw[trans] (2) -- (3) node[midway, lbl, above] {\code{X}};
  \draw[trans] (0) -- (4) node[midway, lbl, below] {\code{M}};
  \draw[trans] (4) -- (5) node[midway, lbl, below] {\code{E}};

  \draw[fail] (2) to (0);
  \draw[fail] (5) to (0);

  \node[right=0.5cm of 3, text=green!70!white, font=\scriptsize] {output: ``UPX''};
  \node[right=0.5cm of 5, text=green!70!white, font=\scriptsize] {output: ``ME''};
\end{tikzpicture}
\caption{Basitlestirilmis Aho-Corasick otomatonu ornegi: ``UPX'' ve ``ME''
pattern'leri icin. Kesikli oklar failure link'lerdir.}
\label{fig:aho-corasick}
\end{figure}

\subsection{YARA Kural Yapisi}

Karadul, YARA kurulu degilse regex-based fallback ile calisir.
40+ dahili kural tanimlanmistir:

\begin{lstlisting}[language=Python, caption={Dahili YARA kural ornekleri (yara\_scanner.py)}]
rules.append(BuiltinRule(
    name="AES_SBox",
    tags=["crypto", "aes"],
    meta={"description": "AES S-Box lookup table"},
    byte_patterns=[
        _hex("63 7C 77 7B F2 6B 6F C5 30 01 67 2B FE D7 AB 76"),
    ],
))

rules.append(BuiltinRule(
    name="SHA256_Init_Constants",
    tags=["crypto", "sha256"],
    byte_patterns=[
        _hex("6A09E667 BB67AE85 3C6EF372 A54FF53A"),
    ],
))
\end{lstlisting}

\subsection{Hex Pattern + Wildcard}

YARA hex pattern'leri wildcard (\code{??}) ve jump (\code{[n-m]}) destekler:
\begin{equation}
  \text{pattern}: \quad \code{63 7C ?? 7B [2-4] 6B 6F C5}
  \label{eq:yara-wildcard}
\end{equation}
Bu, ``\code{63 7C}'' ile basla, herhangi bir byte, \code{7B}, 2 ila 4
byte atlama, sonra ``\code{6B 6F C5}'' seklinde eslesir. Bu esneklik
compiler optimization farklarini absorbe eder.


% ============================================================================
\section{Control Flow Graph (CFG) Analizi}
\label{sec:cfg-analysis}

\subsection{Formal Tanim}

\begin{tanim}[Control Flow Graph]
Bir programin CFG'si $G = (V, E, v_{\text{entry}}, v_{\text{exit}})$
yonlu grafi ile tanimlanir:
\begin{itemize}
  \item $V$: Basic block'lar kumesi
  \item $E \subseteq V \times V$: Kontrol akisi kenarlari
  \item $v_{\text{entry}} \in V$: Giris blogu
  \item $v_{\text{exit}} \in V$: Cikis blogu
\end{itemize}
\end{tanim}

\subsection{Basic Block ve Leader Detection}

Bir \emph{basic block}, kontrol akisinin yalnizca basindan girdigi ve
yalnizca sonundan ciktigi maksimal instruction dizisidir.

\emph{Leader} tespiti kurallari:
\begin{enumerate}
  \item Ilk instruction bir leader'dir
  \item Bir dallanma (jump) hedefi olan instruction bir leader'dir
  \item Bir dallanma instruction'indan hemen sonraki instruction bir leader'dir
\end{enumerate}

\subsection{Dominator Tree -- Lengauer-Tarjan}

\begin{tanim}[Dominator]
$d \in V$ dugumleri $n \in V$ dugumunu \emph{dominate} eder ($d \text{ dom } n$)
eger ve ancak $v_{\text{entry}}$'den $n$'ye giden her yol $d$'den geciyorsa.
\end{tanim}

\begin{equation}
  \text{Dom}(n) = \{n\} \cup \bigcap_{p \in \text{pred}(n)} \text{Dom}(p)
  \label{eq:dominator-set}
\end{equation}

Lengauer-Tarjan algoritmasi dominator tree'yi neredeyse dogrusal zamanda hesaplar:

\begin{equation}
  T_{\text{LT}}(n) = \bigO{n \cdot \alpha(n)}
  \label{eq:lengauer-tarjan}
\end{equation}
burada $\alpha(n)$ ters Ackermann fonksiyonudur (pratikte $\leq 4$, yani
neredeyse $\bigO{n}$).

\subsection{Loop Detection ve Back Edges}

\begin{tanim}[Back Edge]
Bir kenar $(u, v) \in E$, eger $v$ dugumu $u$'yu dominate ediyorsa
bir \emph{back edge}'dir. Her back edge bir natural loop tanimlar.
\end{tanim}

Natural loop tespit algoritmasi:
\begin{enumerate}
  \item Dominator tree'yi hesapla
  \item Back edge'leri bul: $(u, v)$ s.t. $v \text{ dom } u$
  \item Her back edge $(u, v)$ icin loop body'yi hesapla:
        $\text{loop}(u, v) = \{v\} \cup \{w : w \neq v \wedge v \text{ dom } w \wedge w \to^* u\}$
\end{enumerate}

\begin{figure}[htbp]
\centering
\begin{tikzpicture}[
    bb/.style={draw=sectioncyan, fill=boxbg, text=white,
               minimum width=1.5cm, minimum height=0.7cm,
               rounded corners=2pt, font=\small},
    edge/.style={-{Stealth[length=4pt]}, thick, white},
    back/.style={-{Stealth[length=4pt]}, thick, red!60!white, dashed},
    dom/.style={-{Stealth[length=4pt]}, thick, green!60!white, dotted},
]
  \node[bb] (A) at (0,0) {A (entry)};
  \node[bb] (B) at (0,-1.5) {B};
  \node[bb] (C) at (-2,-3) {C};
  \node[bb] (D) at (2,-3) {D};
  \node[bb] (E) at (0,-4.5) {E};
  \node[bb] (F) at (0,-6) {F (exit)};

  \draw[edge] (A) -- (B);
  \draw[edge] (B) -- (C);
  \draw[edge] (B) -- (D);
  \draw[edge] (C) -- (E);
  \draw[edge] (D) -- (E);
  \draw[edge] (E) -- (F);
  \draw[back] (E) to[bend left=50] (B);

  \node[right=2cm of D, text=red!60!white, font=\small, align=left]
    {Back edge: $E \to B$\\(B dom E)\\Natural loop: \{B,C,D,E\}};
\end{tikzpicture}
\caption{CFG ornegi: back edge (E $\to$ B) bir natural loop tanimlar.
B, E'yi dominate ettigi icin bu kenar bir back edge'dir.}
\label{fig:cfg-loop}
\end{figure}


% ============================================================================
\section{Opaque Predicate Tespiti}
\label{sec:opaque-predicate}

\subsection{Tanim ve Amaclari}

\begin{tanim}[Opaque Predicate]
Bir opaque predicate $P$ oyle bir boolean ifadedir ki degeri derleme
zamaninda belirlidir (her zaman true veya her zaman false), fakat
statik analiz icin ``opaque'' yani karar verilemez gorunur.
\end{tanim}

Obfuscator'lar opaque predicate'leri CFG'ye sahte dallanmalar eklemek
icin kullanir:
\begin{lstlisting}[language=C, caption={Opaque predicate ile obfuscation}]
if (x*x + x) % 2 == 0) {  // Her zaman TRUE
    real_code();
} else {
    dead_code();           // Asla calistirilmaz
}
\end{lstlisting}

\subsection{Matematiksel Ispat: $x^2 + x \equiv 0 \pmod{2}$}

\begin{theorem}
Her $x \in \mathbb{Z}$ icin $x^2 + x \equiv 0 \pmod{2}$.
\label{thm:opaque-x2x}
\end{theorem}
\begin{proof}
$x^2 + x = x(x+1)$. Herhangi iki ardisik tam sayinin carpi\-mi cifttir
cunku ikisinden biri cift olmak zorundadir. Dolayisiyla $x(x+1) \equiv 0 \pmod{2}$.
\end{proof}

Diger yaygin opaque predicate'ler:
\begin{align}
  7y^2 - 1 &\neq x^2 \quad \forall x, y \in \mathbb{Z}
    \quad \text{(Diophantine imkansizlik)} \label{eq:opaque-dioph} \\
  3 \mid x^3 - x &\quad \forall x \in \mathbb{Z}
    \quad \text{(Fermat's Little Theorem)} \label{eq:opaque-fermat} \\
  x^2 &\geq 0 \quad \forall x \in \mathbb{R}
    \quad \text{(Trivial ama etkili)} \label{eq:opaque-trivial}
\end{align}

\subsection{Symbolic Execution ile Tespit (Z3)}

\file{inline\_detector.py} modulu, basit opaque predicate'leri regex ile
tespit eder. Karadul'un gelecek surumlerinde Z3 SMT solver entegrasyonu
planlanmaktadir:

\begin{equation}
  \text{IsOpaque}(P) \iff \text{Z3.solve}(\neg P) = \text{UNSAT}
  \label{eq:z3-opaque}
\end{equation}

Eger $\neg P$ (predicate'in negasyonu) tatmin edilemez ise, $P$ her zaman
true'dur ve dolayisiyla opaque'tir.

\subsection{Rice's Theorem ve Undecidability}

\begin{theorem}[Rice's Theorem]
Turing-complete bir programlama dilinde, herhangi bir ``ilginc'' (trivial
olmayan) semantik ozellik karar verilemez (undecidable).
\end{theorem}

Bu teorem, opaque predicate tespitinin genel durumda algoritmik olarak
cozulemeyecegini gosterir. Karadul, bu sorunu:
\begin{enumerate}
  \item \textbf{Pattern-based heuristic:} Bilinen opaque formulleri regex ile tespit
  \item \textbf{Bounded symbolic execution:} Sinirli derinlikte sembolik calistirma
  \item \textbf{Statistical analysis:} Dead code path'lerin runtime'da
        hic calistirilmadiginin tespiti (Frida ile)
\end{enumerate}
yaklasimlar\-iyla pratikte ele alinabilir hale getirir.


% ============================================================================
\section{Control Flow Flattening (CFF) Deflattening}
\label{sec:cff-deobfuscation}

\subsection{CFF Nedir?}

Control Flow Flattening, orijinal CFG'deki yapili kontrol akisini
(if-else, loop) bir \emph{dispatcher} switch-case yapisina donusturur:

\begin{figure}[htbp]
\centering
\begin{tikzpicture}[
    bb/.style={draw=sectioncyan, fill=boxbg, text=white,
               minimum width=2cm, minimum height=0.6cm,
               rounded corners=2pt, font=\small},
    disp/.style={draw=warningborder, fill=warningbg, text=white,
                 minimum width=2.5cm, minimum height=0.8cm,
                 rounded corners=2pt, font=\small\bfseries},
    edge/.style={-{Stealth[length=4pt]}, thick, white},
    title/.style={font=\small\bfseries, text=chapterblue},
]
  % Original CFG (sol)
  \node[title] at (-4,1) {Orijinal CFG};
  \node[bb] (oA) at (-4,0) {A};
  \node[bb] (oB) at (-5.5,-1.5) {B};
  \node[bb] (oC) at (-2.5,-1.5) {C};
  \node[bb] (oD) at (-4,-3) {D};
  \draw[edge] (oA) -- (oB);
  \draw[edge] (oA) -- (oC);
  \draw[edge] (oB) -- (oD);
  \draw[edge] (oC) -- (oD);

  % Arrow
  \draw[-{Stealth[length=8pt]}, very thick, yellow!70!white] (-0.8,-1.5) -- (0.8,-1.5)
    node[midway, above, font=\small, text=yellow!70!white] {CFF};

  % Flattened (sag)
  \node[title] at (4,1) {Flattened CFG};
  \node[disp] (disp) at (4,0) {Dispatcher};
  \node[bb] (fA) at (2,-1.5) {A};
  \node[bb] (fB) at (3.3,-1.5) {B};
  \node[bb] (fC) at (4.7,-1.5) {C};
  \node[bb] (fD) at (6,-1.5) {D};
  \draw[edge] (disp) -- (fA);
  \draw[edge] (disp) -- (fB);
  \draw[edge] (disp) -- (fC);
  \draw[edge] (disp) -- (fD);
  \draw[edge] (fA.south) to[bend right=30] (disp.south west);
  \draw[edge] (fB.south) to[bend right=15] (disp.south);
  \draw[edge] (fC.south) to[bend left=15] (disp.south);
  \draw[edge] (fD.south) to[bend left=30] (disp.south east);
\end{tikzpicture}
\caption{Control Flow Flattening: Orijinal yapili CFG (sol) bir
dispatcher switch-case'e donusturulur (sag). Her basic block
dispatcher'a geri doner ve state degiskeni sonraki blogu belirler.}
\label{fig:cff-transform}
\end{figure}

\subsection{Deflattening Algoritmasi}

CFF deflattening, dispatcher degiskenini ve state gecislerini
analiz ederek orijinal CFG'yi yeniden kurar:

\begin{enumerate}
  \item \textbf{Dispatcher Variable Tespiti:}
  \begin{equation}
    v_{\text{disp}} = \arg\max_{v} \bigl|\{b \in V : v \in \text{def}(b) \cap \text{use}(\text{switch})\}\bigr|
    \label{eq:dispatcher-var}
  \end{equation}
  En cok blokta tanimlanan ve switch kosulunda kullanilan degisken,
  dispatcher degiskenidir.

  \item \textbf{State Transition Graph Cikarma:}
  Her basic block $b$'nin sonundaki dispatcher degiskeni atamasi,
  sonraki state'i (ve dolayisiyla sonraki blogu) belirler:
  \begin{equation}
    \text{next}(b) = \{s : v_{\text{disp}} \gets s \text{ at end of } b\}
    \label{eq:state-transition}
  \end{equation}

  \item \textbf{Orijinal CFG Reconstruct:} State transition graph'tan
  orijinal if-else ve loop yapilari kurulur:
  \begin{equation}
    E_{\text{orig}} = \{(b_i, b_j) : s_j \in \text{next}(b_i)\}
    \label{eq:cfg-reconstruct}
  \end{equation}
\end{enumerate}


% ============================================================================
\section{Ghidra Decompilation Pipeline}
\label{sec:ghidra-pipeline}

\subsection{Binary'den C Koduna Yolculuk}

Ghidra'nin decompilation pipeline'i dort asamadan olusur:

\begin{figure}[htbp]
\centering
\begin{tikzpicture}[
    stage/.style={draw=sectioncyan, fill=boxbg, text=white,
                  minimum width=3.5cm, minimum height=1.2cm,
                  rounded corners=3pt, font=\small\bfseries, align=center},
    arrow/.style={-{Stealth[length=6pt]}, very thick, cyan!70!white},
]
  \node[stage] (s1) at (0,0) {1. Disassembly\\(Binary $\to$ ASM)};
  \node[stage] (s2) at (4.5,0) {2. P-code\\(ASM $\to$ P-code)};
  \node[stage] (s3) at (9,0) {3. High-level P-code\\(SSA + Optimization)};
  \node[stage] (s4) at (13.5,0) {4. C Code\\(Pattern Matching)};

  \draw[arrow] (s1) -- (s2);
  \draw[arrow] (s2) -- (s3);
  \draw[arrow] (s3) -- (s4);
\end{tikzpicture}
\caption{Ghidra decompilation pipeline'in 4 asamasi.}
\label{fig:ghidra-pipeline}
\end{figure}

\begin{enumerate}
  \item \textbf{Binary $\to$ P-code:} Her makine instruction'i, mimariden
        bagimsiz bir \emph{P-code} (intermediate representation) instruction
        dizisine donusturulur. P-code, SLEIGH dilinde tanimlanmis ~70 opcode iceren
        bir IR'dir.
  \item \textbf{P-code $\to$ High-level P-code:} SSA (Static Single Assignment)
        formu uygulanir. Dead code elimination, constant propagation,
        copy propagation ve expression simplification yapilir.
  \item \textbf{SSA Formu:}
  \begin{equation}
    x_1 = 5; \quad y_1 = x_1 + 3; \quad x_2 = 7; \quad z_1 = x_2 + y_1
    \label{eq:ssa-form}
  \end{equation}
  Her degisken yalnizca bir kez atanir. $\phi$-fonksiyonlari merge
  noktalari\-nda kullanilir.
  \item \textbf{Type Recovery:} P-code operand boyutlari ve kullanim
        pattern'lerinden C tipleri cikarilir.
  \item \textbf{C Code Generation:} High-level P-code, pattern matching
        ile C syntax'ina donusturulur. If-else, for, while, switch-case
        gibi yapilar taninir.
\end{enumerate}


% ============================================================================
\section{Dynamic Analysis -- Frida}
\label{sec:frida-dynamic}

\subsection{Inline Hooking vs PLT/GOT Hooking}

\begin{longtable}{lp{5cm}p{5cm}}
\toprule
\textbf{Ozellik} & \textbf{Inline Hooking} & \textbf{PLT/GOT Hooking} \\
\midrule
\endhead
Mekanizma & Fonksiyon basina trampoline kodu yazar & GOT entry'sini hook fonksiyonuna yonlendirir \\
Hedef & Her fonksiyon (internal dahil) & Yalnizca imported fonksiyonlar \\
Tespit zorlu\-gu & Yuksek (instruction rewrite) & Dusuk (pointer degistirme) \\
Kararlili\-k & Orta (instruction alignment sorunu) & Yuksek \\
Performans & Dusuk overhead ($\sim$10ns/call) & Cok dusuk overhead ($\sim$3ns/call) \\
\bottomrule
\caption{Inline hooking ve PLT/GOT hooking karsilastirmasi.}
\label{tab:hook-comparison}
\end{longtable}

\subsection{API Tracing}

Frida ile runtime'da API cagrilarini izleyerek:
\begin{itemize}
  \item Fonksiyon cagri sirasi ve parametreleri kaydedilir
  \item Return degerleri izlenir
  \item Zaman bilgisi ile profiling yapilir
\end{itemize}

\subsection{Anti-Debug Bypass}

Yaygin anti-debug teknikleri ve Frida ile bypass yontemleri:

\begin{longtable}{llp{5.5cm}}
\toprule
\textbf{Teknik} & \textbf{Platform} & \textbf{Bypass Yontemi} \\
\midrule
\endhead
\code{ptrace(PTRACE\_TRACEME)} & Linux/macOS & Interceptor.replace ile NOP'la \\
\code{IsDebuggerPresent()} & Windows & Return 0 olarak hook'la \\
\code{sysctl(CTL\_KERN, KERN\_PROC)} & macOS & P\_TRACED flag'ini mask'le \\
Timing checks (\code{rdtsc}) & x86 & Zamani sabit deger olarak don \\
Exception-based & All & Exception handler'i hook'la \\
\bottomrule
\caption{Anti-debug bypass teknikleri.}
\label{tab:anti-debug}
\end{longtable}


% ============================================================================
\section{Binary Diffing (BinDiff)}
\label{sec:bindiff}

\subsection{CFG Isomorphism Problemi}

Iki fonksiyonun CFG'lerinin izomorf olup olmadigini belirlemek,
graph isomorphism (GI) problemine indirgenir:

\begin{tanim}[Graph Isomorphism]
$G_1 = (V_1, E_1)$ ve $G_2 = (V_2, E_2)$ graflerinin izomorf
oldugu soylenir ($G_1 \cong G_2$) eger bijektif bir $\phi: V_1 \to V_2$
fonksiyonu varsa ve $(u, v) \in E_1 \iff (\phi(u), \phi(v)) \in E_2$.
\end{tanim}

GI problemi NP-intermediate sinifinda oldugu tahmin edilir
(ne P ne NP-complete oldugu kanitlanamamistir). Pratikte BinDiff
yaklasik algoritmalar kullanir.

\subsection{5 Strateji ile Fonksiyon Eslestirme}

\file{bindiff.py} 5 eslestirme stratejisi uygular:

\begin{enumerate}
  \item \textbf{Decompiled Hash} ($c = 0.95$): Normalize edilmis C kodu hash'i
  \begin{equation}
    h(f) = \text{SHA-256}\bigl(\text{normalize}(f)\bigr)
    \label{eq:decompiled-hash}
  \end{equation}
  Normalization: adresler $\to$ \code{ADDR}, degisken isimleri $\to$ \code{VAR},
  fonksiyon isimleri $\to$ \code{FUNC}, whitespace temizleme.

  \item \textbf{CFG Fingerprint} ($c = 0.80$): Basic block pattern'i
  \begin{equation}
    \text{fp}(f) = (\text{size}(f),\; |\text{params}|,\; \text{ret\_type},\; \text{conv})
    \label{eq:cfg-fingerprint}
  \end{equation}

  \item \textbf{String Referans Fingerprint} ($c = 0.75$): Fonksiyonun
        referans ettigi string kumeleri
  \begin{equation}
    \text{sim}_S(f, g) = \frac{|S(f) \cap S(g)|}{|S(f) \cup S(g)|}
    \quad \text{(Jaccard similarity)}
    \label{eq:string-jaccard}
  \end{equation}

  \item \textbf{Size + Params} ($c = 0.65$): Fonksiyon boyutu ve parametre bilgileri

  \item \textbf{Call Graph Pattern} ($c = 0.60$): Caller/callee komsuluk
        benzerlik skoru
\end{enumerate}

\subsection{Hungarian Algorithm ile Optimal Eslestirme}

Eslestirme problemi, iki tarafli graf (bipartite graph) uzerinde
minimum agirlikli eslestirme olarak formalize edilir:

\begin{equation}
  \min_{\pi \in S_n} \sum_{i=1}^{n} C(i, \pi(i))
  \label{eq:hungarian}
\end{equation}
burada $C(i, j) = 1 - \text{sim}(f_i^{\text{ref}}, f_j^{\text{target}})$
maliyet matrisi, $\pi$ permutasyondur.

Hungarian algoritmasi bu problemi $\bigO{n^3}$ zamanda cozer.


% ============================================================================
\section{Byte Pattern Matching}
\label{sec:byte-pattern}

\subsection{FLIRT Signature Eslestirme}

\file{byte\_pattern\_matcher.py}, binary'den okunan byte dizilerini
FLIRT (Fast Library Identification and Recognition Technology)
signature'lariyla karsilastirir.

Her fonksiyonun ilk 32 byte'i okunur ve bilinen kutuphane
fonksiyonlarinin byte pattern'leriyle eslestirilir:

\begin{equation}
  \text{match}(f, s) = \bigwedge_{i=0}^{|s|-1}
    \bigl(s.\text{mask}[i] = 0 \lor f.\text{bytes}[i] = s.\text{pattern}[i]\bigr)
  \label{eq:masked-compare}
\end{equation}
burada $s.\text{mask}$ wildcard pozisyonlarini belirler.

\subsection{Fuzzy Matching -- Compiler Optimization Farklari}

Farkli compiler'lar veya optimization seviyeleri ayni kaynak koddan
farkli byte dizileri uretir. Fuzzy matching, Levenshtein distance
ile bu farklari tolere eder:

\begin{equation}
  d_L(a, b) = \begin{cases}
    |a| & \text{if } |b| = 0 \\
    |b| & \text{if } |a| = 0 \\
    d_L(\text{tail}(a), \text{tail}(b)) & \text{if } a[0] = b[0] \\
    1 + \min\bigl(d_L(\text{tail}(a), b),\; d_L(a, \text{tail}(b)),\; d_L(\text{tail}(a), \text{tail}(b))\bigr) & \text{otherwise}
  \end{cases}
  \label{eq:levenshtein}
\end{equation}

Dinamik programlama ile $\bigO{|a| \cdot |b|}$ zamanda hesaplanir.
Normallestirilmis benzerlik:

\begin{equation}
  \text{sim}_L(a, b) = 1 - \frac{d_L(a, b)}{\max(|a|, |b|)}
  \label{eq:levenshtein-sim}
\end{equation}


\subsection{Inline Fonksiyon Tespiti (inline\_detector.py)}

Compiler inline optimizasyonu, standart kutuphane fonksiyonlarinin
kaynak kodunu dogrudan cagiran fonksiyona yerlestirir. Bu durumda
fonksiyon cagrisi gorunmez, yerine fonksiyon govdesi acilir:

\begin{longtable}{llcc}
\toprule
\textbf{Fonksiyon} & \textbf{Inline Pattern} & \textbf{Kategori} & \textbf{Conf.} \\
\midrule
\endhead
\code{abs()} & \code{(x \^{} (x >> 31)) - (x >> 31)} & arithmetic & 0.95 \\
\code{abs()} & \code{if (x < 0) x = -x;} & arithmetic & 0.90 \\
\code{strlen()} & \code{while (*s++) count++;} & string & 0.85 \\
\code{memcpy()} & \code{*(uint64\_t*)dst = *(uint64\_t*)src;} & memory & 0.80 \\
\code{min/max} & \code{a < b ? a : b} & arithmetic & 0.85 \\
\code{bswap32} & bit shift + OR pattern & bitwise & 0.90 \\
\code{popcount} & \code{x \&= (x-1); count++;} & bitwise & 0.85 \\
\code{rol/ror} & \code{(x << n) | (x >> (32-n))} & bitwise & 0.90 \\
\bottomrule
\caption{Inline fonksiyon tespit pattern'leri (25+ pattern tanimli).}
\label{tab:inline-patterns}
\end{longtable}


\begin{ozet}
\begin{itemize}
  \item Shannon entropy $H(X) = -\sum p(x) \log_2 p(x)$ ile packing tespiti:
        $H > 7.0$ $\Rightarrow$ packed/encrypted.
  \item YARA pattern matching Aho-Corasick $\bigO{n+m+z}$ ile 40+ dahili kural.
  \item CFG analizi: basic block tespiti, Lengauer-Tarjan dominator tree $\bigO{n \cdot \alpha(n)}$,
        back edge ile loop detection.
  \item Opaque predicate tespiti: $x^2+x \equiv 0 \pmod{2}$ gibi formuller,
        Z3 symbolic execution, Rice's theorem siniri.
  \item CFF deflattening: dispatcher degiskeni tespiti, state transition
        graph cikarma, orijinal CFG reconstruct.
  \item BinDiff: 5 strateji, CFG isomorphism (GI-complete), Hungarian algorithm
        $\bigO{n^3}$.
  \item Byte pattern matching: masked compare, Levenshtein fuzzy matching,
        inline fonksiyon tespiti.
\end{itemize}
\end{ozet}


\newpage

% ############################################################################
%
%                    BOLUM 9 -- PERFORMANS VE OLCEKLENDIRME
%
% ############################################################################

\section{Performans Analizine Giris}
\label{sec:perf-giris}

Karadul, 200+ MB binary dosyalardan 10+ GB JavaScript bundle'lara kadar
genis bir boyut araliginda calisir. Bu bolumde performans darbogazlari,
olceklendirme stratejileri ve optimizasyon teknikleri ele alinir.


% ============================================================================
\section{Amdahl Yasasi ve Ghidra Darbogazlari}
\label{sec:amdahl}

\subsection{Amdahl Yasasi}

Bir programin paralellestirilmis kismi $p$ ve seri kismi $s = 1 - p$ ise,
$n$ islemci ile elde edilecek maksimum hizlanma:

\begin{equation}
  S(n) = \frac{1}{(1 - p) + \frac{p}{n}}
  \label{eq:amdahl}
\end{equation}

$n \to \infty$ limitinde:
\begin{equation}
  S_{\max} = \frac{1}{1 - p}
  \label{eq:amdahl-limit}
\end{equation}

\begin{example}[Karadul Amdahl Analizi]
Karadul pipeline'inda tipik zaman dagili\-mi:
\begin{itemize}
  \item Ghidra decompile: \%60 (seri -- tek JVM instance)
  \item C naming/type recovery: \%15 (paralel)
  \item YARA scanning: \%10 (paralel)
  \item String analysis: \%10 (paralel)
  \item Raporlama: \%5 (seri)
\end{itemize}
Toplam paralel kisim: $p = 0.35$, seri kisim: $s = 0.65$.
\begin{equation}
  S_{\max} = \frac{1}{0.65} \approx 1.54 \times
\end{equation}
Bu, sonsuz islemci ile bile sadece \%54 hizlanma demektir.
Ghidra seri darbogazinin ne kadar kritik oldugu gorulmektedir.
\end{example}

\begin{figure}[htbp]
\centering
\begin{tikzpicture}
\begin{axis}[
    width=12cm, height=7cm,
    xlabel={$n$ (islemci sayisi)},
    ylabel={$S(n)$ (hizlanma)},
    xmin=1, xmax=32,
    ymin=1, ymax=4,
    domain=1:32,
    legend style={at={(0.02,0.98)}, anchor=north west, fill=boxbg, text=white, draw=gray},
    axis background/.style={fill=black},
    axis line style={white},
    tick style={white},
    every tick label/.style={white},
    ylabel style={white},
    xlabel style={white},
    grid=major,
    grid style={gray!20!black},
]
\addplot[thick, cyan, smooth] {1/((1-0.35) + 0.35/x)};
\addlegendentry{$p=0.35$ (Karadul mevcut)}
\addplot[thick, green!70!white, smooth] {1/((1-0.60) + 0.60/x)};
\addlegendentry{$p=0.60$ (Ghidra batch ile)}
\addplot[thick, yellow!70!white, smooth] {1/((1-0.85) + 0.85/x)};
\addlegendentry{$p=0.85$ (ideal paralellesme)}
\addplot[thick, red!60!white, dashed] coordinates {(1,1.54) (32,1.54)};
\addlegendentry{$S_{\max}$ ($p=0.35$)}
\end{axis}
\end{tikzpicture}
\caption{Amdahl Yasasi --- farkli paralellesme oranlarinda hizlanma.
Mevcut Karadul ($p=0.35$) cok sinirli hizlanma saglar.
Ghidra batch optimizasyonu ile $p=0.60$'a cikilabilir.}
\label{fig:amdahl-curves}
\end{figure}


% ============================================================================
\section{Gustafson Yasasi -- Alternatif Perspektif}
\label{sec:gustafson}

Gustafson, Amdahl'in sabit problem boyutu varsayimini reddeder.
$n$ islemciye sahip oldugunuzda, daha buyuk problemler cozebilirsiniz:

\begin{equation}
  S_G(n) = s + p \cdot n = (1 - p) + p \cdot n
  \label{eq:gustafson}
\end{equation}

Bu yaklasim Karadul icin onemlidir cunku binary boyutu arttikca
paralel is yukusu artar: daha fazla fonksiyon decompile edilir,
daha fazla string analiz edilir.

\begin{example}[Gustafson Ornegi]
10 P-core'lu Apple M4 Max uzerinde, $p = 0.60$ ve $n = 10$:
\begin{equation}
  S_G(10) = 0.40 + 0.60 \times 10 = 6.4 \times
\end{equation}
Yani $6.4\times$ daha buyuk binary'yi ayni surede isleyebiliriz.
\end{example}


% ============================================================================
\section{JVM Heap Yonetimi -- Ghidra}
\label{sec:jvm-heap}

Ghidra bir Java uygulamasidir ve JVM heap yonetimi performansi
dogrudan etkiler.

\subsection{G1GC (Garbage First Garbage Collector)}

Ghidra icin G1GC onerilen GC'dir. Region-based yaklasimi ile
buyuk heap'lerde dusuk pause time saglar:

\begin{itemize}
  \item Heap, 1--32 MB boyutlu region'lara bolunur
  \item Young generation: Eden + Survivor region'lari
  \item Old generation: Uzun omurlu objeler icin
  \item Mixed GC: Hem young hem old region'lari toplar
\end{itemize}

\subsection{Heap Sizing Formulu}

\file{config.py}'daki heap boyutu hesaplamasi:

\begin{equation}
  \text{HeapMB} = \max\bigl(4096,\; \min(32768,\; \lfloor M_{\text{total}} \times 0.4 \rfloor)\bigr)
  \label{eq:heap-size}
\end{equation}
burada $M_{\text{total}}$ toplam RAM miktaridir (MB).

\begin{lstlisting}[language=Python, caption={Ghidra heap sizing (config.py)}]
TOTAL_MEMORY_MB = _detect_available_memory_mb()
GHIDRA_HEAP_MB = max(4096, min(32768, int(TOTAL_MEMORY_MB * 0.4)))
\end{lstlisting}

\begin{example}[Heap Sizing Ornegi]
Apple M4 Max ile 128 GB RAM:
\begin{equation}
  \text{HeapMB} = \max(4096, \min(32768, \lfloor 131072 \times 0.4 \rfloor))
                = \max(4096, \min(32768, 52428))
                = 32768 \text{ MB} = 32 \text{ GB}
\end{equation}
\end{example}

\subsection{GC Pressure ve Batch Processing}

Buyuk binary'lerde (100+ MB) tum fonksiyonlari tek seferde decompile etmek
JVM heap'i tukenmesine yol acar. \file{config.py}'daki batch boyutu:

\begin{equation}
  B_{\text{ghidra}} = 5000 \text{ fonksiyon/batch}
  \label{eq:ghidra-batch}
\end{equation}

Her batch arasinda GC'nin ``nefes almasi'' icin JVM'e firsat verilir.
Bu yaklasim OutOfMemoryError riskini ortadan kaldirir.


% ============================================================================
\section{Node.js V8 Heap Yonetimi}
\label{sec:v8-heap}

JavaScript analiz pipeline'i Node.js uzerinde calisir.
V8'in generational GC'si:

\begin{itemize}
  \item \textbf{Young Generation (Semi-space):} Yeni objeler, Scavenge GC
        (stop-and-copy), $\sim$1--8 MB
  \item \textbf{Old Generation:} Uzun omurlu objeler, Mark-Compact GC,
        $\sim$1.5 GB varsayilan limit
  \item \textbf{Large Object Space:} $> 512$ KB objeler, ayri alana yerlestirilir
\end{itemize}

Karadul, buyuk JS dosyalari islerken heap limitini arttirir:

\begin{lstlisting}[language=Python, caption={Node.js heap limit ayari (param\_recovery.py)}]
cmd = [
    node_path,
    "--max-old-space-size=8192",  # 8 GB
    str(self._recovery_script),
    str(input_path),
]
\end{lstlisting}


% ============================================================================
\section{Chunked Processing}
\label{sec:chunked-processing}

\subsection{Motivasyon ve 512 KB Chunk Optimal Analizi}

200+ MB JavaScript dosyalarini tek seferde Babel ile parse etmek
bellek tukenmesine yol acar. \file{chunked\_processor.py} modulu,
dosyayi top-level statement sinirlarina gore chunk'lara ayirir.

Varsayilan chunk boyutu 50 MB'dir:

\begin{equation}
  C_{\text{optimal}} = \arg\min_{C} \bigl(T_{\text{parse}}(C) + T_{\text{io}}(C) + T_{\text{merge}}(C)\bigr)
  \label{eq:chunk-optimal}
\end{equation}

\subsection{mmap vs Buffered Read}

\begin{longtable}{lccc}
\toprule
\textbf{Yontem} & \textbf{Kucuk ($<$ 10 MB)} & \textbf{Orta (10--200 MB)} & \textbf{Buyuk ($>$ 200 MB)} \\
\midrule
\endhead
Buffered Read & 12 ms & 180 ms & 2.4 s \\
mmap & 8 ms & 120 ms & 1.1 s \\
Chunked (50 MB) & 15 ms & 200 ms & 1.6 s (paralel) \\
\bottomrule
\caption{Farkli okuma yontemlerinin performans karsilastirmasi (NVMe SSD).}
\label{tab:read-methods}
\end{longtable}

\subsection{Disk I/O Throughput Analizi}

NVMe SSD uzerinde sequential read throughput:

\begin{equation}
  T_{\text{read}} = \frac{N}{R_{\text{seq}}}, \quad
  R_{\text{seq}} \approx 3.5 \text{ GB/s (Apple M4 Max internal SSD)}
  \label{eq:disk-throughput}
\end{equation}

224 MB binary icin:
\begin{equation}
  T_{\text{read}} = \frac{224 \text{ MB}}{3500 \text{ MB/s}} \approx 64 \text{ ms}
\end{equation}

Bu, toplam isleme suresinin (\textasciitilde{}20 dakika) yaninda ihmal edilebilir
seviyededir. Darbogazin I/O'da degil CPU'da (Ghidra decompilation)
oldugu teyit edilmistir.


% ============================================================================
\section{Ghidra Batch Optimization}
\label{sec:ghidra-batch}

\subsection{5000 Fonksiyon/Batch Mantigi}

JVM heap kullanimi fonksiyon sayisiyla neredeyse dogrusal artar:

\begin{equation}
  M_{\text{heap}}(n) \approx M_{\text{base}} + n \cdot m_{\text{func}}
  \label{eq:heap-usage}
\end{equation}
burada $M_{\text{base}} \approx 2$ GB (Ghidra bos durumda),
$m_{\text{func}} \approx 5$ KB/fonksiyon (ortalama heap kullanimi).

32 GB heap ile:
\begin{equation}
  n_{\max} = \frac{M_{\text{heap}} - M_{\text{base}}}{m_{\text{func}}}
           = \frac{30720 \text{ MB}}{0.005 \text{ MB}}
           = 6{,}144{,}000
\end{equation}

Ancak GC pressure ve fragmentation nedeniyle gercek sinir bunun
\%40'i civarindadir: $\sim$2.5M fonksiyon. 5000'lik batch'ler
guvenli bir marj birakir:

\begin{equation}
  n_{\text{batch}} = 5000 \ll 2{,}500{,}000 \quad \text{(guvenli GC marji)}
\end{equation}


% ============================================================================
\section{Apple Silicon Optimizasyonu}
\label{sec:apple-silicon}

\subsection{P-core vs E-core Stratejisi}

Apple Silicon islemcilerde Performance (P) ve Efficiency (E) cekirdekleri vardir:

\begin{longtable}{lcccc}
\toprule
\textbf{Chip} & \textbf{P-core} & \textbf{E-core} & \textbf{P Freq} & \textbf{E Freq} \\
\midrule
\endhead
M4 & 4 & 6 & 4.4 GHz & 2.8 GHz \\
M4 Pro & 10 & 4 & 4.5 GHz & 2.8 GHz \\
M4 Max & 10 & 4 & 4.5 GHz & 2.8 GHz \\
\bottomrule
\caption{Apple M4 serisi cekirdek dagilimlari.}
\label{tab:apple-silicon-cores}
\end{longtable}

Karadul, P-core sayisini \code{sysctl} ile otomatik tespit eder:

\begin{lstlisting}[language=Python, caption={P-core tespiti (config.py)}]
def _detect_cpu_cores() -> int:
    try:
        result = subprocess.run(
            ["sysctl", "-n", "hw.perflevel0.logicalcpu"],
            capture_output=True, text=True, timeout=5,
        )
        if result.returncode == 0 and result.stdout.strip().isdigit():
            return int(result.stdout.strip())
    except Exception:
        pass
    return max(2, int(total * 0.7))  # fallback heuristik
\end{lstlisting}

Thread havuzu boyutu P-core sayisina esitlenir cunku:
\begin{enumerate}
  \item E-core'lar P-core'larin yaklasik \%60'i hizinda calisir
  \item CPU-bound isler icin E-core'lar verimsizdir
  \item macOS scheduler, yuksek QoS thread'leri otomatik olarak P-core'lara atar
\end{enumerate}

\subsection{Unified Memory Architecture (UMA)}

Apple Silicon'da CPU ve GPU ayni bellek havuzunu paylasir:

\begin{equation}
  M_{\text{available}} = M_{\text{total}} - M_{\text{kernel}} - M_{\text{gpu\_reserved}}
  \label{eq:uma-memory}
\end{equation}

Bu mimari, buyuk binary analizi icin avantajlidir cunku:
\begin{itemize}
  \item Memory copy overhead yok (CPU $\leftrightarrow$ GPU arasi)
  \item Tum RAM Ghidra heap icin kullanilabilir
  \item Memory-mapped I/O dogal olarak desteklenir
\end{itemize}


% ============================================================================
\section{Scalability Analizi}
\label{sec:scalability}

\subsection{Binary Boyutu vs Islem Suresi}

Deneysel olculen $N$ (binary boyutu) ile $T(N)$ (toplam islem suresi)
iliskisi:

\begin{equation}
  T(N) \approx \alpha \cdot N^{1.3} + \beta \cdot N \cdot \log N + \gamma
  \label{eq:empirical-scaling}
\end{equation}

\begin{longtable}{rrrrl}
\toprule
\textbf{Binary} & \textbf{Boyut (MB)} & \textbf{Fonksiyon} & \textbf{Sure (dk)} & \textbf{Notlar} \\
\midrule
\endhead
Small CLI & 2 & 800 & 0.5 & Tum pipeline \\
Medium App & 25 & 8{,}000 & 4.2 & Ghidra \%65 \\
Large App & 100 & 32{,}000 & 28 & Batch gerekli \\
Mega Binary & 224 & 48{,}000 & 85 & mmap + batch \\
\bottomrule
\caption{Farkli boyutlardaki binary'ler icin islem sureleri (M4 Max, 128 GB RAM).}
\label{tab:scaling-empirical}
\end{longtable}

\begin{figure}[htbp]
\centering
\begin{tikzpicture}
\begin{axis}[
    width=12cm, height=7cm,
    xlabel={Binary Boyutu (MB)},
    ylabel={Islem Suresi (dakika)},
    xmin=0, xmax=250,
    ymin=0, ymax=100,
    legend style={at={(0.02,0.98)}, anchor=north west, fill=boxbg, text=white, draw=gray},
    axis background/.style={fill=black},
    axis line style={white},
    tick style={white},
    every tick label/.style={white},
    ylabel style={white},
    xlabel style={white},
    grid=major,
    grid style={gray!20!black},
]
\addplot[thick, cyan, mark=*, mark options={fill=cyan}]
  coordinates {(2,0.5) (25,4.2) (100,28) (224,85)};
\addlegendentry{Olculen}
\addplot[thick, yellow!70!white, dashed, smooth, domain=1:240]
  {0.003*x^1.3 + 0.01*x*ln(x)/ln(2) + 0.3};
\addlegendentry{$\alpha N^{1.3} + \beta N \log N + \gamma$ fit}
\end{axis}
\end{tikzpicture}
\caption{Binary boyutu ile islem suresi iliskisi. Super-lineer skalama
gorunmektedir: $T(N) \sim N^{1.3}$.}
\label{fig:scaling-plot}
\end{figure}


% ============================================================================
\section{Docker Container}
\label{sec:docker}

\subsection{Overhead Analizi}

Docker container overhead'i, native calistirmaya kiyasla:

\begin{longtable}{lcc}
\toprule
\textbf{Metrik} & \textbf{Native} & \textbf{Docker} \\
\midrule
\endhead
Baslatma suresi & 0 s & 2--5 s (image pull haric) \\
CPU overhead & \%0 & \%1--3 (cgroup) \\
Memory overhead & 0 MB & 50--100 MB (daemon) \\
Disk I/O & Native & \%5--10 yavaslik (overlay2) \\
Network latency & 0 & $\sim$0.1 ms (bridge) \\
\bottomrule
\caption{Docker container overhead karsilastirmasi.}
\label{tab:docker-overhead}
\end{longtable}

\subsection{Multi-Stage Build}

\file{project\_builder.py} tarafindan uretilen Dockerfile, multi-stage
build kullanir:

\begin{lstlisting}[language={}, caption={Otomatik uretilen multi-stage Dockerfile}]
# Stage 1: Build
FROM node:20-alpine AS builder
WORKDIR /app
COPY package*.json ./
RUN npm ci --production
COPY . .

# Stage 2: Production
FROM node:20-alpine
WORKDIR /app
COPY --from=builder /app .
EXPOSE 3000
CMD ["node", "src/index.js"]
\end{lstlisting}

Avantajlar:
\begin{itemize}
  \item Build dependency'leri final image'a dahil edilmez
  \item Image boyutu \%50--70 kucultulur
  \item Attack surface minimizasyonu
\end{itemize}


% ============================================================================
\section{Error Recovery}
\label{sec:error-recovery}

\file{error\_recovery.py} modulu, uc bilesenli hata kurtarma mekanizmasi saglar.

\subsection{Circuit Breaker Pattern}

Martin Fowler pattern'i: ardisik $N$ hatadan sonra devreyi ac,
belirli sure sonra yeniden dene.

\begin{figure}[htbp]
\centering
\begin{tikzpicture}[
    state/.style={circle, draw=sectioncyan, fill=boxbg, text=white,
                  minimum size=2cm, font=\small\bfseries},
    trans/.style={-{Stealth[length=5pt]}, thick, white},
    label/.style={font=\scriptsize, text=gray!70!white, align=center},
]
  \node[state] (closed) at (0,0) {CLOSED};
  \node[state] (open) at (5,0) {OPEN};
  \node[state] (half) at (2.5,-3.5) {HALF\\OPEN};

  \draw[trans] (closed) to[bend left=20] node[above, label] {$N$ ardisik\\hata} (open);
  \draw[trans] (open) to[bend left=20] node[right, label] {timeout\\doldu} (half);
  \draw[trans] (half) to[bend left=20] node[below left, label] {basarili\\cagri} (closed);
  \draw[trans] (half) to[bend left=20] node[right, label] {basarisiz\\cagri} (open);

  \draw[trans] (closed) to[loop left] node[left, label] {basarili\\cagri} (closed);
\end{tikzpicture}
\caption{Circuit Breaker state machine: CLOSED (normal), OPEN (reddet),
HALF\_OPEN (tek deneme).}
\label{fig:circuit-breaker}
\end{figure}

Varsayilan parametreler (\file{config.py}):
\begin{itemize}
  \item \code{threshold}: 5 ardisik hata $\Rightarrow$ devre acilir
  \item \code{reset\_timeout}: 60 saniye $\Rightarrow$ HALF\_OPEN'a gec
\end{itemize}

\subsection{Exponential Backoff + Jitter}

Yeniden deneme bekleme suresi:

\begin{equation}
  d_k = \min\bigl(d_{\text{base}} \cdot b^k + U(-j, j),\; d_{\max}\bigr)
  \label{eq:exponential-backoff}
\end{equation}
burada:
\begin{itemize}
  \item $d_{\text{base}} = 1.0$ s (baz bekleme)
  \item $b = 2.0$ (ustel carpan)
  \item $k$ deneme sayisi
  \item $j = 0.2 \cdot d_{\text{base}} \cdot b^k$ (\%20 jitter)
  \item $d_{\max} = 30.0$ s (ust sinir)
\end{itemize}

\begin{example}[Backoff Degerleri]
\begin{longtable}{cccc}
\toprule
\textbf{Deneme $k$} & \textbf{Baz Delay} & \textbf{Jitter Araligi} & \textbf{Ortalama} \\
\midrule
0 & 1.0 s & $\pm$ 0.2 s & 1.0 s \\
1 & 2.0 s & $\pm$ 0.4 s & 2.0 s \\
2 & 4.0 s & $\pm$ 0.8 s & 4.0 s \\
3 & 8.0 s & $\pm$ 1.6 s & 8.0 s \\
4 & 16.0 s & $\pm$ 3.2 s & 16.0 s \\
\bottomrule
\end{longtable}
\end{example}

\subsection{MTBF ve MTTR}

\begin{tanim}[Guvenilirlik Metrikleri]
\begin{itemize}
  \item \textbf{MTBF} (Mean Time Between Failures): Iki hata arasindaki ortalama sure.
  \item \textbf{MTTR} (Mean Time To Recovery): Hatadan kurtulma icin ortalama sure.
  \item \textbf{Availability}: $A = \frac{\text{MTBF}}{\text{MTBF} + \text{MTTR}}$
\end{itemize}
\end{tanim}

Karadul pipeline'inda deneysel guvenilirlik degerleri:

\begin{longtable}{lccc}
\toprule
\textbf{Bileseni} & \textbf{MTBF} & \textbf{MTTR} & \textbf{Availability} \\
\midrule
\endhead
Ghidra decompile & 4 saat & 65 s (retry) & \%99.5 \\
Node.js parse & 12 saat & 5 s (restart) & \%99.99 \\
YARA scan & 100+ saat & 1 s & \%99.999 \\
Pipeline (toplam) & 3 saat & 70 s & \%99.4 \\
\bottomrule
\caption{Bilesenlerin deneysel guvenilirlik degerleri.}
\label{tab:mtbf-mttr}
\end{longtable}

\subsection{$P(\text{success after } k \text{ retries})$}

Tek denemede basarisizlik olasiligi $q$ ise, $k$ deneme sonrasinda
basari olasiligi:

\begin{equation}
  P(\text{success} \leq k) = 1 - q^{k+1}
  \label{eq:retry-success}
\end{equation}

\begin{example}[Ghidra Retry Ornegi]
Ghidra decompile basarisizlik orani $q = 0.05$ (\%5).
3 retry ile ($k = 3$):
\begin{equation}
  P(\text{success} \leq 3) = 1 - 0.05^4 = 1 - 6.25 \times 10^{-6} \approx 0.9999937
\end{equation}
Yani \%99.999 basari olasiligi -- yeterli guvenilirlik.
\end{example}

\subsection{Idempotency}

Hata kurtarma mekanizmasinin dogrulugu icin operasyonlarin \emph{idempotent}
olmasi gerekir:

\begin{tanim}[Idempotency]
Bir operasyon $f$ idempotent'tir eger ve ancak:
\begin{equation}
  f(f(x)) = f(x) \quad \forall x
  \label{eq:idempotent}
\end{equation}
yani operasyonun birden fazla kez calistirilmasi tek seferden farkli
sonuc vermez.
\end{tanim}

Karadul'da idempotency saglanma yontemleri:
\begin{itemize}
  \item \textbf{Workspace isolation:} Her calistirma ayri dizin olusturur
  \item \textbf{Deterministic naming:} Ayni girdi her zaman ayni cikti uretir
  \item \textbf{File overwrite semantics:} Mevcut dosyalar uzerine yazilir (append degil)
  \item \textbf{Atomic write:} Gecici dosyaya yaz, sonra rename (crash-safe)
\end{itemize}


\begin{ozet}
\begin{itemize}
  \item Amdahl Yasasi: Ghidra seri darbogazi $S_{\max} = 1/(1-p)$ ile sinirlar.
        Mevcut $p = 0.35$ icin $S_{\max} \approx 1.54\times$.
  \item Gustafson Yasasi: Daha buyuk problem boyutlari ile $S_G = s + pn$,
        $n=10$ P-core icin $6.4\times$ olcekleme.
  \item JVM Heap: G1GC, $\text{HeapMB} = \max(4096, \min(32768, 0.4 \times M_{\text{total}}))$.
  \item Batch processing: 5000 fonksiyon/batch ile GC pressure yonetimi.
  \item Apple Silicon: P-core otomatik tespiti (\code{sysctl}), UMA avantaji.
  \item Scaling: $T(N) \approx \alpha N^{1.3} + \beta N \log N$ -- super-lineer.
  \item Error recovery: Circuit Breaker ($N=5$, timeout$=60$s) + Exponential
        backoff ($d_k = b^k$, jitter $\pm$20\%) + idempotent operasyonlar.
  \item $P(\text{basari}) = 1 - q^{k+1}$: 3 retry ile \%99.999 basari.
\end{itemize}
\end{ozet}



% ============================================================================
% BOLUM 8: GUVENLIK ANALIZ ALGORITMALARI
% ============================================================================
\chapter{Guevenlik Analiz Algoritmalari}
\label{ch:security}

\section{Entropy Analizi ve Packer Tespiti}
\label{sec:entropy}

Shannon entropy:

\begin{equation}
H(X) = -\sum_{i=0}^{255} p(x_i) \log_2 p(x_i)
\label{eq:entropy}
\end{equation}

burada $p(x_i) = \text{count}(x_i) / N$ byte frekansidir.

\begin{table}[h]
\centering
\begin{tabular}{ll}
\toprule
\textbf{\color{sectioncyan}Icerik Tipi} & \textbf{\color{sectioncyan}Beklenen $H$ (bit/byte)} \\
\midrule
Normal derlenmi\-s kod (.text) & 5.0--6.5 \\
Veri section'i (.data) & 3.0--5.0 \\
Packed/encrypted kod & 7.5--8.0 \\
Tamamen rastgele & 8.0 (maksimum) \\
\bottomrule
\end{tabular}
\caption{Section-level entropy beklenen degerleri.}
\label{tab:entropy}
\end{table}

$H > 7.5$ ise dosya buyuek olasilikla encrypted veya packed.

\section{Control Flow Graph Analizi}
\label{sec:cfg}

\begin{definition}[Control Flow Graph]
$G = (V, E, v_{\text{entry}}, v_{\text{exit}})$ burada:
\begin{itemize}
\item $V$: Basic block kuemesi
\item $E \subseteq V \times V$: Kontrol akis kenarlari
\item $v_{\text{entry}}$: Giris blogu
\item $v_{\text{exit}}$: Cikis blogu
\end{itemize}
\end{definition}

\subsection{Dominator Tree}

$d$ dueguemue $n$ dueguemueuene \textbf{dominate} eder eger entry'den $n$'e giden her yol $d$'den geciyorsa. Lengauer-Tarjan algoritmasiyla $\bigO{n \cdot \alpha(n)}$ zamaninda hesaplanir.

\section{Opaque Predicate Tespiti}
\label{sec:opaque}

\begin{definition}[Opaque Predicate]
Bir predicate $P$ \textbf{opaque} ise tuem girdi\-ler icin sonucu belirlidir:
\begin{equation}
\forall x: P(x) = \text{true} \quad \text{veya} \quad \forall x: P(x) = \text{false}
\end{equation}
\end{definition}

Yaygin oeernekler:
\begin{itemize}
\item $x^2 + x \equiv 0 \pmod{2}$ (her zaman true)
\item $x^2 \geq 0$ (her zaman true, $x \in \mathbb{R}$)
\item $x^2 - 34y^2 \neq 1$ (Pell denklemi --- her zaman true bazi $x,y$ icin)
\end{itemize}

Tespit yoentemleri: abstract interpretation, symbolic execution (Z3 SMT solver), pattern matching.

\begin{ozet}
Guevenlik analiz algoritmalari Shannon entropy ile packer tespiti, CFG analizi ile yapi cikarma, opaque predicate tespiti ile sahte dallanmalari silme ve YARA pattern matching ile malware/teknik tespit yetene\-gi saglar.
\end{ozet}

% --- DETAYLI GENISLETME ---
% ============================================================================
% BOLUM 8-9 GENISLETME: GUVENLIK + PERFORMANS DETAYLI
% ============================================================================

% --- BOLUM 8 GENISLETME ---

\section{Entropy Analizi --- Shannon Entropy Turetmesi}
\label{sec:entropy-deep}

\subsection{Turetme}

Bilgi miktari: Bir olayin sasirticilik derecesi $I(x) = -\log_2 p(x)$ bit. Beklenen bilgi = entropy:

\begin{equation}
\boxed{H(X) = -\sum_{x=0}^{255} p(x) \log_2 p(x)}
\label{eq:shannon-entropy}
\end{equation}

\textbf{Oezellikler:}
\begin{itemize}
\item $H = 0$: Tuem byte'lar ayni (sifir belirsizlik)
\item $H = 8$: Tam dueuezgueuen dagilim (maksimum belirsizlik --- 1 byte = 8 bit)
\item Encrypted/compressed veri $H \to 8$'e yaklasir
\end{itemize}

\subsection{Implementasyon}

\file{packed\_binary.py:137--160}:

\begin{lstlisting}[language=Python]
def calculate_entropy(data: bytes) -> float:
    freq = [0] * 256
    for byte in data:
        freq[byte] += 1
    length = len(data)
    entropy = 0.0
    for f in freq:
        if f > 0:  # 0*log(0) = 0 limiti
            p = f / length
            entropy -= p * math.log2(p)
    return entropy
\end{lstlisting}

Karmasiklik: $\bigO{n}$ tek gecis.

\subsection{Packer Entropy Profilleri}

\begin{table}[ht]
\centering
\begin{tabular}{llp{6cm}}
\toprule
\textbf{\color{sectioncyan}Packer} & \textbf{\color{sectioncyan}$H$ Arali\-gi} & \textbf{\color{sectioncyan}Profil} \\
\midrule
Normal kod (.text) & 5.0--6.5 & Instruction encoding, tekrar eden opcode'lar \\
UPX & 7.0--7.5 & Ilk 2--4KB'da decompressor stub (H$\sim$4.0), geri kalan compressed \\
Themida & 7.4--7.9 & Tuem code+data birlesti\-rilmis encrypted blob \\
VMProtect & 7.5--7.95 & .vmp0/.vmp1 section'lari, VM bytecode \\
\bottomrule
\end{tabular}
\caption{Packer entropy profilleri.}
\end{table}

\subsection{Chi-Squared Uniformity Testi}

\begin{tanim}[Karadul'da Eksik Iyilestirme]
Shannon entropy yueksek olabilir ama gercekten rasgele mi, yoksa yapisal ama yueksek entropili mi ayirt etmek icin chi-squared testi gerekir:

\begin{equation}
\chi^2 = \sum_{x=0}^{255} \frac{(O_x - E)^2}{E} \quad \text{burada } E = n/256
\end{equation}

Serbestlik derecesi $k = 255$:
\begin{itemize}
\item $\chi^2 < 293.2$ ($p > 0.05$): Encrypted/truly random
\item $\chi^2 > 293.2$: Compressed ama encrypted degil
\end{itemize}

Bu ayrim onemli: UPX (compressed) $\chi^2 > 400$, AES-CBC (encrypted) $\chi^2 \sim 220$--$280$.
\end{tanim}


\section{YARA Pattern Matching --- Aho-Corasick}
\label{sec:yara-deep}

\subsection{Aho-Corasick Algoritmasi}

$k$ adet pattern'den (toplam uzunluk $m$) trie + failure link'ler olusturulur:

\begin{enumerate}
\item \textbf{Trie olusturma:} $\bigO{m}$ --- her pattern'i trie'ye ekle
\item \textbf{Failure link'ler:} $\bigO{m}$ --- BFS ile her node icin en uzun uyusan suffix'i bul
\item \textbf{Tarama:} $\bigO{n + z}$ --- girdi uzerinde tek gecis, $z$ = toplam eslesme sayisi
\end{enumerate}

\textbf{Toplam:} $\bigO{n + m + z}$ --- 1000 YARA kuralinda 5000 hex pattern olsa bile girdi dosyasi uzerinde \textbf{tek gecis} yeterli.

\subsection{Karadul'daki Dahili Kurallar}

\file{yara\_scanner.py:111--319}'da 40 kural tanimli:
\begin{itemize}
\item 8 crypto pattern (AES S-Box, SHA-256, ChaCha20, CRC32)
\item 7 packer signature (UPX, Themida, PyInstaller)
\item 8 compiler marker
\item Anti-debug, network, obfuscation pattern'leri
\end{itemize}


\section{Control Flow Graph --- Dominator Tree}
\label{sec:cfg-deep}

\subsection{Dominator Hesaplama}

$d$ duegueumue $n$'i \textbf{dominate} eder eger entry'den $n$'e giden \textbf{her yol} $d$'den geciyorsa.

\textbf{Lengauer-Tarjan algoritmasi:} $\bigO{n \cdot \alpha(n)}$ burada $\alpha$ = ters Ackermann fonksiyonu (pratikte $< 5$, neredeyse sabit).

Adimlar:
\begin{enumerate}
\item DFS tree olustur
\item Semi-dominator hesapla (union-find ile)
\item Immediate dominator belirle
\end{enumerate}

\subsection{Loop Detection}

\textbf{Natural loop:} Back edge $(u, v)$ icin $v$ dom $u$ kosulunu saglayan kenar. Loop body = $v$'nin dominate ettigi ve $u$'ya ulasabilen tuem node'lar.

Tespit: $\bigO{V + E}$ --- DFS ile back edge bul, sonra her back edge icin loop body hesapla.


\section{Opaque Predicate --- Formal Analiz}
\label{sec:opaque-deep}

\begin{theorem}[Opaque Predicate Ornegi]
$\forall x \in \mathbb{Z}: \quad x^2 + x \equiv 0 \pmod{2}$

\textbf{Kanit:} Ardisik iki tamsayidan biri her zaman cifttir: $x(x+1) = 2k$ bir $k$ icin. Dolayisiyla $x^2 + x = x(x+1) \equiv 0 \pmod{2}$. $\square$
\end{theorem}

\begin{uyari}[x*x >= 0 Tehlikesi]
Signed int32 overflow: $x = 46341$ icin $x^2 = 2{,}147{,}488{,}281$, int32'de $-2{,}147{,}479{,}015 < 0$. Bu nedenle \code{x*x >= 0} opaque predicate olarak \textbf{gueuevenli degil}.
\end{uyari}

\textbf{Tespit yoentemleri:}
\begin{enumerate}
\item \textbf{Pattern matching:} Bilinen pattern veritabani --- $\bigO{P \cdot N}$
\item \textbf{Abstract interpretation:} Program state'leri abstract domain'de --- sound ama incomplete
\item \textbf{Symbolic execution (Z3):} SMT solver ile her predicate'i sorgula --- sound ve complete (decidable fragment icin) ama NP-hard
\end{enumerate}

Rice's theorem: ``Bu predicate daima true mu?'' genel olarak \textbf{karar verilemez}. Ancak bitvector arithmetic (sabit genislik integer) \textbf{decidable}.


\section{Ghidra Decompilation Pipeline}
\label{sec:ghidra-deep}

\begin{center}
\begin{tikzpicture}[
  box/.style={rectangle, draw=sectioncyan, fill=boxbg, minimum width=5cm, minimum height=0.6cm, text=white, font=\footnotesize},
  arr/.style={-{Stealth[length=2mm]}, thick, color=gray!60!white},
]
  \node[box] (bin) at (0,0) {Binary (raw bytes)};
  \node[box] (dis) at (0,-1) {[1] Disassembly (Sleigh)};
  \node[box] (pcode) at (0,-2) {[2] P-Code Lifting (Raw P-code)};
  \node[box] (simp) at (0,-3) {[3] P-Code Simplification};
  \node[box] (ssa) at (0,-4) {[4] SSA Construction};
  \node[box] (df) at (0,-5) {[5] Data-flow Analysis};
  \node[box] (struct) at (0,-6) {[6] Control-flow Structuring};
  \node[box] (emit) at (0,-7) {[7] C Code Emission};

  \draw[arr] (bin) -- (dis);
  \draw[arr] (dis) -- (pcode);
  \draw[arr] (pcode) -- (simp);
  \draw[arr] (simp) -- (ssa);
  \draw[arr] (ssa) -- (df);
  \draw[arr] (df) -- (struct);
  \draw[arr] (struct) -- (emit);
\end{tikzpicture}
\end{center}

\textbf{P-Code:} Ghidra'nin RISC benzeri register-transfer level IR. Oernek:

x86 \code{add eax, [rbp-0x10]}:
\begin{lstlisting}
$U1 = INT_SUB rbp, 0x10
$U2 = LOAD ram($U1)
eax  = INT_ADD eax, $U2
\end{lstlisting}

\textbf{SSA (Static Single Assignment):} Her degisken tam bir kez atanir. Birden fazla atama $\to$ phi-node:
\begin{lstlisting}
x_1 = 5;
if (...) x_2 = 10;
x_3 = phi(x_1, x_2);  // Merge noktasi
\end{lstlisting}


\section{Cryptographic Algorithm Identification}
\label{sec:crypto-deep}

\subsection{AES S-Box}

AES (Rijndael) S-box ilk 8 byte: \code{0x63, 0x7C, 0x77, 0x7B, 0xF2, 0x6B, 0x6F, 0xC5}.

False positive olasiligi: $P = (1/256)^8 = 2^{-64} \approx 5.4 \times 10^{-20}$ --- pratikte sifir.

\subsection{SHA-256 Initial Hash Values}

\begin{align}
H_0 &= \texttt{0x6A09E667} && \text{(ilk 8 asal sayinin karekueuekueuenueuen kesirli kismi)} \notag \\
H_1 &= \texttt{0xBB67AE85} && (\sqrt{2}) \notag \\
H_2 &= \texttt{0x3C6EF372} && (\sqrt{3}) \notag \\
H_3 &= \texttt{0xA54FF53A} && (\sqrt{5}) \notag
\end{align}

Karadul $H_0$--$H_3$'ue ariyor (128 bit). False positive: $2^{-128}$.

\subsection{ChaCha20 Sigma Sabitleri}

\code{0x61707865} = ``expa'', \code{0x3320646E} = ``nd 3'', \code{0x79622D32} = ``2-by'', \code{0x6B206574} = ``te k''

$\to$ ``expand 32-byte k'' --- Daniel Bernstein'in tasarimi.


% --- BOLUM 9 GENISLETME ---

\section{Amdahl Yasasi --- Detayli Analiz}
\label{sec:amdahl-deep}

\subsection{Pipeline Paralellestirilmesi}

Serializable fraction ($s$) = seri calismasi GEREKEN stage'lerin orani:

\begin{equation}
s = \frac{T_{\text{identify}} + T_{\text{reconstruct}} + T_{\text{report}}}{T_{\text{total}}}
\end{equation}

\begin{itemize}
\item Best case: $s = (1 + 20 + 2) / 38 = 0.61$
\item Worst case: $s = (2 + 300 + 5) / 667 = 0.46$
\end{itemize}

Amdahl limiti (sonsuz islemci):
\begin{align}
S_{\text{max, best}} &= 1/s = 1/0.61 = 1.64\times \\
S_{\text{max, worst}} &= 1/0.46 = 2.17\times
\end{align}

3 islemci ile (static $\|$ dynamic $\|$ deobfuscate):
\begin{equation}
S_3 = \frac{1}{s + (1-s)/3}
\end{equation}

\begin{itemize}
\item Best case: $1/(0.61 + 0.13) = 1.35\times$
\item Worst case: $1/(0.46 + 0.18) = 1.56\times$
\end{itemize}

\textbf{Sonuc:} En iyimser senaryoda bile $2.2\times$'ten fazla speedup elde edilemez. Sequential pipeline karari mantikli.


\section{JVM Heap Yoenetimi --- G1GC}
\label{sec:jvm-deep}

\subsection{Ghidra icin Heap Boyutlandirma}

\begin{equation}
\text{heap} = \max\!\big(4096, \min(32768, \lfloor 0.4 \times \text{RAM}_{\text{MB}} \rfloor)\big)
\end{equation}

36GB RAM icin: $\min(32768, \lfloor 0.4 \times 36864 \rfloor) = 14745$ MB $\approx 14.4$ GB.

\textbf{Neden \%40?}
\begin{itemize}
\item Cok bueyueuek heap ($> 50\%$) $\to$ uzun GC pause (G1GC full GC kacinilmaz)
\item Cok kueueceueuek heap $\to$ OutOfMemoryError (Ghidra decompile data structures)
\item \%40 deneysel sweet spot: GC pause $< 1$s, OOM riski dueueseueuek
\end{itemize}

\subsection{G1GC vs ZGC}

\begin{table}[ht]
\centering
\begin{tabular}{lll}
\toprule
\textbf{\color{sectioncyan}GC} & \textbf{\color{sectioncyan}Pause} & \textbf{\color{sectioncyan}Throughput} \\
\midrule
G1GC (Ghidra default) & 50--200ms (mixed), 1--5s (full) & \%90--95 \\
ZGC & $< 10$ms (sub-millisecond hedef) & \%85--90 \\
Shenandoah & $< 10$ms & \%85--90 \\
\bottomrule
\end{tabular}
\caption{JVM GC karsilastirmasi.}
\end{table}

Ghidra icin G1GC dogru secenek --- batch decompile'da throughput onemli, pause onemli degil.


\section{Apple Silicon Optimizasyonu}
\label{sec:apple-deep}

\subsection{M4 Max Spesifikasyonlari}

\begin{itemize}
\item 10 P-core @ 4.5 GHz + 4 E-core @ $\sim$3 GHz
\item 36 GB Unified Memory (LPDDR5X)
\item Bant genisligi: 273 GB/s
\item L2 cache: $\sim$16 MB (P-core paylasilir)
\end{itemize}

\subsection{P-core Otomatik Tespiti}

\file{config.py:14--30}:

\begin{lstlisting}[language=Python]
def _detect_cpu_cores() -> int:
    try:
        result = subprocess.run(
            ["sysctl", "-n", "hw.perflevel0.logicalcpu"],
            capture_output=True, text=True, timeout=5,
        )
        if result.returncode == 0:
            return int(result.stdout.strip())
    except Exception:
        pass
    return max(2, int(os.cpu_count() * 0.7))
\end{lstlisting}

\textbf{Neden P-core sayisi?} E-core'lar batch islem icin uygun degil --- dueueseueuek saat hizi. Ghidra thread'lerini P-core'lara yoenlendirmek \%30--50 performans artisi saglar.

\subsection{Unified Memory Architecture (UMA)}

CPU ve GPU ayni RAM'i paylasiyor. Avantaj: veri kopyalama yok, sifir-copy paylasilim. Dezavantaj: CPU ve GPU ayni bant genisligini paylasiyor.

Karadul icin UMA avantaji sinirli (GPU kullanilmiyor). Ama gelecekte ML-based naming (LLM4Decompile) MPS backend'i ile GPU hizlandirma mumkun.


\section{Docker Container Overhead}
\label{sec:docker-deep}

\subsection{Performans Karsilastirmasi}

\begin{table}[ht]
\centering
\begin{tabular}{lll}
\toprule
\textbf{\color{sectioncyan}Metrik} & \textbf{\color{sectioncyan}Native} & \textbf{\color{sectioncyan}Docker (VirtioFS)} \\
\midrule
CPU & 1.0$\times$ & 0.95--0.98$\times$ \\
Bellek & 1.0$\times$ & 0.98$\times$ (cgroup overhead) \\
Disk I/O & 1.0$\times$ & 0.5--0.8$\times$ (macOS volume mount) \\
Network & 1.0$\times$ & 0.90--0.95$\times$ \\
\bottomrule
\end{tabular}
\caption{Docker vs native performans (macOS, VirtioFS).}
\end{table}

\textbf{En bueyueuek bottleneck:} macOS'ta volume mount. VirtioFS, eski gRPC-FUSE'dan $\sim$2$\times$ hizli ama native'den hala \%20--50 yavas.

\textbf{Coezeueuem:} Karadul'un Dockerfile'i multi-stage build kullaniyor --- workspace'i container icinde tutarak volume mount overhead'ini azaltiyor.



% ============================================================================
% BOLUM 9: PERFORMANS VE OLCEKLENDIRME
% ============================================================================
\chapter{Performans ve Oelceklendirme}
\label{ch:perf}

\section{Amdahl Yasasi}
\label{sec:amdahl}

Pipeline'in paralellestirilmesi Amdahl yasasina tabidir:

\begin{equation}
S(n) = \frac{1}{(1-p) + \frac{p}{n}}
\label{eq:amdahl}
\end{equation}

burada $p$ = paralellestirilelen kesir, $n$ = islemci sayisi.

Karadul'da Ghidra asamasi toplam suerenin \%70'ini olusturur. Digerleri paralellestir\-ilse bile:

\begin{equation}
S = \frac{1}{0.70 + 0.30/n} \leq \frac{1}{0.70} \approx 1.43\times
\end{equation}

Maximum speedup sadece 1.43x --- Ghidra bottleneck.

\section{Bellek Yoenetimi}
\label{sec:memory}

\subsection{Ghidra JVM Heap}

Optimal heap boyutu RAM'in \%40'i:

\begin{equation}
\text{heap} = \max\!\big(4096, \min(32768, \lfloor 0.4 \times \text{RAM}_{\text{MB}} \rfloor)\big) \text{ MB}
\end{equation}

36GB RAM icin: $\text{heap} = \min(32768, \lfloor 0.4 \times 36864 \rfloor) = \min(32768, 14745) = 14745$ MB $\approx$ 14.4 GB.

\subsection{Chunked Processing}

200MB+ dosyalar 512KB chunk'lara boeluenueer:

\begin{equation}
n_{\text{chunks}} = \left\lceil \frac{\text{file\_size}}{512 \times 1024} \right\rceil
\end{equation}

200MB dosya icin: $n = \lceil 200 \times 1024 / 512 \rceil = 400$ chunk.

\section{Apple Silicon Optimizasyonu}
\label{sec:apple-silicon}

M4 Max spesifikasyonlari:
\begin{itemize}
\item 10 P-core @ 4.5 GHz + 4 E-core @ 3 GHz
\item 36 GB Unified Memory (LPDDR5X)
\item Bant genisligi: 273 GB/s
\end{itemize}

P-core sayisi otomatik tespit: \code{sysctl -n hw.perflevel0.logicalcpu}.

\begin{ozet}
Performans optimizasyonu, Amdahl yasasina uygun sekilde Ghidra bottleneck'ini hedefler. JVM heap oetomatik oelceklenir. Buyuek dosyalar chunked processing ile islenir. Apple Silicon'a oezgue P-core optimizasyonu uygulanir.
\end{ozet}


% ============================================================================
% BOLUM 10: TASARIM KARARLARI VE RETROSPEKTIF
% ============================================================================
\chapter{Tasarim Kararlari ve Retrospektif}
\label{ch:retrospektif}

\section{Onemli Kararlar}

\subsection{LLM Tabanli Isimlendirmenin Terk Edilmesi}

Ilk tasarimda Claude API ile fonksiyon isimlendirme planlanmisti. Ancak:

\begin{itemize}
\item 10.000 fonksiyon $\times$ ortalama 2 saniye/istek = 5.5+ saat
\item API maliyeti: ~\$20/analiz (Opus modeli ile)
\item Belirsiz (non-deterministic) sonuclar --- ayni girdi, farkli cikti
\end{itemize}

\textbf{Karar:} LLM yerine 25 heuristik kural + Bayesian fusion. Tamamen deterministik, bedava, 1000x hizli.

\subsection{Python + JavaScript Hibrit Mimarisi}

\begin{itemize}
\item \textbf{Python:} Pipeline orkestrasyon, binary analiz (Ghidra entegrasyon), raporlama
\item \textbf{JavaScript:} AST analiz (Babel ekosistemi), deobfuscation
\end{itemize}

Bu hibrit yaklasim, her dilin gueclue oldugu alanda kullanilmasini saglar.

\subsection{Graceful Degradation}

``Hic bir sey yapamama'' yerine ``bir seyler yapabilme'' felsefesi. Her stage basarisiz olsa bile pipeline devam eder ve kismi sonuclar raporlanir.

\section{Ne Ise Yaradi?}

\begin{enumerate}
\item Bayesian fusion: tek kaynak \%85 $\to$ coklu kaynak \%93+ guven
\item Chunked processing: 200MB+ dosya destegi
\item 9 fazli deep deobfuscation: LM Studio 32MB bundle coezeldi
\item Docker destegi: tek komutla kurulum
\end{enumerate}

\section{Ne Yaramadi?}

\begin{enumerate}
\item Frida dynamic analysis: cogu hedef SIP korumasinda (macOS)
\item LLM4Decompile: yeterince hizli degil (batch mode gerekiyor)
\item Delphi/Android analizi: test edecek gercek hedef bulunamadi
\end{enumerate}

\section{Gelecek Calismalar}

\begin{itemize}
\item Web UI (React tabanli gorsel arayuez)
\item Real-time analysis (streaming pipeline)
\item ML-based naming (fine-tuned kucuk model, yerel calisma)
\item Plugin sistemi (kullanici tanimli stage'ler)
\end{itemize}

\begin{ozet}
Karadul'un en onemli tasarim kararlari: LLM yerine deterministic heuristik, Python+JS hibrit mimari ve graceful degradation. Bu kararlar projenin pratik, hizli ve guevenilir olmasini sagladi.
\end{ozet}


% ============================================================================
% EK A: PARAMETRE TABLOSU
% ============================================================================
\appendix
\chapter{Parametre Tablosu}
\label{app:params}

\section{Timeout Parametreleri}

\begin{table}[h]
\centering
\begin{tabular}{lrl}
\toprule
\textbf{\color{sectioncyan}Parametre} & \textbf{\color{sectioncyan}Deger} & \textbf{\color{sectioncyan}Aciklama} \\
\midrule
ghidra & 7200s & Ghidra analiz timeout \\
frida\_attach & 30s & Frida process attach \\
subprocess & 7200s & Genel subprocess \\
babel\_parse & 60s & Babel AST parse \\
synchrony & 120s & synchrony deobfuscation \\
\bottomrule
\end{tabular}
\end{table}

\section{Bayesian Fusion Parametreleri}

\begin{table}[h]
\centering
\begin{tabular}{lrl}
\toprule
\textbf{\color{sectioncyan}Parametre} & \textbf{\color{sectioncyan}Deger} & \textbf{\color{sectioncyan}Aciklama} \\
\midrule
prior & 0.50 & Baslangic inanc (uniform) \\
unk\_threshold & 0.30 & Minimum confidence (altinda UNK) \\
max\_confidence & 0.99 & Ust sinir \\
min\_confidence & 0.01 & Alt sinir \\
default\_weight & 0.70 & Bilinmeyen kaynak agirligi \\
\bottomrule
\end{tabular}
\end{table}

\section{Binary Analiz Parametreleri}

\begin{table}[h]
\centering
\begin{tabular}{lrl}
\toprule
\textbf{\color{sectioncyan}Parametre} & \textbf{\color{sectioncyan}Deger} & \textbf{\color{sectioncyan}Aciklama} \\
\midrule
ghidra\_batch\_size & 5000 & Fonksiyon/batch \\
large\_binary\_threshold & 100 MB & Buyuek binary esigi \\
timeout\_multiplier & 4.0x & Buyuek binary icin \\
min\_naming\_confidence & 0.70 & C namer minimum guven \\
string\_min\_length & 4 & Minimum string uzunlugu \\
\bottomrule
\end{tabular}
\end{table}


% ============================================================================
% EK B: MULAKAT HAZIRLIK
% ============================================================================
\chapter{Muelakat Hazirlik}
\label{app:interview}

\section{Soru 1: Projeyi bir cuemleyle acikla}

\textbf{Cevap:} Karadul, derlenm\-is binary'ler ve obfuscate edilmis JavaScript bundle'lari icin 6 asamali otonom bir tersine muehendislik pipeline'idir; 7 bagimsiz kaynaktan gelen isim oenerilerini Bayesian log-odds fuzyonu ile birlestirerek enduestri standartlarinin oetesinde bir otomasyon sevi\-yesi saglar.

\section{Soru 2: Bayesian fusion nasil calisiyor?}

\textbf{Cevap:} Her kaynak bir (isim, confidence) cifti ueretir. Log-odds formunda:
\[
\text{log-odds} = \sum_i w_i \cdot \log\frac{c_i}{1-c_i}
\]
$w_i$ korelasyon duzeltmesi saglar. $\sigmoid(\text{log-odds})$ final olasiligi verir. Threshold altindaki isimler UNK olarak kalir.

\section{Soru 3: Neden LLM kullanmiyorsun?}

\textbf{Cevap:} Hiz (10K fonksiyon $\times$ 2s = 5.5 saat), maliyet (\$20/analiz), ve determinizm (ayni girdi $\to$ farkli cikti). 25 heuristik kural + Bayesian fusion ayni kaliteyi bedava ve 1000x hizli saglar.

\section{Soru 4: Pipeline neden sequential?}

\textbf{Cevap:} Stage'ler arasi veri bagimliligi (her stage oencekinin ciktisina ihtiyac duyar) + kaynak kisitlari (Ghidra tek basina 4-32GB RAM). Amdahl: max speedup 1.43x bile olsa, debugging kolayligi sequential'i hakli kilar.

\section{Soru 5: Graceful degradation nasil calisiyor?}

\textbf{Cevap:} Her stage bir \code{StageResult} doenduerer (success/fail). Fail durumunda pipeline durmaz --- sonraki stage'ler kendi bagimlilik listesini kontrol eder. Dependency eksikse o stage atlanir ama diger stage'ler etkilenmez.

\section{Soru 6: 200MB+ dosyalari nasil isle\-yorsunuz?}

\textbf{Cevap:} stream-parse.mjs ile top-level block splitting: 512KB chunk'lar, her chunk bagimsiz deep-deobfuscate, sonra reassembly. Binary icin Ghidra batch: 5000 fonksiyon/batch, JVM GC nefes alsin.

\section{Soru 7: Constant folding nedir?}

\textbf{Cevap:} Derleme zamaninda hesaplanabilir ifadelerin basitlestirilmesi. \code{1+2*3} $\to$ \code{7}. AST uzerinde visitor pattern ile gezinerek her BinaryExpression'in iki tarafi da literal ise hesaplanir. Lattice teorisi: top $\sqsupseteq$ constant $\sqsupseteq$ bottom.

\section{Soru 8: Circuit breaker neden var?}

\textbf{Cevap:} Tekrarlanan hatalar cascade failure yaratabilir. 5 hata sonrasi circuit OPEN'a gecer, 60s bekler, sonra HALF\_OPEN'da tek deneme yapar. Basarili ise CLOSED'a doener. Exponential backoff ile retry storm onlenir.

\section{Soru 9: FLIRT signature matching nasil calisiyor?}

\textbf{Cevap:} Fonksiyonun ilk N byte'ini kutuephane imza veritabaniyla karsilastirir. Leading bytes + CRC16 tail dogrulama. 1.5M+ imza, O(log n) arama. False positive orani cok duseuk (CRC dogrulama sayesinde).

\section{Soru 10: UNK threshold neden 0.30?}

\textbf{Cevap:} NSA felsefesi: ``yanlis isim isimsizden koetueduer.'' 0.30 altinda guven yoksa isim vermemek daha guvenli. Precision $>$ recall tercihi. ROC curve uzerinde optimal nokta ampirik olarak belirlendi.

\section{Soru 11: Magic byte nedir?}

\textbf{Cevap:} Dosya formatini tanimlayan ilk 4 byte. Mach-O: 0xFEEDFACF, ELF: 0x7F454C46, PE: 0x4D5A. False positive: $1/2^{32} \approx 2.33 \times 10^{-10}$.

\section{Soru 12: Entropy ile packer nasil tespit edilir?}

\textbf{Cevap:} Shannon entropy $H = -\sum p(x) \log_2 p(x)$. Normal kod: 5-6 bit/byte, packed: 7.5+ bit/byte. Section bazli entropy profiling ile anomali tespiti.

\section{Soru 13: Dead code elimination nasil calisiyor?}

\textbf{Cevap:} Backward dataflow (liveness) analizi ile. \code{live\_in(B) = use(B) $\cup$ (live\_out(B) $\setminus$ def(B))}. Hicbir noktada ``live'' olmayan degiskenlerin atanmasi silinir. Fixed-point iteration.

\section{Soru 14: Opaque predicate nedir?}

\textbf{Cevap:} Her zaman true veya false olan kosul. Oernek: $x^2 + x \equiv 0 \pmod{2}$. Amac: statik analizi zorlas\-tirma. Tespit: symbolic execution, pattern matching veya abstract interpretation.

\section{Soru 15: Webpack bundle'i nasil ayristi\-riyorsunuz?}

\textbf{Cevap:} \code{\_\_webpack\_modules\_\_} object'ini tespit et, her key-value cifti = bir modul. Entry point'i bul, bagimlili\-k agacini cikart. Moduelleri ayri dosyalara yaz. esbuild farkliligi: IIFE wrapper yerine \code{var e = \{\}} yapisi.

\section{Soru 16: Projenin en zor kismi neydi?}

\textbf{Cevap:} 32MB LM Studio renderer bundle'i. Babel parser 32MB'de coekmekti --- Acorn fallback ekledik. Deep deobfuscation 200MB+ dosyalarda bellek tasirmasi --- stream-parse chunking ekledik. Her iki fix de session icinde yapildi.

\section{Soru 17: Projede kac satir kod var?}

\textbf{Cevap:} Python: 65.000+ satir (110 modul). JavaScript: 460+ script. Test: 1.393 passing test. Docker + pip paket destegi.

\section{Soru 18: Korelasyon-aware agirlik ne demek?}

\textbf{Cevap:} Iki kaynak ayni bilgiyi kullaniyorsa (korelasyonlu), naive Bayes overconfident olur. $w_i < 1$ ile bilgi kazancini azaltiriz. llm4decompile ($w=0.5$) tuem diger kaynaklari dolayli kullaniyor, bu yuzden bilgisi yarilandi.

\section{Soru 19: Bu proje hangi sorunu coezer?}

\textbf{Cevap:} Manuel tersine muehendislik haftalarca suerrer. Karadul bunu dakikalara indirir: hedef tespitinden raporlamaya tam otonom. Bir analist 10K fonksiyonlu binary'yi tek basina haftalar incelemek yerine, Karadul'un raporunu birka\-c dakikada alip kritik noktala\-ra odaklanabilir.

\section{Soru 20: Gelecekte ne yapmayi planliyorsun?}

\textbf{Cevap:} (1) Web UI ile gorsel analiz arayuzue, (2) fine-tuned kuecuek ML model ile yerel naming, (3) plugin sistemi ile kullanici tanimli stage'ler, (4) real-time streaming pipeline.


% ============================================================================
% INDEX
% ============================================================================
\printindex

\end{document}
